\bookpart{共同基础与图像处理}{Common Foundations and Image Processing}
\label{part:image-processing}

\BookPartGuide{%
一张照片可以用于判断对象类别、确定空间范围，也可以用于恢复受损内容或制作新图像。不同任务对输出和允许修改的内容有不同要求。本部分从图像观测与任务目标出发，说明这些要求怎样决定计算过程与评价依据。

概述建立本卷共用的观测、条件与成果地图，并说明图像、视频和多模态任务的联系。视觉与多模态模型章接着说明编码、连接、学习和输出接口，为后续任务建立计算依据。图像识别与检索章研究类别、区域、测量量和匹配关系；图像复原、增强与压缩章研究内容保真、呈现质量与比特成本；图像生成与编辑章研究条件创建、转换、局部修改和新视角合成。各章保留输入、结果和来源之间的对应，区分教学算例与模型实验。}

\chapter{图像和多模态处理概述}
\label{chap:vision-overview}

同一段街道录像可以支持许多不同的请求。交通统计关心经过多少辆车，图像检索希望找回拍到某辆车的其他照片，夜景增强希望看清路边招牌，视频编辑则可能要求把车身改成另一种颜色。这些任务都读取视觉信息，但需要保留的证据、允许改变的内容和应当返回的结果并不相同。仅仅选择一个能够接收图像的模型，还没有完整规定要解决的问题。

前面的模型与范式介绍了表示、预测、生成和学习的通用机制。本卷从应用端继续追问：观测怎样进入计算，模型输出怎样还原成对象、区域、时间区间或媒体成果，怎样判断结果满足了最初的请求？理解本章只需知道训练利用样本调整预测规则，使用时再把这条规则作用于新输入；具体网络尚不是阅读前提。本章先规定观测、任务和依据，紧接着的\ref{chap:vision-multimodal-models}章“视觉与多模态模型”建立共用接口，任务章再逐步加入空间、时间和跨来源的要求。这样的安排使我们先知道需要什么结果，再研究哪些计算能够形成它。

\section{图像和多模态处理的基本概念}
\label{sec:vision-overview-concepts}

描述一个处理问题，需要同时给出观测、成果和约束。观测决定算法实际上看到了什么，成果决定预测变量是什么，约束则规定哪些误差不能接受。把这三项混在一起，容易把“图像变清晰了”误认为“图中文字恢复正确了”，或把“能输出一个框”误认为“找到了请求中的对象”。

以“统计这段街景中转弯的车辆”为例，完整请求还要说明统计哪段时间、哪个路口，以及同一车辆多次出现怎样计数。系统需要取得相关画面，保留拍摄时刻，把局部预测对应到原画面，再判断哪些对象完成了约定的动作，最后按个体而不是按帧计数。这是一条有输入输出依赖的处理过程：没有时刻便难以确定事件区间，没有个体对应便可能重复计数。具体检测和跟踪算法可以不同，这些任务关系却都需要明确。

\subsection{处理对象与观测范围}
\label{sec:vision-overview-observations}

街景照片没有直接存下“汽车”“招牌”这些名称，而是把成像结果记录在离散网格上。\term{像素}（pixel）是网格中的一个位置及其记录值，\term{通道}（channel）则区分该位置记录的不同分量。常见的红绿蓝图像（red, green, blue, RGB）在每个位置记录三个颜色分量，但通道也可以承载深度、光谱或其他测量。为了说明算法接收多少位置、每个位置有多少数值，把图像写为三阶数组
\begin{equation}
  \bm{X}\in\mathbb{R}^{h\times w\times c},
  \qquad x_{i,j,k}=\text{第$i$行、第$j$列、第$k$通道的值。}
  \label{eq:vision-image-array}
\end{equation}
这里$h,w,c$分别是高度、宽度与通道数。数组字形沿用本书的粗体约定，元素值是标量。输入范围也属于定义的一部分：8位整数的$0$至$255$与归一化后的$0$至$1$不能直接混用。

例如，把一幅8位RGB照片交给要求$[0,1]$输入的程序，需要按约定将各分量除以255；这改变了数值尺度，不改变采样位置。若程序要求先减去通道均值，还要记录相应常数。训练和使用必须采用相容的输入约定，否则同一照片会形成不同数值，不能把由接口不一致造成的变化都归因于场景不同。

\term{分辨率}（resolution）在数字图像接口中通常指网格尺寸，例如$1920\times1080$；它决定有多少采样位置，却不能独自说明可辨细节。失焦、高噪声或压缩过重都可能使较大的图像缺少有效信息。镜头、曝光、光照、传感器和后续成像处理影响像素值，同一物体在不同照片中的像素未必相同。把一幅小图插值到更大网格，也没有增加新的实拍观测。

像素坐标只是图上的位置。本卷在没有另行说明时使用横向向右、纵向向下的图像坐标；连续框边界用两个坐标之差计算宽高。若讨论离散像素中心、图像缩放或归一化坐标，会明确给出变换。对象身份用于说明不同位置是否属于同一实体，三维坐标用于说明实体位于哪个空间位置。两辆车可在不同时间经过同一像素，一辆车也可以跨越许多像素，因此这三种描述不能互换。

模型看到的网格还可能只是原图的一部分。若先裁出招牌再缩小，预测的文字区域就处于裁剪图坐标中；返回原图位置时，需要撤销缩放并加回裁剪起点。只保存裁剪后的图片，会丢失它原来位于哪幅照片、哪个位置的依据。由二维位置换到三维位置则是另一类问题，需要投影关系、相机条件及尺度信息，不是把横纵坐标再添一个数就能完成。

\term{视频}（video）把画面与时间关联起来，也就是记录不同拍摄时刻的视觉观测。为了区分画面顺序和实际时间，可表示为$\{(\bm{X}_t,\tau_t)\}_{t=1}^{n}$，其中$n$是帧数，帧编号$t$表示顺序，时间戳$\tau_t$以秒为单位表示拍摄或播放时刻。帧率（frame rate）$f$表示每秒记录多少帧；固定帧率时，相邻帧间隔为$1/f$，抽帧步长为$k$时，被保留的相邻画面间隔为$k/f$。这两个间隔分别属于原始记录和实际输入，比较运动或事件时不能混用。

可变帧率或丢帧时，应使用真实时间戳，不能简单把帧编号当秒数。镜头切换会中断局部运动的连续性，相机移动也会让静止物体在图上移动。例如，相机向右平移时，路灯可在画面中向左移动；只凭这个像素位移，无法断定路灯正在运动。视频的空间变化需要结合相机、对象对应和时间条件解释。

文字、图像和声音还可以共同描述街景中的一个事件。\term{模态}（modality）指信息的类型或通道，例如视觉画面、文字表达和声音记录。\term{多模态处理}（multimodal processing）利用不同类型信息之间的联系完成任务，也就是让一种信息帮助选择、补充或核查另一种信息。画面可以显示车辆的运动，声音可以提供鸣笛时刻，文字可以指定应当寻找哪次转弯；这些联系需要学习或建立对应，不能由“同时输入”自动得到。

同一辆车的多张照片增加的是视角和观测次数，图像与文字配对增加的是信息类型及两者之间的联系。视频虽然含有许多帧，在本卷仍按时序视觉输入展开；页面布局说明文字、图像与表格怎样排列，也不必另立为一种模态。这样的约定用于确定章节职责，不宣称各种文献都采用完全相同的模态划分。

\begin{example}[采样之后还能看见什么]
  \label{ex:vision-overview-sampling}
  以下为构造的教学情形。设视频按区间$[0,12)$秒记录，固定帧率为25帧每秒，含300帧，第$t$帧的时间为$(t-1)/25$秒。保留第1、26、51帧并依此每隔25帧取一帧，得到0至11秒处的12幅画面，间隔为1秒。设一辆车只在4.2秒至4.6秒间出现在画面内，则这些输入中完全没有这辆车。后续模型即使正确处理每幅输入，也不能据此证明原视频中没有车辆出现。
\end{example}

这个例子把信息获取和模型计算分开了。错误可能发生在进入模型之前；增加语言推理长度不能补回从未保留的那段观测。长视频、高分辨率页面和细小目标都需要检查同样的问题。

\subsection{处理目标与输出形式}
\label{sec:vision-overview-outputs}

当请求是“图中有没有公交车”时，结果可以是类别判断；请求变为“公交车在哪里”时，还需要空间范围。\term{目标检测}（object detection）通常返回对象的边界框（bounding box）、类别与得分，边界框用矩形圈出对象所在范围。\term{图像分割}（image segmentation）则为像素指定类别或对象归属，也就是在许多位置分别说明那里是什么或属于谁。框包含的背景多少、两个同类对象是否分开，都影响结果含义。

分割结果可以用掩码（mask）表示，即逐位置记录哪些像素属于目标。例如，一个框可以圈住整辆公交车并包含周围背景，车辆掩码则可以沿轮廓排除背景。若只给每个像素“车辆”类别，两辆相邻车辆仍可能连成同一片区域；若要求分别编辑两辆车，就还需区分它们的个体归属。对象跟踪（object tracking）进一步把不同时刻的范围关联为同一对象的轨迹（trajectory），即同一对象在不同时刻的位置序列，不能用一串没有身份对应的框替代。

数量、关节位置和距离属于测量结果。按检测实例计数需要规定怎样处理边界上的对象；从图像估计距离则需要交代相机、尺度和坐标依据。检索返回满足相关性要求的记录及排序。“外观像同类汽车”和“拍到了同一辆汽车”是不同的相关性，不能用同一组正例默认替代。

空间测量还要区分要测哪个量。姿态（pose）可以描述对象的关节位置或朝向，深度（depth）描述沿约定相机方向的远近；两者既不同于类别，也不同于框的宽度。若没有已知尺度或其他标定依据，单幅照片中的“大车”和“近车”可能给出相近的像素尺寸。语言描述“车离路灯很近”也没有直接给出以米计的距离，测量要求不能由自然语言措辞替代。

另一类请求需要返回图像本身。\term{图像复原}（image restoration）依据退化观测估计原内容，例如从带噪照片估计较少噪声的图像；\term{图像增强}（image enhancement）改善指定用途下的呈现，例如调整暗部层次；\term{图像压缩}（image compression）用较少比特表示图像，并规定解码后的允许误差。三者可以组成一条处理链，但分别回答内容估计、呈现和表示成本的问题。

压缩成果首先是码流（bitstream），也就是按规定组织的一串比特；接收端通过解码把它还原为图像。只比较解码图像是否清晰，遗漏了码流大小；只比较文件是否小，遗漏了失真及解码条件。增强成果则需要按用途核验，例如招牌能否读清、暗部层次是否便于观看；这两项要求可能选择不同的亮度变化。

\term{图像生成}（image generation）根据分布或条件创建图像，\term{图像编辑}（image editing）则对已有图像实施指定变化。复原与生成都可能产生原输入中看不到的细节，区别在于任务要求：复原希望估计原场景，创作可以按要求改变场景。一个去模糊模型把招牌上的字改成另一个清晰的字，即使观看效果很好，也没有完成文字保真的复原目标。

理解和生成因此不是按“输出文字还是输出图片”划分。回答“车身是什么颜色”虽然生成了文字，任务仍是依据观测作出判断；按请求把车身改成蓝色，任务则允许改变原内容。共享一个模型并不取消这一区别，因为两个请求允许使用的依据和允许发生的变化不同。

图\ref{fig:vision-task-map}把同一街景可能形成的成果并列起来。请沿各行比较结果需要对应的语义、范围、身份、尺度或内容条件，而不是把各行理解为互斥的模型类别。共享网络可以接不同输出模块，一个模块预测框，另一个模块组织描述；图中也把描述与编辑、配声放在同一成果分支，但描述仍需观测支持，不因此具有任意创作的权限。

\begin{figure}[htbp]
  \centering
  % 教学任务地图；宽约108mm，采用项目语义配色。
\begin{tikzpicture}[
  font=\figurefont\small,
  task/.style={draw=BookRule,fill=BookPaper,minimum width=30mm,
    minimum height=9mm,align=center,inner sep=2mm},
  link/.style={-{Stealth[length=2mm]},draw=BookMuted,line width=0.8pt},
  note/.style={anchor=west,text=BookMuted,font=\figurefont\footnotesize}
]
  \node[task,fill=FigureInput!10,draw=FigureInput,minimum width=20mm] (input) at (0,0) {街景观测};
  \node[task] (recognition) at (4,2.6) {类别、属性};
  \node[task] (region) at (4,1.3) {框、掩码、轨迹};
  \node[task] (measure) at (4,0) {数量、距离、排序};
  \node[task] (restore) at (4,-1.3) {复原、压缩结果};
  \node[task] (create) at (4,-2.6) {编辑、描述、配声};
  \draw[draw=BookMuted,line width=0.8pt] (input.east) -- (1.8,0);
  \draw[draw=BookMuted,line width=0.8pt] (1.8,-2.6) -- (1.8,2.6);
  \foreach \name/\y in {recognition/2.6,region/1.3,measure/0,restore/-1.3,create/-2.6}{
    \draw[link] (1.8,\y) -- (\name.west);
  }
  \node[note] at (5.8,2.6) {对象与语义};
  \node[note] at (5.8,1.3) {范围与身份};
  \node[note] at (5.8,0) {尺度与相关性};
  \node[note] at (5.8,-1.3) {保真与比特};
  \node[note] at (5.8,-2.6) {条件与对应};
\end{tikzpicture}

  \caption{同一街景观测可以服务不同任务。各行表示一种成果及其主要核验关系，均为教学构造，不代表某个模型的实测输出。图像、视频、文字和声音可以分别进入相应分支；输入相同不意味着监督或评价相同。}
  \label{fig:vision-task-map}
\end{figure}

\subsection{输入依据与结果约束}
\label{sec:vision-overview-evidence}

处理问题中的信息来源有不同职责。照片和录音是观测，人工类别、框和掩码是训练或评价标注，查询说明希望找什么，创作条件说明希望生成什么，相机参数与曝光记录则提供辅助测量。一句“红色汽车正在转弯”若是检索查询，不能同时被当作视频中确有这一事件的证据。

训练时，一条样本可以包含街景图像、车辆范围和类别标注，算法据此比较预测与目标，再调整模型中可学习的参数。使用时，人工正确框不再作为输入，系统必须从新观测形成候选并按任务规则选择。相机参数来自标定，查询来自用户，模型权重来自训练；这些信息的取得过程不同，不应都称为“模型自己知道”。具体误差、参数更新及停止条件由共用模型章和各任务章展开。

配对信息也可能不完整。网页文字常只描述图片的一部分，旁白可能先讲下一步操作，字幕可能来自画外人物。需要把图像编号、区域、时间与来源保留下来，再检查是否对应同一实体或事件。模型得分度量的是给定训练与接口下的匹配程度，不是对配对事实的无条件证明。

一条可回查的结果至少要说明使用了哪份原记录、哪段时间和哪个区域，输入经历了怎样的裁剪或采样。若请求要求“转弯时鸣笛的车”，系统既要找到转弯范围，也要取得对应录音，并核对声音是否能归于这辆车。只有同期鸣笛声而没有声源依据时，可以报告观察到的两件事及其时间关系，不能自动补出同一主体，更不能由先后顺序推出原因。

约束必须落实到可以检查的量。车辆计数可以规定在某时间段越过指定线的个体数；文字读取可以要求完整字符串一致；压缩可以给定每像素比特数；交互分割可以限定用户点击次数。分辨率、允许等待的时间和候选选择预算会改变可用信息，也必须进入比较协议。

以12秒的街景计数为教学例子，若任务只统计在$[0,12)$秒内越过路口入口线的车辆，那么停在路边却未越线的车辆不计入，同一车辆连续出现在20帧中也只计一次。边界时刻、遮挡时的处理及无法判定身份时的输出都应预先约定。这样得到的数值才有确定的统计范围和对象；“检测到20个框”还不能直接当作“经过20辆车”。

形成最终成果也要撤销输入组织造成的坐标变化，并保留不能恢复的部分。局部区域的框需要回到原图，抽帧得到的事件边界需要按时间戳转换，码流需要由规定解码器还原，语言答案需要回指支持它的记录。若采样遗漏了关键时刻，或解码约定不匹配，应说明证据不足或接口错误，不能用一个看似完整的输出掩盖缺口。

数据划分应对应实际使用中的独立单位。如果相邻视频帧随机分入训练与测试，两边可能含有几乎相同的背景和对象，测得的结果就不能充分代表面对新视频的表现。医学图像常需按患者或采集中心划分，工业图像可能需要按批次与设备划分，文档可能需要按来源与模板划分。选用哪一种划分取决于希望估计的未来变化。

若部署目标是另一地区的路口，按地区留出数据比混合各地相邻帧更接近这一检验问题；若目标是在同一路口处理新设备，则又需要按设备或采集批次检查。验证数据用于选择设置，最终测试数据用于估计选定系统面对留出条件的结果，两者职责不能混用。这些划分不保证覆盖全部未来情况，但决定当前评价可以支持哪种外推。

ImageNet-C是分类扰动基准，给既有图像施加规定类型和严重度的合成退化；WILDS是分布变化基准集合，把真实数据中的环境差异组织成训练与评价条件，不同数据集采用不同分组。二者提供的证据不同：在合成噪声上类别错误较少，并不能直接推出跨医院或跨地域也有同样结果。这里由像素退化扩展到数据来源变化，原有结论只能保留在已经检验的退化类型内\scite{hendrycks2019corruptions,koh2021wilds}

\section{图像和多模态处理的研究内容}
\label{sec:vision-overview-tasks}

本卷按主要成果及新增信息条件安排应用。“视觉与多模态模型”章先说明图像、文字和声音怎样进入表示，哪些接口返回可比较的分数、语言响应或媒体。图像任务随后在单幅或多视图观测上形成结果，视频任务增加跨帧对应和信息到达条件，多模态任务再组合查询、证据与多种输出。同一模型可以跨任务使用，章节归属取决于当前主要问题，不取决于模型名称。

\subsection{图像应用的研究内容}
\label{sec:vision-overview-image-tasks}

在工业质检中，“这个零件是否异常”和“异常位于哪里”是两个相互联系的任务。正常样本的局部特征可以给出偏离程度，区域结果还需与具体缺陷对应。若照片本身严重模糊，复原可以改善可见性，但生成出的划痕不能作为真实缺陷的新证据。图像识别、区域判断与复原应在相同对象上保留来源关系。

在这组问题中，识别把局部和整体线索变成类别、属性或异常判断，检测和分割把判断落到范围，姿态和空间恢复再形成关节、相机、深度或三维结构。空间恢复（spatial reconstruction）利用图像等观测估计场景的空间组织，不等于任意创建一个好看的三维场景。\ref{chap:vision-image-recognition}章沿这些成果进入测量及匹配，解释每一步需要怎样的监督、坐标和尺度依据。

照片整理提供另一组联系。场景类别可以先缩小候选，再由整图表示找相似照片，局部对应与几何核验进一步检查是否拍到了同一建筑。按“街道”“夜景”整理是类别检索，按“红色车身”整理是属性检索，找回同一辆车或同一建筑则是实例检索（instance retrieval），即把具体实体而不是同类外观作为相关性依据。最后一种请求需要保留能够辨认个体的细节，不能只靠场景名称。

\ref{chap:vision-image-restoration}章分别研究退化观测中的内容估计、指定用途下的呈现及有限比特表示。\ref{chap:vision-image-generation}章研究按条件创建、输入到输出的内容转换、局部或连续编辑，以及已观测场景的新视角呈现。对于工业零件，恢复原有刻痕与创作另一种表面纹理具有不同目标；对于照片整理，压缩后仍能区分个体与只保留场景大意也具有不同要求。质量、测量和生成因此可以协作，但需要分别说明哪些细节可改变。

\subsection{视频应用的研究内容}
\label{sec:vision-overview-video-tasks}

看到每一帧中的车辆还不足以统计车辆经过路口的数量。统计需要区分同一辆车的重复出现，跟踪需要在遮挡前后维持身份，转弯判断又要把轨迹与道路及时间范围联系起来。\ref{chap:vision-video-recognition}章分别讨论这些输出。逐帧检测正确但身份交换，是跟踪错误；身份正确但事件边界偏移，则是活动定位错误。

对同一段车辆转弯录像，轨迹回答“同一辆车在何处”，事件区间回答“何时发生了约定动作”，检索结果回答“哪份录像或哪段片段符合查询”。活动和事件不是仅凭一帧的类别就能确定的量，例如转弯需要观察方向随时间改变，并结合道路关系解释。\ref{chap:vision-video-recognition}章“视频感知、跟踪与检索”据此连接对象和属性、范围延续、活动与事件、视频及片段检索，不把这些成果合并成一个逐帧准确率。

邻帧也可以提供单幅图像缺少的细节。某帧的路牌模糊，另一帧可能清晰；只有把两帧中同一位置的内容对应起来，额外观测才可能帮助恢复。把对应内容放到可比较位置称为对齐（alignment），把有效信息传给其他帧称为传播（propagation）。更多帧同时带来等待、存储与错误传播；若输出时不允许等待后续帧，就不能直接使用双向读取前后帧的算法。

超分辨率（super-resolution）估计更密空间网格上的内容，插帧（frame interpolation）估计已有两帧之间的画面，二者分别扩大空间和时间采样。例如，暗处的车辆若在相邻帧中可见，正确对应可以利用这条补充信息；若误把另一辆车的纹理传过来，单帧虽然更清晰，主体却错了。\ref{chap:vision-video-restoration}章讨论这些跨帧恢复、采样扩展及编码问题，清晰度、身份延续与时间信息分别核验。

视频编码还要规定编码端使用哪些历史、向解码端传递什么，解码端怎样仅凭码流和允许的参考帧恢复画面。离线处理（offline processing）可以读取整段给定录像，在线处理（online processing）按输入到达顺序工作；两种方式都可能维护长历史，但能否读取相对于当前时刻的后续帧不同。“输入帧数更多”“历史更长”和“当前输出更准确”分别描述信息量、保留范围和任务结果，不构成彼此的保证。

\subsection{多模态应用的研究内容}
\label{sec:vision-overview-multimodal-tasks}

“找到车辆转弯的片段”把文字查询与时间证据联系起来；“为什么画面中有人回头”可能还需要听见声音，但可观察到先后关系并不等于识别了因果机制。文档和图表增加了字符串、阅读顺序、行列与单位的对应。\ref{chap:vision-multimodal-understanding}章从检索、定位、读取和结构恢复推进到描述与问答，始终保留答案所依赖的来源。

例如，“结合声音说明转弯时发生了什么”需要先对应画面与录音的时间，区分发动机声、鸣笛声和画外声，再形成有来源的描述。文档中同一个数字若属于不同表头，含义也可能不同；界面上的同一个图标若处于不同面板，指向的操作又可能不同。多图、文档、图表、界面和音画因此不是仅增加几个输入文件，还增加了跨页、跨区域、跨时刻的对应要求。结构化信息（structured information）把读到的内容组织为字段及明确关系，例如把数字与其行名、列名和单位一同返回，而不是只输出一串字符。

“生成同一辆车转弯并带发动机声音的片段”则改变了问题性质。这里的车、运动和声音是创作条件，不是待解释的实拍证据。\ref{chap:vision-multimodal-generation}章讨论视频、配声、联合音视频及空间成果的制作。固定已有视频再生成声音，与画面和声音共同生成，具有不同的可修改范围；单幅文生图虽然也是跨模态任务，其完整应用流程集中在图像生成章。

制作多个成果时，还需要说明它们共享什么条件。车辆参考图规定主体外观，文字规定动作，已有录音可以规定节奏，但参考图并没有规定未见一侧的全部结构。视觉条件响应（visually conditioned response）把画面中的判断表达为文字或语音，也就是根据所见内容组织回答，仍需核查事实，并按播放要求组织输出。纯文本处理和不依赖视觉的语音应用属于卷九的范围；本卷只在视觉证据或媒体协调改变问题时展开相应关系。

\section{研究历史与当前前沿}
\label{sec:vision-overview-history-and-frontiers}

从局部图像匹配到连续音画响应，应用扩展不断增加需要核验的对象和关系。我们先看这些要求怎样随研究发展进入计算，再看已有路线仍未满足哪些约束；历史中的接口变化与当前的证据、控制和预算问题由此联系起来。

\subsection{研究历史}
\label{sec:vision-overview-history}

视觉研究没有一条覆盖所有任务的单线发展史。分类、几何、复原、编码与生成长期并行，也不断借用彼此的表示和计算。以下按应用要求的扩展组织几个转折，阶段相互重叠，代表工作只用于解释新增问题。

\subsubsection{人工设计方法与观测约束}
\label{sec:vision-overview-classical-history}

同一街景在两次拍摄中会改变大小、方向和亮度，恢复或匹配不能只比较原像素是否相等。早期及长期使用的人工设计路线把所需关系写成明确规则：滤波（filtering）按邻域规则组合像素，成像与退化建模说明场景怎样变成带模糊或噪声的记录，多视图几何（multi-view geometry）利用不同相机投影之间的约束恢复空间关系。已知模糊过程时，可以研究哪些变化被抹弱及怎样限制恢复时的噪声放大；多幅照片有可靠对应时，可以利用投影关系估计空间结构。观测模型也由此指出输入无法唯一确定哪些内容。

尺度不变特征变换（scale-invariant feature transform, SIFT）的2004年论文提供了局部匹配的代表入口：它寻找图像中在不同尺度下可重复辨认的位置，为位置确定参考方向，并把附近外观整理为可比较的描述。这样可以从两幅建筑照片提出局部对应，再由几何约束核查。名称中的尺度不变性指设计所针对的尺度变化，不是任意视角、遮挡或材质变化下的匹配保证\scite{vision-lowe2004}

这类方法没有因深度网络出现而失去全部作用。几何一致性仍可检查学习式匹配结果，退化与采样关系仍可约束复原，码流解码仍须遵守编码端与解码端的信息一致性。后来变化较大的是哪些表示与计算由数据学习，物理观测的限制并未随之消失。

\subsubsection{学习式视觉应用的扩展}
\label{sec:vision-overview-learned-history}

若任务是从许多照片判断类别，人工规定每种外观线索会遇到覆盖和组合的困难。学习式路线让训练样本同时约束内部表示和最终输出。AlexNet是2012年大规模图像分类的代表网络，它把像素变换成多层表示，再给出整图类别分数；训练中的图像类别监督说明应当选择什么类别，却没有直接告诉模型对象的准确范围\scite{krizhevsky2012}

要判断路面、车辆和建筑分别占据哪些位置，整图类别就不够了。全卷积网络（fully convolutional network, FCN）的代表工作于2014年提交预印本、2015年发表于CVPR，把预测组织为与空间位置对应的类别图；这种密集输出（dense output）在许多位置分别给出结果，而不是只返回一个整图标签。它使学习到的图像表示进入像素级任务，同时需要与空间范围相对应的监督\scite{vision-1411-4038}

检测进一步返回多个对象范围，实例分割区分同类个体，全景分割（panoptic segmentation）共同组织可数对象的个体区域与路面、天空等背景区域。这里增加的是结果怎样组织、哪种标注能够核验它的问题。视频还要跨帧关联同一对象，不能由单幅分类发展独自概括这些进展，也不能只用“网络变深了”解释输出职责的变化。

与此同时，视频中的时空表示和对象关联、低照度成像、图像转换及学习式压缩也在发展。它们分别处理时间证据、成像条件、输入输出映射和有限比特表示。一个网络能够在多种任务中使用，不代表同一检查点在没有对应监督或条件接口时就具备所有能力。

\subsubsection{可复用能力与开放条件}
\label{sec:vision-overview-open-history}

在约2020至2023年的代表路线中，一项变化是让结果组织和使用条件更加开放。检测Transformer（Detection Transformer, DETR）的2020年工作把检测结果看作无序对象集合：训练时将预测与标注逐一匹配，使集合中的元素对应不同目标。这里的无序指“第一个框”没有固定类别或身份，不指对象之间没有关系；检测的具体训练与输出选择在识别章展开\scite{vision-2005-12872}

对比语言图像预训练（contrastive language-image pre-training, CLIP）的2021年工作让图片和文字形成可比较的表示，用户因此可以在使用时写出候选类别描述，再比较新图像与各描述的分数。开放类别（open-vocabulary categories）在这里表示候选不再只能由固定分类输出表规定，不意味着任意新概念都能识别正确。CLIP没有仅凭这种匹配接口就获得自由句子生成能力\scite{radford2021}

生成式视觉语言模型（generative vision-language model）则把视觉信息接入语言预测，依据图像与请求组织描述或答案。与CLIP选候选不同，它需要逐步形成可变长度文字，也增加了核查每条回答是否由观测支持的要求。两类接口及其训练依据由“视觉与多模态模型”章建立，历史在这里关注它们让用户提出了哪些新的请求。

用户还可以用点或框选择待分割对象，用空间条件约束生成布局，用少量参考图像指定创作主体。提示扩展了交互方式，也引入了目标歧义：落在车门上的一点可以指车门，也可以指整辆车。开放词汇和通用接口扩大了候选范围，结果仍须依照具体目标核验。

新视角合成（novel-view synthesis）形成了另一条支线：给定已拍照片及相机条件，估计从另一位置观看这个场景会得到什么画面。2020年的神经辐射场（neural radiance field, NeRF）用可查询的空间密度和方向相关颜色表示场景，2023年的三维高斯溅射（3D Gaussian Splatting）用一组空间高斯元素表示可渲染内容。两者依据已拍视图拟合场景，再计算新相机下的图像；未观测区域仍可能受先验或表示约束影响，照片呈现准确也不自动证明表面几何适合测量。相关机制在\ref{sec:vision-generation-novel-view}节展开\scite{vision-nerf2020,vision-gaussian2023}

\subsubsection{连续、多源与交互应用}
\label{sec:vision-overview-continuous-history}

用户在第一帧指出一辆车后，往往还希望在后续画面中继续获得它的区域。分割一切模型2（Segment Anything Model 2, SAM 2）于2024年公开，这里依据同年10月的v2说明它怎样将点或框等提示与视频历史联系起来，输出后续帧中的目标掩码。这里的历史记忆保存供当前分割读取的信息，不等于已经认定了现实车辆的永久身份；遮挡、离开和重新出现仍需分别检查\scite{vision-sam2}

长录像中的问题还可能回到早先片段。MovieChat于2023年提交预印本、2024年发表于CVPR，研究把密集视觉信息整理成长短期记忆以支持问答；压缩历史使较长输入可被利用，但压缩后的表示是否保留后来查询所需细节，仍取决于实际保存的内容。它扩展的是历史可用范围，不是给任意过去问题都提供完整证据\scite{vision-moviechat2024}

空间结果和响应形式也在扩展。2025年的视觉几何Transformer（Visual Geometry Grounded Transformer, VGGT）是联合几何预测模型，可从多视图形成相机、深度和点图等量，点图（point map）为图像位置提供对应三维点；Qwen2.5-Omni技术报告则把图像、音画输入与文字和语音响应联系起来。前者需要检查坐标、尺度和几何相容，后者还要检查答案内容、发音和响应时序，不能用一个共同接口分数概括这些不同任务\scite{vision-2503-11651,vision-qwenomni2025}

生成系统也从单幅图片扩展到视频、空间内容与音视频。共享模型和多阶段制作同时存在：一种路线共享部分表示与主干，另一种路线让专用模型通过明确接口传递结果。仅凭一个产品入口可以接收多种输入，无法判断内部是否共同训练、是否共享参数，更无法推出各项输出之间已保持一致。

\subsection{当前前沿与研究方向}
\label{sec:vision-overview-frontiers}

开放、连续和多源接口扩大了应用范围，也使证据保存、主体一致和成果协调成为更具体的要求。以下以截至2026年10月17日核对的代表文献定位这些问题，区分预印本首稿、修订和会议年份；这不是领域材料的穷尽。近期方法只用于说明改变了什么信息条件或状态组织，其局部结果仍限于对应数据、版本与观察预算。

\subsubsection{高分辨率与长时间中的证据保存}
\label{sec:vision-overview-evidence-frontier}

整幅页面压缩到有限网格后，小字可能消失；长视频均匀采样后，短暂动作可能没有进入输入。研究问题是怎样先发现可能有用的时空范围，再读取足够细的原始内容，并保存以后可以回查的依据。网络中的视觉词元（visual token）是代表一个图像块或汇总范围的计算单元，增加这类单元的容量并不自动补回抽帧时删除的时刻。输入容量、证据选择和实际保留内容因此需要一起记录。

在前面的采样例子中，若查询是“那辆车的车牌是什么”，重新读取4.2至4.6秒的画面可以补充证据；但前提是原记录仍可访问，并有线索把查询引向这段时间。只保存“路口曾有车辆”的语言摘要，没有保存车牌区域，就不能从摘要恢复字符。主动再次观察与增加历史摘要解决的是不同缺口，也承担不同的读取、存储和等待成本。

VideoZeroBench的2026年v1为长视频问题标注答案、时间区间及空间框，并通过不同层级分别提供证据提示或要求模型返回证据。由此可以区分“给出证据后能否回答”和“能否自己找到证据并回答”，而不只看一个答案分数。其有限基准用于诊断这些环节，不证明全部长视频任务具有相同难度；证据获取和评价在视频感知章及\ref{sec:vision-understanding-evaluation}节展开\scite{vision-videozerobench2026}

\subsubsection{空间、时间与主体状态的一致性}
\label{sec:vision-overview-consistency-frontier}

相机、深度、点图与轨迹若描述同一场景，应在共同坐标及尺度约定下相容。一个输出把车辆放在路灯前方，另一个输出却给出相反的深度关系，即使各自看起来合理，组合结果也不能直接使用。VGGT提供联合几何预测的路线，但共同预测不等于所有输出自动相容；空间问答又检查模型能否依据观测回答方向、距离或数量问题。几何误差与答案错误需要分别评价，坐标和尺度约定也需要恢复到实际使用的场景。

时间上的一致性也有不同对象。车辆遮挡后再次出现，需要核验身份；动作反向播放，需要核验事件顺序；生成视频中的物体离开画面后重现，需要核验主体与状态保持。这些问题都涉及连续性，却不能由一个通用的“稳定分数”替代。

例如，轨迹每帧只移动少量像素，仍可能在两车交会时交换身份；生成车辆外观始终相同，仍可能在离开镜头期间停在旧位置。研究需要将位置变化、身份延续和状态演化各自对应到可检查的量，再检查它们能否共同描述一个过程。图像空间恢复、视频对应和生成历史分别承担其中的计算，不能由视觉平滑推出物理过程正确。

\subsubsection{不同模态对同一事件的对应}
\label{sec:vision-overview-alignment-frontier}

声音和画面可以互补，也可能冲突。画面中的人嘴唇在动，不表示录音一定来自此人；字幕出现一个名称，也不证明该对象可见。模型需要保留谁提供了什么信息、发生在何时，再决定如何融合。

多源模型提供了联合计算的接口，但评价还要检查删除、替换或错配某种模态后的变化。模态消融（modality ablation）是在约定设置中移除或改变一种输入，再比较结果，以诊断系统使用了哪些信息。若问题单凭字幕就能回答，原本的准确率不能独自证明视觉路径有效；若替换鸣笛声后答案改变，也只能说明系统对这项替换敏感，尚未说明它找对了声源。

替换录音可能同时改变音量、环境噪声和输入分布，删除字幕也可能造成训练中少见的输入格式。核验时应尽量固定问题、画面和其他条件，记录究竟改变了什么，再将结果变化与时空对应是否正确分开检查。跨模态接口由共用模型章解释，来源核验和错配诊断由多模态理解章展开。

\subsubsection{生成内容的控制与跨成果协调}
\label{sec:vision-overview-generation-frontier}

“一辆红色汽车在蓝色路牌左侧”同时要求对象存在、数量正确、属性归属和空间关系成立。局部修改还要求未指定对象保持，视频增加持续身份和动作，声音增加声源与时间对应。即使整幅图与请求的语义相近，也可能把红色错赋给路牌；评价需要核查关系的两端，而不只计算整图与文字的相似程度。

Show-o2的2025年v3研究在共享主干中结合文字预测和视觉生成，以不同输出模块形成文字及图像或视频，具体能力取决于相应变体。它为“解释参考画面后按条件制作”提供了能力组合的路线，不能仅凭共享主干推断解释必然准确或全部修改条件都被遵守\scite{vision-showo22025}

LTX-2的2026年v1则以相互连接的音频和视频分支共同形成媒体，使两种输出可以在生成过程中交换信息。这不同于先固定完整视频再补声音，也不同于Show-o2的文字与视觉输出组合。对于“车辆转弯后鸣笛两次”，验收要同时检查车辆、动作区间、两次声音及其对应，不从自然画面或同步短片段直接推出长叙事、物理接触及全局状态都正确；模型接口由共用模型章解释，条件遵循、未修改内容保持和跨成果协调由图像生成及多模态生成章展开\scite{vision-ltx22026}

\subsubsection{持续观察与有限预算下的应用}
\label{sec:vision-overview-streaming-frontier}

车辆正在经过路口时，在线系统不能先看完未来的转弯再声称即时作出判断。沿用前文的在线输入条件，因果性（causality）在这里约束的是当前输出不能读取尚未到达的观测。为明确允许读取的范围，设第$t$帧的拍摄时间为$\tau_t$、到达系统的时间为$a_t$，两者以同一时钟的秒数表示。在当前时刻$\tau$可使用的观测集合为
\begin{equation}
  \mathcal{O}(\tau)
  =\{(\bm{X}_t,\tau_t,a_t):a_t\leq\tau\}\text{。}
  \label{eq:vision-causal-observations}
\end{equation}
集合条件按到达时间筛选，而不是按帧编号筛选；在忽略传输与采集延迟的理想情形下，才有$a_t=\tau_t$。这个集合说明哪些信息已经到达，不保证系统仍完整保存它们；实际输入还受所保存的历史内容和读取预算限制。

若算法必须等待$\delta$秒后才能形成输出，它在时刻$\tau+\delta$可以使用更晚到达的观测，但这段等待应计入响应延迟（response latency），即从请求发出到约定结果就绪所经历的时间。图\ref{fig:vision-observation-window}在忽略传输延迟的情形下区分这些范围。离线处理只读取已经给定的完整记录，不是能够取得任意现实未来；对于先前的查询时刻，它可以使用记录中更晚的帧。

\begin{figure}[htbp]
  \centering
  % 时间坐标只作范围示意，不表示某个实验的采样频率。
\begin{tikzpicture}[x=1cm,y=1cm,font=\figurefont\small]
  \fill[FigureInput!10] (0,0.15) rectangle (4.6,1.3);
  \fill[FigureModel!10] (4.6,0.15) rectangle (7,1.3);
  \fill[BookPaper] (7,0.15) rectangle (10.2,1.3);
  \draw[-{Stealth[length=2mm]},BookInk,line width=0.8pt] (0,0) -- (10.5,0)
    node[below left] {观测时间};
  \foreach \x in {0.4,1.4,2.4,3.4,4.4,5.4,6.4,7.4,8.4,9.4}{
    \fill[BookInk] (\x,0) circle (1.5pt);
  }
  \draw[BookInk,line width=0.9pt] (4.6,-0.25) -- (4.6,1.5);
  \draw[BookMuted,dashed,line width=0.8pt] (7,-0.25) -- (7,1.5);
  \node[below] at (4.6,-0.25) {$\tau$};
  \node[below] at (7,-0.25) {$\tau+\delta$};
  \node[align=center] at (2.3,0.75) {已到达历史\\即时可用};
  \node[align=center] at (5.8,0.75) {等待期间\\新增观测};
  \node[align=center] at (8.6,0.75) {其余未来\\离线才可用};
\end{tikzpicture}

  \caption{同一查询在不同信息条件下可读取的帧，示意中忽略传输延迟。圆点表示实际观测，竖线是查询发起时刻$\tau$；在线即时处理只读取左侧，允许等待$\delta$时才可加入中间区间，离线处理还可使用右侧已给定记录中的后续帧。颜色之外用文字和线型区分范围。}
  \label{fig:vision-observation-window}
\end{figure}

循环结构用先前内部状态处理当前输入，缓存（cache）保存供后续计算读取的已有信息，长记忆扩大可保留的历史；这些名称都描述计算组织，因果性仍取决于实际读取了哪些时刻。例如，用双向读取得到的特征更新一个循环状态，该状态仍含有未来信息，不能因外层具有循环结构就称整个过程为即时在线。

Flash-VStream的2025年路线分开处理持续到达的帧和用户问题，用概括历史的信息选择需要保留或取回的细节。TempoGround的2026年v1则面向流式语言定位，关联相邻帧对象，预测进入、持续或离开的存在状态，再形成二维框及相机坐标系中的三维框。前者关注历史怎样供问答使用，后者关注语言目标怎样随到达画面保持对应；两者都需要记录采样、历史容量、观察范围及响应时间，具体机制由视频感知和多模态理解章展开\scite{vision-flash-vstream2025,vision-tempoground2026}

响应预算还要明确结束点：第一文字单元出现、完整答案就绪和全部媒体播放完成不是同一时刻。存储更多摘要可能保存较长经历，却仍遗漏小目标；处理更多区域可能取得细节，却增加读取和等待。在线结果用已到达证据判断当前或过去，未来预测则估计尚未观测的状态，两者可以出现在同一系统中，但应分开标明和评价。

\subsubsection{从单次输出到持续、可查询的场景状态}
\label{sec:vision-overview-queryable-state-frontier}

设街景中的一辆车离开镜头，几秒后又从另一个方向返回。“现在画面里有什么”“它上次在哪里”“回来的是不是同一辆车”和“离开期间可能去了哪里”需要不同信息。当前检测回答可见对象，历史记录保留过去观测，对象对应判断身份，离视期间的位置则只能由额外观测或预测约束。保存最后一幅画面不能自动说明这辆车此刻仍停在旧位置。

持续、可查询的场景状态（persistent, queryable scene state）指一种随观测更新、能按对象和时刻读取的内部记录或表示。它把身份、位置、属性、事件及其依据联系起来，使后续问题可以回到同一对象，并在新证据出现时修订原判断。这里的状态可以不完整，也可以保留多个可能解释；它不是要求系统精确存下现实世界的全部物理量。

为了区分“维护记录”和“回答当前问题”，可以用一个教学接口写出两步关系。设$\mathcal{S}_{t-1}$是处理第$t$帧前的状态，$o_t=(\bm{X}_t,\tau_t,a_t)$是已经到达的观测，$q$是查询，$\hat r_t$是返回的结果及证据索引，则
\begin{equation}
  \begin{aligned}
    \mathcal{S}_t&=U(\mathcal{S}_{t-1},o_t),\\
    \hat r_t&=R(\mathcal{S}_t,q)\text{。}
  \end{aligned}
  \label{eq:vision-overview-state-update-query}
\end{equation}
$U$表示根据新观测更新状态，$R$表示按查询读出结果；$\mathcal{S}_t$是组织多项信息的结构，不预设为固定维向量。初始状态可以为空，也可以含有明确来源的既有地图。这是问题分工的示意，不是已经实现的通用算法：每个具体系统还须规定更新哪些记录、如何处理冲突、删去什么，以及无法确定时怎样返回。

沿车辆例子复核这两步，首次看见车辆时可以建立临时身份，保存区域和时刻；离开镜头后，把可见状态改为“当前未见”，保留最后观测及可能运动的预测，但不把预测伪装成新观测。车辆返回时，再用新外观、位置和其他可用证据核查身份；若原先关联错误，应修订身份关系并保留修订依据。用户随后问“上次在何处”应读历史记录，问“现在在哪里”则需要当前证据或明确注明估计，两个查询不能都返回同一旧框。

V-JEPA 2.1的2026年v1提供了特征预测的相关路线：训练时从部分可见图像或视频预测完整输入的内部表示，并把预测约束扩展到可见位置和多个表示层，以支持局部空间结构与时间信息的学习。这里的潜在特征（latent feature）是网络内部对观测的数值表示，不是直接可检查的实体表。该路线为视频感知提供表示与预测依据，但不能仅由特征预测推断系统已能维护现实对象的永久身份或回答任意状态查询\scite{vision-vjepa21-2026}

InfiniteVGGT的2026年v1从另一端研究持续几何：在因果图像流中保留首帧参考信息，并按预算选择滚动保存的历史特征，持续形成相机、深度与空间结果。它改变的是几何计算如何在有限内部缓存下延续，不是无界保存全部历史；有限长序列上的检验也不构成无限时间无漂移的保证。其空间恢复机制归入图像识别与检索章，与视频感知中的实体历史相互补充\scite{vision-infinitevggt-2026}

LiveWorld的2026年v1则研究生成环境中的离视动态（out-of-sight dynamics），即对象不在当前视野内时仍可能继续发生的变化。它将静态背景与动态实体分开维护，为局部活动实体设置虚拟观察位置，生成离视期间的演化，再将更新状态用于返回视角的画面。这个过程模拟未观测内容，不是从真实录像恢复了确定发生的事件；动态区域数目、局部演化和几何条件也限制其范围。其机制由多模态生成章展开，不能与真实历史证据混用\scite{vision-liveworld-2026}

《Generalist Open-World Temporal Perception》的2026年v1是一篇观点论文，提出让感知状态具有持续性、可寻址性、跨输出相容性、来源对应、可修订性及不确定性表达。可寻址性（addressability）在这里指后续查询能够再次指向同一实体、表面或事件，而不是只能读取一段无法定位的摘要。论文将这些要求组织为研究议程，并讨论相关组件证据；它不是已实现一个满足全部要求的通用系统的实验报告\scite{vision-generalist-temporal-perception-2026}

这些路线共同指向如何组织可持续使用的信息，却不处理同一种状态。研究中的核验至少要分清真实观测是否可回查、身份是否可修订、几何是否在共同参考下相容，以及未观测演化是否明确作为预测或创作。保留内部缓存、答对一次问题或生成一段连贯视频，只分别说明其中一部分；要支持持续查询，还需在新证据、遮挡、重访和预算变化下检查这些关系是否仍成立。

\section{本章小结}
\label{sec:vision-overview-summary}

街景上的分类、定位、增强与编辑之所以需要不同解法，是因为它们要求不同的成果，也容许不同的内容变化。明确观测、输出和约束之后，才能判断模型需要保留什么表示、取得什么监督，以及怎样恢复和核验结果。图像建立空间关系，视频增加身份、时间与信息到达条件，多模态处理再连接查询、不同来源及不同输出。

研究范围从专门任务扩展到开放、连续和多源应用，并没有取消证据边界。采样可能遗漏内容，预测可能依赖先验，生成可能补出没有观测支持的细节；持续场景状态还要把过去观测、当前判断和未观测预测分开保存，并允许新证据修订原关联。今后的证据保留、跨模态对应和长期状态研究，需要同时考虑能回答什么、依据是否可查以及预算是否允许，而不只扩大一个输入接口。

后续先由共用模型章建立编码、连接和输出的计算，再由任务章追踪监督、参数来源、逐次计算和结果恢复。本章对开篇街景请求的共同回答是：任务规定需要形成的成果，来源与预算规定可以使用的信息，核验检查成果是否仍指向正确对象、区域、时刻及用途。只有这些关系明确，分类、定位、复原、编辑与持续查询的能力声明才有可比较的含义。

\chapter{视觉与多模态模型}
\label{chap:vision-multimodal-models}

同一张街景照片，可以用来寻找相似图片、回答“右侧停着什么车”，也可以作为生成另一种天气下街景的条件。若再给一段录音，系统还能比较画面与声音是否相关，或把回答说出来。这些功能表面上都接受多种输入，内部需要的计算却不同。一个能比较图文相似度的模型，未必有生成句子的输出层；一个会说话的视觉助手，也未必能生成视频。

概述已经区分观测、查询和创作条件，本章在这些对象上建立后续任务共用的模型接口。我们追踪观测怎样被编码、不同来源的信息在哪里相遇、输出又怎样恢复成分数、语言或媒体。编码和连接决定模型能利用哪些证据，训练目标决定哪些对应关系得到学习，输出模块决定结果能够以什么形式出现。理解和生成可以共享其中一部分，也可以由几个专用模型协作完成。

以下沿用同一个教学场景：街景左侧有红色自行车，右侧有红色汽车，录音中可能出现发动机声。图片只提供静态外观，声音是否来自右车还需额外依据。我们可以要求选择车辆类别、回答车窗状态，也可以要求创建保留两车位置的雨天变体。这样，表示匹配、语言响应和媒体生成能够在相同对象上比较，而不会把生成的内容当作实拍证据。本章所用的小向量和概率均为教学构造，不是所介绍模型的测量结果。

\section{多模态模型的组成与接口}
\label{sec:vision-models-interfaces}

连接图像和文字之前，需要区分原始观测、网络内部表示和最终结果。它们的数值形式可能相似，却承担不同职责。例如，一张照片可以形成许多空间位置上的特征，也可以被汇总为一个整图向量；后者便于匹配，却未必足以恢复前者包含的小字和边界。

\subsection{各模态的编码与组织}
\label{sec:vision-models-encoding}

沿用前文记号，图像为$\bm{X}\in\mathbb{R}^{h\times w\times c}$，视频由带时间戳的帧$\bm{X}_t$组成。图像编码器把像素转换为$\bm{F}\in\mathbb{R}^{n_{\mathrm v}\times d_{\mathrm v}}$，每一行表示一个位置或一个汇总单元。对图像块编码，$n_{\mathrm v}$通常随空间分辨率增加。位置编码使模型能够区分“左侧的车”和“右侧的车”，但位置编码不能补回裁剪时删除的对象。

以不重叠图像块为例，设每块边长为$p$像素，且$h,w$均可被$p$整除，则共有$n_{\mathrm v}=hw/p^2$个块，不计额外汇总位置。把一个$p\times p\times c$块按固定顺序展开为$\bm x_i\in\mathbb R^{p^2c}$，线性映射$\bm A\in\mathbb R^{d_{\mathrm v}\times p^2c}$和偏置$\bm b\in\mathbb R^{d_{\mathrm v}}$形成$\bm A\bm x_i+\bm b$，再加上该块的位置表示。网络后续层在这些单元之间交换信息，得到的特征已不等于原像素。视觉Transformer（vision transformer, ViT）沿这条图像块序列计算，注意力层让不同位置的信息交互；卷积编码器则先用局部卷积逐层计算，再从空间网格取特征。不同网络可以形成相同形状的接口，却未必保留同样内容。

文字先被切成离散词元（token），即分词器按词表规则划出的字符或子词单元，再经嵌入表转为向量。音频则可以由波形转换为时频表示，再划分成局部单元编码。时频表示的两维分别描述时间和频率，和图像的横纵坐标不同。视频还需保留帧的实际时间、帧内空间位置和所属片段；单纯把第1帧、第2帧写入序列，无法区分两帧相隔0.1秒还是10秒。

若文字词表大小为$V$，第$j$个词元编号为$a_j\in\{1,\ldots,V\}$，嵌入表$\bm B\in\mathbb R^{V\times d_{\mathrm t}}$给出的输入就是第$a_j$行。分词器规定哪些字符串对应哪些编号，不能在使用时任意换表。音频波形$s_i$还需采样率$f_{\mathrm s}$，第$i$个采样点对应$i/f_{\mathrm s}$秒；将波形分窗并计算频谱后，每列仍应记录窗口的时间范围。视频按帧编码时，可以先在各帧得到视觉网格，再沿时间组织；采用时空编码器时，单元可以同时覆盖数帧，此时其时间范围也不再等于单帧时间戳。

实际系统常把连续的视觉特征也称为“视觉词元”。这里的词元表示可进入序列模型的一个计算单元，不一定经过离散词表量化。它与文字词元的整数编号、媒体生成器使用的离散编码、原始像素值，都不能直接互换。比较词元数时，必须说明每个单元对应多少空间或时间范围。

设街景原图宽1200像素、高800像素，裁出左上角为$(400,200)$、宽高为$400\times400$的区域，再调整到$200\times200$作为局部输入。以连续边界坐标计，局部结果$(x',y')=(50,80)$对应原图
\[
  (x,y)=(400+2x',\,200+2y')=(500,360)\text{。}
\]
因此，输入记录至少应保留原图编号、裁剪边界和缩放关系。若视频只选取0、2、4秒三帧，还应保留真实时间戳，而不是让下游把三帧当成连续拍摄的相邻帧。编码前的选择已经限定了模型可能看到的证据。

同一段视频还可以抽取1、3、5秒三帧。两组输入的帧数相同，却可能包含不同车辆或动作；只有保留采样位置，后续才能知道某个答案依据哪组观测。裁剪与插值也应区分：从原始高分辨率照片重新裁取车窗，可能增加相对于缩小整图的可见细节；把已经缩小的车窗插值放大，仅改变数组尺寸，不增加原始像素证据。后面的主动观察接口使用的是前一种信息获取。

\subsection{跨模态信息的连接}
\label{sec:vision-models-connection}

\term{视觉语言模型}（vision-language model, VLM）把视觉输入与语言联系起来，可以输出可比较表示，也可以生成语言。本书采用这一较宽的用法。连接音频、深度或惯性信号的模型还需相应编码和训练依据，不能仅凭“多模态”名称确定支持范围。

一种安排是两侧独立编码，再把表示映射到共同空间计算相似度。图像向量不依赖当前文字查询，适合预先建立图库。另一种安排让文字计算读取视觉特征：语言隐状态提供查询，也就是当前希望取用信息的向量；视觉特征提供用于比较的键和承载内容的值，再按匹配权重读取。此时视觉底层特征可以缓存，但读取结果随问题和已生成文字改变。还有一种安排把多源表示放入共享主干，通过位置、模态标记和注意力范围组织交互。

把第二种安排写成一次可检查的读取，有助于明确哪里依赖当前问题。设语言隐状态为$\bm H\in\mathbb R^{m\times d_{\mathrm l}}$，视觉特征为$\bm F\in\mathbb R^{n_{\mathrm v}\times d_{\mathrm v}}$，分别投影得到$\bm Q\in\mathbb R^{m\times d_{\mathrm a}}$、$\bm K\in\mathbb R^{n_{\mathrm v}\times d_{\mathrm a}}$与$\bm V_{\mathrm a}\in\mathbb R^{n_{\mathrm v}\times d_{\mathrm o}}$。交叉注意力（cross-attention）让每个文字位置按与视觉位置的相关程度加权读取视觉值：
\begin{equation}
  \bm R=\operatorname{softmax}_{\mathrm{row}}
     \left(\frac{\bm Q\bm K^\top}{\sqrt{d_{\mathrm a}}}+\bm M\right)
     \bm V_{\mathrm a}.
  \label{eq:vision-models-cross-read}
\end{equation}
这里$\bm M\in\mathbb R^{m\times n_{\mathrm v}}$在允许读取的位置取0，禁止读取的位置取$-\infty$，每行至少要有一个可读位置。行softmax使一个文字位置的权重和为1，$\bm R$的每一行是相应视觉值的加权和。问题从“车是什么颜色”变成“车窗是否打开”时，$\bm H$及查询改变，即使$\bm F$未变，读取权重也可改变。注意力数值只表示内部加权，不自动构成对象定位的标注。

这三种安排描述计算发生在哪里，不是互斥的完整模型类别。一个系统可以先用独立表示召回候选，再让联合模型逐个检查。表示对齐指训练使不同来源的内容在比较或交互中形成有用联系，参数共享指不同输入实际使用相同的权重，连接模块则负责把形状或信息路径接起来。两套编码器即使输出同样维数，也没有因此完成语义对齐。

以图库查询为例，独立编码可以保存$M$幅图的向量，新请求只需编码一次文字并比较候选。若每幅图都要与当前问题共同通过一个大型主干，通常就要执行$M$次相关联合计算。前者能否找到关系复杂的目标、后者能否在预算内处理全部候选，都需要按任务确认；模块安排只说明这些成本和能力从何而来。

\subsection{输出表示与解码}
\label{sec:vision-models-output}

输出分数、输出文字和输出媒体需要不同的预测对象。匹配接口可以计算一个向量内积；语言接口在每一步给出词表上的分布；图像或音频接口则预测像素、波形或能够被解码的中间表示。中间表示必须与解码器的训练约定相容，任意同维向量都不能送入同一个解码器。

表\ref{tab:vision-model-interfaces}固定同一街景输入，比较不同请求所需的结果。若给一个图文匹配模型附加图像生成器，匹配模型可以提供条件，但生成能力属于完整连接后的系统。反过来，视觉问答虽然输出文字，评价的仍是回答是否依据画面作出了正确判断。

\begin{table}[htbp]
  \centering
  \small
  \begin{tabularx}{\linewidth}{@{}p{22mm}XX@{}}
    \toprule
    请求 & 模型直接预测的对象 & 恢复最终结果还需什么\\
    \midrule
    选择车辆类别 & 候选描述的匹配分数 & 给定候选集与选择规则\\
    回答停车位置 & 下一文字词元的分布 & 逐词解码与来源核验\\
    生成雨天变体 & 图像潜表示的更新量 & 生成迭代与图像解码器\\
    说出判断 & 语音编码序列 & 声学解码与波形合成\\
    \bottomrule
  \end{tabularx}
  \caption{输入相同而输出接口不同的教学比较。每行描述一种完整路径，并不意味着同一个检查点同时具有四种能力。}
  \label{tab:vision-model-interfaces}
\end{table}
\FloatBarrier

\section{多模态理解模型}
\label{sec:vision-models-understanding}

如果任务是从给定候选中选择，模型可以用可比较表示形成判断。如果任务要求自由描述或回答，则还需利用语言接口组织结果。本节先展开匹配如何得到训练，再说明这种接口的用途与边界；语言响应的具体连接方式在\ref{sec:vision-models-language-generation}节展开。

\subsection{CLIP的图文双编码与相似度判定}
\label{sec:vision-models-contrastive}

为了用文字指定视觉概念，需要让“红色汽车”这样的描述与包含相应内容的图片能够比较。对比语言图像预训练（contrastive language-image pre-training, CLIP）以真实图文配对为监督，分别编码两种输入，再把正确配对与同一批次中的其他组合区分开。其图像编码器可以采用经修改的残差卷积网络ResNet或ViT，文字编码器采用Transformer，双编码结构不要求两边网络相同\scite{radford2021}

设一批有$n$对图文$(\bm{X}_i,q_i)$。图像和文字编码器分别输出$f_{\mathrm v}(\bm{X}_i)\in\mathbb R^{d_{\mathrm v}}$与$f_{\mathrm t}(q_i)\in\mathbb R^{d_{\mathrm t}}$，线性投影矩阵$\bm{W}_{\mathrm v}\in\mathbb R^{d\times d_{\mathrm v}}$、$\bm{W}_{\mathrm t}\in\mathbb R^{d\times d_{\mathrm t}}$将它们映射到$\mathbb{R}^d$。对于非零投影结果，归一化表示为
\begin{equation}
  \bm{v}_i=\frac{\bm{W}_{\mathrm v}f_{\mathrm v}(\bm{X}_i)}
                    {\lVert\bm{W}_{\mathrm v}f_{\mathrm v}(\bm{X}_i)\rVert_2},
  \qquad
  \bm{u}_i=\frac{\bm{W}_{\mathrm t}f_{\mathrm t}(q_i)}
                    {\lVert\bm{W}_{\mathrm t}f_{\mathrm t}(q_i)\rVert_2}\text{。}
  \label{eq:vision-clip-embedding}
\end{equation}
内积$s_{i,j}=\bm{v}_i^\top\bm{u}_j$就是余弦相似度。归一化使向量长度不能单独抬高分数，模型必须改变方向来改善配对。原CLIP共同学习编码器、线性投影以及控制分数尺度的温度参数。

温度$\tau>0$把相似度转换为$\ell_{i,j}=s_{i,j}/\tau$，这些未归一化分类分数称为logit。温度越小，同样的相似度差会形成越大的logit差。固定一幅图，在本批次$n$段文字中识别其配对文本，可以使用行方向交叉熵，即为正确配对的负对数选择概率计分；固定一段文字，反过来在图像中识别配对，则使用列方向交叉熵：
\begin{equation}
  \begin{aligned}
    L_{\mathrm{v\to t}}
      &=-\frac1n\sum_{i=1}^n
        \log\frac{\exp(\ell_{i,i})}{\sum_{j=1}^n\exp(\ell_{i,j})},\\
    L_{\mathrm{t\to v}}
      &=-\frac1n\sum_{i=1}^n
        \log\frac{\exp(\ell_{i,i})}{\sum_{j=1}^n\exp(\ell_{j,i})},\\
    L_{\mathrm{CLIP}}&=\tfrac12
        \bigl(L_{\mathrm{v\to t}}+L_{\mathrm{t\to v}}\bigr)\text{。}
  \end{aligned}
  \label{eq:vision-clip-loss}
\end{equation}
对角线表示数据记录中的正确配对，其他单元在该目标中充当负例。两项损失约束相同的相似度矩阵，但分母分别沿行和列求和。若批次中另有一幅也符合该描述的图，单一配对标签会把它当负例，这就是配对监督和真实语义关系之间的差别。

交叉熵为何会改变相对得分，可以从一行的目标直接推导。令$p_{i,j}=\exp(\ell_{i,j})/\sum_k\exp(\ell_{i,k})$，第$i$行的误差为$-\ell_{i,i}+\log\sum_k\exp(\ell_{i,k})$。第一项只在$j=i$时导数为$-1$，第二项的导数为$p_{i,j}$，因此两者相加得到$p_{i,j}-\mathbb I[j=i]$。若再令$\widetilde p_{i,j}$为第$j$列沿图像方向归一化的概率，完整双向平均的梯度为
\begin{equation}
  \frac{\partial L_{\mathrm{CLIP}}}{\partial\ell_{i,j}}
   =\frac{p_{i,j}+\widetilde p_{i,j}
                 -2\mathbb I[i=j]}{2n}.
  \label{eq:vision-clip-logit-gradient}
\end{equation}
正配对概率还小于1时，对角线的导数为负；负配对的导数为正。沿负梯度改变参数，便倾向于增大前者、减小后者。这里的logit不是独立可调参数，实际更新还须通过内积、归一化和两个编码器传回权重。

归一化也参与这个过程。对投影结果$\bm a\ne\bm0$与$\bm v=\bm a/\lVert\bm a\rVert_2$，其导数矩阵为$(\bm I-\bm v\bm v^\top)/\lVert\bm a\rVert_2$。这一步把沿$\bm v$方向、只改变长度的梯度分量去掉，保留能改变方向的分量。图像编码器收到同一行所有文字比较的梯度，文字编码器收到同一列所有图像比较的梯度，所以一次真实配对同时约束了两侧。零向量处公式未定义，实现需避免或用数值稳定措施处理，不能把非零条件省掉。

取一个可复算的二维教学批次：两幅图的表示分别为$(1,0)^\top$、$(0,1)^\top$，两段配对文字的表示分别为$(0.8,0.6)^\top$、$(0.6,0.8)^\top$。四个向量均为单位向量。设两组标签为“红色汽车”和“红色自行车”，并取$\tau=0.2$，就得到
\[
  \bm{S}=\begin{bmatrix}0.8&0.6\\0.6&0.8\end{bmatrix},
  \qquad
  \bm{\ell}=\begin{bmatrix}4&3\\3&4\end{bmatrix}\text{。}
\]
每一行的正确配对概率为$e^4/(e^4+e^3)\approx0.7311$，交叉熵为$\log(1+e^{-1})\approx0.3133$；列方向在这个对称例子中相同。行损失对第一行两个logit的梯度，暂不计批次平均系数，为$(-0.2689,0.2689)$。梯度下降会提高正确配对相对错误配对的分数。图\ref{fig:vision-clip-matrix}标出了两个归一化方向。

完整双向平均中$n=2$，式\eqref{eq:vision-clip-logit-gradient}给出第一行导数$(-0.1345,0.1345)$。若只为观察分数变化，暂把logit当独立变量，并取教学步长0.1，则$(4,3)$更新为约$(4.0134,2.9866)$，正确选择概率增至约0.7363。这不是真实网络的一次训练结果，因为真实参数共用、单位向量受约束；它只复算目标希望把两个分数差推向哪个方向。

训练时，先把数据记录中的图片按模型约定裁剪、缩放和归一化，把对应文本按固定分词器编码，再抽取一批记录构造整个$\bm S$。原CLIP从随机初始化的编码器开始学习，用Adam优化器调整编码器、投影和分数尺度，并按预定训练日程停止，而不是等到全部配对概率变成1。温度通常通过学习对数尺度保持为正；它改变softmax的尖锐程度，不改变未缩放余弦分数的排序。训练后保留图像处理方式、分词器、两侧权重、投影和尺度，丢弃当前批次的负例矩阵。面对新数据的正确率仍需另行评价，下降的训练目标没有给出泛化保证。

\begin{figure}[htbp]
  \centering
  % 教学数据由四个单位向量的内积构造，温度为0.2。
\begin{tikzpicture}[x=1cm,y=1cm,font=\figurefont\small,
  arr/.style={-{Stealth[length=2mm]},BookInk,line width=0.8pt}]
  \node[anchor=west] at (0,4.5) {缩放后分数$\ell_{i,j}=s_{i,j}/\tau$};
  \node[align=center] at (3.8,3.55) {$q_1$\\红色汽车};
  \node[align=center] at (6.0,3.55) {$q_2$\\红色自行车};
  \node[anchor=east] at (2.6,2.5) {图像1};
  \node[anchor=east] at (2.6,1.2) {图像2};
  \foreach \x/\y/\val in {3.8/2.5/4,6.0/2.5/3,3.8/1.2/3,6.0/1.2/4}{
    \draw[BookRule,line width=0.5pt,fill=BookPaper]
      (\x-1,\y-0.55) rectangle (\x+1,\y+0.55);
    \node at (\x,\y) {$\val$};
  }
  \draw[BookTeal,line width=1.2pt] (2.8,1.95) rectangle (4.8,3.05);
  \draw[BookTeal,line width=1.2pt] (5.0,0.65) rectangle (7.0,1.75);
  \draw[arr] (7.2,2.5) -- (8.15,2.5);
  \node[anchor=west,align=left] at (8.2,2.5) {行softmax\\$0.7311$\\$0.2689$};
  \draw[arr] (3.8,0.4) -- (3.8,-0.2);
  \node[align=center,anchor=north] at (3.8,-0.3)
    {列softmax\\$(0.7311,\;0.2689)^\top$};
  \node[anchor=west] at (0,-1.35) {对角线：数据配对\qquad 非对角线：批次负例};
\end{tikzpicture}

  \caption{CLIP双向目标的两对图文教学示例。行固定图像、列固定文字，粗框标出训练配对。行、列softmax分别形成两个分类问题，图中的0.7311是批次内选择概率，不是“图中确有该物体”的校准概率。}
  \label{fig:vision-clip-matrix}
\end{figure}

使用时，可以把类别$c$写成描述模板$q_c$，缓存其表示$\bm{u}_c$，再对新图像选取$\argmax_c\bm{v}^\top\bm{u}_c$。写“这是一张红色汽车的照片”和只写“汽车”，提供给文字编码器的上下文不同，得到的方向也可能不同。若任务只有“汽车”“自行车”两个候选，模型一定会从中选一个，除非另加拒绝规则；加入“公交车”后，归一化分数的分母又会改变。

这种零样本分类不使用当前类别任务的标注样本去拟合分类器，但仍依赖大规模图文预训练。比较结果时应记录候选集合、模板、图像处理和模型版本。整图相似度没有直接给出车辆的框、实例掩码或自由答案，这些输出需要其他接口。对比学习用正对应与候选参照构成相对判定，候选来源和数量会改变训练信号；通用目标与采样问题可参考\BookChapterRef{chap:data-derived-knowledge-acquisition}{卷八“从数据中获得知识”章}，这里的关键是图文两侧如何共同构成配对判定。

\subsection{ImageBind的图像配对与多模态表示联系}
\label{sec:vision-models-imagebind}

图文配对之外，视频帧和同时录到的声音、彩色图与对应深度、佩戴设备的画面与惯性信号，也提供了跨来源联系。若必须为所有模态两两采集配对数据，模态种类增加后组织成本会迅速上升。ImageBind用图像作为配对中心，把多个编码器接到已有视觉语义空间中\scite{vision-imagebind2023}

原实现分别编码图像、文字、音频、深度、热成像和惯性信号。图像与视频共用视觉编码器，视频输入由短时间范围内的帧构成。音频转换为梅尔频谱（mel-spectrogram），即按接近听觉尺度的频率分组记录各时间窗的能量，再划块编码。深度预处理为表示距离关系的单通道视差图，热成像也按单通道图组织；惯性信号则把加速度计和陀螺仪的时间序列通过一维投影送入Transformer。各编码器末端有独立线性投影，输出归一化到同一维数。

以声音为例，一批训练数据包括画面$\bm{X}_i$与对应声音$a_i$。模型比较$\bm{v}_i$和$\bm{u}^{\mathrm a}_j$，沿图像到音频、音频到图像两个方向使用与式\eqref{eq:vision-clip-loss}同类的配对目标。深度、热成像和惯性数据各自与其配对画面计算目标，不要求一条记录同时拥有全部模态。原论文实现冻结已预训练的图像、文字编码器，学习音频、深度、热成像和惯性编码路径。

更具体地说，设模态$m$的数据集为$\mathcal D_m=\{(\bm X_i,a_i^{m})\}$，其编码与投影参数为$\bm\theta_m$。每次在一个已存在的配对数据集中抽取批次，计算图像向量与该模态向量的双向目标$L_m$；误差经归一化和投影反传到$\bm\theta_m$。冻结图像端使误差只推动新增模态去适配图像空间，文字端也不因音频批次改变。原论文按模态分别给出AdamW日程，音频、深度和热成像设置为64轮、惯性设置为8轮，并非每条记录同时联合全部模态。训练后保留各自的预处理、编码器与投影；推理不再需要配对标签和批次内负例。

继续使用街景例子，假设发动机录音形成单位向量$(0.9,\sqrt{0.19})^\top$，图库中两幅候选的图像表示为$(1,0)^\top$和$(0,1)^\top$。两个内积约为0.9和0.4359，排序先返回前者。若两幅候选都含车辆但只有后者是这段录音的真实来源，几何排序仍可能选错；训练配对学到的是所覆盖数据中的联系，并非每次录音的声源认证。检索应返回候选编号与分数，才能继续核查内容和来源，而不是仅输出“找到了发动机”。

固定图像和文字空间的意义在于，新增模态需要向已经形成的坐标关系适配，而不是让所有编码器同时任意移动。训练仍需真实对应和足够数据覆盖。画面中既有车又有树而录音只有风声时，整段音画配对也没有显式指出声音来自哪个区域。

使用声音检索画面时，只需编码声音，再按它与图库向量的内积排序。使用声音选择文字类别时，可以比较声音向量与候选文字向量，即使新增音频编码器没有直接用音文对训练。论文在所测数据中观察到了这种未直接配对模态之间的联系；共同空间本身不是保证联系正确的定理。

一个几何事实有助于理解这种联系的可能性。若某段声音和某段文字都与同一个图像向量很接近，即$\lVert\bm{u}^{\mathrm a}-\bm v\rVert_2\leq\epsilon_{\mathrm a}$、$\lVert\bm{u}^{\mathrm t}-\bm v\rVert_2\leq\epsilon_{\mathrm t}$，三角不等式给出
\[
  \lVert\bm{u}^{\mathrm a}-\bm{u}^{\mathrm t}\rVert_2
  \leq\epsilon_{\mathrm a}+\epsilon_{\mathrm t}\text{。}
\]
但实际训练不能让每个语义一致的样本都围绕同一个准确图像锚点，也不能消除背景、视角和声源歧义，因此这只是有条件的几何解释。ImageBind的音频到图像生成展示还外接了图像生成系统，用声音表示替代相应条件；ImageBind本体仍输出表示。

\subsection{理解结果的形成与接口选择}
\label{sec:vision-models-understanding-interface}

街景查询“找相似画面”可以直接返回向量排序。“描述发生了什么”没有固定候选句表，需要根据视觉内容预测一串文字。“指出右侧车辆”则要求把表达对应到区域，仅有一句“右边有一辆车”尚未给出可检查的位置。三者都涉及理解，但需要的结果不同。

相似度适合在已定义候选中选择，语言生成适合组织可变长度的描述和答案，定位接口则要输出框、点或掩码。语言模型也可以按约定输出坐标，但仍需训练依据、坐标转换和定位评价。具体的语言指代与区域对应见\ref{sec:vision-understanding-grounding}节。

同一表示也可服务生成。图像语义向量可以作为生成器的条件，图文相似度的梯度可以调整采样方向，生成结束后的相似度还可以筛选候选。这些用途分别改变生成输入、生成过程和选择规则，不能混为“CLIP生成图像”。下一节从预测对象和训练模块出发，解释这些不同连接。

\section{多模态生成模型}
\label{sec:vision-models-generation}

生成描述、答案和语音响应，主要需要把已观察到的内容表达出来；生成图像、视频和环境声音，还要预测完整媒体中的大量细节。两类输出可以在一个系统中结合，所需目标和解码器却不同。本节分别追踪语言响应和媒体内容的生成路径。

\subsection{语言响应的生成}
\label{sec:vision-models-language-generation}

对街景问题“右侧停着什么车”，需要从同一视觉证据形成长度可变的答案，而不是从两个类别描述中选一个。设观测经过编码得到视觉表示$\bm{F}$，请求为$q$，回答的词元为$y_1,\ldots,y_m$。条件语言生成通常按此前已经产生的前缀逐词展开：
\begin{equation}
  p_{\bm\theta}(y_1,\ldots,y_m\mid\bm F,q)
  =\prod_{j=1}^m p_{\bm\theta}(y_j\mid y_{<j},\bm F,q)\text{。}
  \label{eq:vision-vlm-language}
\end{equation}
训练时，把真实回答前缀提供给模型，以目标词元的负对数概率监督；使用时，前缀逐步由模型自己产生。式中的条件依赖只有通过实际连接才有可能成立。下面几种方法改变的重点，就是视觉信息怎样到达预测$y_j$所用的语言隐状态，以及这条路径上哪些权重被训练。

一条训练记录应包含图片、请求和参考回答，必要时再含先前轮次及其角色。令$\mathcal T_i$为第$i$条记录中需要监督的回答位置，用户问题和作为条件的图片不属于这组目标位置。因为参考序列的概率是逐词条件概率的乘积，取负对数后得到可相加的误差；对$n$条记录形成经验目标
\begin{equation}
  \widehat L_{\mathrm{text}}
   =-\frac1n\sum_{i=1}^n\sum_{j\in\mathcal T_i}
      \log p_{\bm\theta}(y_{i,j}\mid
          y_{i,<j},\bm F_i,q_i).
  \label{eq:vision-vlm-text-objective}
\end{equation}
可再按目标词元数归一化，但归一化口径会改变长回答的权重，训练记录应明确采用哪一种。若参考回答中的“汽车”在当前位置预测概率为0.2，这个词元贡献$-\log0.2\approx1.6094$；若概率增到0.5，贡献降为约0.6931。负对数概率促使模型提高标注词元的概率，不直接判断标注句子的事实真假。

使用时，语言头把当前位置隐向量映射为$V$个logit，softmax后得到词元分布。选择最大概率词元是贪心解码；按分布抽样可以产生不同措辞；同时保留若干前缀、按累计对数概率选择则是束搜索。选出的词元接回上下文，再计算下一位置，直到生成结束标记或达到长度上限，最后由分词器恢复字符串。教学上若“右侧”“是”“汽车”的条件概率分别为0.8、0.9、0.7，这条三词元路径的概率为0.504，尚须乘上结束词元概率才是完整回答概率。实际分词未必把这些词各切成一个单元，例子只展示从局部分布到完整序列的计算。

\subsubsection{BLIP-2的查询提取与视觉前缀}
\label{sec:vision-models-blip2}

图像编码器输出数百个特征，已有语言模型却没有学过怎样读取它们。BLIP-2在二者之间训练查询Transformer（querying transformer, Q-Former），即用一组可学习查询从图像特征中抽取信息的网络。查询在这里是模型参数，不是用户的问题；其输出是一组随图片改变的连续表示\scite{vision-blip22023}

设图像特征为$\bm F\in\mathbb R^{n_{\mathrm v}\times d_{\mathrm v}}$，查询矩阵为$\bm Q_0\in\mathbb R^{r\times d_{\mathrm q}}$。查询通过自注意力交换信息，并以交叉注意力读取$\bm F$，形成$\bm Z\in\mathbb R^{r\times d_{\mathrm q}}$。原论文使用32个查询，无论图像编码器输出多少特征，交给下一模块的查询输出数量保持固定。压缩减少后续语言计算所读的视觉单元，也要求查询选择并保留与任务相关的信息。

这里可以用式\eqref{eq:vision-models-cross-read}跟踪一次查询读取。把$32\times768$查询隐表示投影为查询，把冻结图像特征投影为键和值，就得到$32\times n_{\mathrm v}$的权重矩阵。每行读取图像网格的一组加权值，再通过残差与后续层形成查询输出。以两个视觉值$(1,0)^\top$、$(0,1)^\top$及单个教学查询为例，若读取权重为$(0.25,0.75)$，输出贡献就是$(0.25,0.75)^\top$；这是信息汇聚，不是自动裁出右车。Q-Former的共享自注意力层从预训练BERT初始化，新增交叉注意力随机初始化，查询本身也作为参数学习。已有图像编码器来自CLIP或EVA-CLIP，后接的语言权重来自OPT或FlanT5，不是所有模块都从零开始。

第一阶段连接冻结图像编码器和Q-Former，使用图文对同时学习三种目标。对比目标让图像查询和文字分别编码，二者在自注意力中不能互看，以免把待匹配文本直接泄漏给图像表示；每段文字与多个查询输出比较，取最高相似度作为图文分数。配对判定目标则允许查询和文字双向交互，由各查询输出的二分类logit取平均，判断图文是否匹配，并使用困难负例。

第三种目标是图像条件文字生成。查询之间可以互看，但不能读取目标文字；每个文字词元可以读取全部查询和此前文字。由于文字不能直接读取冻结图像特征，生成目标迫使Q-Former经过查询路径传递有用信息。三种目标共用Q-Former参数，但注意力掩码不同，分别约束独立匹配、细致核验和信息表达。

设归一化查询输出为$\bm z_{i,k}$、文字表示为$\bm t_j$，独立匹配的分数是$\max_{1\leq k\leq r}\bm z_{i,k}^{\top}\bm t_j$，再按式\eqref{eq:vision-clip-loss}形成批次目标。配对判定则对一条图文组合给出两类概率$(p_0,p_1)$，真实配对标签$b=1$、采样负例标签$b=0$，误差为$-b\log p_1-(1-b)\log p_0$。条件文字生成用式\eqref{eq:vision-vlm-text-objective}，但此时预测者是Q-Former的文字分支。三个可计算误差相加，再把梯度传回共享层、查询和各目标的输出层，图像编码器权重不变。最大相似度让最相关查询承担该配对的直接比较梯度，并不强制32个查询分别对应32类对象。

第二阶段用$\bm W\in\mathbb R^{d_{\mathrm l}\times d_{\mathrm q}}$把$\bm Z$线性投影到语言模型的嵌入维数$d_{\mathrm l}$，得到$\bm P=\bm Z\bm W^\top\in\mathbb R^{r\times d_{\mathrm l}}$。这些向量放在文字输入前面，称为\term{视觉前缀}（visual prefix）：它们占据语言模型的输入位置，用连续向量提供图像条件，却不对应可直接打印的文字。冻结图像编码器和大型语言模型后，生成损失训练Q-Former与投影层。对于解码器式语言模型，目标是根据视觉前缀生成图像描述；对于编码器—解码器式语言模型，论文还使用前缀语言建模，把部分文字与视觉前缀交给编码器，预测剩余文字。

第一阶段保留训练好的Q-Former作为第二阶段初始化，另加语言宽度的投影。每个批次重算图像查询、视觉前缀和参考文字的概率，误差经冻结语言层反传到$\bm P$，再更新投影与Q-Former。原论文用AdamW和预定学习率日程，第一阶段250K步、第二阶段80K步；这些停止点是训练配置，不是证明目标已全局最小。图文描述可以来自原始网页，也可以由论文的描述生成与筛选流程获得，合成文字不应被当作无误标注。最终保留图像编码器、查询模块、投影、语言模型与分词器，预测不再运行第一阶段的配对二分类头。

在街景问答中，模型先计算图像特征，再取得32个查询输出，接上“右侧停着什么车”及相应提示，由语言模型逐词生成答案。固定图片的视觉前缀可用于不同文字请求，但压缩时未保留的细节无法仅靠换一种问法恢复。这里的冻结说明针对两阶段预训练，不能推广到原论文所有下游微调设置，更不能作为所有VLM的固定训练顺序。

\subsubsection{Flamingo的视觉重采样与交错上下文}
\label{sec:vision-models-flamingo}

若输入包含“示例图像—示例回答—待答图像—问题”，视觉信息还需要与文字在序列中的位置对应。Flamingo先把每幅图像或视频压缩为固定数量的视觉单元，再让语言模型在中间层读取这些单元。原模型冻结已有视觉编码器与语言模型，训练新增连接部分\scite{vision-flamingo2022}

视觉编码器输出空间特征；视频则对采样帧分别编码，并加入时间嵌入。感知重采样器（Perceiver Resampler）使用可学习查询，通过注意力把数量可变的特征重采样为64个输出。这里的“重采样”是学习式信息汇聚，不是从原图均匀抽取64个像素。其注意力的键和值还包含查询隐表示本身，使这些潜在单元能够参与后续汇聚。

与BLIP-2把视觉表示放在语言输入前面不同，Flamingo在冻结的语言层之间插入门控交叉注意力与前馈模块。交叉注意力的查询来自当前语言隐状态，键和值来自重采样后的视觉表示。若用$C(\bm H,\bm Z)$概括一个新增交叉注意力分支，残差更新的核心形式为
\begin{equation}
  \bm H'=\bm H+\tanh(\alpha)\,C(\bm H,\bm Z)\text{，}
  \label{eq:vision-flamingo-gate}
\end{equation}
其中$\alpha$是从0开始学习的标量门控。初始时新增分支对输出的贡献为0，语言模型保留原来的计算；门控和连接权重随后从文字预测损失中学习。完整模块还包含相应的前馈分支，不应把式\eqref{eq:vision-flamingo-gate}当成整层结构。

初始贡献为0不意味着门控也学不到信号。若把后续误差对$\bm H'$的导数记为$\bm G$，则$\partial L/\partial\alpha=(1-\tanh^2\alpha)\langle\bm G,C(\bm H,\bm Z)\rangle$，其中内积对矩阵全部元素求和。$\alpha=0$时这个导数仍可非零；而该分支权重的直接梯度先乘以$\tanh\alpha$，在门控打开后才逐渐获得信号。教学上若原隐向量为$(1,0)^\top$、分支贡献为$(0,2)^\top$，门控从0变为0.2后，更新结果为$(1,0.3948)^\top$。这说明新视觉信息怎样进入冻结语言计算，不是说每个门控实际都朝正方向变化。

交错输入还需要控制可见范围。原Flamingo的每个文字位置直接交叉注意到它之前最近的一幅图像或一个视频的视觉表示，而不是同时直接读取全部历史图片。历史图像仍可通过此前文字隐状态和语言自注意力间接影响结果。因此，示例中的红车和待答图中的红车虽然类别相同，也必须保持来源位置，不能把示例属性无条件移到目标上。

训练数据包含网页中的交错图文、图文对及视频文字对，目标仍是对可见视觉上下文条件下的文字进行预测。使用时，把几个已知示例和新请求放入上下文即可，示例本身不引起梯度更新。它们帮助说明当前要做什么任务，但示例数、输入处理和上下文压缩仍限制能保存的细节。图\ref{fig:vision-language-connectors}把前缀、层间读取与直接投影三条路径放在一起，便于区分视觉信息实际进入语言计算的位置。

一条交错训练记录还要保存视觉单元插入文字的位置，才能形成上述可见掩码。对每个数据来源分别计算文字负对数概率，再按来源权重累积梯度，更新重采样器、新增门控连接层及相应新增词元；原NFNet视觉权重和Chinchilla语言权重保持冻结。训练到预定日程后保留这些模块。推理把示例图$A$、示例答案“车窗关闭”、目标图$B$和问题依次放入上下文，目标答案位置直接读取$B$的64个视觉单元，再逐词解码；$A$的答案提供格式示范，却不能代替观察$B$的车窗。闭集任务还可对每个候选答案累积条件对数概率，按相同长度处理规则选择，而不是必须自由生成一句长回答。

\subsubsection{LLaVA的特征投影与视觉指令适配}
\label{sec:vision-models-llava}

若保留图像的网格特征，可以直接把每个视觉单元映射到语言嵌入空间，再让语言主干学习怎样使用这些输入。原始LLaVA采用这一连接方式，并用视觉指令数据训练回答行为。本节依据原论文v2中的线性投影设置，后来的多层感知机连接器不属于这里的模型\scite{vision-llava2023}

CLIP视觉编码器从图片取得网格特征$\bm F$，投影矩阵把各位置的特征转为与文字嵌入同宽的$\bm P$。这些视觉单元与问题词元一起进入自回归语言模型。投影并没有把每个图像块翻译成一个汉字或单词，而是调整连续特征的坐标，使语言层可以对它们计算注意力。

设$\bm W\in\mathbb R^{d_{\mathrm l}\times d_{\mathrm v}}$，则采用行存特征时$\bm P=\bm F\bm W^\top$。网格顺序和图像输入边界帮助模型知道哪些位置属于图片；语言自注意力随后在视觉位置与文字位置之间计算，和Flamingo另设一条交叉注意力分支不同。以两个二维教学网格$(1,0)$、$(0,1)$及$\bm W=\left[\begin{smallmatrix}2&0\\0&1\end{smallmatrix}\right]$为例，投影得到$(2,0)$、$(0,1)$。投影改变坐标宽度或尺度，但没有额外读取像素，也没有建立“第一行一定是自行车”的物体语义。

训练分为特征对齐和视觉指令适配。对齐阶段用图像描述构造“请简要描述图片”一类请求，冻结视觉编码器与语言模型，只更新投影层。语言模型的误差仍需沿计算图传回投影输入，冻结权重并不等于可以截断这条反向传播路径。指令适配阶段继续冻结视觉编码器，同时更新投影与语言模型，数据则包括对话、详细描述和复杂问答。

对多轮数据，输入保留图片、用户请求和前轮回答，损失监督助手应生成的回答词元。若请求是“比较两辆车的状态”，模型依赖同一视觉网格形成两个对象的描述，并逐词预测比较句。训练示例帮助模型学会“比较”这一请求形式，能否正确读取车的位置、朝向和遮挡仍由视觉证据及训练覆盖决定。

原论文的数据构造用图片已有描述和对象框向文字教师提供视觉内容的符号表示，再生成对话、详细描述和复杂问答；教师在这一流程中没有直接看原像素。训练LLaVA时仍把真实图片交给CLIP视觉编码器，不能把教师的输入方式误写成学生只读文字。第一阶段从预训练CLIP与Vicuna加载权重，随机初始化投影，用描述监督提高回答概率；第二阶段从第一阶段投影继续，用助手位置掩码同时更新投影与Vicuna，视觉权重不变。原论文使用AdamW及有限轮次日程，最终保留更新后的语言模型、投影和固定视觉编码器。

冻结视觉端时，$\bm F$可作为常量；若误差对$\bm P$的导数为$\bm G\in\mathbb R^{n_{\mathrm v}\times d_{\mathrm l}}$，投影的梯度就是$\bm G^\top\bm F$，形状与$\bm W$相同。这给出“语言误差怎样训练连接器”的直接依据。第二阶段再允许语言权重改变，模型便同时学习读取视觉前缀和组织请求答案。使用时先得到网格投影，接入问题与角色标记，再按前面的逐词规则生成回答。预测出的“右侧汽车车窗关闭”应回查目标图；语言损失只对训练回答负责，不能保证任意新图片的细节判断正确。

原模型可以生成自然语言，但没有因此得到任意精度的坐标、实例掩码或图像生成能力。若回答写了“右车正在行驶”，单张静态图片只提供姿态和场景线索，通常不足以确定实际速度；模型必须避免把语言中常见的情境补全当成观测。

\begin{figure*}[htbp]
  \centering
  % 每行仅保留用于区分接入位置的模块，不展开完整Transformer层。
\begin{tikzpicture}[x=1cm,y=1cm,font=\figurefont\small,
  block/.style={draw=FigureModel,fill=FigureModel!10,line width=0.8pt,
    minimum height=9mm,minimum width=21mm,align=center},
  input/.style={block,draw=FigureInput,fill=FigureInput!10},
  output/.style={block,draw=FigureOutput,fill=FigureOutput!10},
  arr/.style={-{Stealth[length=2mm]},BookInk,line width=0.8pt}]
  \node[anchor=west,font=\figurefont\small\bfseries] at (0,6.7) {BLIP-2：查询压缩后作为输入前缀};
  \node[input] (bf) at (1.1,5.7) {视觉特征};
  \node[block] (bq) at (4.1,5.7) {Q-Former};
  \node[block] (bp) at (7.2,5.7) {投影后的\\视觉前缀};
  \node[block] (bl) at (10.7,5.7) {语言模型};
  \node[output] (bo) at (14.1,5.7) {回答词元};
  \draw[arr] (bf) -- (bq);
  \draw[arr] (bq) -- (bp);
  \draw[arr] (bp) -- (bl);
  \draw[arr] (bl) -- (bo);
  \node[anchor=north] at (4.1,5.15) {可学习查询读取图像};
  \node (btext) at (10.7,4.7) {文字问题};
  \draw[arr] (btext) -- (bl);

  \node[anchor=west,font=\figurefont\small\bfseries] at (0,3.9) {Flamingo：在语言层之间读取视觉表示};
  \node[input] (ff) at (1.1,2.45) {视觉特征};
  \node[block] (fr) at (4.1,2.45) {Perceiver\\Resampler};
  \node[block] (fc) at (7.2,2.45) {门控交叉\\注意力};
  \node[block] (fl) at (10.7,2.45) {后续语言层};
  \node[output] (fo) at (14.1,2.45) {回答词元};
  \draw[arr] (ff) -- (fr);
  \draw[arr] (fr) -- node[above] {$K,V$} (fc);
  \draw[arr] (fc) -- (fl);
  \draw[arr] (fl) -- (fo);
  \node (fh) at (7.2,3.35) {语言隐状态$Q$};
  \draw[arr] (fh) -- (fc);
  \node[anchor=north] at (4.1,1.9) {每个图像／视频64个单元};

  \node[anchor=west,font=\figurefont\small\bfseries] at (0,0.7) {原始LLaVA：网格特征投影后进入共享序列};
  \node[input] (lf) at (1.1,-0.25) {视觉网格};
  \node[block] (lp) at (4.1,-0.25) {线性投影};
  \node[block] (ls) at (7.2,-0.25) {视觉单元\\与问题序列};
  \node[block] (ll) at (10.7,-0.25) {语言模型};
  \node[output] (lo) at (14.1,-0.25) {回答词元};
  \draw[arr] (lf) -- (lp);
  \draw[arr] (lp) -- (ls);
  \draw[arr] (ls) -- (ll);
  \draw[arr] (ll) -- (lo);
\end{tikzpicture}

  \caption{三种视觉接入方式的结构示意，按各原论文的预训练路径绘制。BLIP-2由少量查询形成视觉前缀，Flamingo在语言层中新增读取视觉表示的交叉注意力，原始LLaVA把视觉网格线性投影后交给语言主干。图中省略各模块内部层，框旁文字指出视觉信息进入的位置；冻结范围由正文所述阶段决定。}
  \label{fig:vision-language-connectors}
\end{figure*}
\FloatBarrier

\subsubsection{Qwen3-VL的多层视觉接入与视频时间组织}
\label{sec:vision-models-qwen3vl}

多页文档和视频包含分散证据，只读取视觉编码器的最末层、再把所有帧当成无时间含义的图片，可能丢失细节与先后关系。Qwen3-VL技术报告v1在动态分辨率编码、视觉合并、跨层接入和时间表示上组织这些信息。报告区分思考版与非思考版，以下机制不意味着两类后训练版本具有相同响应行为。

视觉编码器接受动态分辨率输入，输出长度随图像大小变化的视觉特征。视觉语言合并模块使用多层感知机（multilayer perceptron, MLP），即在逐位置的线性映射之间加入非线性运算；这里用两层映射把相邻$2\times2$特征压缩为一个语言宽度的视觉单元\scite{vision-qwen3vl2025}
空间合并降低上下文长度，同时也把四个位置的信息汇集到一个位置表示中；实际细节保留需结合输入分辨率和训练检查。

多层视觉接入DeepStack进一步利用视觉编码器三个不同深度的中间特征。每组特征经独立合并模块映射后，加到语言主干前三层相应的视觉位置隐状态上。视觉信息由多个深度进入语言计算，而不是全部仅依赖初始输入。残差接入不增加额外序列位置，但仍增加了特征投影和存储等计算。

位置表示有两项不同安排。交错多模态旋转位置编码（interleaved multimodal rotary position embedding, interleaved MRoPE）按各坐标旋转查询与键的部分分量，使比较分数能够反映相对位置。它把时间、高度、宽度对应的旋转位置维度分散到不同频率范围，避免把某个坐标维度集中在单一频段。视频时间戳则以文字形式放在每个时间块之前，例如“3.0秒”，提供明确的时间读数。两者共同使用，但文字时间戳不是从MRoPE隐向量中自动打印出来的数值。

在一个教学视频中，采样块对应1.0、3.0和8.0秒。保留这些时间戳，模型才有条件判断后两个观测相隔5秒；只给顺序编号会把这种间隔信息删掉。不过，如果敲门动作发生在4.2秒且没有进入任何采样块，时间戳再准确也无法补回动作。多图输入同样要保留图号，避免把第二幅图的车窗状态归到第一幅图。

报告的预训练先只更新合并模块，冻结视觉编码器和语言主干，随后转为全部模块训练，并逐步扩展上下文窗口。后训练包含指令适配、蒸馏与强化学习等安排，思考版使用相应的推理形式数据，非思考版使用标准回答形式。使用时应选定具体检查点、格式和观察预算。报告支持的256K词元上下文是容量约定，不是对任意长文档小字和短暂视频事件的保存保证。

这些阶段的样本仍需要明确输入与目标。图像描述样本以图片为条件、描述为目标，视频问答样本以带时间戳的帧组和问题为条件、答案为目标，交错文档样本按顺序保留图片与文字。它们用文字预测误差训练，但来源、目标位置及长度不同，不能把每张图都当作一个类别标签。视觉端由SigLIP-2检查点继续训练，语言端由Qwen3初始化，合并模块将空间与层间特征接入语言计算；预热阶段误差只改变合并权重，全参数阶段再共同改变三个模块。扩大上下文改变的是一次可计算序列的允许长度，并不替代更密采帧或读取原图。

多图比较时，模型先分别编码两幅街景，以图号和输入边界标记组织视觉块，再把问题接入；答案中的“第二张图车窗打开”需要依赖第二图的实际特征。视频步骤问答则在每个帧组前接入实际时间文字，模型生成事件与时间描述时使用这些序列条件。语言头形成候选词元后，仍按式\eqref{eq:vision-vlm-language}解码，时间字符串不是额外测量器自动返回的真值。最终保留具体思考或非思考检查点及其处理器；后训练产生的响应形式不能反推输入采样一定保留了所问动作。

\subsubsection{Qwen2.5-Omni的音画输入与文字语音响应}
\label{sec:vision-models-qwen-omni}

视频中的敲门动作和敲击声可能互相补充，也可能错位或来自画外。Qwen2.5-Omni技术报告v1通过Thinker--Talker结构处理这类输入并输出文字与语音。Thinker读取多模态上下文、生成文字及隐表示，Talker利用这些信息生成语音编码，再恢复成声音\scite{vision-qwenomni2025}

输入音频被转换为梅尔频谱，经音频编码器形成时间特征；图像和视频由视觉编码器处理。时间对齐多模态旋转位置编码（time-aligned multimodal rotary position embedding, TMRoPE）给音视频提供共同时间尺度，报告中一个时间位置单位对应40毫秒。音频单元沿时间排列，视频单元同时保留帧时间和空间位置。输入还按实际时间划分为2秒块，每块内先放视觉表示、后放音频表示，再按块交错。

因此，序列中相邻的视觉和声音单元不一定只相差一个物理采样点，实际时间由位置记录共同表达。这个安排也不证明两段内容来自同一个声源。若门外另有人敲击金属，画面中的手部动作与声音可能在时间上接近，却缺乏确定来源。回答应分别依据看见的动作、听见的声音和额外文字线索。

Thinker按式\eqref{eq:vision-vlm-language}的同类方式预测回答文字。Talker接收Thinker的高层隐表示以及已采样文字词元的嵌入，以双轨自回归方式组织文字和语音编码。隐表示提供语义、语气等上下文，离散文字则帮助确定具体词语，避免语义接近但读音不同的内容仅由连续表示区分。

语音编码还不是可以播放的波形。报告的流式解码路径用流匹配模型将编码转为梅尔频谱，再由声码器合成波形；\ref{sec:vision-models-flow-matching}节将解释所需的速度预测目标。解码采用有限块上下文，包含回看和一定前看，因此流式表示分块产生结果，并不意味着没有等待。

训练职责也有区别。多模态预训练起步时固定语言模型，先适配连接再训练相关编码器，随后展开全参数与长上下文训练；Thinker再用多模态指令数据适配回答。Talker单独经历语音续接、偏好优化和说话人适配等训练安排，以学习从语义上下文到语音的生成。理解监督不能替代语音发音监督，音画时间对齐也不能替代输出声学解码。

以文字与对应朗读音频为训练记录，先用语音编码器把真实声音转为离散编码，再以此前编码、文字与Thinker隐表示预测下一个语音编码；目标编码的负对数概率提供Talker的监督。它与回答文字的目标可以相互联系，但真实语音决定音色、节奏与发音的参考，不能只用一条文字回答推出全部声学目标。波形解码器另以编码和真实频谱为配对学习频谱恢复，声码器学习从频谱到波形。推理时Talker使用自己此前生成的编码，流式解码只读取约定的有限块，最终得到可播放的采样序列，而不是把Thinker的隐藏向量直接保存成音频文件。

在敲门例子中，输入路径保留动作和声音，Thinker形成“画面中的人抬手敲门，并伴有敲击声”这样的候选回答，Talker把确定的回答组织为语音。该版本支持文字与语音输出，不因此具有图像或视频生成模块。首段声音的到达时间还包含输入编码、文字生成、语音编码预测及波形解码等开销。

\subsubsection{观察决策、视觉工具返回与后续语言计算}
\label{sec:vision-models-observation}
\label{sec:vision-models-countercredit}

如果整图缩小后右车的车窗只剩几像素，语言模型继续生成更长解释仍没有得到新的观测。另一种计算是暂停回答，指定原图中的车窗区域，由工具返回该区域的像素，再继续语言预测。观察动作（visual observation action）指这次取回视觉内容的决定，不是答案中的一句“我再仔细看看”。它改变后续计算能够读取的信息，而仅延长原上下文中的文字生成不改变像素来源。

一次调用可以记录为$(h_i,c_i,o_i)$：$h_i$是调用前已有的完整状态，包括请求、先前图片与返回记录、已生成文字；$c_i$是此次动作及参数；$o_i$是实际返回图像。裁剪参数还需原始图片或先前观察的来源编号、坐标范围和输出尺寸。若在先前裁剪中再次裁剪，应组合两次坐标变换才能回到原图，而不能把第二次的局部坐标直接当作原图坐标。工具实际取出像素后，模型重新编码$o_i$，把其表示放入相应图像轮次，随后才预测后续词元。越界、空区域或调用格式无效时，没有可用的新观察，系统应保留失败记录，而不是默默把失败当成已经看到了对象。

有了这个接口，还要学习什么时候值得调用。最终答对只能说明整条路径的结果，不能说明这次裁剪有必要，也不能说明答案真的依赖返回像素。2026年预印本CounterCredit在已经发生的调用前状态上，分别检查“继续这条视觉分支比立即回答增加了多少正确答案得分”和“真实像素比同格式随机返回增加了多少得分”。本节依据其v1方法、训练配置和返回格式，范围限于返回单幅裁剪图像的工具\scite{vision-countercredit-2026}。

训练记录包含原图$\bm X$、问题$q$和参考答案$y$。策略$\pi_{\bm\theta}$从已有视觉语言检查点生成交错文字与裁剪的轨迹$\xi$，最后按任务规则取得正确性回报$r_{\mathrm{out}}(\xi)\in\{0,1\}$。为了比较同一调用的不同返回，论文不重新生成一段不同的调用前解释，而是固定已实现的$h_i$。真实分支接入原动作$c_i$与原返回$o_i$；立即回答分支从$h_i$直接要求回答；随机分支固定$h_i$与$c_i$及其文字参数，只把返回像素替换成从原图均匀采样的同尺寸图块$\widetilde o_{i,k}$。随机图块不事先筛成“保证无关”，因而可能偶然含有答案证据，这种随机参照的限制必须保留。

先明确比较用的可计算量。对于多选问题，设$O$为题目给出的选项字母集合，$e$为固定的立即作答提示，则参考答案得分为
\begin{equation}
  U_{\bm\theta}(p)
   =\frac{\pi_{\bm\theta}(y\mid p,e)}
          {\sum_{a\in O}\pi_{\bm\theta}(a\mid p,e)}.
  \label{eq:vision-countercredit-score}
\end{equation}
$p$是待探测的状态前缀，分母只在选项字母间归一化，三个分支使用相同提示与答案起始格式。得分不是最终抽样答对的频率，也不是外部判断器给整段推理打分。自由回答没有这个有限选项分母，论文改用参考答案在真实前缀监督下的每词元平均对数概率，单位为每词元奈特；它与$[0,1]$概率差的尺度不同，论文分别校准多选死区0.05与自由回答死区0.44奈特每词元。探测得分停止梯度，只用来计算策略训练信号。

令$U_i^{\mathrm{real}}=U_{\bm\theta}(h_i,c_i,o_i)$、$U_i^{\mathrm{now}}=U_{\bm\theta}(h_i)$，并用$K$个随机返回估计参照均值。两项差值为
\begin{equation}
  \begin{aligned}
    D_i&=U_i^{\mathrm{real}}-U_i^{\mathrm{now}},\\
    \overline U_i^{\mathrm{rand}}
      &=\frac1K\sum_{k=1}^{K}
          U_{\bm\theta}(h_i,c_i,\widetilde o_{i,k}),\\
    E_i&=U_i^{\mathrm{real}}-\overline U_i^{\mathrm{rand}}.
  \end{aligned}
  \label{eq:vision-countercredit-values}
\end{equation}
$D_i$称为决策价值（decision value），比较实际视觉分支与立即作答；$E_i$称为证据价值（evidence value），比较实际像素与同动作下的随机返回。前者还包含动作格式或延后回答的影响，后者尽量保持这些因素不变。它们检查的是当前策略的参考答案得分变化；即使原来已经能答对，增加置信度仍可能产生正值，因此论文的“需要”不能直接改写成逻辑上“不看必错”。

只奖励其中一项会遗漏不同失败情形。若$D_i$大而$E_i$接近0，随机图块也带来同样收益，这次调用没有显示真实像素的额外作用。若$E_i$大而$D_i$接近0，像素相对于随机返回有信息，但立即作答已同样充分。论文因此取$a_i$为探测资格指示，包括有效单图返回、存在不同随机图块和得分有限，取$\varepsilon$为预先校准的误差死区，并定义
\begin{equation}
  \begin{aligned}
    b_i&=\min\{E_i,D_i\}
        =U_i^{\mathrm{real}}
          -\max\{U_i^{\mathrm{now}},\overline U_i^{\mathrm{rand}}\},\\
    v_i&=r_{\mathrm{out}}(\xi)\,a_i\,\mathbb I[b_i>\varepsilon].
  \end{aligned}
  \label{eq:vision-countercredit-verification}
\end{equation}
这里$v_i=1$才表示此次调用获得双重核验。第二个等式来自同时减去两个参照时取较小差值，等价于比较较高参照得分；所以任何一边没有超过死区都不能通过。正确性因子还要求整条轨迹最终答对。$\varepsilon$用训练前的随机图块对照校准后固定，避免把有限$K$导致的小波动全当成收益，且不同得分单位使用不同值。

仍以车窗问题作教学计算。假设立即回答的正确选项得分为0.40，实际裁剪后为0.80，三个随机返回的得分为0.35、0.45、0.40，则$\overline U^{\mathrm{rand}}=0.40$，$D=E=b=0.40$。取多选配置$\varepsilon=0.05$，有效调用且最终答对时即可通过。若实际与随机返回都为0.80，则$D=0.40$而$E=0$，不能通过；若立即回答为0.82、实际为0.80、随机均值为0.40，则$E=0.40$而$D=-0.02$，也不能通过。三种情形固定了同一个问题及像素来源，展示两个信号为何不能互相替代。

调用收费还要防止模型仅靠减少成本而偏好错误答案。设$n(\xi)$为已执行调用数，$\mathcal V(\xi)=\{i:v_i=1\}$，核验通过的调用获得超过门槛部分的返还，未通过的调用付费。论文的轨迹价格与总回报为
\begin{equation}
  \begin{aligned}
    P(\xi)
      &=\min\left\{C,\,
         \gamma\sum_{i\in\mathcal V(\xi)}(b_i-\varepsilon)\right\}\\
      &\quad-\rho\bigl(n(\xi)-|\mathcal V(\xi)|\bigr),\\
    R(\xi)&=r_{\mathrm{out}}(\xi)+P(\xi).
  \end{aligned}
  \label{eq:vision-countercredit-price}
\end{equation}
$\gamma$把得分收益换成回报，$C$是每条轨迹返还上限，$\rho>0$是每个未核验调用的费用。失败调用和错误答案轨迹中的调用也计费，不通过便不能免费。若每条轨迹最多$n_{\max}$次调用且$\rho n_{\max}<1$，正确轨迹至少获得$1-\rho n_{\max}>0$，错误轨迹没有返还、总回报至多0，因此费用不改变正确与错误轨迹的排序。这只证明给定回报的界，不保证优化必然学会正确策略。

用论文多选配置$\gamma=0.5$、$\rho=0.05$、$C=0.7$、$n_{\max}=6$继续复算。上述$b=0.40$的调用返还$0.5(0.40-0.05)=0.175$；若随后还有一个未通过调用，轨迹价格为$0.175-0.05=0.125$，答对时总回报为1.125。六个调用全部未通过但最终答对，回报仍为0.70；零调用但答错则为0。这里费用单位是训练回报，不是货币或实测运行秒数，部署成本还需另外记录。

策略更新使用组相对策略优化（group relative policy optimization, GRPO），即对同一个问题抽样一组回答轨迹，用组内相对回报决定哪些生成动作应提高概率。CounterCredit把正确性和价格分开处理。若第$j$条轨迹的正确性为$O_j$、价格为$P_j$，组均值为$\overline O,\overline P$，$\sigma_{\min}>0$为避免零分母的标准差下限，论文取0.15，其优势为
\begin{equation}
  \widehat A_j=
    \frac{O_j-\overline O}
         {\max\{\operatorname{std}(O),\sigma_{\min}\}}
       +(P_j-\overline P).
  \label{eq:vision-countercredit-advantage}
\end{equation}
正确性项除以组内标准差，价格项只减均值、不除标准差，所以多一次未核验调用造成的$\rho$差值仍保持原回报尺度。组内全错时，论文把价格优势设为0，避免只根据错误轨迹中的调用数学习。根据优势提高或降低动作概率的策略学习原理可参考\BookChapterRef{chap:rl-learning}{卷七“学习”章}；这里新增的是价格如何进入已有策略更新，而不是另一种视觉编码器。

为了把这个优势接回参数，令$r_{j,k}=\pi_{\bm\theta}(a_{j,k}\mid h_{j,k})/\pi_{\mathrm{old}}(a_{j,k}\mid h_{j,k})$为轨迹$j$在生成位置$k$的概率比。对位置取平均的目标项为$\min\{r_{j,k}\widehat A_j,\operatorname{clip}(r_{j,k},1-\delta,1+\delta)\widehat A_j\}$，$\delta$限制概率比的有效更新幅度；优化器沿增大该项的方向调整策略权重。参考答案、随机图块及探测得分不通过这个目标反向传播，它们先形成停止梯度的优势，真正的梯度来自已抽样动作的对数概率。新的策略再生成下一批轨迹，如此重复到预定训练日程，而不是一直调用到探测得分不再变化。

原实验从Qwen2.5-VL-7B-Instruct的裁剪轨迹监督微调开始，固定视觉编码器与投影器，随后在同一格式和调用预算下做三轮GRPO训练。最终保留学到的策略及裁剪接口，推理时不提供参考答案、不运行这些对照探测；只根据当前状态选择立即回答或实际观察。训练探测每个合格调用还需$K+2$次得分前向，不能把它计作推理免费取得的真值。本文接口可以由后续问题驱动的证据获取回用，但论文当前信号不直接覆盖搜索、OCR文字返回或任意外部工具；随机参照偏差、策略得分误校准和跨模态分布变化也会限制核验的解释。

\subsection{媒体内容的生成}
\label{sec:vision-models-media-generation}

生成一幅雨天街景时，仅有“雨天”这个语义判断还不够，模型需要决定每处颜色、纹理与空间关系。直接在像素中生成可以保留细粒度自由度，但计算量大；在压缩表示中生成可以减少计算，却必须有相应解码器。条件信息怎样进入生成过程，又形成另一项独立选择。

\subsubsection{LDM的空间潜表示与文字条件接入}
\label{sec:vision-models-ldm}

潜空间扩散模型（latent diffusion model, LDM）先学习图像压缩与重建，再在压缩空间中学习生成。这里的潜表示保留二维网格结构，与CLIP把整图汇总成一个语义向量的空间不同。原论文的文字条件LDM通过交叉注意力把文字表示送入去噪网络\scite{pgm-generative-rombach2022}

图像自编码器由编码器$E$与解码器$D$组成。编码得到$\bm Z=E(\bm X)\in\mathbb R^{h'\times w'\times d_z}$，解码得到$D(\bm Z)$。自编码器（autoencoder）通过还原输入来学习压缩坐标，不等于按类别汇总一幅图片。自编码阶段使用重建相关的感知与对抗目标，并配合潜空间正则化，学习在压缩后保留可重建的信息。该阶段完成后固定自编码器，在其潜空间中训练扩散网络，不把图像重建器和每个生成任务从头共同重训。

重建训练的输入与目标都是同一张图：输入先经$E$压缩，再经$D$还原，像素或已有感知网络特征的差异衡量保留信息的误差。感知误差比较的是两图特征，不要求逐像素相同；局部判别器则根据真实与重建图块学习区分，编码器和解码器反过来提高重建图的真实程度。KL正则版本限制潜分布偏离标准高斯的程度，向量量化版本则学习有限代码及其对应，这两种约束不能当作同一个目标。原论文给出不同自编码器设置，选择哪一套便同时确定了后续潜变量尺度和解码器。这个阶段提高重建保真度，不直接训练“根据雨天文字创造新场景”。

为说明计算对象，把潜数组按固定次序展开为向量$\bm z$。取噪声步$k$，给它加入标准高斯噪声$\bm\epsilon$：
\[
  \bm z^{(k)}
  =\sqrt{\bar\alpha_k}\bm z+
    \sqrt{1-\bar\alpha_k}\bm\epsilon\text{，}
\]
其中$\bar\alpha_k$由噪声日程指定，$(k)$表示加噪或去噪过程的步，不是视频帧。文字请求$c$经条件编码器得到$\bm C=E_{\bm\phi}(c)$。去噪网络以当前潜表示、噪声步和文字为条件预测所加噪声，目标为
\begin{equation}
  L_{\mathrm{LDM}}
  =\mathbb E_{\bm X,c,k,\bm\epsilon}
    \left[
      \left\lVert\bm\epsilon-
      \epsilon_{\bm\theta}
      \bigl(\bm z^{(k)},k,E_{\bm\phi}(c)\bigr)
      \right\rVert_2^2
    \right]\text{。}
  \label{eq:vision-ldm-objective}
\end{equation}
原论文的条件设置共同优化$\bm\theta$与$\bm\phi$；后来使用冻结文字编码器的实现，需要另行说明。

这里$\bar\alpha_k=\prod_{j=1}^{k}\alpha_j$，$\bar\alpha_0=1$，$0<\alpha_j<1$，每个$\alpha_j$控制一次保留原状态的比例。独立高斯加噪的合成仍为高斯，所以可以从干净潜表示一次构造任意$k$的带噪输入，而不必真的先算完前$k$步。每次训练先抽图文对，再抽$k$和同维标准噪声，编码干净潜网格后直接计算$\bm z^{(k)}$。网络输出与噪声同形，平方误差对输出的导数为$2(\widehat{\bm\epsilon}-\bm\epsilon)$，优化器把这个差传到去噪网络、文字编码器和连接投影。固定的图像自编码器提供目标坐标，不参与这次更新。以有限批次均值估计式\eqref{eq:vision-ldm-objective}，按所选训练日程和验证结果保留检查点，而不是把抽样期望当成已精确求解的总体目标。

交叉注意力中，查询来自当前去噪网络的空间特征，键和值来自文字序列。假设当前有$n_z$个空间单元、文字有$m$个词元，注意力权重是$n_z\times m$的矩阵，每个空间位置可以按当前状态读取不同文字信息。它没有直接为“红色”指定一个真实对象框，词语与空间的正确对应仍需通过数据和目标学习。

使用时如何从预测噪声得到更干净状态，还需要一步。把带噪关系中的噪声替换成预测值后，先估计干净潜向量：
\begin{equation}
  \widehat{\bm z}
   =\frac{\bm z^{(k)}
       -\sqrt{1-\bar\alpha_k}\,
           \bm\epsilon_{\bm\theta}(\bm z^{(k)},k,\bm C)}
          {\sqrt{\bar\alpha_k}}.
  \label{eq:vision-ldm-clean-estimate}
\end{equation}
这来自对前向加噪等式移项，并非使用时已知真实$\bm z$。若采用标准高斯反向采样，把$\beta_k=1-\alpha_k$和
$\widetilde\beta_k=\beta_k(1-\bar\alpha_{k-1})/(1-\bar\alpha_k)$记为噪声方差系数，反向均值与更新可写为
\begin{equation}
  \begin{aligned}
    \bm\mu_k
      &=\frac1{\sqrt{\alpha_k}}
        \left(\bm z^{(k)}
          -\frac{\beta_k}{\sqrt{1-\bar\alpha_k}}\,
             \bm\epsilon_{\bm\theta}(\bm z^{(k)},k,\bm C)\right),\\
    \bm z^{(k-1)}
      &=\bm\mu_k+\sqrt{\widetilde\beta_k}\,\bm\eta_k,
         \qquad \bm\eta_k\sim\mathcal N(\bm0,\bm I).
  \end{aligned}
  \label{eq:vision-ldm-reverse-interface}
\end{equation}
均值形式由前向高斯链的后验及噪声参数化得到；最后一步不再加随机噪声。不同采样器可以改用确定性或跳步更新，所以这个接口说明基本依赖，不规定所有LDM检查点的实际步数。令$K$为所选噪声日程的末步，且$\bar\alpha_K$足够小，生成便从近似标准高斯的$\bm z^{(K)}$开始，按$k=K,\ldots,1$反复预测，再把末态恢复为二维潜网格，用同一$D$解码，反转训练时的像素尺度处理才得到显示图片。

沿街景例子，假设只查看二维潜网格中的一个两维单元，干净值为$(2,-1)^\top$，抽样噪声为$(0,1)^\top$，$\bar\alpha_k=0.64$。训练输入为$(1.6,-0.2)^\top$；若网络预测噪声$(0.1,0.8)^\top$，误差平方和为$0.01+0.04=0.05$，输出导数为$(0.2,-0.4)^\top$。式\eqref{eq:vision-ldm-clean-estimate}恢复的干净估计为$(1.925,-0.85)^\top$，而不是原来的$(2,-1)^\top$。这条计算明确显示噪声误差怎样改变潜表示；它还不是完整街景像素，因为其他网格位置和解码器共同决定最终图像。

使用时先编码“雨天街道，右侧停着红色汽车”，再从潜空间噪声出发多步去噪，最终用$D$恢复图片。若高宽都压缩为原来的八分之一，潜网格位置数变为$1/64$；这并不意味着全部计算或内存恰好减少64倍，因为通道数、网络宽度和注意力实现仍然不同。压缩还可能损失小字等细节，生成后的事实与控制要求由\ref{chap:vision-image-generation}章的任务协议检查。

\subsubsection{unCLIP的图像表示先验与图像解码}
\label{sec:vision-models-unclip}

LDM直接用文字条件预测空间潜表示，unCLIP则在文字与图像之间加入一个整图语义表示。冻结CLIP后，先验模型根据文字预测CLIP图像表示，另一个生成解码器再依据该表示形成图片。这是DALL-E 2报告中的unCLIP路线，生成器并非CLIP本身\scite{vision-unclip2022}

设文字为$c$，CLIP图像表示为$\bm z_{\mathrm i}$，图像为$\bm X$，模型可写成
\begin{equation}
  p(\bm X\mid c)
  =\int p_{\mathrm{dec}}(\bm X\mid\bm z_{\mathrm i},c)\,
          p_{\mathrm{prior}}(\bm z_{\mathrm i}\mid c)\,
          \mathrm d\bm z_{\mathrm i}\text{。}
  \label{eq:vision-unclip-factorization}
\end{equation}
这里表达模型的分层生成分布，不宣称训练后已等于真实数据分布。先验处理“这段文字可能对应哪些图像语义表示”，解码器处理“这个语义表示可能对应哪些具体像素”。后者还可以读取文字，不能把完整条件一律删成只有图像向量。

训练先验时，图文对中的图像经冻结CLIP形成目标向量。论文比较了两条路线：自回归先验先对图像向量作主成分降维与分量量化，再逐项预测离散编码；扩散先验直接对连续图像向量加噪，并在文字条件下回归干净向量。扩散先验采用干净表示预测目标，不能直接照搬LDM的噪声预测符号而不作转换。

自回归先验先在训练图像表示上估计主成分方向，按方差从大到小保留分量，再把每个分量的值分到离散桶。原报告用319个主成分、每分量1024桶。训练以文字、CLIP文字表示及图文点积的量化条件为前缀，逐分量预测正确桶编号，负对数概率之和训练先验Transformer。使用时先选点积条件，再逐项抽桶、取对应数值、逆主成分变换，恢复候选图像表示。量化与截断产生额外误差，不能把所得向量当作原CLIP图像的精确逆。

扩散先验保留连续向量。令$\bm z_{\mathrm i}^{(k)}$为按前节高斯规则加噪的图像表示，文字及CLIP文字表示合写为条件$c$，先验网络$f_{\bm\psi}$的目标为
\begin{equation}
  L_{\mathrm{prior}}
   =\mathbb E_{\bm X,c,k,\bm\epsilon}
      \left[\left\lVert
         f_{\bm\psi}(\bm z_{\mathrm i}^{(k)},k,c)
           -\bm z_{\mathrm i}\right\rVert_2^2\right].
  \label{eq:vision-unclip-prior-objective}
\end{equation}
目标$\bm z_{\mathrm i}$由冻结CLIP计算；先验预测干净向量，误差传回$\bm\psi$，不改变CLIP坐标。采样时，若需要噪声形式，可用
$\widehat{\bm\epsilon}=(\bm z_{\mathrm i}^{(k)}-\sqrt{\bar\alpha_k}\widehat{\bm z}_{\mathrm i})/\sqrt{1-\bar\alpha_k}$换算，再实施反向更新。论文在使用时采两个候选图像表示，按与CLIP文字表示的点积选择其中一个；这是表示阶段的选择，尚没有生成两幅完整图片。

训练生成解码器时，以真实图像的CLIP表示和相应文字为条件，对带噪图像学习去噪。图像表示被投影到时间步嵌入及额外上下文单元，连接到由GLIDE结构改造的图像生成网络；随后使用图像上采样模块恢复更高分辨率。CLIP在先验与解码器训练期间保持冻结，两个生成模块学习不同的映射。

解码训练的目标现在是像素而非整图向量：图文对给出$\bm X$，CLIP给出$\bm z_{\mathrm i}$，对$\bm X$加噪后用$(\bm z_{\mathrm i},c)$条件预测噪声，学习反向图像分布。先验和解码器分别抽样、计算各自误差并按各自日程更新，最后保留两套权重及上采样模型。生成器先形成$64\times64$图片，再经$256\times256$、$1024\times1024$上采样路径恢复高分辨率；后两级收到较低分辨率图像条件，而不是重新从一段文字独立生成。原文的模块形式改自GLIDE，不表示必须把一个现成GLIDE模型不经训练直接接上即可工作。

对“右侧红车、左侧红色自行车的雨天街景”，先验产生的$\widehat{\bm z}_{\mathrm i}$只表达一组整体条件。固定它和文字、换不同图像噪声，解码器可以产生多幅雨天街景。若两候选表示与文字点积为0.72和0.68，选择第一项只说明这一表示比较规则，不能提前判定最终图片一定有两辆车或位置正确。参考图直接提供图像表示时跳过先验，但仍需像素生成与上采样；要保留原图细节，报告的反演路线还记录扩散末端状态等残余信息，而不是仅保存一个CLIP向量。

这一分层路径把语义条件与像素细节分开：相近图像表示可以对应纹理、布局细节不同的结果，表示匹配并不决定完整重建。参考图片的变化可以展示哪些内容在生成中较易保持，但“语义相似”仍不能替代对原图每处细节的检查。

\subsubsection{GLIDE的文字条件与语义引导}
\label{sec:vision-models-glide}

条件输入规定去噪网络能够读取什么，引导规定使用时怎样组合或修正预测，生成后重排则只决定从已有候选中选谁。GLIDE在像素空间学习文字条件扩散，提供了比较这三个环节的具体入口。本节建立生成接口，后续图像创作在这个接口上选择控制方式与评价协议\scite{vision-glide2022}。

GLIDE先用文字Transformer编码请求，将文字信息接入图像去噪网络。训练根据图文配对构造带噪图像，让去噪预测依赖文字。为了实现无分类器引导，训练时还以一定概率去掉文字条件，使同一个网络能给出条件预测和无条件预测。使用时对同一噪声状态分别计算两者，再按选定强度组合。

一条训练记录仍是图像与描述，但带噪对象改为像素数组。文字网络的末端表示接入噪声步条件，文字序列又作为去噪网络注意力层的键和值，当前图像空间特征提供查询。抽样噪声步后，以所加噪声为目标，用平方误差训练文字编码、去噪网络及连接权重，论文还训练反向方差相关项。文字条件随机置为空序列时，网络不能再读取原描述；所以无条件预测是同一网络在真实训练过的空条件下运行，并非把预测中的某个词向量事后删掉。

令$\bm\epsilon_{\mathrm c}$与$\bm\epsilon_{\varnothing}$分别表示同一带噪图片、同一噪声步上的条件与空条件预测，无分类器引导（classifier-free guidance, CFG）用
\begin{equation}
  \widehat{\bm\epsilon}
     =\bm\epsilon_{\varnothing}
        +g(\bm\epsilon_{\mathrm c}-\bm\epsilon_{\varnothing})
  \label{eq:vision-glide-cfg-interface}
\end{equation}
替代反向均值计算中的噪声预测，$g$是使用时设定的无量纲强度。$g=1$恢复条件预测，$g>1$沿条件与空条件预测的差再向外推；这里的线性组合作用于噪声预测，不是把两幅完成图片混合。通过式\eqref{eq:vision-ldm-reverse-interface}同类的像素反向更新形成下一状态，重复到无噪声末态，再由上采样网络扩大分辨率。训练的步数和优化停止由原模型日程指定，推理的采样步数则由所选采样器指定，二者不是同一个计数。

固定街景文字和一个二维教学噪声状态，假设$\bm\epsilon_{\varnothing}=(0.2,0.4)^\top$、$\bm\epsilon_{\mathrm c}=(0.1,0.8)^\top$。取$g=2$得到$(0,1.2)^\top$。这个结果不是条件预测的概率，而是要送入下一反向更新的数值；第二维已经超过两种原预测的范围，说明引导是外推。增加强度可能改变条件遵循与样本多样性，也可能引入异常，不能由组合形式推出越大越好。

噪声感知CLIP引导需要另一套图文比较模型。它必须读取当前噪声水平的图像，把文字相似度对图像的梯度用于修正生成方向。原GLIDE实验专门训练了适应噪声输入的CLIP，不能把只看清晰图片的原CLIP不加说明地放在同一位置。固定请求与初始噪声时，改变引导强度会改变每一步的预测或更新，最终像素也可能随之改变。

设噪声适应图像编码与文字编码形成匹配分数$s_{\bm\phi}(\bm x^{(k)},c)$，反向图像分布的原均值为$\bm\mu_k$、协方差为$\bm\Sigma_k$，CLIP梯度引导写成
\begin{equation}
  \widehat{\bm\mu}_k
   =\bm\mu_k+g_{\mathrm{clip}}\bm\Sigma_k
           \nabla_{\bm x^{(k)}}s_{\bm\phi}(\bm x^{(k)},c).
  \label{eq:vision-glide-clip-interface}
\end{equation}
梯度指出当前像素怎样局部变化会提高匹配得分，协方差把它转成与反向更新相容的均值修正。CLIP权重在采样中固定，求导对象是当前图像状态，不是临时微调编码器。教学上若$\bm\mu_k=(0.5,0.5)^\top$、$\bm\Sigma_k=0.1\bm I$、梯度为$(1,-2)^\top$、$g_{\mathrm{clip}}=2$，修正均值就是$(0.7,0.1)^\top$。它发生在本步抽样之前，因而能改变未来轨迹；生成后的CLIP排序没有这条像素梯度路径。

生成后重排则先完成$K$个候选，之后才计算各自的匹配分数并选择。重排不修改任何候选，也不能从全部错误的候选中创造正确答案。三种环节可以组合，但应把额外模型调用、梯度和候选数量计入预算。图文相似度仍需与对象数量、属性归属和空间关系的独立核验配合。

\subsubsection{流匹配的速度目标与连续生成}
\label{sec:vision-models-flow-matching}

媒体生成不一定把网络输出解释为“应去掉多少噪声”，也可以直接学习“当前表示应向哪个方向变化”。\term{流匹配}（flow matching）就是通过回归预先构造路径上的速度来训练连续生成模型。它把训练时可计算的局部目标，与使用时求解整条生成轨迹联系起来\scite{vision-flowmatching2023}

用一个明确的教学构造说明。设数据潜向量为$\bm z_1\in\mathbb R^d$，噪声为$\bm z_0\sim\mathcal N(\bm0,\bm I)$，路径参数$s\in[0,1]$。在两者之间定义直线路径
\begin{equation}
  \bm z(s)=(1-s)\bm z_0+s\bm z_1,
  \qquad
  \bm u=\frac{\mathrm d\bm z(s)}{\mathrm ds}
       =\bm z_1-\bm z_0\text{。}
  \label{eq:vision-flow-path}
\end{equation}
$s=0$对应噪声，$s=1$对应数据，$s$是无量纲的生成进度，速度的单位是每单位生成进度的潜坐标变化，不是米每秒或视频帧速度。这里采用无残余噪声的理想化端点以便手算；原流匹配论文也讨论终点保留微小方差等路径，不能把所有实现的日程等同于这一个式子。

训练时取一对数据与条件$(\bm z_1,c)$，抽样噪声$\bm z_0$和时刻$s$，就能直接算出中间状态$\bm z(s)$与目标速度$\bm u$。速度网络$\bm v_{\bm\theta}$的平方误差目标为
\begin{equation}
  L_{\mathrm{FM}}
  =\mathbb E_{\bm z_1,c,\bm z_0,s}
     \left[
       \left\lVert
         \bm v_{\bm\theta}(\bm z(s),s,c)
           -(\bm z_1-\bm z_0)
       \right\rVert_2^2
     \right]\text{。}
  \label{eq:vision-flow-loss}
\end{equation}
计算一次训练目标不必先求解从0到$s$的微分方程，因为中间状态有显式构造。更新的是速度网络参数，样本端点和抽样噪声给出监督。

有限训练集上，把$n$个样本的随机路径误差取平均，便得到一次批次目标。若某个输出速度的残差为$\bm r=\bm v_{\bm\theta}-\bm u$，平方误差导数为$2\bm r$，再经速度网络反传到可训练权重。可从已有生成检查点初始化，也可随机初始化，具体依实现而定；媒体编码器若固定，目标端点在训练中不因这次速度误差改变。优化器按有限训练日程更新，保留速度网络、条件编码和对应媒体解码器。流匹配提供目标构造，不给任意网络和数据配置一个自动收敛到真实分布的保证。

使用时却不知道目标$\bm z_1$，所以不能直接取$\bm z_1-\bm z_0$走向答案。模型需要从一个新噪声出发，根据当前状态预测速度，求解
\begin{equation}
  \frac{\mathrm d\widehat{\bm z}(s)}{\mathrm ds}
    =\bm v_{\bm\theta}(\widehat{\bm z}(s),s,c),
  \qquad
  \widehat{\bm z}(0)\sim\mathcal N(\bm0,\bm I)\text{，}
  \label{eq:vision-flow-ode}
\end{equation}
再把$\widehat{\bm z}(1)$送入媒体解码器。简单的欧拉法在$s_k$处用
$\widehat{\bm z}^{(k+1)}
=\widehat{\bm z}^{(k)}+(s_{k+1}-s_k)
\bm v_{\bm\theta}(\widehat{\bm z}^{(k)},s_k,c)$
更新。有限步长有数值误差，网络也有预测误差，两者不能混为同一项。

为什么从许多不同样本路径学习，仍能得到一个只依赖当前状态的速度？固定输入$(\bm z,s,c)$时，平方损失的最优预测是目标速度的条件均值。令$\bar{\bm u}=\mathbb E[\bm u\mid\bm z,s,c]$，则
\[
  \mathbb E[\lVert\bm v-\bm u\rVert_2^2\mid\bm z,s,c]
  =\lVert\bm v-\bar{\bm u}\rVert_2^2+
    \mathbb E[\lVert\bm u-\bar{\bm u}\rVert_2^2\mid\bm z,s,c]\text{。}
\]
展开平方后，交叉项因$\mathbb E[\bm u-\bar{\bm u}\mid\bm z,s,c]=\bm0$消失。右侧第二项不依赖$\bm v$，因此网络学习的是穿过当前状态的样本路径速度的条件平均。在适当正则条件下，这一边缘速度场可以推动对应的概率分布路径；它不要求网络辨认训练时某条唯一的端点连线。

这个分布结论还需要速度场及密度足够规则，使样本运动可由连续性方程描述。对条件路径密度$p_s(\bm z\mid c)$，质量守恒写为$\partial_s p_s=-\nabla_{\bm z}\cdot(p_s\bar{\bm u})$；局部流入流出的差决定该位置密度如何改变。理论上的条件平均场满足这条对应关系，而有限数据和有限容量训练得到的$\bm v_{\bm\theta}$只是近似。因此，训练时知道直线端点，推理时只知道当前状态；两者不应被画成同一条已知答案路径。

取二维教学样本$\bm z_0=(0,1)^\top$、$\bm z_1=(2,-1)^\top$，目标速度为$(2,-2)^\top$，$s=0.25$时状态为$(0.5,0.5)^\top$。假如使用时在该位置预测速度$(1.2,-2.8)^\top$，步长0.25给出下一状态$(0.8,-0.2)^\top$，而该训练直线在$s=0.5$处为$(1,0)^\top$。这个差异说明一次速度预测误差怎样传到位置，不能用单条训练直线宣称网络推理必然沿直线运行。图\ref{fig:vision-flow-step}给出几何对应。

\begin{figure}[htbp]
  \centering
  % 直线路径为教学构造；粗箭头使用特意偏离目标的速度预测。
\begin{tikzpicture}[x=3cm,y=1.65cm,font=\figurefont\small,
  arr/.style={-{Stealth[length=2mm]},BookInk,line width=0.8pt}]
  \draw[arr] (-0.15,-1.4) -- (2.55,-1.4) node[right] {$z_1$};
  \draw[arr] (-0.15,-1.4) -- (-0.15,1.5) node[above] {$z_2$};
  \foreach \x in {0,1,2}{
    \draw[BookMuted,line width=0.5pt] (\x,-1.4) -- (\x,-1.46);
    \node[below] at (\x,-1.46) {$\x$};
  }
  \foreach \y in {-1,0,1}{
    \draw[BookMuted,line width=0.5pt] (-0.15,\y) -- (-0.19,\y);
    \node[left] at (-0.19,\y) {$\y$};
  }
  \draw[BookBlue,line width=0.9pt] (0,1) -- (2,-1);
  \fill[BookBlue] (0,1) circle[radius=1.1mm];
  \fill[BookTeal] (2,-1) circle[radius=1.1mm];
  \node[anchor=south west] at (0,1.13) {$s=0,\ (0,1)$};
  \node[anchor=west] at (2.02,-0.96) {$s=1$};
  \fill[BookInk] (0.5,0.5) circle[radius=1mm];
  \node[anchor=south west] at (0.52,0.62) {$s=0.25,\ (0.5,0.5)$};
  \draw[BookTeal,line width=1pt,fill=white] (1,0) circle[radius=1.1mm];
  \node[anchor=west] (ideal) at (1.35,0.1) {路径上$(1,0)$};
  \draw[BookTeal,line width=0.6pt] (1.04,0) -- (1.3,0.1);
  \draw[-{Stealth[length=2.2mm]},BookAmber,line width=1.3pt]
    (0.5,0.5) -- (0.8,-0.2);
  \fill[BookAmber] (0.7825,-0.23) rectangle (0.8175,-0.17);
  \node[anchor=east] at (0.76,-0.35) {预测$(0.8,-0.2)$};
  \draw[BookMuted,dashed,line width=0.8pt] (0.82,-0.17) -- (0.97,-0.03);
  \node[anchor=west] at (0,-2.05) {粗箭头：$0.25(1.2,-2.8)^\top$};
\end{tikzpicture}

  \caption{流匹配教学样本的路径与一次欧拉更新。实线连接已知训练端点，粗箭头为预测速度乘以步长后的位移，虚线标示与训练路径上$s=0.5$位置的差别。横纵轴是两个潜变量坐标；路径参数$s$不是视频中的事件时间。}
  \label{fig:vision-flow-step}
\end{figure}

即便每条条件路径都很简单，整体数据分布对应的轨迹也未必是直线，更不能仅据此声称得到了全局最优传输。规范化流（normalizing flow）描述通过可逆变换建立分布的模型，流匹配描述训练速度场的一类目标，两者相关但不是同义词。若$\bm z_1=T(\bm z_0)$为可微可逆变换，密度还须补上体积变化：
\begin{equation}
  \log p_1(\bm z_1)
    =\log p_0(\bm z_0)
       -\log\left|\det
          \frac{\partial T(\bm z_0)}{\partial\bm z_0}\right|.
  \label{eq:vision-flow-density-transform}
\end{equation}
例如$T(z)=2z$把同样概率质量铺在两倍长度上，所以密度减半，对数密度减去$\log2$。这解释了雅可比行列式为何不能从概率计算中省略。

连续情形把许多微小变换相接。若速度场对空间可微且产生唯一可逆轨迹，沿轨迹的密度变化为
\begin{equation}
  \frac{\mathrm d}{\mathrm ds}
        \log p_s(\widehat{\bm z}(s))
    =-\nabla_{\bm z}\cdot
          \bm v_{\bm\theta}(\widehat{\bm z}(s),s,c).
  \label{eq:vision-flow-density-divergence}
\end{equation}
散度是速度对各同名坐标导数之和，度量局部体积的扩张或收缩；扩张会降低单位体积上的密度。该式来自连续性方程和沿轨迹的链式求导，不是流匹配平方误差中的另一项必选监督。只需生成样本的系统可以不计算似然，要求密度的系统还要积分散度并计入求解代价。相关条件与变换原理可参考\BookChapterRef{chap:pgm-generative}{卷六“生成模型”章}。

\subsubsection{LTX-2的音视频双分支与联合解码}
\label{sec:vision-models-ltx2}

若只生成敲门画面，再由另一个模型补声音，后一步只能适应已固定的动作。联合生成可以让画面和声音在生成过程中互相影响。LTX-2技术报告v1采用分别表示、相互连接的音视频分支，学习文字条件下的联合媒体输出。

视频经因果变分自编码器（variational autoencoder, VAE）编码为空间与时间潜网格，音频经频谱转换和独立音频VAE编码为时间潜序列。VAE把媒体压缩到受分布约束的连续表示，再由对应解码器还原；这里的因果约束限制编码使用未来时间的信息，不等于自由问答的推理因果性。两个分支的长度、通道宽度和压缩方式可以不同，因此不需要把音频潜向量误解释成视频像素。文字经编码与连接模块形成条件，分别进入两种媒体的生成分支。

每个分支具有自己的注意力和前馈计算，又通过双向交叉注意力交换信息。视频查询读取音频的键和值，音频查询反过来读取视频的键和值。视频分支内部同时使用时空位置，跨音视频连接则着重使用时间位置，因为音频时间单元没有图像平面上的横纵坐标。时间编码使对应关系可被表达，但正确的敲击同步仍要从配对数据和生成结果检验。

训练把配对音画编码为各自目标潜表示，构造噪声状态，用流匹配监督两个分支的速度\scite{vision-ltx22026}
沿用前节的教学记号，可以把联合状态写成$(\bm z^{\mathrm v}(s),\bm z^{\mathrm a}(s))$，两边的速度预测都依赖当前音画状态及文字条件。各项损失的尺度需要按维度和具体训练设置处理，不能不计音视频长度差异而把元素误差总和直接比较。

报告说明结构与条件路径，公开训练源码进一步给出噪声与速度的符号方向：它以$\sigma=1$表示噪声、$\sigma=0$表示干净数据，构造$(1-\sigma)\bm z_{\mathrm{clean}}+\sigma\bm\epsilon$，目标为$\bm\epsilon-\bm z_{\mathrm{clean}}$。令本章$s=1-\sigma$，便恢复式\eqref{eq:vision-flow-path}的噪声到数据方向，同时速度需变号；不能在减小$\sigma$的采样器中误用本章正向速度。源码中的音视频批次共用$\sigma$，各自抽同形噪声，各自按有效元素平均平方误差，最后把两项相加；作为条件的已知视频位置可由掩码排除监督。这个可核查实现不代表已公开全部基础预训练日程。

文字条件来自冻结的Gemma3多层隐表示，经特征提取投影及每模态的文字连接器处理。报告的初期阶段共同训练特征投影与生成模型，之后固定该投影；文字连接器仍与音视频生成模块共同训练。两个媒体自编码器和声码器提供已定义的解码空间，生成训练保留它们的输入输出约定。学习后的速度、连接模块、两套媒体解码器及文字处理器一起构成使用接口，不能只保留一个视频分支就宣称仍有联合声音输出。

若把一个敲门教学片段的状态简化为视频和音频各一个标量，噪声端为$(0,1)$、数据端为$(2,-1)$，目标正向速度为$(2,-2)$。$s=0.25$时联合状态为$(0.5,0.5)$；若两分支预测为$(1.5,-1.5)$，总平方残差为$0.25+0.25=0.5$。步长0.25更新到$(0.875,0.125)$。真实系统有长时间潜序列与双向读取，不能从这两个标量直接判断敲门同步；例子只显示两种媒体在同一次生成步中各有目标、误差和更新。

使用时，从两个潜空间的噪声开始共同更新。视频潜表示由视频VAE解码为帧，音频潜表示先恢复梅尔频谱，再由声码器恢复波形。若文字要求“人物敲门两次，随后停下”，验收至少要对应两次接触、两次声音以及停顿，不能只听到类似敲门声就认定通过。联合生成不同于为固定视频配声，接口本身也不提供视觉问答；LTX-2中的文字编码器服务生成条件，不等于系统具有自由问答输出头。

\section{统一理解与生成模型}
\label{sec:vision-models-unified}

前面的方法可以形成可比较表示、语言或媒体。如果用户希望“解释这张图，再按要求画一个变体”，可以在一个主干中共享计算，也可以把理解模型和生成器连接起来。判断“统一”的含义，需要看哪些参数、表示和训练目标实际共用，不能只看界面是否只有一个输入框。

\subsection{Show-o2的共享视觉表示与双输出头}
\label{sec:vision-models-showo2}

理解图像需要语义特征，生成图像还需要纹理和局部结构。Show-o2在同一视觉潜空间中提取这两类信息，送入共享语言主干，再用不同输出头分别处理文字和视觉生成。本节依据论文v3的结构及训练说明\scite{vision-showo22025}

图像和视频由3D因果VAE编码。因果性约束编码中未来信息的使用，图像可以作为相应的单帧情形处理。视觉潜表示经过两条路径：语义层从潜网格中提取高层信息，投影层保留便于后续处理的低层特征。语义层在主要两阶段训练之前，利用预训练SigLIP的图像特征进行蒸馏，使其能够从干净和带噪潜表示中提取语义。

两个路径按对应的空间及时间位置对齐，沿特征维拼接，再经归一化与MLP融合为共享视觉表示。这里的共享表示保留多个位置，并非把整幅图压成一个CLIP向量。干净参考图以路径终点$s=1$的状态进入，生成中的视觉块则带有当前路径时刻。相同的视觉接入模块因此需要处理不同噪声水平。

设语义路径输出为$\bm S\in\mathbb R^{n_z\times d_{\mathrm s}}$、低层路径输出为$\bm P\in\mathbb R^{n_z\times d_{\mathrm p}}$，融合输入为$[\bm S,\bm P]\in\mathbb R^{n_z\times(d_{\mathrm s}+d_{\mathrm p})}$。同一行须来自相同潜网格范围，再经均方根归一化和两层映射得到语言宽度的$\bm U$。归一化调整每行的数值尺度，拼接则保留两类特征供映射共同读取；行数不匹配时不能只因同属一张图就直接拼接。3D VAE在这里把时空像素压缩为可连续生成的潜网格，实际使用的权重来自Wan2.1，语言主干来自Qwen2.5-Instruct。语义路径由预训练SigLIP视觉层与新的潜块嵌入组成，先接受干净图片特征的蒸馏，再接入主要训练。

文字嵌入和视觉块按交错顺序进入主干。序列块之间遵守因果顺序，使后续输出可以依赖先前条件；同一视觉块内部允许完整注意力，使一个待生成图像或视频块中的位置共同参与当前速度预测。这里的“因果”描述输出块的条件顺序，不等于每个视频帧只能读取此前帧。

语言头把隐状态转换为下一词元分布，流头则预测视觉潜表示的速度。后者包含带生成时刻调制的Transformer层，其输出经过数值更新和VAE解码才成为媒体。两种训练目标可以概括为
\begin{equation}
  L=\lambda_{\mathrm{text}}L_{\mathrm{NTP}}+L_{\mathrm{FM}}\text{，}
  \label{eq:vision-showo-loss}
\end{equation}
其中$L_{\mathrm{NTP}}$是文字的下一词元预测损失，$L_{\mathrm{FM}}$是视觉速度误差，$\lambda_{\mathrm{text}}$控制文字项权重。两项可以通过共享主干联系起来，但输出维数和误差对象不同。

一条理解样本把干净图像块和问题作为条件，监督后续答案文字；一条生成样本把文字描述作为条件，对目标图像潜表示抽噪声和路径时刻，监督流头速度。交错样本还可以让先前干净图像与文字成为后来图像的条件，不应对作为条件的图片强行添加“预测下一词元”的目标。训练以各样本允许的目标位置分别计算式\eqref{eq:vision-vlm-text-objective}和式\eqref{eq:vision-flow-loss}，再按式\eqref{eq:vision-showo-loss}合并。目标权重改变共享参数受到两种梯度的比例，不能通过一个标量权重把词元误差和潜变量误差解释成同一种单位。

流头每次给出与视觉潜网格相同形状的速度，但语言头仍给出词表维数的logit。比如一个$27\times27$的视觉网格有729个空间单元，生成头对各潜块预测若干连续分量；语言头在一个回答位置对每个词元编号计分。这一差别解释了“统一”为什么不要求相同输出头。对共享权重$\bm\theta_{\mathrm{shared}}$，一次联合更新的梯度是$\lambda_{\mathrm{text}}\nabla L_{\mathrm{NTP}}+\nabla L_{\mathrm{FM}}$；两个梯度可能指向不同方向，因此共享可以传递信息，也可能产生任务冲突，不能仅据共享结构保证两项能力同时改善。

原论文的第一阶段只训练投影、时空融合和流头，保留已有语言知识，并学习视觉生成接口；第二阶段用高质量理解与生成数据适配除VAE外的整体模型。图\ref{fig:vision-showo-heads}把共享范围与两个输出头分开表示。VAE不随这两个阶段更新，统一训练也不意味着每个部件都在每个阶段改变。

具体更新从前述检查点与蒸馏后的语义层开始。第一阶段对图文批次重算文字预测及潜速度误差，只改变新接入的投影、融合与流头；冻结语言层仍需要把误差传到连接输入。第二阶段再放开除VAE外的模型，用高质量理解和生成样本共同更新，并按有限步数训练日程保留最终检查点。原论文1.5B设置以AdamW训练，并在第一阶段和第二阶段改变文字项权重；视频与交错版本另外加入相应样本。已有权重来源、目标位置、开放参数及数据混合共同决定实际得到的模型，不能只说“统一训练”而省略这四项。

\begin{figure*}[htbp]
  \centering
  \begin{tikzpicture}[x=1cm,y=1cm,font=\figurefont\small,
  block/.style={draw=FigureModel,fill=FigureModel!10,line width=0.8pt,
    minimum height=9mm,minimum width=20mm,align=center},
  input/.style={block,draw=FigureInput,fill=FigureInput!10},
  output/.style={block,draw=FigureOutput,fill=FigureOutput!10},
  arr/.style={-{Stealth[length=2mm]},BookInk,line width=0.8pt,
    rounded corners=1.2mm}]
  \node[input] (z) at (1.1,1.8) {干净／带噪\\视觉潜表示};
  \node[block] (sem) at (4.0,2.7) {语义层};
  \node[block] (low) at (4.0,0.9) {低层投影};
  \node[block] (fuse) at (6.9,1.8) {按位置拼接\\与融合};
  \node[block] (trunk) at (9.8,1.8) {共享主干};
  \node[block] (text) at (12.7,2.7) {语言头};
  \node[block] (flow) at (12.7,0.9) {流头};
  \node[output] (words) at (15.4,2.7) {回答词元};
  \node[output] (media) at (15.4,0.9) {更新潜状态\\VAE解码};
  \draw[BookInk,line width=0.8pt] (z.east) -- (2.55,1.8);
  \draw[arr] (2.55,1.8) |- (sem.west);
  \draw[arr] (2.55,1.8) |- (low.west);
  \draw[arr] (sem.east) -- (5.45,2.7) |- (fuse.west);
  \draw[arr] (low.east) -- (5.45,0.9) |- (fuse.west);
  \draw[arr] (fuse) -- (trunk);
  \draw[BookInk,line width=0.8pt] (trunk.east) -- (11.25,1.8);
  \draw[arr] (11.25,1.8) |- (text.west);
  \draw[arr] (11.25,1.8) |- (flow.west);
  \draw[arr] (text) -- (words);
  \draw[arr] (flow) -- (media);
  \node (prompt) at (9.8,0.05) {文字上下文};
  \draw[arr] (prompt) -- (trunk);
  \node[anchor=west] at (0,-0.75) {视觉块之间保持因果顺序，块内位置允许完整注意力。};
\end{tikzpicture}

  \caption{Show-o2的共享视觉接入与双输出头示意，依据论文v3重绘。视觉潜表示的语义路径和低层路径按位置融合，再与文字一起进入共享主干。语言头给出词元，流头给出视觉潜变量的速度；生成分支需迭代更新潜状态并重新经过模型，图中省略迭代回边以突出共享和专用模块的边界。}
  \label{fig:vision-showo-heads}
\end{figure*}
\FloatBarrier

对于“解释参考图，再生成雨天变体”，参考图先经过干净视觉路径，语言头生成解释；随后的视觉块依据此前上下文，通过流头逐步形成媒体。解释正确并不自动保证生成正确，生成中仍可能丢失刚才已识别的对象关系。论文包含文字、图像和视频路线，但视频与交错数据的训练安排随变体而异，报告中的较大7B变体没有沿用同样的视频生成训练数据。选择检查点时必须确认实际训练范围，不能把模型族的全部展示合并成每个版本的能力，更不能额外推断音频生成。

两条输出路径可以继续在街景上复算。理解时，参考图按$s=1$送入，文字头逐词给出“右侧是红色汽车”；生成开始标记出现后，系统分配新的视觉噪声块，保留此前图片和文字作为条件。假设新块的一个二维单元在$s=0.25$为$(0.5,0.5)^\top$，流头预测$(1.2,-2.8)^\top$，步长0.25更新为$(0.8,-0.2)^\top$。更新后的块必须重新经过双路径融合、共享主干与流头，才得到下一时刻速度，不能把第一次隐表示重复使用到终点。完整视觉块到$s=1$后经VAE解码成图片，再作为后续交错输出的上下文；词元生成和潜空间求解由不同停止条件控制。

\subsection{理解模型与专用生成器的接口组合}
\label{sec:vision-models-composition}

另一条路线是由视觉语言模型读取图片和请求，形成明确的内容说明，再调用专用生成器。以“把玩具狗旁的花瓶改成蓝色，其余尽量保持”为教学请求，理解模型可以输出目标短语“花瓶”、保留对象“玩具狗”和修改属性“蓝色”。若生成器支持局部编辑，还需要把目标短语转换成原图坐标中的区域或掩码，并把原图作为独立条件传递。

只传一句“蓝色花瓶旁有玩具狗”时，生成器没有收到原图的精确纹理和位置，因而只能依据语义重建场景。传原图与掩码可以保留更多空间条件，但若理解模型把目标区域选成玩具狗，生成器仍可能准确执行了错误区域上的请求。组合系统应保存中间说明、区域来源、输入图片和最终结果，使错误能定位到读取、传递还是生成环节。

连续特征也可以通过适配模块传给生成器，减少必须把全部细节说成文字的限制。然而，不同编码器的特征坐标和尺度未必兼容，适配模块通常需要自己的训练依据。分别训练后直接连接、固定两端学习适配器、共同更新多个模块，是不同的训练安排，不能仅因推理时连续调用就称为共同学习。

与Show-o2相比，这类组合通常不共享理解与生成主干，可以独立替换生成器或读取模型，却增加了中间格式转换和多次调用。共享主干减少某些重复计算，但也需要在同一组参数中协调任务。两种路线的选择取决于能传递哪些证据、哪一步容易出错，以及最终成果是否满足请求。

\subsection{理解、图像生成与反思编辑的共同训练}
\label{sec:vision-models-unireason}

共享理解与生成接口以后，还可能遇到两个不同缺口。用户要求“从对面看同一条街，两辆车的左右关系随视角改变”，字面条件需要转成明确视角；初始图片即使依据正确文字生成，仍可能把两车画在错误位置。前一个问题发生在生成前的条件形成，后一个问题需要看见已生成图片之后才能判断。2026年预印本UniReason 1.0把生成前的知识推理和生成后的反思编辑放入共同训练，所核查版本为v1，不把其有限一轮修订推广成任意次数自我改进\scite{vision-unireason-2026}。

UniReason基于Bagel已有检查点，而不是将Show-o2直接追加一个反思头。Bagel的Transformer混合结构（mixture of transformers, MoT）使理解与生成专家在共同的交错上下文中计算：理解路径读取视觉编码特征并预测文字，生成路径在VAE潜空间预测连续速度。这里“专家”指面向不同输入与目标的参数路径，不是外部人工审稿者。模型有统一上下文和可交互表示，但并非理解与生成全部权重相同；语言预测和图像速度也仍有不同输出目标。已有ViT提供观测特征，VAE提供生成坐标，本节集中解释它们怎样在交错样本中接受监督，不另把Bagel展开成一套任务算法。

生成前的知识推理样本应包含请求$c$、目标文字推理$T$与目标图像$\bm I$；编辑样本还包含原图$\bm I_{\mathrm{src}}$。例如请求涉及相反视角时，$T$把场景对象、观察方向和相应左右关系写成可用的生成条件，而不是仅重复“画得正确”。原论文的文生图数据从人工选择的知识种子和类别说明开始，由Gemini-2.5 Pro扩充请求并生成文字推理，再由Qwen-Image形成图片；编辑数据以UniREdit-Data-100K的原图、指令、目标结果为基础补充推理。Gemini另按指令符合、视觉质量和推理正确性筛选样本。人工种子、已有记录、合成教师和筛选器承担不同职责，筛选通过仍不等于每一条合成知识都已成为事实真值。

反思编辑样本又有不同的监督边界。初始生成器先给出$T^{(0)}$与草图$\bm I^{(0)}$；视觉核验器比较图片与请求，形成可执行的偏差说明$T^{(1)}$；编辑教师根据原草图和偏差说明形成$\bm I^{(1)}$。论文用Gemini-2.5 Pro核验对象是否存在、属性是否正确以及风格、真实程度和审美方面的偏差，由Qwen-Image-Edit实施编辑，再用比较判断保留被认为改善且符合反馈的结果。一条记录因此保存请求、初始文字、初始图片、反思文字及修订图片，不能只存最终图片后声称模型学过“观察错误”的过程。

训练分两阶段。第一阶段固定基础模型的理解分支，只更新生成分支，用没有推理文字的已有文生图和编辑数据加强图像条件落实；它不监督模型生成解释。第二阶段从这一检查点继续，开放模型的理解与生成参数，用知识推理和反思编辑样本共同微调。单轮知识样本同时监督$T$和$\bm I$；反思样本把$T^{(0)},\bm I^{(0)}$作为上下文，不对这段草图过程计算目标，而只监督$T^{(1)},\bm I^{(1)}$。这避免把待纠正草图也当成应完整模仿的输出。开放参数只说明哪些权重允许变化，真正的梯度仍由受监督的位置及可见条件决定。

把文字目标位置记为$\mathcal T$、被监督图像位置记为$\mathcal G$，第$g$幅目标图片的潜向量为$\bm z_{g,1}$，抽样噪声为$\bm z_{g,0}$，其可见上下文为$C_g$。按本章噪声到数据的方向记号，两个误差可以写成
\begin{equation}
  \begin{aligned}
    L_{\mathrm{text}}
      &=-\sum_{j\in\mathcal T}
           \log p_{\bm\theta}(y_j\mid C_j),\\
    L_{\mathrm{img}}
      &=\sum_{g\in\mathcal G}\mathbb E_{s,\bm z_{g,0}}
        \left\lVert
          \bm v_{\bm\theta}(\bm z_g(s),s,C_g)
              -(\bm z_{g,1}-\bm z_{g,0})
        \right\rVert_2^2,\\
    L_{\mathrm{joint}}
      &=\lambda_{\mathrm{text}}L_{\mathrm{text}}
         +\lambda_{\mathrm{img}}L_{\mathrm{img}}.
  \end{aligned}
  \label{eq:vision-unireason-joint-objective}
\end{equation}
这里$C_j$只含当前目标位置之前允许读取的文字与图片，图像项则使用式\eqref{eq:vision-flow-path}的可构造中间状态。若实现用反向噪声参数，需如前节同时变换时间和速度符号。上式是按本章记号整理的监督关系，不声称论文给每个像素单独定义了知识正确性损失。文字负对数概率训练参考反思措辞，图像速度误差训练目标媒体，数据筛选才提供两者是否符合知识及请求的外部依据。

实际批次把不同类型记录打包，按监督掩码累积两种误差，再反传到各自专家和参与条件计算的模块。原论文第二阶段取$\lambda_{\mathrm{text}}=2$、$\lambda_{\mathrm{img}}=1$，用Adam和学习率预热、余弦日程；第一阶段30K次迭代，第二阶段10K次迭代。这里开放的是统一模型的理解与生成分支，公开训练说明仍将已有VAE设为冻结，不能把“全模型微调”解释成重训媒体坐标。迭代次数规定何时停止该实验，不保证推理已经正确或每幅图都会越改越好。最终保留经过共同微调的统一模型及对应视觉处理、VAE解码接口；合成教师和筛选器是训练数据构造者，不是每次推理必须另调用的隐藏部件。

沿街景例子作一次可检查的教学推演。先固定几何简化：只有观察方向改变，画面横向坐标映射为$x'=w-x$，不涉及遮挡及透视。原图汽车中心横坐标为900、自行车为300，宽度$w=1200$，新视角条件应是汽车在300、自行车在900。模型先生成这段明确条件$T^{(0)}$，再从噪声经流更新得到$\bm I^{(0)}$。假设草图仍把汽车放在右侧，模型重新编码草图，反思文字指出“汽车与自行车的左右位置仍按原视角，应交换”，然后把草图和这条修订条件送入编辑路径。公式中的文字目标是这条反思，图像目标是符合新左右关系的$\bm I^{(1)}$；初始错误草图只供读取，不再被标成应模仿的正确图像。

再看这个样本的一次目标计算。若反思中的某目标词元概率为0.5，其负对数概率为0.6931；假设被监督修订图像的一个二维潜单元残差为$(0.1,-0.2)^\top$，平方误差为0.05。仅计这两个教学单元，权重组合为$2\times0.6931+0.05\approx1.4363$。实际目标还要累积其他词元、潜单元及批次，不可把这个数冒充论文训练日志。它说明文字与图像误差怎样在同一次参数更新中共同起作用，也说明“共同训练”不是仅在推理时先后调用两个模型。

推理没有教师目标图像。模型先根据请求及可选原图形成文字条件，再生成草图；若启用反思路径，就以草图为新的可读视觉内容形成偏差说明，并执行一轮编辑，保留初始与修订结果供比较。修订图像不是将原噪声继续多走一步，而是带有草图条件的另一次图像生成过程，需额外编码、文字预测、流求解和解码。论文实现的一轮细化既可能纠正位置，也可能损坏原来正确的颜色或对象；验收应同时检查原条件保持与新增修改。真正相反视角的场景还会产生遮挡与透视变化，前面的$x'=w-x$只用于教学复算，不能当作一般三维视角转换公式。后续图像创作章可回用本节的训练与推理接口，另行记录候选比较和修订成本。

\section{模型能力与使用约束}
\label{sec:vision-models-constraints}

一个模型支持的输入和输出，是开始选型的依据。实际使用还需确认信息经过编码、压缩和生成后，是否保留了当前任务所需的内容，以及达到这一结果用了多少资源。下面从信息、训练依据和成本三个方面回到同一街景与音画请求。

\subsection{信息保留与输出能力}
\label{sec:vision-models-information}

整图缩小会让车牌字符变成几个像素，局部裁剪会删除周围关系，视频抽帧会漏掉短动作，固定查询数会把许多视觉位置汇总成少量表示。投影和共享主干可以重新组织已经保留的信息，却不能保证恢复这些阶段丢弃的证据。

可以固定问题，分别提供整图、局部高分辨率图，以及带原图位置的二者组合。如果只有局部图能正确读取车牌，就应进一步判断错误来自原始分辨率、输入缩放还是后续压缩，而不是立即归为语言推理失败。组合输入还要确认模型知道局部图来自哪里，否则可能把两个视图当成两辆不同的车。

输出也需逐项验收。匹配分数正确只说明当前候选中的选择，语言流畅只说明形式自然，视频画面逼真只说明感知质量的一部分。数量、位置、时间与来源需要对应的检查。精确区域和答案证据见\ref{chap:vision-multimodal-understanding}章，媒体连贯性和控制见\ref{chap:vision-multimodal-generation}章。

\subsection{模态贡献与训练依据}
\label{sec:vision-models-modality-contribution}

模型接受声音，并不说明一次回答确实使用了声音。可以在保持问题不变时移除音频，或换成时间错开的声音，检查结果与正确答案怎样变化。若问题只问车的颜色，移除声音后答案不变是合理的；若问题问敲击发生了几次，错误声音改变了答案，则需要判断模型是否把它与画面进行了适当核验。

同样，替换街景图却保留文字问题，可以检查模型是否依赖常见语言模式猜测。替换条件必须使正确答案或证据关系发生可解释的变化；只观察回答变了，无法证明模型使用了正确模态。多模态评价还需要覆盖缺失模态、冲突模态和训练中未直接出现的组合。

训练依据决定了可以期待哪一层联系。整图配对监督适合建立全局对应，框或时间区间标注提供更细的位置联系，指令示例说明请求和回答怎样组织。它们可以联合使用，但不能互相冒充。某个新组合的成功应在同类样本与清楚的任务协议中复核，不能由一次演示推断任意模态都已对齐。

\subsection{计算资源与使用成本}
\label{sec:vision-models-cost}

多模态成本可以按输入编码、联合上下文计算、输出生成和媒体解码记录。固定图像大小还不够，裁剪数、视频采样、音频长度、历史轮数和输出时长都可能改变实际输入与输出规模。一个只生成短答案的调用，不能与生成数秒音视频的调用仅比较总延迟。

在标准稠密注意力中，$n$个序列位置的成对分数规模为$n^2$，隐藏宽度为$d$时，相关乘加量含$O(n^2d)$项。若有256个文字单元和1024个视觉单元，总长为1280；视觉单元增加到4096时，总长4352，成对位置数约变为原来的11.56倍。这个教学比例只描述该计算项，不能直接当作实际延迟倍率，局部注意力、缓存与分块实现都会改变资源安排。

冻结模块减少需要保存的参数梯度和优化器状态，却不省去推理中的前向计算。若可训练连接器位于冻结语言模型之前，训练还可能需要经过冻结层计算输入梯度。媒体生成则通常反复调用速度或去噪网络，成本取决于步数和每步条件计算；候选重排还增加整条生成路径的重复次数。

共享主干和模型组合也需按实际调用比较。前者可能复用上下文，后者可以独立缓存编码或选择不同规模的专用模块。部署时还需记录编码、缓存和生成各环节的资源消耗，以及等待和输出长度的限制。

\section{本章小结}
\label{sec:vision-models-summary}

街景照片、请求和声音进入同一系统，需要经过有来源与位置约定的编码，再由匹配空间、视觉读取或共享主干建立联系。CLIP与ImageBind说明可比较表示怎样形成；视觉前缀、层间交叉注意力和多层接入说明语言生成怎样利用视觉；潜空间生成与速度预测则为媒体输出提供了不同的计算路径。

这些路径可以在同一模型中共享，也可以通过专用模块组合。语言回答依靠下一词元预测表达判断，图像和音视频输出还需相应潜表示、生成更新与解码器。输入支持范围、训练过的对应关系和实际输出能力应分别确认。

当一次输入不足以回答时，新的裁剪应从原始来源取回像素，而不是只延长解释。CounterCredit通过固定调用前状态的两个对照，把视觉分支的收益与真实返回像素的收益分开，适用范围仍受裁剪参照与评分方式限制。生成中的反思编辑则读取已经形成的图片：UniReason用选择后的文字与图像目标共同训练一轮修订，但修订费用和条件保持还需单独检查。二者都把后续语言计算与新图像联系起来，前者取得真实观测，后者处理创作结果，不能把这两类来源混为证据。

编码和连接中的证据丢失、时间对应不可靠及不同输出目标之间的取舍，仍需要结合具体任务研究。重点是怎样在固定预算中保留必需的局部证据、建立可靠时间对应，并协调不同输出目标；扩大名义容量只是其中一种手段。选好接口后，仍需回到任务检查答案与证据、内容与条件是否一致，这也是后续理解与生成两章的重点。

\chapter{图像识别与检索}
\label{chap:vision-image-recognition}

一张路口照片中有两辆车。给出“街道”这个场景名，尚未说明车辆在哪里；给出两个框，尚未说明哪辆是红色、是否存在遮挡、距离相机多远；找到另一张相似街景，也未必找到了同一辆车。本章围绕这些逐渐具体的请求，建立从像素到语义、范围、测量和对应关系的计算。

通用的卷积、注意力和表示学习已经提供了特征提取工具。在这里，需要进一步说明特征对应整幅图、一个对象还是一个局部位置，监督约束哪个输出，以及模型预测如何还原成用户需要的结果。街景贯穿对象与空间任务，正常零件与划痕则用于说明没有完整缺陷类别表时怎样判断异常。以下算例均为教学构造，数值用于复核机制。

\section{语义与属性判定}
\label{sec:vision-recognition-semantics}

类别说明“属于什么”，属性说明“具有什么特征”，异常判断说明“是否偏离规定的正常状态”。三种输出可以来自同一表示，但需要不同的标签和判定依据。场景分类把整幅画面的环境概括为街道或厨房，对象分类则需要明确被判断的对象；厨房这一场景名不能展开为一份完整的物品清单。

\subsection{对象类别与场景判定}
\label{sec:vision-recognition-classification}

在单标签任务中，一个样本只对应类别表中的一项；多标签任务允许多项同时成立。例如车辆裁剪图可按主要车型给一个标签，整张街景却可以同时含公交车、轿车和行人。损失函数必须遵守这一区别：互斥类别可用softmax，多种可共存对象通常分别输出是否存在的得分。

\subsubsection{已有视觉表示与分类头：DINOv3的任务接入}
\label{sec:vision-recognition-dinov3}

已有编码器可以把图像转成全局向量及保留位置的特征。DINOv3提供了这类可复用表示；本节关心的是拿到表示后怎样完成有标签的任务，其完整自监督训练另属表示学习。全局向量压缩整图信息，空间特征仍与图像块位置相连，二者适合接入不同的预测头\scite{vision-2508-10104}。

一条分类训练记录包含图像$\bm{X}_i$及类别$y_i$。先按所选检查点约定缩放、裁剪和归一化，保持标签仍对应被保留内容，再从冻结编码器提取全局向量；局部部位任务应换成部位裁剪或已对齐的空间特征。编码器权重来自已有预训练，新增线性头可以从小随机权重和零偏置开始。训练、调参和测试图像应按车辆、拍摄批次或物种采集条件分开，避免同一对象的近重复图同时提供学习和测试依据。

设冻结编码器$f_{\bm{\theta}}$得到$\bm{z}=f_{\bm{\theta}}(\bm{X})\in\mathbb{R}^{d}$，分类头包含$\bm{W}\in\mathbb{R}^{k\times d}$与$\bm{b}\in\mathbb{R}^{k}$。对$k$个互斥类别，计算
\begin{equation}
 p(y=r\mid\bm{X})=
 \frac{\exp(\bm{w}_r^\top\bm{z}+b_r)}
 {\sum_{j=1}^{k}\exp(\bm{w}_j^\top\bm{z}+b_j)}\text{，}
 \label{eq:vision-linear-head}
\end{equation}
这里$\bm{w}_r^\top$是$\bm{W}$的第$r$行。负对数给真实类别低概率时较大代价；在样本独立的条件模型中，连乘似然取负对数便成为逐样本和，经验目标为$L=-n^{-1}\sum_i\log p(y_i\mid\bm X_i)$，平均使目标不随批次大小直接放大。冻结训练只更新$\bm{W},\bm{b}$；全量适配还更新$\bm{\theta}$，所需梯度存储和训练数据条件随之变化。

对一个样本，记类别得分为$a_r=\bm{w}_r^\top\bm{z}+b_r$、概率为$p_r$，则$\partial(-\log p_y)/\partial a_r=p_r-\mathbf{1}\{r=y\}$。链式求导给出头部梯度$(p_r-\mathbf{1}\{r=y\})\bm{z}$及偏置梯度。每批求平均后，最简单的梯度更新为$\bm{W}^{(t+1)}=\bm{W}^{(t)}-\eta_t\nabla_{\bm W}L$，偏置同样更新，$\eta_t>0$是学习率；带动量的优化器另外保存梯度统计。按事先规定的更新预算训练，依据不参与梯度的验证集选择头部参数，不声称神经网络已经找到全局最优。后文的任务网络沿用“构造监督、前向计算误差、反向更新、验证选择”的关系，但可训练范围、误差项和输出恢复分别说明。

假设三类车型的得分为$(2,1,0)$。经过softmax计算，概率依次约为0.665、0.245和0.090。它们是在给定类别集合内归一化的结果，不能证明图像中一定存在可辨车辆。若编码器读取整幅街景，得分还可能依赖道路背景；若任务要判定左车的车型，应先取得左车区域，或使用明确的区域监督。

若这条样本真实类别为第二类，得分梯度约为$(0.665,-0.755,0.090)$。下降更新提高第二类相对于其他类的得分；这一方向仅说明当前监督如何作用，并不保证下一批或新领域的分类正确。使用时保留编码器、最终分类头、类别顺序和相同预处理，计算得分后取最大项；需要未知状态时再使用验证确定的拒绝规则。场景分类保留整图，逐车分类则保留对象编号及裁剪来源，两者不能混用同一输出记录。

外观接近的物种分类揭示了另一种困难：全局背景足以区分拍摄地点，却可能不足以区分鸟喙或叶缘。iNaturalist包含细粒度类别和不均衡样本，评价时应同时检查小部位分辨率和类别覆盖\scite{vision-1707-06642}
例如常见类别有90张且全对，少见类别有10张且全错，样本平均准确率为90\%，两个类别的平均准确率却为50\%。这两项分别强调观测频率与类别公平覆盖，不能互相替代。

\subsubsection{文字指定类别与零样本识别：CLIP接口}
\label{sec:vision-recognition-clip}

如果使用时才给出“轿车、货车、公交车”，固定输出维度的分类头不能直接容纳新类别。图文表示把类别描述也编码为向量，使图像可以与使用时给定的候选比较。CLIP的双编码器训练已在\ref{sec:vision-models-contrastive}节展开；这里使用已经学到图文对应的两侧编码器及投影，不为当前类别表再训练一个分类头。

令$\bar{\bm{z}}$与$\bar{\bm{t}}_r$分别为单位长度的图像向量和第$r$条描述向量，按$s_r=\bar{\bm{z}}^\top\bar{\bm{t}}_r$排序。使用温度$\tau>0$可计算候选内归一化得分
\begin{equation}
 q_r=\frac{\exp(s_r/\tau)}
 {\sum_{j\in C}\exp(s_j/\tau)}\text{，}
 \label{eq:vision-clip-classification}
\end{equation}
其中$C$是这次输入的非空候选索引集合，图文归一化要求原向量非零。分母随候选变化，$q_r$不是不依赖候选的存在概率。可以对“这是一张公交车照片”等多个模板的单位文字向量求平均后再归一化，但模板集合须在评价前固定，并排除平均后为零的退化情形。

设四项相似度为$0.31,0.28,0.19,0.12$，最后一项是不在图中的“帆船”。最大值只使第一项在这四项中领先。把前三项全部换成不相关词，模型仍会返回一个最大值。开放使用要另设含未知对象的验证协议、拒绝阈值或目标存在判断；单凭相似度最大不能完成拒识。整图相似也无法保证“红色”绑定到货车而非旁边的轿车。

\subsection{对象属性判定}
\label{sec:vision-recognition-attributes}

要记录“红色货车”和“蓝色轿车”，仅输出整图属性集合$\{\text{红色},\text{蓝色}\}$会丢掉归属。可以先定位对象，再对各区域提取表示，形成$(\text{对象编号},\text{属性},\text{取值})$记录。位置与属性共享编号，使后来计数、搜索或人工核验仍指向同一个对象。

属性之间是否互斥应由标签定义决定。“红色车身”和“车灯开启”可以同时成立，“完整车身的主色”则可以规定为一个单标签字段。对于可共存的$k$个二值属性，令$p_{i,r}$为预测，$a_{i,r}\in\{0,1\}$为标注，$m_{i,r}$表示该属性是否有可靠标注，则可用
\begin{equation}
 L_{\mathrm{attr}}=
 -\frac{\sum_{i,r}m_{i,r}
 [a_{i,r}\log p_{i,r}+(1-a_{i,r})\log(1-p_{i,r})]}
 {\sum_{i,r}m_{i,r}}\text{。}
 \label{eq:vision-attribute-loss}
\end{equation}
分母要求至少有一个有效标注；全无标注的批次不计算此项。把“未标注”直接置为“不存在”，会训练出系统性的假负例。

具体接入可令区域向量$\bm{z}_i$经过每属性一个线性得分$h_{i,r}$，再用sigmoid得到$p_{i,r}=1/(1+e^{-h_{i,r}})$。单个有效标注的未平均二元目标对得分的梯度是$p_{i,r}-a_{i,r}$，其中$a_{i,r}$仍指二值标注。冻结区域编码器时只训练属性头，也可按数据条件适配编码器。各属性阈值由独立验证记录选择，输出保留每车的属性分数和可见性；若主色被定义为互斥字段，则另用softmax头，不把二元输出强行当成唯一主色。

裁剪区域过小会切掉车灯，过大又可能纳入旁车颜色。应检查属性所需部位是否可见，并在不可判定时保留未知状态。完整对象编号并不能自行修正错误的视觉证据，但它使属性错绑能够被单独定位。

\subsection{正常状态与异常判定}
\label{sec:vision-recognition-anomaly}

工业零件的异常类型可能不断出现，训练集却主要含正常样本。此时可学习或保存正常外观，再度量新观测与正常外观的偏离。重建误差、分布密度和局部近邻距离是不同的偏离依据，得分都需要结合独立验证集解释。

\subsubsection{正常局部特征记忆与近邻评分：PatchCore}
\label{sec:vision-recognition-patchcore}

PatchCore从已有编码器的中间层提取并局部聚合特征，将正常图像中各位置的向量存入记忆。相邻尺度的特征对齐后保留局部上下文，使用时不需要先列出所有划痕和破损类别\scite{vision-2106-08265}

全部正常片段可能很大，因而从特征集合$\bm{M}$中保留代表子集$\bm{M}_{\mathrm{c}}$。其覆盖目标可写为
\begin{equation}
 \min_{\substack{\bm{M}_{\mathrm{c}}\subseteq\bm{M}\\
 |\bm{M}_{\mathrm{c}}|=m}}
 \max_{\bm{z}\in\bm{M}}\min_{\bm{u}\in\bm{M}_{\mathrm{c}}}
 \lVert\bm{z}-\bm{u}\rVert_2\text{。}
 \label{eq:vision-patchcore-coreset}
\end{equation}
原方法使用近似贪心选择，而非精确穷举全部子集。已选一个起点后，反复加入离当前子集最远的特征，可以保留正常模式中的不同部分；这个过程在建立记忆，不是在用缺陷标签重新训练编码器。

选择时为每个正常特征保存到当前子集的最近距离，新加入$\bm{u}$后更新为原距离与$\lVert\bm{z}-\bm{u}\rVert_2$的较小者，达到预算$m$就停止。原工作用随机投影辅助近似选择，保留的片段仍用于原特征空间评分。最内层最小值寻找一个正常参照，外层最大值检查最难覆盖的正常模式；若只最小化平均距离，少量正常模式可能在记忆中完全缺席。完成后保存编码器、所用中间层及聚合设置、代表记忆和验证阈值。建立记忆的数据必须符合“正常”定义，异常混入记忆会把某些缺陷变成近邻参照。

查询图像位置$p$的基本局部分数为$d_p=\min_{\bm{u}\in\bm{M}_{\mathrm{c}}}\lVert\bm{z}_p-\bm{u}\rVert_2$。将这些分数按空间排列、插值和平滑可形成异常热图。最大局部距离$s_*=\max_p d_p$是整图评分的起点；原方法还根据最异常片段$\bm{z}_*$所对应正常近邻$\bm{u}_*$周围的支持情况重新加权。令$\bm{N}_b(\bm{u}_*)$包含$\bm{u}_*$及其在记忆中的近邻，则
\begin{equation}
 s=\left(1-
 \frac{\exp\lVert\bm{z}_*-\bm{u}_*\rVert_2}
 {\sum_{\bm{u}\in\bm{N}_b(\bm{u}_*)}
  \exp\lVert\bm{z}_*-\bm{u}\rVert_2}\right)s_*\text{。}
 \label{eq:vision-patchcore-score}
\end{equation}
邻域大小$b$是需固定的设置，实际应大于1；否则分子分母相同，分数退化为零。若最近正常片段较孤立，其他邻居的距离更大，括号中的系数更接近1，最大距离受到的压低更少。这个重加权改变整图分数，不改变局部距离图的来源。

假设一维正常记忆为$\{0,1,4\}$，查询局部特征为$\{0.2,0.9,6\}$，对应距离是$0.2,0.1,2$，最大值由第三个片段给出。若$b=2$，取最近正常片段4及其邻居1，整图分数为$2[1-\exp(2)/(\exp(2)+\exp(5))]\approx1.905$。区域线索保留了小范围的偏离，但若新设备使所有特征整体平移，正常零件也可能远离记忆。高分只能表示对当前正常参照的偏离，不能直接命名缺陷或判定业务不合格。

\section{对象定位与区域划分}
\label{sec:vision-recognition-regions}

定位要求把判断接回图上的位置。边界框给出大致范围，掩码给出像素归属，透明或混合边界还需要连续贡献比例。希望处理哪个目标，可以由固定类别、文字或点框提示指定；目标选择与范围形成需要分别解释。

\subsection{对象范围的定位}
\label{sec:vision-recognition-detection}

\term{边界框}（bounding box）用一个矩形包围目标。以下采用连续边界$\bm{b}=(x_1,y_1,x_2,y_2)$，宽高分别为$x_2-x_1$和$y_2-y_1$，不加离散像素计数中的1。检测结果是多个$(\bm{b},c,s)$组成的集合，其中$c$是类别，$s$是得分。标签应说明框包围可见部分还是完整对象；旋转目标也可改用中心、宽高和角度表示。

比较两个框的重叠程度时，\term{交并比}（intersection over union, IoU）将交集面积除以并集面积：
\begin{equation}
 \operatorname{IoU}(\bm{b},\bm{b}')
 =\frac{|\bm{b}\cap\bm{b}'|}
 {|\bm{b}|+|\bm{b}'|-|\bm{b}\cap\bm{b}'|}\text{。}
 \label{eq:vision-iou}
\end{equation}
这里矩形应有正面积，无交集时分子为零。取两个面积均为800的框，边界分别为$(10,10,50,30)$和$(20,10,60,30)$。它们的交集面积为600，IoU为$600/1000=0.6$。若它们是对同一辆车的两个预测，应去重；若属于两辆互相遮挡的车，直接去掉一个则会漏检。

\subsubsection{整图网格预测：YOLO}
\label{sec:vision-recognition-yolo}

原始YOLO把图像划为$S\times S$网格，每个单元预测$B$个框和一组共享的条件类别概率。目标中心落在哪个单元，就由该单元承担监督。框中心相对于单元定位，宽高相对于整图归一化；每框置信度学习目标存在与框质量的联合信号\scite{vision-1506-02640}

若单元列行为$(j,i)$，预测中心偏移为$(u,v)$，输入宽高为$(w,h)$，则中心还原为$((j+u)w/S,(i+v)h/S)$。训练在负责单元的多个预测框中选择当前IoU最高者承担该对象的坐标监督。原始平方误差目标分别约束中心、宽高的平方根、置信度和类别；没有负责对象的候选也接受置信度监督，但权重较低，以免大量背景压倒对象信号。

为看清同一目标怎样进入误差，记单元为$a$、框槽为$b$，$I^{\mathrm{resp}}_{a,b}$标出当前负责框，$I^{\mathrm{other}}_{a,b}$标出其他框槽，$I^{\mathrm{obj}}_a$标出含目标中心的单元，三者均为零一指示。预测中心、宽高和置信度分别为$(\hat x,\hat y)$、$(\hat w,\hat h)$和$\hat C$，真实值去掉帽子；宽高为正并相对整图归一化。$\lambda_{\mathrm c},\lambda_{\mathrm n}>0$分别是坐标和非负责框置信误差的权重。先把位置误差分为中心与尺寸两部分：
\begin{equation}
 \begin{aligned}
 e_{a,b}^{\mathrm{pos}}&=(\hat x_{a,b}-x_a)^2
                       +(\hat y_{a,b}-y_a)^2,\\
 e_{a,b}^{\mathrm{size}}&=(\sqrt{\hat w_{a,b}}-\sqrt{w_a})^2
                        +(\sqrt{\hat h_{a,b}}-\sqrt{h_a})^2,\\
 L_{\mathrm{coord}}&=\lambda_{\mathrm c}\sum_{a,b}
 I^{\mathrm{resp}}_{a,b}(e_{a,b}^{\mathrm{pos}}+e_{a,b}^{\mathrm{size}})\text{。}
 \end{aligned}
 \label{eq:vision-yolo-training-process}
\end{equation}
类别和置信误差则按同一职责选择监督：
\[
 \begin{aligned}
 L_{\mathrm{conf}}&=\sum_{a,b}I^{\mathrm{resp}}_{a,b}
 (\hat C_{a,b}-\operatorname{IoU}_{a,b})^2\\
 &\quad+\lambda_{\mathrm n}\sum_{a,b}
 I^{\mathrm{other}}_{a,b}\hat C_{a,b}^{\,2},\\
 L_{\mathrm{class}}&=\sum_a I^{\mathrm{obj}}_a\sum_r
 (\hat p_a(r)-\mathbf{1}\{r=c_a\})^2\text{。}
 \end{aligned}
\]
平方根使相同绝对宽高误差在较小框上占更大相对权重；它不改变中心仍由单元承担的分配。对训练图求平均，将三项相加，依当前预测重新选择负责框，再反向更新网络。原工作先做分类预训练再训练检测网络，新增检测层接受框监督；推理只保留最后的网络，不保留逐批负责关系。

教学中$S=7$、输入$448\times448$，单元列行$(2,3)$的中心偏移$(0.5,0.25)$还原为$(160,208)$。若归一化宽高为$(0.2,0.1)$，像素宽高是$(89.6,44.8)$，再由中心加减半宽高得到四边。置信度0.8与该单元货车条件得分0.75相乘，货车框分数为0.6。输入曾缩放或加边时，最后还须去掉加边并除缩放比例，才是原图候选。

每框置信度写为$\mathbb{P}(\mathrm{obj})\operatorname{IoU}$，乘上共享的$\mathbb{P}(c\mid\mathrm{obj})$得到类别相关得分。它将语义与定位质量合在一起，不能只解释为类别概率。使用时解码全部框，按阈值过滤，再用\term{非极大值抑制}（non-maximum suppression, NMS）去掉同类中与当前高分框过度重叠的低分框。

NMS按得分从高到低取一个框，保留它并删除重叠超过阈值的候选，再对余下框重复，直到没有候选。对于前述IoU为0.6的两个框，若得分为0.9与0.7，阈值为0.5，第二个会被删除。这一规则处理重复预测，却可能伤害拥挤场景。一个单元共享类别预测也限制了原始版本处理同单元多对象的能力；后来的YOLO版本改变了预测、分配与去重设计，不能沿用此处全部细节。

\subsubsection{多尺度特征与密集训练权重：FPN、RetinaNet}
\label{sec:vision-recognition-fpn}

远处车辆在图上只占几个像素，深层特征的空间步幅可能使它消失。浅层保留更多位置细节，深层又具有较大的上下文。\term{特征金字塔网络}（feature pyramid network, FPN）把深层特征上采样，与相应浅层经通道变换后的特征相加，再用卷积整理，形成不同空间尺度的预测特征\scite{vision-1612-03144}

第$\ell$层的步幅$s_\ell$连接特征坐标与原图坐标。训练需要按使用该金字塔的检测器规定，将目标或锚框分配到尺度层；FPN本身只生成特征，并不规定一种通用的最终检测集合。较密特征有助于保留小目标，但若目标在输入缩小时已经丢失，特征融合不能补回观测。

记主干第$\ell$层为$\bm{C}_\ell$，自顶向下的中间特征为$\bm{M}_\ell$，则相同空间网格上可以写$\bm{M}_\ell=\operatorname{Conv}_{1\times1}(\bm{C}_\ell)+\operatorname{Up}(\bm{M}_{\ell+1})$，再用$3\times3$卷积得到$\bm{P}_\ell$。横向卷积先统一通道数，上采样先统一位置数，否则两项无法相加。对步幅$s_\ell$的输出位置$(i,j)$，检测器按其网格约定放置原图参考中心；不同层即使特征形状不同，最终框也需回到同一原图坐标再比较。

RetinaNet在金字塔各位置放置预定大小和长宽比的\term{锚框}（anchor box），即坐标回归的参考框。网络预测类别以及相对于参考框的中心偏移、宽高对数修正。训练按锚框与真实框的IoU分配前景、背景和忽略项，推理则解码修正、筛选并执行NMS。

设锚框中心和尺寸为$(x_a,y_a,w_a,h_a)$，目标为$(x_*,y_*,w_*,h_*)$。训练回归标签为$((x_*-x_a)/w_a,(y_*-y_a)/h_a,\log(w_*/w_a),\log(h_*/h_a))$；给定预测修正$\bm{\delta}$，逆变换为
\begin{equation}
 \begin{aligned}
 \hat x&=x_a+w_a\delta_x,\qquad \hat y=y_a+h_a\delta_y,\\
 \hat w&=w_a e^{\delta_w},\qquad \hat h=h_a e^{\delta_h}\text{。}
 \end{aligned}
 \label{eq:vision-retinanet-decode}
\end{equation}
中心差除参考尺寸，宽高比取对数，使标签无像素单位且能表达乘性缩放。教学锚框$(50,30,40,20)$配目标$(60,30,80,20)$，标签为$(0.25,0,\log2,0)$，逆算恢复原目标。负锚框没有对应框，不计算坐标回归项；边界分配的忽略项不参与分类。

大量背景锚框在训练后很容易正确分类，但数量仍可能主导总损失。令$p_t$为模型分给真实二值标签的概率，\term{焦点损失}（focal loss）为
\begin{equation}
 L_{\mathrm{focal}}=-\alpha_t(1-p_t)^\gamma\log p_t\text{，}
 \qquad\gamma\geq0\text{。}
 \label{eq:vision-focal-loss}
\end{equation}
其中$\alpha_t$调整正负样本权重。$\gamma=0$时退回加权交叉熵；$\gamma=2$时，$p_t=0.99$的容易样本乘上$10^{-4}$，而$p_t=0.5$的困难样本乘上$0.25$。这个差别让训练更关注尚未学好的候选\scite{vision-1708-02002}
改变容易背景的权重没有增加少见类别的观测，因而前景背景失衡与类别长尾仍须分别检查。

各锚框按类输出独立sigmoid得分，分类采用焦点项，正锚框回归采用平滑$L_1$项，训练按正例数量归一化；没有正例的批次应按实现保留非零归一化下界。主干可从分类预训练初始化，金字塔与任务头由检测监督学习。对同一批的前景和背景计算误差后更新共享参数，使用时从所有层解码参考框、合并同类候选再NMS，最后逆变换到原图。焦点因子只是改变每个候选的训练贡献，没有从后处理集合中直接删除背景。

\subsubsection{逐位置边界距离预测：FCOS}
\label{sec:vision-recognition-fcos}

预定义参考框带来了大小与比例的设计。FCOS改为让特征位置直接预测到对象左、上、右、下边界的距离$(l,t,r,b)$。若位置对应原图$(x,y)$，预测框就是$(x-l,y-t,x+r,y+b)$，不再以锚框为坐标起点\scite{vision-1904-01355}

监督仍然需要分配。位于真实框内且满足尺度范围的位置可以成为正例；原方法用不同层的回归范围分摊对象，并在多个框竞争一个位置时选择面积较小的框。负例只学习分类，正例还学习距离及中心度。这里描述的是原始FCOS，后续采用中心区域采样等改动时应另记配置。

每层特征位置先按步幅对应原图采样中心，再由目标四边减该中心形成$(l,t,r,b)$标签。训练分类采用焦点目标，正例的恢复框与真实框用IoU误差比较，中心度使用二元交叉熵；三项按权重组成批次目标，梯度同时到达特征主干及三个预测分支。距离分支按实现约束为正，推理由距离和采样中心直接得到四边，类别分数与中心度组合后跨层筛选。对于无目标位置，只学“不是这些类别”，不能给它虚造四边距离。

\term{中心度}（centerness）衡量位置是否接近框中心，训练目标为
\begin{equation}
 c_*=\sqrt{\frac{\min(l,r)}{\max(l,r)}
                 \frac{\min(t,b)}{\max(t,b)}}\text{。}
 \label{eq:vision-fcos-centerness}
\end{equation}
式中距离来自分配的真实框，框内部距离为正。中心点四侧相对均衡，$c_*=1$；靠近边界时其中一个比例趋于零。推理把预测中心度与类别得分组合，降低边缘候选的排序，再用NMS形成集合。

对于框$(20,10,80,50)$，中心$(50,30)$的目标距离为$(30,20,30,20)$。位置$(25,30)$的距离变成$(5,20,55,20)$，中心度为$\sqrt{1/11}\approx0.302$。即使这两个位置都能精确回归同一框，中心度也不同，说明它不是预测框IoU的直接测量。图\ref{fig:vision-box-points}把位置、距离与框的关系画在同一坐标系中。

\begin{figure}[htbp]
 \centering
 % 连续像素坐标下的FCOS教学例；所有长度与正文算例一致。
\begin{tikzpicture}[x=1mm,y=-1mm,font=\figurefont\small,
  dist/.style={-{Stealth[length=2mm]},draw=BookBlue,line width=0.9pt}]
 \draw[draw=BookMuted,line width=0.7pt,->] (10,0)--(90,0) node[right] {$x$};
 \draw[draw=BookMuted,line width=0.7pt,->] (10,0)--(10,58) node[below] {$y$};
 \draw[draw=BookTeal,fill=BookTeal!7,line width=1pt] (20,10) rectangle (80,50);
 \draw[draw=BookMuted,dashed,line width=0.7pt] (50,10)--(50,50);
 \draw[draw=BookMuted,dashed,line width=0.7pt] (20,30)--(80,30);
 \fill[BookInk] (50,30) circle (1mm);
 \fill[BookBlue] (25,30) circle (1mm);
 \node[anchor=north west] at (51,33) {中心：$c_*=1$};
 \node[anchor=south west] at (26,27) {边缘：$c_*\approx0.302$};
 \draw[dist] (25,18)--(20,18);
 \draw[dist] (25,18)--(80,18);
 \node[anchor=south] at (22.5,17) {$5$};
 \node[anchor=south] at (57,17) {$55$};
 \draw[draw=BookBlue,line width=0.7pt] (25,18)--(25,29);
 \node[anchor=north] at (20,51) {$(20,50)$};
 \node[anchor=south] at (80,9) {$(80,10)$};
\end{tikzpicture}

 \caption{同一真实框内的中心位置与边缘位置。实线箭头标出边缘位置到左右边界的距离，虚线标出中心轴；两点都可以回归同一矩形，但中心度目标分别为1与约0.302。坐标采用向右、向下为正的连续像素边界。}
 \label{fig:vision-box-points}
\end{figure}

\subsubsection{无序对象集合与唯一匹配：DETR}
\label{sec:vision-recognition-detr}

密集候选需要在输出端去重。DETR把目标直接视为无序集合：固定$n$个对象查询读取图像特征，每个查询预测一个框和类别，类别还包含“无对象”。查询是学习得到的向量，并非预先指定给某辆车的身份\scite{vision-2005-12872}

训练标注有$m\leq n$个对象，记第$i$个真实类别与框为$(c_i,\bm{b}_i)$，第$j$个预测为$(\hat{\bm{p}}_j,\hat{\bm{b}}_j)$。寻找从真实对象到不同预测位置的单射$\sigma$：
\begin{equation}
 \hat{\sigma}=\argmin_\sigma
 \sum_{i=1}^{m}\left[
 -\hat p_{\sigma(i)}(c_i)
 +\lambda_1\lVert\bm{b}_i-\hat{\bm{b}}_{\sigma(i)}\rVert_1
 +\lambda_{\mathrm{g}}L_{\mathrm{GIoU}}
    (\bm{b}_i,\hat{\bm{b}}_{\sigma(i)})\right]\text{。}
 \label{eq:vision-detr-matching}
\end{equation}
框在共同归一化坐标中比较。广义IoU还惩罚两个框最小包围矩形中未被并集覆盖的区域，使无交集时也有几何信号。匹配代价用于选对应；选定后，分类损失采用负对数概率，定位项只计算匹配对象，未匹配查询学习“无对象”，并降低该类的权重。

具体计算中，主干从缩放图像产生二维特征，展平后加位置编码，编码器在图像位置间交换上下文。每个可学习查询经解码器读取位置特征，再由共享类别头和框头产生结果。框头输出归一化中心、宽高；预测序列没有必须对应目标编号的顺序。匹配可以用匈牙利算法求当前代价矩阵的最小单射，选定后视对应为本轮固定监督，不对离散匹配本身求梯度。

若$E$为两框最小包围矩形，广义交并比为$\operatorname{GIoU}=\operatorname{IoU}-|E\setminus(\bm b\cup\bm b')|/|E|$，定位损失使用$1-\operatorname{GIoU}$。未相交框的IoU均为零，这个额外空白比例仍能区分隔得近和隔得远。训练目标是匹配查询的负对数分类及两项框误差，加未匹配查询的无对象分类；多图求平均后更新主干、编码解码器、查询和任务头。原始实现可用已有分类主干初始化，其余部分由检测任务学习。

假设两个真实对象对三个查询的总代价矩阵为
\[
 \bm{C}=\begin{pmatrix}1&2&9\\1.2&5&8\end{pmatrix}\text{。}
\]
逐行独立取最小都会选第一个查询，违反唯一对应。集合匹配选择第一对象到第二查询、第二对象到第一查询，总代价为$3.2$，小于反过来分配的$6$。三辆车的情形同样要求三个不同查询；改变标注排列只改变匹配的表示，不应改变最佳集合代价。

使用时没有真实框参与匹配，直接过滤“无对象”和低分输出。原始DETR不执行常规NMS，重复预测主要由唯一匹配训练约束。查询数限制一次可输出的对象数，原始全局注意力、训练周期与小目标效果也各有条件；这些限制不能概括后来所有Transformer检测器。

\Needspace{65mm}
\subsubsection{文字条件下的对象预测：Grounding DINO}
\label{sec:vision-recognition-grounding-dino}

当目标变成“红色货车”时，需要把类别与属性文字接入检测计算。Grounding DINO分别编码图像和文字，通过跨模态特征增强交换信息，再用语言相关性选择用于初始化的对象查询；解码器继续读取两侧特征并更新框与文字对应得分\scite{vision-2303-05499}

文字经分词形成若干词元，候选图像特征与词元比较，训练用对象范围以及对象对应的词语监督匹配。推理时不仅要还原框，还要说明得分对应“货车”、包含属性的短语，还是更长表达中的一部分。词元阈值、框阈值和短语聚合规则都会影响最终输出。

令增强后的候选图像特征为$\bm{v}_i$、文字词元为$\bm{t}_r$，语言引导查询选择依据候选与文字的相似度，例如从$\max_r\bm{v}_i^\top\bm{t}_r$较高的位置选初始化查询。此时选中的仍是候选位置，不能当成已经确认的货车。跨模态解码继续更新查询和归一化框，并形成各查询对各词元的匹配logit。训练记录需包含图像、文本、目标框及文本中的目标词元关联，关联可以跨多个词元；分类误差约束这些关联，$L_1$和广义IoU约束匹配框，集合分配避免同一监督对象任意占多个查询。

可使用已有视觉和语言编码器初始化，经过有图文范围对应的预训练或任务适配获得连接权重；使用已发布检查点时保留其分词、图像尺度及阈值协议，不为每次文字查询更新权重。以“红色货车”查两辆车时，先解码每个候选框，再保留满足框阈值的项，从相应词元得分恢复短语，最后将框中心宽高乘原图尺寸并转换为四边。若模型只在“货车”词元高分而“红色”得分弱，最终记录不能悄悄把整个属性短语写成已核实事实。

“红色货车”要求属性绑定，“公交车旁的货车”又要求利用参照物消歧。能接收这两种字符串只说明接口容纳了条件，不能证明关系已经正确解析。可分别遮去参照物、更换颜色词或保留多个同类目标，检查预测是否按所需证据改变。完整的表达定位及证据协议在\ref{chap:vision-multimodal-understanding}章讨论。

\subsubsection{开放词汇定位之外的未知对象发现}
\label{sec:vision-recognition-unknown-discovery}

仓库画面中出现一个尚未加入类别表的搬运装置。固定分类器可能把它认作货车，也可能把它当背景；文字定位允许用户写出装置名称，但用户此时未必知道应当查询什么。\emph{未知对象发现}（unknown object discovery）要求在没有名字时先提出“这里可能有一个对象”，保留区域供后续核验。它与开放词汇定位的差别在于名称是否已作为输入，也与拒识不同：拒识已有裁剪图只回答是否属于已知类，发现还要确定未知对象在哪里。

用集合$C_{\mathrm{k}}$表示训练已标注类别，输出可以是已知记录$(\bm{b},c,s)$，也可以是未命名记录$(\bm{b},\mathrm{unknown},s_{\mathrm{o}})$。对象性（objectness）$s_{\mathrm{o}}$衡量一个候选是否应被当作独立对象，而不是它属于哪个已知类。未匹配标注的候选中同时含真实背景和未标注对象，把它们全部用零监督，会使发现能力向已知类别偏移。

OWOBJ把对象性建模接入已有检测器的训练，而不另造一套推理接口。以其基于可变形DETR的实现为例，图像特征经编码、查询解码得到$\bm{q}_i\in\mathbb{R}^d$；已知类别头、框头和对象性分支读取这些查询。训练先按已知框和类别作唯一匹配，匹配查询确定属于对象，其他查询暂不能判成背景。原论文由查询特征的移动平均估计均值$\bm{\mu}$与正定对角协方差$\bm{\Sigma}$，用
\begin{equation}
 \begin{aligned}
 d_{\mathrm{M}}(\bm{q})=
 (\bm{q}-\bm{\mu})^\top\bm{\Sigma}^{-1}(\bm{q}-\bm{\mu}),\\
 s_{\mathrm{o}}(\bm{q})=\exp[-d_{\mathrm{M}}(\bm{q})]\text{。}
 \end{aligned}
 \label{eq:vision-owobj-objectness}
\end{equation}
生成对象性软标签。马氏距离（Mahalanobis distance）按各特征方向的正常变动尺度比较偏离，避免原始数值尺度大的通道独占距离。$s_{\mathrm{o}}$在距离为零时为1，距离增大时下降；它是该训练设计中的对象性信号，不是已经校准的真实对象概率\scite{vision-owobj-2025}。

匹配查询继续接受已知类别和框监督，并通过减小$d_{\mathrm M}$提高对象特征的密度；未匹配查询用$s_{\mathrm{o}}$而非一律为零的标签监督额外对象性输出。原工作还用变分近似的先验匹配项约束分布：高斯先验的方差随估计后验调整，并加入方差偏置，以处理少标注条件下的分布收缩。已知与未匹配查询之间的能量间隔项则改变分类得分的相对约束。这里的能量来自分类logit的对数指数和，不是像素亮度或物理能量。训练更新对象性相关特征及检测器参数，移动统计随批次更新；使用阶段沿用所接检测器的框和对象性输出，不再用真实标注给新图匹配。

例如批次的经验均值为$\bar{\bm q}$，移动均值可按$\bm\mu^{(t+1)}=\rho\bm\mu^{(t)}+(1-\rho)\bar{\bm q}$更新，$0\leq\rho<1$控制旧批次保留程度，对角方差同样利用当前统计更新。方差需要保持正数才能求逆，初始化和数值下界必须随实现固定。整体训练在已有检测框与分类误差外加入匹配对象的距离项、先验匹配及能量间隔，仍从所接检测器的预训练权重和规定预算开始。未知输出筛选沿用该检测器的对象性分类器，而式\eqref{eq:vision-owobj-objectness}主要提供训练软标签；不能在使用时拿新图真标签重新估计对象分布。

作为一维教学例子，若$\mu=0$、$\Sigma=1$，两个未匹配查询为0.2和2，则软对象性分别为$e^{-0.04}\approx0.961$和$e^{-4}\approx0.018$。给前者较强对象监督可以避免直接把未知候选抹掉，但若背景也靠近这个分布，它同样会被提出。这说明对象性需要同时检查未知召回和背景误检，不能由“保留更多区域”推断发现已经正确。

发现后可暂存原图编号、框、得分和当前类别表，再由人工或后续证据命名新类，形成新增训练记录。加入新类别是类别增量（class-incremental learning）条件，即后续学习扩展输出集合还要保留旧类能力；它不是这一次未命名检测已经完成的步骤。增量更新可参考\BookChapterRef{chap:ket-update-maintenance}{卷八“如何更新并维持已有知识”章}，当前必须保留未知区域及其来源。仓库装置被误报为已知货车、墙面被误报未知对象、加入装置类后货车退化，是三个不同错误。概念提示可在\ref{sec:vision-recognition-sam3}节进一步产生实例掩码，但只查询已有名称仍不能替代未知发现协议。

\subsection{像素归属与对象区分}
\label{sec:vision-recognition-segmentation}

框内常含背景，也可能含另一辆车。需要准确范围时，应进一步指定像素归属。\term{语义分割}（semantic segmentation）给像素分配类别，同类车辆可以连成一片；\term{实例分割}（instance segmentation）给不同车辆各自的掩码；\term{全景分割}（panoptic segmentation）还共同覆盖道路、天空等非实例区域，使每个像素获得类别，可数对象的像素再获得实例编号。

这一区别决定结果的组织方式。两辆车的并集可以是正确的车辆语义区域，却是错误的单个实例。道路通常按区域类别处理，不需要把每块路面任意编号。检测、语义与实例标签不能未经说明相互转换。

图\ref{fig:vision-segmentation-outputs}沿用同一个离散街景，显示增加实例编号与背景类别各带来什么信息。这里仅画出可见区域，框状图元表示教学掩码，不是算法检测框。

\begin{figure}[htbp]
 \centering
 % 同一离散场景的三种分割输出；字符在灰度下保留类别和身份区别。
\begin{tikzpicture}[font=\figurefont\small,x=4.5mm,y=-4.5mm]
 \foreach \offset/\title in {0/语义分割,8/实例分割,16/全景分割}{
  \begin{scope}[xshift=\offset*4.5mm]
   \node[anchor=south] at (3,0) {\title};
   \fill[BookPaper] (0,0.4) rectangle (6,4.4);
   \draw[BookMuted,line width=0.7pt] (0,0.4) rectangle (6,4.4);
   \draw[BookRule,line width=0.5pt] (0,1.4)--(6,1.4);
   \ifnum\offset=8
    \node[text=BookMuted] at (3,3.6) {不标背景};
   \else
    \node[text=BookMuted] at (3,3.6) {道路};
    \node[text=BookMuted] at (3,0.9) {天空};
   \fi
   \filldraw[fill=BookTeal!18,draw=BookTeal,line width=0.8pt]
      (0.5,1.8) rectangle (2.8,2.8);
   \ifnum\offset=0
    \filldraw[fill=BookTeal!18,draw=BookTeal,line width=0.8pt]
      (3.2,1.8) rectangle (5.5,2.8);
    \node at (1.65,2.3) {车};
    \node at (4.35,2.3) {车};
   \else
    \filldraw[fill=BookAmber!16,draw=BookAmber,dashed,line width=0.8pt]
      (3.2,1.8) rectangle (5.5,2.8);
    \node at (1.65,2.3) {车1};
    \node at (4.35,2.3) {车2};
   \fi
  \end{scope}
 }
\end{tikzpicture}

 \caption{三种分割结果的区别。语义分割保留类别，实例分割分开可数对象，全景分割共同保留背景类别与对象编号；文字编号和边界线型使身份区别在灰度打印中仍可辨认。图中区域为离散教学构造。}
 \label{fig:vision-segmentation-outputs}
\end{figure}

\subsubsection{从空间特征形成像素判断}
\label{sec:vision-recognition-dense-pixels}

密集特征仍比原图稀疏。把它们简单放大只能增加采样位置，细边界需要由保留位置的信息支持。单幅语义图和双时相变化图都以像素为输出，但前者判断类别，后者判断同一区域的状态是否改变。

\Needspace{60mm}
\paragraph{多尺度上下文与边界恢复：DeepLabv3+}
\label{sec:vision-recognition-deeplab}

若减小空间步幅，计算成本随特征尺寸上升；若持续下采样，路杆和车窗边缘可能难以恢复。DeepLabv3+用空洞卷积扩大特征读取范围，结合不同采样间隔的上下文，再以解码器融合较浅层细节\scite{vision-1802-02611}

空洞率$r$表示卷积核相邻采样位置间隔为$r$，例如一维三点核读取$x_{i-r},x_i,x_{i+r}$，参数个数仍为三。多尺度分支与全局汇总使位置可以同时利用邻域和较大范围。解码阶段上采样语义特征，对浅层特征先调整通道，再拼接和卷积，最后恢复输入网格上的类别得分。

训练对有标注像素计算交叉熵，忽略无效标签；推理逐位置选择类别并按预处理的逆坐标变换恢复原图。道路内部的类别主要依赖上下文，细路杆又需要足够分辨率。改变输出步幅和输入尺寸会改变这两类证据与预算，应共同记录。即使两车边缘都清楚，纯类别图仍未分配实例身份。

对训练图的有效位置集$\Omega_i$，类别标注为$y_{i,p}$，输出logit为$h_{i,p,r}$。按类别softmax得到$p_{i,p,r}$后，像素目标为
\begin{equation}
 L_{\mathrm{pixel}}=
 -\frac{\sum_i\sum_{p\in\Omega_i}\log p_{i,p,y_{i,p}}}
        {\sum_i|\Omega_i|}\text{。}
 \label{eq:vision-dense-pixel-training}
\end{equation}
每个位置的得分梯度是预测类别分布减真实类别指示，再经上采样、解码器和上下文分支传回主干。可从分类主干初始化，新增上下文和解码层由像素监督学习。缩放、裁剪图像时，对离散标签采用保持类别的采样，不能用颜色插值制造新类别。推理宜先把类别得分图插值回原图，再作类别选择；直接插值整数类别编号没有概率含义。该恢复只补采样网格，不补被输入删掉的细杆。

\paragraph{双时相上下文与变化图：BIT}
\label{sec:vision-recognition-bit}

遥感影像的目标可能是发现新增建筑，而不是给单幅影像中的每栋楼分类。BIT输入同一区域的两期图像，以共享的特征提取器形成空间表示，再把每期特征通过学习的空间权重汇总为少量语义词元。两期词元共同编码，使一幅图的局部判断能参考另一时期的上下文\scite{vision-2103-00208}

可将词元汇总理解为$\bm{T}=\bm{A}^{\top}\bm{F}$：$\bm{F}\in\mathbb{R}^{n\times d}$是$n$个位置的特征，$\bm{A}\in\mathbb{R}^{n\times k}$的每列在位置上归一化，得到$k$个加权汇总。编码后的词元再通过交叉注意力返回两期空间位置，恢复上下文增强的特征；对两期特征作差异计算并用预测头输出变化概率。

训练使用配准影像与像素变化标签，模型参数由变化损失更新。新建筑应对应语义变化，树木变绿是否计为变化则取决于标注定义。若两期建筑整体错开几个像素，差异分支也会产生边缘响应，因此配准是任务前提。上下文交换能够利用关联，却不能把传感器差异自动解释成真实建设活动。

每次训练同时裁剪、翻转两期图像及变化标签，保持同一位置仍指同一地表；不能对两期各取互不对应的随机区域。若增强特征分别为$\bm{F}_1',\bm{F}_2'$，变化头从$|\bm{F}_1'-\bm{F}_2'|$形成两类得分，再在有效像素上用式\eqref{eq:vision-dense-pixel-training}的类别目标学习“变化／不变”。主干可以从已有图像权重初始化，词元汇总、双期编码和空间解码由变化数据学习。冻结或适配主干需按实际设置记录。使用时保留相同双期顺序和配准，把分数恢复到共同原网格，阈值选出变化区域；区域只说明标签所定义的状态差异，没有自动解释变化原因。

\subsubsection{从对象或区域形成掩码}
\label{sec:vision-recognition-mask-sets}

为了分开同类对象，可以先取得候选对象，再在对象内预测掩码；也可以直接预测一组掩码并匹配实例。用户点框提示又提供了另一种目标选择依据。以下沿这三种接口展开。

\paragraph{区域对齐与实例掩码：Mask R-CNN}
\label{sec:vision-recognition-mask-rcnn}

Mask R-CNN在候选框上并行预测类别、框修正和实例掩码。为了让框内位置与原特征准确对应，\term{区域对齐}（RoIAlign）保留连续坐标，以双线性插值在各小格中采样，再汇总成固定尺寸特征，避免先把框边界粗略取整造成的位移\scite{vision-1703-06870}

候选区域来自区域建议网络（region proposal network, RPN），即在特征位置的参考框上先预测对象性和框修正，形成可能含对象的矩形。RPN的前景背景分配依据真实实例框；候选解码和去重后，第二阶段再按候选与真实框的重叠分配类别、框修正和实例掩码。训练可以用已有分类主干初始化，并共同优化RPN目标和区域分支目标。一个正候选必须记住匹配实例编号，否则即使“车辆”类别正确，仍可能使用另一辆车的掩码标签。

设四个相邻特征值为$f_{0,0},f_{0,1},f_{1,0},f_{1,1}$，连续位置在水平方向偏移$a$、竖直方向偏移$b$，二者均在$[0,1]$内，则插值为
\begin{equation}
 f(a,b)=(1-a)(1-b)f_{0,0}
 +a(1-b)f_{0,1}+(1-a)bf_{1,0}+abf_{1,1}\text{。}
 \label{eq:vision-roialign}
\end{equation}
用数值$(0,4,2,6)$与$(a,b)=(0.25,0.5)$可得$f=2$，既保留横向的四分之一偏移，也保留纵向中点。若直接取最近格点，这些偏移信息就会消失。

原方法对每个类别预测一个二值掩码通道。训练时只对该正候选匹配到的真实类别通道计算像素二元损失，其他类别通道不在此样本上竞争像素归属。使用时取预测类别的掩码，将区域坐标缩放回原框并放回原图。两个重叠车辆由两个候选各自形成掩码，重叠区域的实例归属与类别判断是不同问题。

区域掩码网格的每个目标值由匹配实例在该区域中的前景归属取得，预测经sigmoid后用二元交叉熵约束。总目标将RPN对象性及框误差、区域类别及框误差、正候选掩码误差相加，再沿共享特征反向更新；负候选没有实例掩码目标。推理先按类别与框分数筛选、去重，再对保留实例取相应掩码通道。若一个$28\times28$掩码对应原框$(20,10,80,50)$，归一化区域位置$(\xi,\eta)$回到$(20+60\xi,10+40\eta)$，放置时还需处理框越界和原图尺寸。恢复后的掩码、类别和框共享实例编号，不能分别排序后重新拼接。

区域对齐解决采样量化误差，不能补回候选阶段漏掉的车辆，也不能修正把候选匹配到错误实例的监督。其关键点扩展改变输出目标为有部位语义的位置，不能由掩码准确直接推断关节准确。

\paragraph{掩码查询与局部注意力：Mask2Former}
\label{sec:vision-recognition-mask2former}

先有框再分割并非唯一方式。Mask2Former从多尺度像素特征与查询直接产生掩码集合。每个查询预测类别以及掩码向量，掩码向量与各位置像素嵌入比较，形成该查询的前景概率；训练用集合匹配连接预测与目标，再计算分类、二元掩码和Dice损失\scite{vision-2112-01527}

Dice项度量软预测与真实前景的重叠。对于预测$p_j$和真实标签$m_j$，一种带平滑量$\epsilon>0$的形式为
\begin{equation}
 L_{\mathrm{Dice}}=
 1-\frac{2\sum_j p_jm_j+\epsilon}
 {\sum_j p_j+\sum_j m_j+\epsilon}\text{。}
 \label{eq:vision-dice}
\end{equation}
它与逐像素损失提供不同信号：一个关注区域重叠，另一个约束各位置的预测。训练中的采样点、归一化及匹配权重还需按具体实现固定。

掩码注意力把上一轮预测的前景区域变为后续跨注意力的可读范围。对于查询$i$与位置$j$，在注意力得分中加入屏蔽项$M_{i,j}$：允许读取为0，不允许为$-\infty$，再作softmax。这样每个查询重点读取自己可能对应的区域；多尺度特征与反复更新使区域可以细化。

设第$i$个查询解码成掩码向量$\bm{e}_i\in\mathbb{R}^d$，位置$p$的像素特征为$\bm{f}_p\in\mathbb{R}^d$，掩码概率为$\sigma(\bm{e}_i^\top\bm{f}_p)$。每次查询更新后重新预测掩码并形成下一层读取范围，训练匹配则比较类别和采样掩码误差，不用检测框替代掩码监督。匹配查询承担分类、二元掩码与Dice项，未匹配查询学无对象；训练数据的裁剪同时作用到目标区域。若屏蔽使所有位置都不可读，注意力的分母将失效，实现需恢复可读位置，而不能直接对全$-\infty$求softmax。

语义输出可将各查询的类别概率与该查询掩码概率相乘，对查询求和，再逐像素选类；实例输出保留查询编号，按类别及掩码质量筛选；全景输出还需处理区域竞争、无对象和不可数类别。三种规则使用相同输入街景，却保存不同信息。例如两个车辆查询可以在语义图中合并，实例图中不能合并。最终将分数图恢复原图后再组织区域，避免把裁剪坐标当作原像素。

原工作按任务与数据分别训练，不能把统一架构理解为一份参数无条件通用。错误的早期掩码也可能限制后续读取，因此要检查小对象和被遮挡对象是否在迭代中持续被漏掉。

\paragraph{图像与提示编码、歧义掩码：SAM}
\label{sec:vision-recognition-sam}

有时用户只希望选中图上的一个区域，甚至尚未给它类别名。SAM把图像编码、提示编码与掩码解码分开，同一图像的昂贵编码可以复用，点或框变化时更新提示并运行较轻的解码器。点提示包含位置与前景、背景标签，框提示提供空间范围；已有掩码也可作为稠密提示\scite{vision-2304-02643}

解码器让提示与图像表示交换信息，输出若干候选掩码及质量估计。训练提供图像、目标掩码与采样提示，歧义输出通过多个候选容纳一个提示的不同合理解释。质量预测估计候选与监督目标的重叠程度，服务选择，并非在使用时知道了真实掩码。

图像编码器从已有视觉权重开始，提示位置先归一化到图像坐标，并附正负点或框角类型的学习嵌入；稠密掩码提示经下采样与图像特征对齐。解码器让提示读取图像，也让图像读取提示，输出掩码向量；向量与上采样图像特征逐位置比较，得到低分辨率掩码logit。多个输出中对监督掩码误差较小的一项承担歧义训练的分割目标，原方法结合焦点及Dice项，并用真实候选IoU监督质量头。训练模拟后续点击，在错误区域补提示，参数由提示条件下的残差更新；不能把真实目标直接放进使用时的质量分支。

使用同一街景时先缓存图像编码，点或框改变只更新提示计算及解码。候选logit上采样到输入图，去除加边后恢复原尺寸，再按阈值变成掩码；质量得分可推荐候选，用户意图仍决定选择车门还是整车。交互停止可以由用户确认或预设点击预算控制，不存在仅凭质量头保证正确意图的终止规则。

落在车门上的正点可能指车门或整车。用户选择接近意图的候选，再在漏掉的车尾补正点，或在被误选的路面补负点，才逐步消除歧义。应报告提示来源、点击次数与纠错策略；用真实标注框作提示与由检测器自动给框，具有不同的信息与错误条件。SAM的区域输出本身不自动附带“货车”名称。

\paragraph{概念存在与实例范围：SAM 3}
\label{sec:vision-recognition-sam3}

点提示通常选择一个局部区域，“找出所有校车”则需要判定概念是否出现，并枚举符合概念的实例。SAM 3接收概念短语和带正负标签的图像区域示例，区域视觉特征参与概念表示；随后由对象查询预测框、掩码和局部匹配得分。原始报告将这一接口称为可提示概念分割（promptable concept segmentation, PCS），即用概念而非单个局部点击指定全部符合条件的实例\scite{vision-2511-16719}。

概念是否在整幅图像中出现，与每个候选是否匹配，承担不同职责。共享存在词元与对象查询一同进入解码器，再由存在预测头得到$p_{\mathrm{pres}}$；查询$i$的局部分数$s_{i\mid\mathrm{pres}}$表达概念存在条件下的匹配，最终使用$s_i=p_{\mathrm{pres}}s_{i\mid\mathrm{pres}}$。训练分别提供存在标签与实例范围；局部分类损失在概念存在的图像上计算，使全局缺席主要由共享分支学习。不能仅用“图里有车”这一标签学出所有车辆的像素边界。

文字先取得词元特征，区域示例则同时编码框位置、示例中的视觉内容及正负标签，与文字共同形成概念条件。图像特征通过融合编码读取条件，对象查询再细化框与掩码；负示例告诉模型哪些相似对象不属于当前概念，不是将该区域的像素全部删除。训练记录包含概念、存在与否、全部匹配实例及提示。存在头用二元交叉熵覆盖正图与难负概念；概念存在时，实例按类别和范围代价唯一匹配，再用局部得分、框和掩码误差训练查询。缺席概念不向局部分类分支施加逐查询负梯度，职责与普通DETR的无对象训练有别。

这些接口来自有阶段区分的参数学习。SAM 3先用图文对预训练视觉和文字编码器，再以图像、短语和掩码三元组训练图像检测器及两侧编码器；三元组既有伪标注也有人工标注。高质量图像阶段加入存在头，并根据漏检采正区域、根据高分误检采负区域，累加示例以模拟纠错。该阶段联合降低框$L_1$、广义IoU、掩码焦点与Dice以及分类、存在误差。图像检查点因而已经学过概念与范围关联；使用时更换“校车”文字只改变条件，不重新学习权重。后续视频跟踪冻结主干的阶段不应倒推为前面的图像训练也一直冻结。

若教学中存在分数为0.9，两个实例局部分数为0.8和0.6，最终分数为0.72和0.54；不存在概念的存在分数为0.05时，同样局部分数只剩0.04和0.03。最终阈值决定保留哪些查询，再将各掩码放回原图并保持独立实例编号。共享分支的一个错误会同时改变所有查询，故不能只看保留掩码的边界而忽略被整体压低的概念。

对于不存在的概念，即使某个候选与“车”这一宽泛外观相似，共享存在判断也可以抑制错误输出。反过来，共享存在分支若把“校车”判为不存在，多个局部查询都可能被压低，形成相关的漏检。评价因此分别检查概念有无、每实例召回和掩码边界。简短概念与区域示例扩大了目标选择方式，但尚不足以推断任意关系嵌套或长指令都已解析；视频身份传播在\ref{chap:vision-video-recognition}章继续讨论。

\subsection{软边界与前景贡献的估计}
\label{sec:vision-recognition-matting}

发丝和半透明边缘的一个像素可能同时收到前景与背景的光。二值掩码只能选择保留或删除，无法表达这种混合。\term{图像抠图}（image matting）估计前景贡献比例$\alpha\in[0,1]$。在线性颜色空间的简化合成模型中，每个像素满足
\begin{equation}
 \bm{x}=\alpha\bm{f}+(1-\alpha)\bm{b}\text{，}
 \label{eq:vision-matting}
\end{equation}
其中$\bm{x},\bm{f},\bm{b}\in\mathbb{R}^3$是观测、前景与背景颜色。只有三个颜色观测，却有前景、背景和$\alpha$共七个未知标量，需要额外条件或先验；二值分割标签并不能唯一解出它们。

\subsubsection{三分图条件与alpha细化：Deep Image Matting}
\label{sec:vision-recognition-deep-matting}

Deep Image Matting使用\term{三分图}（trimap）指定确定前景、确定背景和未知区域，把RGB图与这一提示通道共同输入编码解码网络，再用细化网络改善$\alpha$边界。确定前景约束$\alpha=1$，确定背景约束$\alpha=0$，学习的重点是未知带中的混合比例\scite{vision-1703-03872}

训练可利用已知前景、背景与$\alpha$构造合成图像，监督$\alpha$误差，并按式\eqref{eq:vision-matting}检查预测贡献比例能否重建合成观测。使用时输出$\alpha$及其已知区约束；要把对象放到新背景$\bm{b}'$，仍需估计或取得前景颜色$\bm{f}$，再合成为$\alpha\bm{f}+(1-\alpha)\bm{b}'$。只取得$\alpha$并不意味着前景颜色已无污染。

RGB与三分图组成四通道输入，编码器压缩空间信息、解码器恢复未知带的$\hat\alpha$，细化网络读取图像及初步$\alpha$再预测细节。训练在未知区$U$中比较$|\hat\alpha_p-\alpha_p|$，并比较$\hat\alpha_p\bm f_p+(1-\hat\alpha_p)\bm b_p$与合成观测的颜色误差。两项通过同一混合公式连接：前者约束贡献比例，后者约束该比例的颜色后果。已有视觉权重可初始化编码层，新增输入和解码部分由合成监督学习，细化阶段再利用已训练的初步结果。三分图可以从目标$\alpha$构造未知带以训练，但使用时必须实际取得提示，不能把真实$\alpha$隐藏在输入中。

例如灰度前景为0.8、原背景为0.2、$\alpha=0.25$，观测为0.35。换成0.6的背景，正确合成为$0.25\times0.8+0.75\times0.6=0.65$；若二值化为背景，结果为0.6，发丝的贡献消失。若直接把观测0.35当成前景颜色，又会混入旧背景。提示误把发丝标成确定背景时，网络还可能被错误条件限制。边界、贡献比例和最终合成质量应分别核验。

\section{数量、姿态与空间测量}
\label{sec:vision-recognition-measurement}

“有几辆车”“人的手腕在哪里”“路牌有多远”分别要求数量、具有部位语义的位置以及带坐标或尺度的空间量。它们可以调用前面的对象范围，但监督和误差不能由检测指标替代。

要把手腕或路牌放到三维空间中，先看一个已经选定的表面点如何成像。\emph{针孔相机}（pinhole camera）用相机中心和一个成像平面描述投影：空间点与相机中心的连线穿过图像平面，交点就是它的像素位置。这里忽略镜头畸变，采用位于相机前方的虚拟成像平面，使像素坐标向右、向下增加，相机光轴向前。实际镜头存在畸变时，应先用标定参数校正，或在投影中加入畸变模型。

同一点可以在不同坐标系中表达。\emph{世界坐标}（world coordinates）是为整个场景选定的公共参照，\emph{相机坐标}（camera coordinates）则以当前相机中心为原点。记世界点为$\bm{P}_{\mathrm w}\in\mathbb{R}^3$，相机点为$\bm{P}_{\mathrm c}=(x_{\mathrm c},y_{\mathrm c},z)^\top$。\emph{外参}（extrinsic parameters）$\bm{R}\in\mathbb{R}^{3\times3}$和$\bm{t}\in\mathbb{R}^3$把前者变成后者，即$\bm{P}_{\mathrm c}=\bm{R}\bm{P}_{\mathrm w}+\bm{t}$，其中$\bm{R}^{\top}\bm{R}=\bm{I}$、$\det\bm{R}=1$。这里$\bm{t}$是世界到相机变换的平移，不是相机中心的世界坐标；相机中心为$\bm{C}_{\mathrm w}=-\bm{R}^{\top}\bm{t}$，逆变换为$\bm{P}_{\mathrm w}=\bm{R}^{\top}(\bm{P}_{\mathrm c}-\bm{t})$。

只知道相机点还不能得到像素，因为焦距和图像原点仍需确定。\emph{内参}（intrinsic parameters）描述这一步映射。忽略像素轴倾斜时，内参矩阵为
\[
 \bm{K}=\begin{pmatrix}
 f_x&0&c_x\\0&f_y&c_y\\0&0&1
 \end{pmatrix}\text{，}
\]
其中$f_x,f_y>0$是以像素为单位的焦距，$(c_x,c_y)$是主点，即光轴落到图像平面的位置。沿光轴深度$z>0$的点与投影平面构成相似三角形，横纵方向分别给出$u-c_x=f_xx_{\mathrm c}/z$、$v-c_y=f_yy_{\mathrm c}/z$。把这两个关系写在同一个齐次式中，得到
\begin{equation}
 z\begin{pmatrix}u\\v\\1\end{pmatrix}
 =\bm{K}(\bm{R}\bm{P}_{\mathrm w}+\bm{t})\text{。}
 \label{eq:vision-camera-projection}
\end{equation}
齐次坐标的末项让平移可以与投影共同表达，并不增加一个物理空间维度。$z=0$时除法未定义，$z<0$的点在相机后方，不能按当前可见图像点使用。

式\eqref{eq:vision-camera-projection}也说明像素为何不能单独确定深度。令$\tilde{\bm{p}}=(u,v,1)^\top$，则沿这条像素射线的相机点都是$\bm{P}_{\mathrm c}(z)=z\bm{K}^{-1}\tilde{\bm{p}}$，其中$z>0$待定。光轴深度与到相机中心的直线距离$\rho$不是同一个量：
\begin{equation}
 \rho=\lVert\bm{P}_{\mathrm c}\rVert_2
 =z\sqrt{1+\left(\frac{u-c_x}{f_x}\right)^2
            +\left(\frac{v-c_y}{f_y}\right)^2}\text{。}
 \label{eq:vision-depth-range-measurement}
\end{equation}
教学中取$f_x=f_y=800$像素、主点$(320,240)$和像素$(520,240)$。若$z=4$米，反投影得到$(1,0,4)$米，直线距离为$\sqrt{17}\approx4.123$米；若$z=8$米，点$(2,0,8)$仍投到同一像素。两次投影相同，测量量却不同。

多视图增加了射线约束，但还需区分\emph{相对尺度}（relative scale）与\emph{度量尺度}（metric scale）。前者只确定长度比例，后者还确定一单位对应多少米。若把全部世界点和相机平移同时乘$a>0$，相机点与$z$也乘$a$，式\eqref{eq:vision-camera-projection}的像素不变。已知相机基线、已知物体长度或带单位的深度观测可以确定这个比例；只把首相机设为原点，只固定了参照位置和朝向，没有测出尺度。网络从人体或街景训练分布预测的尺度是附加先验，也须在当前设备与场景中验证。

同一个世界点进入第二幅图时，要先回到世界坐标，再进入第二相机。对两相机外参$(\bm{R}_1,\bm{t}_1)$与$(\bm{R}_2,\bm{t}_2)$，消去世界点可得
\begin{equation}
 \begin{aligned}
 \bm{P}_{\mathrm c,2}
 &=\bm{R}_{2,1}\bm{P}_{\mathrm c,1}+\bm{t}_{2,1},\\
 \bm{R}_{2,1}&=\bm{R}_2\bm{R}_1^\top,\qquad
 \bm{t}_{2,1}=\bm{t}_2-\bm{R}_2\bm{R}_1^\top\bm{t}_1\text{。}
 \end{aligned}
 \label{eq:vision-relative-camera-measurement}
\end{equation}
将第一幅图的反投影点代入，再用第二幅图的$\bm{K}_2$投影，就得到一个可检查的跨图位置。裁剪和缩放也必须进入内参：若横纵缩放为$a_x,a_y$，裁剪原点为$(u_0,v_0)$，则新焦距为$(a_xf_x,a_yf_y)$，新主点为$(a_x(c_x-u_0),a_y(c_y-v_0))$。把新像素配旧内参会直接改变射线方向。

落在图像内且位于相机前方，只是可见的必要条件。\emph{遮挡}（occlusion）指同一射线上更近的表面挡住更远的表面；不透明场景中的像素通常记录最先遇到的表面。两幅图之间比较一个投影点时，还要检查源图深度是否与该点相容，否则取到的可能是挡在前面的车辆。人体被挡住的手腕可以由结构先验推断，但不是已经观察到的关节点。以下姿态、深度和场景恢复都沿用这些坐标、尺度与可见性约定；计数只需规定个体与区域规则，不必为此先重建整个场景。

\subsection{对象数量的估计}
\label{sec:vision-recognition-counting}

稀疏对象可以先逐个识别再计数，密集遮挡场景则可以直接估计数量的空间分布。前者保留对象清单，后者强调总量，两者取得正确总数的原因可能不同。

\subsubsection{检测实例汇总与重复计数控制}
\label{sec:vision-recognition-instance-counting}

给定去重后的检测集合，可按类别、区域与可见条件筛选，再把满足条件的指示量相加。若区域为$R$，以框中心$\bm{c}_i$是否在$R$内为计入规则，则
\begin{equation}
 \hat n_R=\sum_i
 \mathbf{1}\{c_i\in C_{\mathrm{target}}\}
 \mathbf{1}\{s_i\geq\eta\}
 \mathbf{1}\{\bm{c}_i\in R\}\text{。}
 \label{eq:vision-detection-count}
\end{equation}
这里$c_i$表示类别，$\bm{c}_i$表示二维中心，两者字形不同。中心落入与“任意可见部分落入”是两种不同规则，边界目标会得到不同计数。

例如真实有5辆车，漏掉1辆，另1辆被重复预测，总数仍可能为5。只看计数误差会掩盖两种错误；应同时检查对象匹配，尤其是遮挡和边界处的实例。跨帧统计还要避免同一辆车被多次计入，需由跟踪维持身份。

\subsubsection{多尺度密度预测与积分：MCNN}
\label{sec:vision-recognition-mcnn}

在人群过密、单个身体范围难以标注时，可用头部点标注学习\term{密度图}（density map），即每个位置分摊多少计数贡献。MCNN以多个不同感受范围的卷积分支处理尺度变化，连接各列特征后回归一张密度图\scite{vision-mcnn2016}

若标注点为$\bm{p}_i$，可用归一化核构造监督
\begin{equation}
 D(\bm{p})=\sum_{i=1}^{n}G_{\sigma_i}(\bm{p}-\bm{p}_i),
 \qquad
 \sum_{\bm{p}}G_{\sigma_i}(\bm{p}-\bm{p}_i)=1\text{。}
 \label{eq:vision-density-target}
\end{equation}
原工作令核宽与邻近头部点的平均距离相关，用局部拥挤程度近似尺度变化。卷积列和汇总层以预测密度与目标密度的平方误差训练。使用时对预测网格求和得到人数；密度可以是非整数分布，最终总数也可保留实值用于评价。

更具体地，若第$i$点最近$k$个邻点的平均距离为$\bar d_i$，可设$\sigma_i=\beta\bar d_i$，其中$\beta>0$及邻点数由标签构造约定。它只是从点间距得到尺度近似，孤立点或点数不足时需回退核宽。每列卷积读取不同邻域，共同输出单位为“该密度格分摊人数”的$\hat D_i$。经验目标可写为$L=(2n)^{-1}\sum_{i=1}^n\sum_p(\hat D_i(p)-D_i(p))^2$，每格梯度与$\hat D_i-D_i$成正比，训练将它传回汇总层及各列。原MCNN分别预训练卷积列后再联合适配；使用已训练参数时没有逐对象分类步骤。

若连续密度单位为人／像素面积，缩放$(x',y')=a(x,y)$使面积扩大$a^2$，需取$D'(x',y')=D(x'/a,y'/a)/a^2$，变量代换才给出$\int D'=\int D$。离散网格应相应重分配计数质量；输出步幅变大时，标签下采样也应保存总和。裁剪训练时先按规定中心规则确定哪些点计入，再为这些点在裁剪域构造归一化核；把原密度图任意裁一块，会把边外对象的核尾也计入，并丢掉边内对象的一部分质量。

三个点各贡献1，整图密度之和便为3。若一个边缘点的核有一半落到图外，直接截断后该点只贡献约0.5；有限网格构造标签时需按计入规则归一化有效核。图像缩放也要保持总和：将每个密度单元分成四个，四个新单元应分享原来的质量，不能每个都复制原值。

总和正确仍可能把本属左半边的人数移到右半边。应同时看局部计数、空间误差与场景变化。透视改变对象尺度，新的拍摄高度和密度范围可能超出训练条件；密度图没有自动恢复每个人的身份或完整位置。

\subsection{关节与姿态的估计}
\label{sec:vision-recognition-pose}

人体关节不同于图像中任意可匹配的角点：左腕、右肘具有固定部位含义。多人场景还必须知道这些部位属于谁。二维位置位于图像平面，三维姿态则要增加深度与坐标约定。

\subsubsection{二维部位定位与关联：OpenPose}
\label{sec:vision-recognition-openpose}

自顶向下路线先检测人，再在各人区域内定位关节；多人重叠时，候选框和裁剪质量影响后续姿态。OpenPose采用自下而上路线，在整图预测关节热图以及连接部位的\term{部位亲和场}（part affinity field, PAF），再组装人体\scite{vision-1611-08050}

热图在对应关节附近取得高响应，峰值形成候选位置。PAF是每种肢段的一张二维向量场：真实肢段附近写入从一种部位指向另一种部位的单位向量，其他位置为零，重叠肢段按规则合并。训练从关键点标注构造两类目标，对有效位置计算热图与向量场误差，并逐阶段细化。

设候选肘部、腕部为$\bm{a},\bm{b}$，方向$\bm{u}=(\bm{b}-\bm{a})/\lVert\bm{b}-\bm{a}\rVert_2$，对应PAF为$\bm{L}$。沿候选肢段平均方向一致程度：
\begin{equation}
 s(\bm{a},\bm{b})=
 \int_0^1
 \bm{L}((1-\lambda)\bm{a}+\lambda\bm{b})^\top
 \bm{u}\,\mathrm{d}\lambda\text{。}
 \label{eq:vision-paf-score}
\end{equation}
实现中在有限个点上采样求平均，再按部位连接类型作匹配并组合人体。$\bm{a}=\bm{b}$时方向未定义，应排除这一退化候选。

每类关节的目标热图可由标注位置周围的高斯响应构造，多人同类关节按规定合并；每种肢段在其附近带状区域写入方向，重叠区域平均有效方向。各阶段对有标注位置计算热图及PAF平方误差，未标注身体区域以有效掩码排除，梯度共同更新特征及分支。推理先从热图峰值取点，再对每种肢段建立候选两端的评分矩阵，禁止同一部位候选在该连接类型中任意占多条连接，按匹配结果逐步拼成身体。人体组合还需利用共享关节编号合并连接，不能把所有高分线段简单相加。

假设五个采样点的方向内积为$(1,1,0.8,1,0.7)$，平均为0.9；连到另一人的手腕时变为$(0.1,0,-0.2,0,0.1)$，平均为0。两个手腕都定位正确，也仍需这一关联步骤才可能把手臂接对。投影重合、严重遮挡和交叉肢段会使证据混杂，因此关节位置与人体组装应分开评价。

\subsubsection{二维关节到三维姿态：Martinez基线}
\label{sec:vision-recognition-lifting}

取得二维关节后，可以把按固定部位顺序排列的$2j$维坐标输入回归网络，输出$3j$维三维坐标。Martinez等人的基线使用全连接层与残差模块学习这一映射，使二维检测和三维恢复的误差来源能够分开研究\scite{vision-1705-03098}

训练规范化输入坐标，以对应的三维姿态作回归监督。三维目标通常相对于根关节表示，并遵守指定相机坐标、尺度和关节集合。使用时必须复现训练的规范化，再将输出还原为该协议中的坐标；关节顺序错位会使一个部位的输入被当成另一个部位，即使维度仍然匹配。

记按训练集统计规范化的二维坐标为$\bm{x}_i\in\mathbb{R}^{2j}$，相机坐标中的三维关节减去根关节后为$\bm{y}_i\in\mathbb{R}^{3j}$。提升网络先把$\bm{x}_i$投到较宽隐层，残差块以$\bm{h}'=\bm{h}+F_{\bm{\theta}}(\bm{h})$更新，其中$F$包含全连接、归一化、非线性及训练时随机屏蔽；最后线性层恢复三维输出。以$(1/n)\sum_i\lVert g_{\bm{\theta}}(\bm{x}_i)-\bm{y}_i\rVert_2^2$训练，从初始化权重逐批更新，保留网络和训练集规范化统计。输出逆规范化后得到相对根关节坐标，不含被减去的根位置；要得到人体到相机的距离，还需要另估根平移，不能直接由这组相对坐标读取。

节首教学像素$(520,240)$可对应$(1,0,4)$或$(2,0,8)$。人体长度先验能够在这些解释中形成偏好，但真正尺度来自哪个标注协议仍需说明。真实二维输入隔离关节检测误差，检测二维输入则要先逆人体裁剪并按同一规范化处理。已知相机的多视图关节可以通过射线求交或最小重投影误差估计三维点，增加的是观测，不是把单目回归网络的输出自动变成多视图证据。

向前或向后弯曲的肢段也可能有相近二维投影，网络用训练中的人体结构先验选择一种三维解释。评价要说明是否允许旋转、平移及尺度对齐：相似变换对齐后的关节误差较小，不意味着原始输出已具有正确朝向或距离。

\subsubsection{从关节坐标到可提示的人体表面恢复}
\label{sec:vision-recognition-sam3d-body}

两个三维骨架可以具有相同关节位置，却对应不同肩宽、躯干厚度和手指表面。若需要把人体轮廓投回图像，关节点还不够。\emph{人体网格恢复}（human mesh recovery, HMR）估计由顶点和连接面片组成的人体表面，同时解释其姿态、体形与相机。表面顶点数只是表示容量，不能证明被遮挡的背部已经恢复真实形状。

参数化人体表示用少量身体参数生成骨架和网格，使输出遵守同一人体连接结构。SAM 3D Body使用Momentum Human Rig（MHR），将骨架结构、关节姿态和表面形状分开控制。骨架结构规定连接和骨骼长度，姿态规定各关节如何相对旋转，表面形状规定骨架周围的身体外形。记这些参数为$\bm{\beta}_{\mathrm{sk}},\bm{\theta}_{\mathrm{pose}},\bm{\beta}_{\mathrm{surf}}$，已提供的人体模型分别产生关节和顶点：
\[
 \bm{J}=\mathcal{J}(\bm{\beta}_{\mathrm{sk}},
                   \bm{\theta}_{\mathrm{pose}}),\qquad
 \bm{V}=\mathcal{V}(\bm{\beta}_{\mathrm{sk}},
                   \bm{\theta}_{\mathrm{pose}},
                   \bm{\beta}_{\mathrm{surf}})\text{。}
\]
这里$\bm{J}\in\mathbb{R}^{j\times3}$、$\bm{V}\in\mathbb{R}^{v\times3}$，$j,v$由模型资产规定；这两个表达式只说明输出分工，不是把MHR的全部变形计算替换成一个线性式。把关节或顶点经相机投影，才得到可与原图比对的位置和轮廓\scite{vision-sam3d-body-2026}。

模型首先取得人体裁剪及裁剪到原图的变换，归一化后由共享图像编码器产生空间特征。已发布版本可使用DINOv3视觉主干，身体和手部解码器由人体任务监督学习；推理还需加载相应MHR模型资产，只有网络权重不足以确定顶点连接和变形规则。身体解码器的查询中包含初始MHR和相机估计、辅助二维和三维关节点，以及可选的用户提示。二维关键点提示写成$(x,y,\ell)$，其中$\ell$指出关节身份；位置编码与学习的关节嵌入使“左腕在这里”不同于一个无部位含义的点。掩码提示经卷积编码后加到图像特征，提供人体占据区域。查询读取图像和提示，输出姿态、表面形状、骨架及相机参数，再通过MHR资产恢复完整网格。

手部占据整人裁剪中的像素可能很少，模型另设可选手部裁剪和独立手解码器，仍共用图像编码器权重。身体解码器本已能输出包含手足的完整人体，手解码器提供手部专门训练形成的补充。原方法不把局部手参数直接替换进去后就结束：直接替换可能连带改变肘部，因而用手解码得到的腕位置和身体解码得到的肘位置再次提示身体分支，细化后按人体连接树合并。遮挡手部若没有可靠检测，应保留身体输出或关闭局部细化，而不是把不存在的裁剪当成有效观察。

训练样本的三维依据来自多阶段标注：人工关键点、可微拟合、多视图几何和密集关键点检测以不同组合产生监督，还包含合成和多视图资料。原论文首稿为2026年2月，并不据此认定系统首次发布也发生在该月。训练将二维关节放在相应裁剪坐标中监督，将身体三维关节减去骨盆、手关节减去手腕后比较；已标注手部才参与对应目标。关节使用带学习不确定性权重的$L_1$项，姿态与形状参数使用$L_2$项，另加关节限位和手框误差。随机多轮提示模拟纠错，对用户给定关节提高监督权重，部分三维项逐步升权。训练得到的是图像到人体参数的映射，不是为每张新图重新收集真实三维网格。

教学中设一个人左腕被遮住，无提示预测投回原图$(120,150)$，而用户根据额外依据给点$(128,146)$。两者像素残差为$\sqrt{8^2+4^2}\approx8.94$像素。加入关节提示后应重跑解码、恢复网格并重新投影，检查该残差以及肘腕连接、手部形状和身体轮廓是否共同改变；加入人体区域掩码时则主要检查轮廓条件。这个计算只是定义需要检查的量，没有假造提示后的模型误差或宣称误差必然下降。

使用时人体框可由外部检测器取得，框误差属于完整流程的一部分。用户点需先变到人体裁剪坐标，预测关节和网格在相机条件下投影后，再用裁剪逆变换回原图。提示轮次可由用户确认或预算停止，最终保存人体编号、MHR参数、相机、网格及提示来源。检查点的量纲与首相机约定需要保留；未知真实身高时，不能把模板的身体尺寸当成该人的测量结果。

一个轮廓相符的网格仍可能把手腕置于不同深度，节首的尺度歧义也仍存在。应分别检查关节、顶点表面、手足细节、提示一致性和相机，而不只比较一张渲染图。相机坐标中的合理人体不能直接给出世界坐标的真实身高、动作意图或生物力学量；后两类还需时序、力学与采集条件。连续姿态及身份关联由\ref{chap:vision-video-recognition}章承接。

\subsection{深度与距离的估计}
\label{sec:vision-recognition-depth}

现在把测量对象从关节换成每个可见表面像素的光轴深度。节首的式\eqref{eq:vision-camera-projection}已经给出像素射线，尚未确定的是射线上的$z$。单幅外观可以提供训练先验，多幅已知相机的图像可以提供跨射线约束；两条路线都需说明输出采用深度、逆深度还是其他相对量，以及是否已经确定度量尺度。

\subsubsection{单幅相对深度与度量适配：Depth Anything V2}
\label{sec:vision-recognition-depth-anything}

单目模型利用物体形状、透视、遮挡和训练先验，从一张图预测密集深度。Depth Anything V2以合成数据的精确深度训练较强教师，再为大量真实图像产生深度伪标签，供不同规模学生学习；图像编码器与稠密解码器把表示接回各空间位置\scite{vision-2406-09414}

这一训练桥接应区分两个信息来源：合成图像提供几何标签，真实图像提供目标外观，伪标签质量依赖教师。原工作以已有DINOv2编码器和密集预测变换器（dense prediction transformer, DPT）解码器作为起点；后者把不同层空间特征逐步融合、上采样成逐位置数值，不把全图汇总向量当作一张深度图。较大教师先用合成RGB和精确深度训练，教师为真实RGB产生伪深度后，学生只从这些真实图与伪标签对继续学习，教师标签不随学生每次梯度更新而变成真实测量。

基础检查点预测的是尺度、平移不变的相对逆深度。为看清“允许尺度和平移”如何作用，设预测相对量为$\hat q_p$、目标为$q_p$，一个教学用的平方对齐目标可以写成
\begin{equation}
 \min_{a,b}\frac1{|\Omega|}
       \sum_{p\in\Omega}(a\hat q_p+b-q_p)^2\text{。}
 \label{eq:vision-relative-depth-alignment}
\end{equation}
这里$\Omega$是当前图的有效监督位置，$a,b$是比较时的对齐量，不是已知米制相机参数。对$a,b$求导为零得到$\sum\hat q_p(a\hat q_p+b-q_p)=0$、$\sum(a\hat q_p+b-q_p)=0$；若预测有非零方差，便有$a=\operatorname{Cov}(\hat q,q)/\operatorname{Var}(\hat q)$、$b=\bar q-a\bar{\hat q}$。若全图预测为常数，尺度无从估计。这个平方例子解释对齐消去什么信息，不冒称原论文尺度平移不变目标的逐项实现。原训练还比较相邻位置残差的梯度，防止整体对齐虽合适、局部边界却过于平滑。

例如预测$(1,2,3)$与目标$(3,5,7)$在$a=2,b=1$下完全重合；对齐误差为零没有证明原预测的数值已是正确逆米单位。对伪标签样本，原工作忽略损失最大的部分区域以降低噪声影响，并加入特征对齐以保持已有语义信息；这不保证被忽略位置都是真噪声。训练的编码器和解码器分别按设定学习率逐批更新，验证选择可用检查点，不在使用时为每幅图再求真实深度标签。伪标签学习的一般依据可参考\BookChapterRef{chap:data-derived-knowledge-acquisition}{卷八“从数据中获得知识”章}，本节的新增条件是深度约定与空间有效区域。

使用相对检查点时，将原图按模型尺寸和归一化处理，编码成空间特征，解码出$\hat q$，再把浮点图插值回原图有效区域并保留相对逆深度身份。颜色图只是显示映射，不能从色值直接读米。度量适配改用带单位的深度标签更新已有网络，使输出在该数据、最大深度及相机条件下对应实际距离；模型版本与量纲需随结果一同保留。

若三个位置的相对逆深度为$(0.8,0.4,0.2)$，可以按该约定判断第一个较近，却不能未经校准说深度为$(1.25,2.5,5)$米。用带实际单位的深度标签适配后，模型才学习特定数据条件下的度量映射；即便如此，相机和场景变化仍可能带来尺度误差。边界位置、前后排序和实际距离应分别评价。

\subsubsection{相机条件下的代价体与深度推断：MVSNet}
\label{sec:vision-recognition-mvsnet}

如果有多张重叠照片和已知相机，同一表面在不同图像中的对应可以提供深度约束。MVSNet选定参考图，对每个候选深度$d_k$，把参考像素反投影成空间点，再投到其他相机，在源图特征中插值取值。整个平面上的映射可写为依赖$d_k$的可微单应变换\scite{vision-1804-02505}

这一步可由节首投影直接算出。参考像素$\tilde{\bm p}_r$在假设深度$d_k$下为$d_k\bm K_r^{-1}\tilde{\bm p}_r$，进入源相机得到$\bm R_{s,r}d_k\bm K_r^{-1}\tilde{\bm p}_r+\bm t_{s,r}$。由于$\bm e_3^\top\bm K_r^{-1}\tilde{\bm p}_r=1$，其中$\bm e_3=(0,0,1)^\top$，除去齐次比例$d_k$便得到
\begin{equation}
 \begin{aligned}
 \tilde{\bm p}_s&\sim\bm H_{s,r}(d_k)\tilde{\bm p}_r,\\
 \bm H_{s,r}(d_k)=
 \bm K_s\left(\bm R_{s,r}
 +\frac{\bm t_{s,r}\bm e_3^\top}{d_k}\right)\bm K_r^{-1}\text{。}
 \end{aligned}
 \label{eq:vision-mvsnet-homography-process}
\end{equation}
该式只针对参考相机中$z=d_k$的平面，不能用一个深度平面代替整幅起伏场景。特征图尺度不同于图像时，先按步幅调整内参和采样坐标，再用双线性插值读取源特征；越界或后方投影不得当成有效观测。

正确深度应使同一表面的特征比较一致。将参考特征和变换后的各源特征记为$\bm{f}_r$，共有$n_{\mathrm{v}}$个视图，可以用方差形成匹配代价：
\begin{equation}
 \begin{aligned}
 \bar{\bm{f}}&=\frac1{n_{\mathrm{v}}}\sum_{r=1}^{n_{\mathrm{v}}}\bm{f}_r,\\
 \bm{c}(d_k,u,v)&=\frac1{n_{\mathrm{v}}}\sum_{r=1}^{n_{\mathrm{v}}}
   (\bm{f}_r-\bar{\bm{f}})\odot(\bm{f}_r-\bar{\bm{f}})\text{。}
 \end{aligned}
 \label{eq:vision-mvs-cost}
\end{equation}
这里$\odot$表示逐元素相乘。所有像素和候选深度的代价共同形成\term{代价体}（cost volume），它保留“每个位置在各个深度假设下多不一致”的信息。三维卷积在空间与深度轴上整理代价，经softmax得到深度权重$\pi_k$，用$\hat d=\sum_k\pi_kd_k$回归深度，并可作后续细化。训练只对具有有效参考深度的位置计算误差。

同一像素的教学标量特征在正确假设下为$(1,1.1,0.9)$，均值1、方差约0.0067；错误假设读到$(1,2,3)$，方差为$2/3$。真实网络在多通道、多位置及深度轴上整理代价，不能仅靠这一个方差保证深度正确。若候选深度为2、4、6米，权重为0.1、0.7、0.2，期望回归得到4.2米，不必等于某个离散候选。双峰权重也会把结果放在两个表面之间，因而概率集中度和有效区域需要另行检查。

训练记录含参考图、若干重叠源图、已知相机及有效参考深度。先确定深度范围和候选间隔，形成特征代价体和预测深度，在有效像素上以$L_1$比较初步及细化结果，再更新二维特征提取、三维代价整理及细化网络。使用时保留这些训练参数，但相机是当前图集提供的条件，不是从训练集复制的外参。恢复深度到参考网格后，可按置信和多视图几何一致性过滤，再反投影合并点云；每幅图看到的表面不同，直接叠加所有点仍可能留下遮挡冲突。若候选数为$k$、特征网格为$h'\times w'$、通道为$d$，仅代价体就需$O(kh'w'd)$存储。

在平行、已校正的双目简化情形中，焦距$f$、相机基线$b$和水平视差$\Delta u$满足$d=fb/\Delta u$。取$f=800$像素、$b=0.2$米、$\Delta u=40$像素，得到$d=4$米。视差误差1像素在此处约引起0.1米深度误差；同样误差在更远处会被放大。这一关系说明已知基线如何提供实际尺度。

图\ref{fig:vision-depth-hypotheses}展示两个深度假设如何查询不同源图位置。弱纹理可能让多个候选都具有相近代价，遮挡使源图看到另一表面，反光破坏外观一致性，相机误差又使正确深度也无法对齐。需要同时报告深度误差、可恢复面积与候选数带来的存储成本。

\begin{figure}[htbp]
 \centering
 % 针孔投影的几何关系；示意深度候选，不按物理长度绘制。
\begin{tikzpicture}[font=\figurefont\small,x=1cm,y=1cm,
 ray/.style={draw=BookBlue,line width=0.9pt},
 alt/.style={draw=BookAmber,dashed,line width=0.9pt}]
 \coordinate (a) at (0,0);
 \coordinate (b) at (7,0);
 \coordinate (p) at (2,3);
 \coordinate (q) at (3,4.5);
 \draw[ray] (a)--(q);
 \draw[ray] (b)--(p);
 \draw[alt] (b)--(q);
 \draw[draw=BookInk,line width=1.2pt] (-0.3,1)--(2.2,1);
 \draw[draw=BookInk,line width=1.2pt] (4.5,1)--(7.5,1);
 \fill[BookInk] (a) circle(1.1mm);
 \fill[BookInk] (b) circle(1.1mm);
 \fill[BookTeal] (p) circle(1.1mm);
 \draw[BookAmber,fill=white,line width=0.9pt] (q) circle(1.1mm);
 \fill[BookBlue] (0.667,1) circle(1mm);
 \fill[BookBlue] (5.333,1) circle(1mm);
 \draw[BookAmber,fill=white,line width=0.9pt] (6.111,1) circle(1mm);
 \node[below] at (a) {参考相机};
 \node[below] at (b) {源相机};
 \node[above right] at (0.667,1) {同一参考像素};
 \node[below] at (5.333,0.9) {$u'_1$};
 \node[below] at (6.111,0.9) {$u'_2$};
 \node[left] at (p) {深度$d_1$};
 \node[right] at (q) {深度$d_2$};
 \node[right] at (7.5,1) {源图平面};
 \draw[draw=BookMuted,<->,line width=0.7pt] (0,-0.8)--(7,-0.8)
   node[midway,below] {已知相机相对位置};
\end{tikzpicture}

 \caption{已知相机下，参考像素沿射线形成不同深度假设，分别投到源图的不同位置。图为针孔几何教学示意，只表示对应关系，不按物理长度绘制。跨图比较可排除不一致的假设，但遮挡或重复纹理会使判断不唯一。}
 \label{fig:vision-depth-hypotheses}
\end{figure}

\subsection{相机、场景几何与空间语义的恢复}
\label{sec:vision-recognition-geometry}

前一方法把相机作为已知条件。若相机未知，需要共同解释视图关系、表面点和投影。经典流程先寻找对应，再估计相机和三维点；学习式方法可以直接预测这些几何量，仍需处理共同坐标与尺度。

\Needspace{65mm}
\subsubsection{成对点图与全局对齐：DUSt3R}
\label{sec:vision-recognition-dust3r}

DUSt3R输入图像对，以两侧编码及跨图解码交换信息，为两幅图的像素分别输出三维\term{点图}（point map）：每个有效像素对应一个三维向量。两张点图都用第一幅相机的坐标表达，并预测每点置信度，使无需预先给出相机的成对几何成为直接输出\scite{vision-2312-14132}

训练记录含同一场景的图像对，以及由深度、内参和外参得到的真实点图与有效位置。两幅图共享视觉编码器，两个解码分支反复交换跨图信息，再分别输出各图像素的点图；第二张点图也变换到第一相机。编码解码权重可以从CroCo跨视图预训练初始化，输出头由三维监督学习，不把“使用时不输入相机”误写成“训练不曾使用几何标签”。

原工作分别用有效三维点到原点的平均距离规范化预测和目标，得到尺度无关点误差$e_i$。置信目标为$\sum_i(c_i e_i-\alpha\log c_i)$，其中$\alpha>0$，原实现以$c_i=1+\exp h_i>1$产生置信度。对固定$e_i>0$，该项导数为$e_i-\alpha/c_i$，较大误差倾向于较小置信，但受$c_i>1$限制；对数项使提高置信也需承担误差代价。以$\alpha=1$、教学误差$e=0.1$为例，无约束驻点为$c=10$；$e=2$的驻点0.5不在允许范围，便倾向于下边界。这是可靠性权重的训练关系，不是校准正确率。所有有效点误差反向更新网络，无效天空或透射区域按标签条件排除。

三张建筑照片形成多对输入后，每一对以自己的首图为参考，不能直接把全部点拼成场景。全局对齐优化各对点图到共同坐标的旋转、平移和尺度，使共享视图及可信点彼此相容。相机、深度与对应再从已对齐的几何联系恢复；一次成对前馈与多图优化是两个不同计算阶段。

先将有足够重叠的图对组成连接图，边$e$保存两张在同一对坐标中的点图$\bm X^{v,e}$及置信。对每幅视图求共同点图$\bm\chi^v$，对每条边求尺度$s_e>0$与刚体变换$(\bm R_e,\bm t_e)$，使$\sum_{e,v,p}c_p^{v,e}\lVert\bm\chi_p^v-s_e(\bm R_e\bm X_p^{v,e}+\bm t_e)\rVert_2$较小。相同边上的两图共用一个变换；对重复视图的每个像素共享同一$\bm\chi$，才把不同边联系起来。固定公共参照及尺度规范防止所有点任意缩到零。由初步连接初始化这些场景变量，按残差迭代优化至预设步数或改变量足够小，网络权重不随该场景对齐重新训练。

相机恢复再检查点图是否服从针孔关系。例如假设主点在中心、像素轴正交且焦距相同，可由$(u-c_x,v-c_y)\approx f(x/z,y/z)$对高置信点拟合$f$；对应三维点与二维像素可进一步估计相机姿态。将共同点图改写为相机、深度的反投影，还可在对齐目标中同时优化这些量。点图头本身没有硬性保证每个点满足同一个相机，所以相机恢复与跨图核验不能省略；连接图不连通时，也没有依据决定两个子场景的相对位置。

若甲乙图预测建筑宽度为2个单位，乙丙图为5个单位，需要先估计约2.5的相对尺度并对齐朝向位置。单对看似合理不保证多对闭环一致；弱重叠、动态行人和错误高置信点都会干扰对齐。没有额外度量依据时，共同场景仍只具有约定尺度。

\subsubsection{多视图共享特征与联合几何输出：VGGT}
\label{sec:vision-recognition-vggt}

逐对预测再优化可能重复处理同一视图。VGGT让多幅图的词元共同参与计算，交替进行帧内注意力与跨全部视图的全局注意力。帧内计算保留各图的结构，全局计算交换视角证据，随后由不同预测头输出相机、深度和点图；点跟踪头还接收参考图中的二维查询点\scite{vision-2503-11651}

训练使用对应几何量监督各预测头，把相关任务接在共享表示上。输出以首相机为参考，并沿用训练的尺度规范；这使多视图结果可在同一约定下比较，却没有自动赋予米单位。相机头给出的内外参和深度头的输出可以反投影成点，再与直接点图比较，从而检查联合输出是否几何一致。

一条训练记录是同一场景的多张RGB、每图相机、有效深度和共同坐标点图，跟踪监督还包含查询点的跨图位置及可见性。这些几何来自传感器、合成引擎或运动恢复结构，而不是只有图像类别。训练先把全部量变到首相机坐标，再以有效三维点到原点的平均距离$s$除相机平移、点坐标和深度。$s$是该场景的规范化长度，不作为使用阶段已知的真实米数。与DUSt3R在目标比较时再归一化预测不同，VGGT要求网络直接学习这套规范输出。

编码器用已有DINOv2视觉权重初始化，将每图变成带位置的图像块词元，并附一个相机词元和若干辅助寄存词元；后者是网络中暂存汇总信息的计算位置，不是实际相机或空间点。首图使用与其他图不同的学习词元标记参照。帧内层让图像词元与自己的相机词元交换信息，全局层再跨全部视图读取，因此相机词元始终与所属图像关联。最终相机头读取相机词元，稠密头将多层图像特征恢复为空间网格，分别产生深度、点图、置信及跟踪特征。

相机头的九维编码包含四维旋转四元数、三维平移和横纵视场角。四元数是表示旋转的四个归一化数，解码成满足正交条件的$\bm R_i$；视场角$\varphi_x,\varphi_y$在针孔与中心主点条件下给出$f_x=w/[2\tan(\varphi_x/2)]$、$f_y=h/[2\tan(\varphi_y/2)]$。因此得到的是当前网络图像网格的内参，而非不依赖缩放的固定焦距。深度$D_i$仍是第$i$相机自身光轴深度；只有点图和相机外参共同用首相机公共参照。恢复公共点可计算
\begin{equation}
 \bm{P}_{1,i}(u,v)=
 \bm R_i^\top\left[
 D_i(u,v)\bm K_i^{-1}(u,v,1)^\top-\bm t_i\right]\text{。}
 \label{eq:vision-vggt-unprojection-process}
\end{equation}
这里$\bm R_i,\bm t_i$沿用世界到第$i$相机约定，公共世界就是首相机。先将相机、像素和深度放在同一处理网格，才能逆投影；若改回原图坐标，也须相应改回$\bm K_i$。

训练以Huber误差比较相机编码，以置信加权位置误差和空间梯度误差监督有效深度与点图，以点位置误差和二元可见性目标监督跨图查询。Huber项在小残差区使用平方、大残差区使用线性代价，减轻极大相机残差的单项贡献，但不是自动排除错误标签。原模型联合更新视觉主干、交替注意力和任务头，按场景采样2至24幅图，图像与几何网格共同缩放；主干权重来自预训练，几何输出仍需这一步联合监督。原文使用有限迭代预算与AdamW，没有给任意图集恢复正确几何的保证。

使用三张建筑图时，加载既有检查点，按相同图像处理保留缩放和图号，一次联合前向得到每图相机和稠密量。可用式\eqref{eq:vision-vggt-unprojection-process}从深度构造公共点，也可读取直接点图；两条输出不应未经核验混拼。教学中首相机为单位旋转、零平移，主点$(320,240)$、焦距800、像素$(520,240)$与深度4个规范单位恢复点$(1,0,4)$，仍不能将4个单位直接写成4米。按置信和可见条件保留点，投回其他相机检查像素和深度是否相容。

跟踪接口需先把原图查询点变到当前处理网格，在查询图特征中插值取其描述，再与各图跟踪特征相关，估计跨图二维位置和可见性。输出位置逆变换回各自原图，并保留查询编号、源图及不可见状态。它可用于无时序图集中的对应，不是未给查询时自动生成一份车辆身份轨迹。训练中随机点的监督与使用时指定的窗角查询对应同一种位置输出，新增动态运动条件则由视频章解释。

多图前馈输出之后仍可进行几何优化，但应把这部分额外时间单独记录。图像数和每图词元数共同增加全局注意力成本，动态对象、超出训练范围的相机以及缺少重叠的图像也会削弱预测依据。新视角渲染将在\ref{sec:vision-generation-novel-view}节使用这些几何条件，生成外观与恢复可测量表面仍是不同目标。

\subsubsection{大规模离线几何与测试时场景压缩}
\label{sec:vision-recognition-vgg-t3}

建筑的视图从三张增到数百张，完整图集已可用，但同时交换全部词元变得昂贵。若每图有$p$个词元、共有$n_{\mathrm v}$幅图，则全局注意力比较$(n_{\mathrm v}p)^2$个查询键组合，键和值本身占$O(n_{\mathrm v}pd)$空间。高效注意力实现可以避免完整保存得分矩阵，却没有消除这些组合的计算。离线条件允许共同利用全部视图，需要改变的是场景信息的读取方式。

VGG-T$^3$把每个全局层中长度随视图增加的键值集合压缩到一个固定大小的多层感知机（multi-layer perceptron, MLP）。这里MLP将一个键向量映射成对应的值向量，用其参数保存当前场景的映射；它不是一个固定长度的三维点清单。原方法以VGGT权重初始化，保留图像编码、帧内计算和预测头，适配全局层的投影及新增测试时训练模块；查询键改用$L_2$归一化。值中的图像块部分先在各图网格上经过短二维卷积混合邻域，得到$\bm{v}'_i$，相机和辅助寄存词元不经过这项空间卷积。邻域混合使拟合目标带有局部上下文，避免仅复现同一词元的两个线性投影\scite{vision-vgg-t3-2026}。

对一个全局层，设当前图集产生$N=n_{\mathrm v}p$个键$\bm{k}_i$、查询$\bm{q}_i$和目标值$\bm{v}'_i$，维数均为$d$。测试时训练（test-time training, TTT）指在使用当前输入时临时更新一小组场景权重$\bm{\phi}$；预训练主干并未因这个名称变成未训练网络。为使键读出相应值，内层目标及查询计算可写为
\begin{equation}
 \begin{aligned}
 \bm{\phi}_*&\approx\argmin_{\bm{\phi}}
   \frac1N\sum_{i=1}^{N}
   \lVert T_{\bm{\phi}}(\bm{k}_i)-\bm{v}'_i\rVert_2^2,\\
 \bm{o}_i&=T_{\bm{\phi}_*}(\bm{q}_i)\text{。}
 \end{aligned}
 \label{eq:vision-vgg-t3-fast-weights}
\end{equation}
近似号强调实际只做规定次数的更新，不求一个已证明存在的精确最优映射。每层先从训练得到的初始化取得快速权重，按键值残差计算内层梯度，再将查询送入更新后的MLP；所得词元进入下一层，全部层完成后才由相机、深度和点图头解码。原文内层采用Muon优化，训练内层用一步，通常测试用两步；这属于所报告配置，不是任意场景的收敛结论。

平方残差目标还允许分批拟合全集。固定一次更新开始的$\bm{\phi}$，把词元划成批次$B_s$，则
\[
 \nabla_{\bm{\phi}}L=
 \frac1N\sum_s\sum_{i\in B_s}
 \nabla_{\bm{\phi}}
 \lVert T_{\bm{\phi}}(\bm{k}_i)-\bm{v}'_i\rVert_2^2\text{。}
\]
这一等式来自有限和的微分线性性，要求各批都在同一组权重上算梯度，累加后再更新；逐批改权重不是同一次全集梯度。可把未处理批放在主机内存，或跨设备分摊，降低单卡同时保存图集特征的需求。拟合和读取在固定MLP容量、分辨率及更新步数下随词元数线性增长，原图、编码特征、输出点图和数据传输仍随视图增多。

用标量线性模型$T_\phi(k)=\phi k$作教学简化，两个键值样本为$(1,2)$、$(2,4)$。在$\phi=0$处，平均平方目标的梯度为$-10$，取教学步长0.1更新到$\phi=1$，再更新一次到1.5；查询1.5读出2.25，还没有达到目标映射给出的3。真实方法用非线性MLP和不同优化器，这个例子仅说明“有限次拟合”与“精确保存全部键值”有差别。新视图若作为定位查询，可冻结已有各层快速权重只做读取，但原文相机头还保留建图图像的相机词元，不能把全部定位存储都说成与图集无关。

比较同一建筑逐步增加的视图，应同时记录内层步数、点图误差、相机误差、重投影、宽基线失败、总显存及总时间。原文在其基准中还指出相机预测较不精确，并采用直接点图头而非相机深度反投影评估点图。固定容量映射产生近似误差，不能由资源增长方式推断几何精度保持不变。这里完整图集先到齐；下一小节改变的是信息到达条件。

\subsubsection{流式几何状态、地图记忆与长期漂移}
\label{sec:vision-recognition-infinitevggt}

沿走廊前进时，第$t$幅图已经到达，返回走廊的画面却尚未发生。当前输出不能借未来视图纠正自身。\emph{因果几何预测}（causal geometry prediction）限制当前帧只读取自身与已到达历史。直接保留每层的历史键值，可以复用旧特征，但缓存随$t$增长；只保留最近几帧，又可能删去初始参照与回访线索。

InfiniteVGGT在已有因果几何模型上采用免训练的滚动缓存策略。图像仍经帧内和因果跨帧计算，输出相机、深度、点图及跟踪特征；变化在于哪些历史键值继续可读，不是为当前流重新拟合一个VGG-T$^3$式MLP。各层、各注意力头把第一帧的键值固定为锚点，因为几何输出以首相机为公共参照；其余历史与新到词元构成可裁减候选集\scite{vision-infinitevggt-2026}。

相邻画面含大量重复特征，原方法用归一化键的方向差异选择保留内容，避免读取完整注意力矩阵。对一个层和头，令候选非零键的单位向量为$\hat{\bm{k}}_i$，平均向量为$\bm{\mu}=m^{-1}\sum_i\hat{\bm{k}}_i$，多样性评分为
\begin{equation}
 s_i=-\frac{\hat{\bm{k}}_i^\top\bm{\mu}}
                 {\lVert\bm{\mu}\rVert_2}\text{。}
 \label{eq:vision-infinitevggt-diversity}
\end{equation}
与平均方向差别较大的键取得较大分数，并与自己的值一起保留。这个评分近似识别特征冗余，不是直接测量三维点的真实信息量；平均向量为零时余弦无定义，复现需有数值稳定和退化处理，不能从理想公式推出每组输入都可直接计算。

不同层的候选多样性不同。原文将层平均评分$\bar s_\ell$经softmax变成预算比例$b_\ell=\exp(\bar s_\ell/\tau)/\sum_j\exp(\bar s_j/\tau)$，用总可变预算$B$分配各层保留数，再按分数选前若干键值，最后加回不可裁减的首帧锚点。实现需把预算变为整数并核对总量；预算$B$、温度$\tau>0$和锚点成本均须记录。追加、评分、裁减、保留是每次到帧后的状态更新，权重本身不因裁减而接受梯度训练。

教学中令可变候选键方向为$(1,0)$、$(1,0)$、$(0,1)$，平均为$(2/3,1/3)$。前两者与平均方向余弦为$2/\sqrt5\approx0.894$，第三者为$1/\sqrt5\approx0.447$，第三者的负余弦评分更高；预算只容两项时，它和一个重复方向可以留下。示例说明了为何裁减可以保留不同方向，却没有证明被删方向在之后回访时不重要。

固定缓存控制的是工作状态，不是全部地图输出的永久存储。每帧仍产生点图，若全部写盘，总输出量仍随时间增加；当前帧到缓存的读取及评分也占计算。走廊返回起点时，应比较短程相机误差、累计轨迹漂移、回访位置重合与缓存成本。保留首帧和部分旧键有助于维持参照，却不是显式回环检测与全图校正。Long3D约万帧的有限序列可提供长程测试，不能证明无限时间稳定；地图参照与视频中的对象身份也不是同一个状态。

\subsubsection{三维几何与开放词汇语义的共同恢复}
\label{sec:vision-recognition-ff3r}

房间两幅照片都标出“椅子”，还不能知道它们是否指向同一张椅子的同一表面。二维语义可借投影进入三维，但几何错位会把椅面特征放到墙上；直接合并近邻三维点又可能把相邻椅子和背景混在一起。共同恢复需要分别保留位置、外观和可查询语义，并检查三者指向相同表面。

FF3R输入至少两幅多视图图像，沿用VGGT式几何词元，同时以冻结的CLIP-LSeg编码器取得逐位置语义特征。末层几何词元作为查询读取语义词元的键和值，使几何解码获得语义上下文。解码器输出深度和相机，并预测像素对齐的三维高斯：其中心由深度反投影确定，协方差描述空间扩展，颜色参数描述外观，不透明度描述遮挡贡献，另附语义向量$\bm{f}_i$。高斯是可投影、可聚合的场景表示，不应把其中心直接当作精确物体表面\scite{vision-ff3r-2026}。

为用图像监督这些三维属性，需要把它们投回观察相机。每个高斯中心用节首投影得到二维位置，空间协方差经投影的局部线性近似成为像平面的扩展区域；按深度先后计算每项对像素的贡献。若第$i$项在该像素的覆盖与不透明度合成系数为$a_i\in[0,1]$，前方透过比例为$T_i=\prod_{j<i}(1-a_j)$，则颜色和特征分别按$\sum_iT_i a_i\bm c_i$与$\sum_iT_i a_i\bm f_i$汇总，并按规定处理未覆盖背景。两者共用贡献权重，因此特征的监督也会反向影响所预测几何；完整新视角渲染在生成章展开，此处只用投回观测图的残差解释训练依据。

训练先从同一场景采样上下文视图，保留图号并采用相同裁剪尺度，初始化几何主干和深度头于已有VGGT权重，新增融合与输出部分随机初始化；几何交替注意力参与适配，CLIP-LSeg保持冻结。将预测三维表示投回观测视图，比较渲染RGB与真实RGB，以及渲染特征与CLIP-LSeg特征。原论文还明确使用VGGT给出的相机、深度伪标签作为蒸馏约束，深度只在高置信位置比较。因此“无需人工相机、深度和语义标注”不等于“只有RGB残差”，也不等于没有预训练监督。

若两视图看到同一表面，它们的语义特征还应在几何对应下相容。使用式\eqref{eq:vision-relative-camera-measurement}，将目标图像像素$\tilde{\bm{p}}_t$在预测深度$D_t$下反投影，再投到上下文图$c$：
\begin{equation}
 \bm{p}_c=
 \Pi\!\left[
 \bm{K}_c\left(
 \bm{R}_{c,t}D_t(\bm{p}_t)\bm{K}_t^{-1}\tilde{\bm{p}}_t
 +\bm{t}_{c,t}\right)\right]\text{，}
 \label{eq:vision-ff3r-warp}
\end{equation}
其中$\Pi$用第三坐标除前两坐标。按$\bm{p}_c$插值采样源特征，计算$1-\cos(\bm{f}_t,\bm{f}_{c\to t})$，只在投影落入图像且深度相容的位置求平均，再交换两图求另一方向。边界和遮挡筛选让“不曾看见同一表面”的位置不被强行拉到一起；特征不为零是余弦计算条件。

视图密集时，可按三维位置把高斯分到有限体素，即固定大小的立方区域。FF3R先求同一体素的平均语义向量$\bar{\bm{f}}_s$，以余弦距离$d_i^{\mathrm{sem}}$度量每项偏离，再把几何置信$C_i$与语义一致性共同用于聚合：
\begin{equation}
 \begin{aligned}
 w_{i\to s}=
 \frac{\exp(C_i-\lambda d_i^{\mathrm{sem}})}
 {\sum_{j:\,v(j)=s}\exp(C_j-\lambda d_j^{\mathrm{sem}})},\\
 \bar{\bm{a}}_s=\sum_{i:\,v(i)=s}w_{i\to s}\bm{a}_i\text{。}
 \end{aligned}
 \label{eq:vision-ff3r-voxel}
\end{equation}
$v(i)$是体素编号，$\bm{a}_i$可以是位置、颜色或语义属性，$\lambda$控制语义距离权重。权重和为1，所以它是体素内的加权汇总，不是把全部属性相加。若教学中两项置信均为1、语义距离分别为0与1、$\lambda=1$，权重为$e/(e+1)\approx0.731$与0.269；偏离项贡献下降但仍未消失。跨物体体素过大或平均语义本身错误，汇总仍会污染边界；平均语义向量为零时还需余弦的数值退化处理。

完整目标包含RGB像素与感知误差、特征余弦误差、双向几何特征误差以及深度和相机蒸馏项。按采样图集计算这些可微残差，更新可训练部分，训练完成后保留融合和解码参数；新房间使用前馈预测获得几何、外观与语义表示。查询“椅子”时，把文字编码为相同语义空间向量，与三维或投影后的特征比较，再按规定阈值和区域规则恢复结果。仅颜色渲染自然，不能证明椅子语义或尺度正确。

同一椅面的两次观察应落到共同三维位置并具有相容语义，未见背面则仍只有先验预测。检查可见表面的深度、跨视角语义归属、遮挡筛选和体素边界，才能判断查询结果是否真的指向同一对象。作者项目页标注CVPR 2026 Findings，这一身份不同于把预印本泛称为CVPR主会论文。新视角外观在\ref{sec:vision-generation-novel-view}节展开，未观测内容的创作归\ref{chap:vision-multimodal-generation}章。

\section{图像匹配与检索}
\label{sec:vision-recognition-retrieval}

识别已经把图像变成类别、范围和测量。检索还需要依据查询，从图库找出相关图像并排序。紧凑描述适合筛选候选，局部对应保留细节，几何约束再检查对应是否能共同解释视角变化。

\subsection{查询与相关性的确定}
\label{sec:vision-recognition-relevance}

“找类似的红色汽车”和“找同一辆车的其他照片”使用不同正例。前者容许不同实例，只要外观或语义相近；后者需要车身细节、车牌等区分身份的证据，并可能允许颜色因光照变化。若没有预先规定相关性，高相似度排序甚至可能朝错误目标优化。

查询可以是整图或裁剪区域。整幅照片包含建筑、树木和天空，局部查询只指定一块招牌；两者对应的容许背景、遮挡和视角变化不同。图库中的近重复图也可能使任务过于容易。评价前应固定查询范围、正例定义、忽略项、图库版本和干扰图规模，文字查询则由\ref{chap:vision-multimodal-understanding}章继续连接。

\subsection{候选图像的形成}
\label{sec:vision-recognition-candidates}

在大图库中逐一做精细匹配成本很高。全局向量先快速保留较少候选，局部特征再进入后续核验。近似索引可以加速向量搜索，但它引入的召回损失与表示本身的区分能力应分开检查。

\subsubsection{全局特征聚合与检索学习：GeM}
\label{sec:vision-recognition-gem}

空间特征需要汇总为固定维向量。平均池化平等计入各位置，最大池化只保留最强响应。\term{广义均值池化}（generalized-mean pooling, GeM）在两者之间调节：对第$k$通道的非负特征$x_{p,k}$，计算
\begin{equation}
 g_k=\left(\frac1n\sum_{p=1}^{n}x_{p,k}^{\,q_k}\right)^{1/q_k},
 \qquad q_k>0\text{。}
 \label{eq:vision-gem}
\end{equation}
指数$q_k=1$为平均，趋于无穷时趋近最大值。实现通常在幂运算前给特征加下界，避免零点数值问题；指数可以固定，也可与编码器参数一同学习\scite{vision-1711-02512}

对响应$(1,1,1,4)$，平均为1.75，$q=3$时为$(67/4)^{1/3}\approx2.559$，最大值为4。提高指数让局部强响应更影响全图描述，但并不保证该响应一定属于目标。聚合后将$\bm{g}$归一化为单位向量，以欧氏距离或内积排序；单位向量满足$\lVert\bm{u}-\bm{v}\rVert_2^2=2-2\bm{u}^{\top}\bm{v}$，两种排序因此等价。

检索训练需要知道哪些图拍到了同一实例。原工作利用三维重建与相机几何选择具有重叠内容的正例和困难负例，故“不需人工检索标注”并不等于没有监督依据。对比目标使正图对接近，对距离小于间隔$m$的负图对施加惩罚；训练后保留编码器与池化参数，为图库和查询使用相同的特征处理。

一条训练图对$(\bm X,\bm X')$的同实例标签为$y\in\{0,1\}$，两图经过共享编码、GeM和单位归一化后，距离为$d=\lVert\bar{\bm g}-\bar{\bm g}'\rVert_2$。对比目标的一种标准写法为
\begin{equation}
 \ell=\tfrac12 y d^2+
 \tfrac12(1-y)[\max(0,m-d)]^2\text{。}
 \label{eq:vision-gem-contrastive-process}
\end{equation}
正对的距离导数为$d$，下降更新倾向于缩短距离；负对在$d<m$时导数为$d-m<0$，倾向于增大距离，超过间隔则不给这项梯度。教学中$m=1$，正对距离0.4的目标为0.08，负对距离0.4的目标为0.18，负对距离1.2为零。平方距离是训练约束，不是与实际世界中建筑距离同量纲的测量。

原工作由重建得到的共同三维轨迹及相机共视挑选正对，从不同实例中挑选当前描述相近的困难负对。以已有分类编码器初始化，训练共同更新视觉特征与可学习池化指数，重新挖掘困难负例时使用当前参数，而测试图库不能参与形成监督。使用时为每图保存归一化描述与图号，对查询同样处理，按内积取前若干候选。指数、可能的描述变换和归一化须一致；只编码图库而忘记查询的同样变换，会改变距离定义。裁剪查询需保留裁剪来源，后续局部验证仍针对该查询范围。

同一建筑的两张照片可能因视角差异较远，另一栋同样玻璃幕墙的建筑却可能很近。全局描述提供候选，不能证明局部细节一致。已有DINOv3向量可以作为不同参数来源的筛选基线，但比较应固定输入、归一化、图库与训练信息。

\Needspace{55mm}
\subsubsection{局部关键点与描述：SuperPoint}
\label{sec:vision-recognition-superpoint}

若希望保留窗角或招牌角等可重复的位置，可以为每幅图输出稀疏点坐标和描述向量。SuperPoint使用共享编码器连接兴趣点分支与描述分支；兴趣点没有“左腕”那样的对象部位语义，而是适合跨图重复辨认的位置\scite{vision-1712-07629}

原始检测头在较粗网格上预测每个$8\times8$区域内的64个位置及一个“没有点”的通道。经过通道归一化、去掉无点通道和空间重排，得到原图分辨率的点分数，再阈值筛选并抑制过近峰值。描述头产生较稠密的低分辨率特征，在保留点位置插值采样，并将描述归一化。最终保存$(x,y,\text{得分},\bm{d})$，而不仅是一张响应图。

训练先以具有已知角点的合成几何图形学习检测，再对真实图像施加多次已知单应变换，将检测结果逆变换并汇总，形成更稳定的伪标签。已知变换也提供两图描述应对应的位置，用于联合训练点与描述。单应适应提供的是这些变换下的监督，不能覆盖任意三维视差。

对一个粗网格单元，64个点位置与无点通道组成65类，标签按变换后的可见兴趣点生成，交叉熵使点分数回到正确子位置。描述监督由已知单应把一个位置映到另一图，保留落在有效范围的对应；对应描述需内积较高，其他位置按间隔降低相似度。多次单应适应是产生真实图伪标签的步骤，联合训练再用这些伪标签及变换约束更新共享编码与两个分支。它与使用阶段在一对新照片上找对应不同，没有从新图真实相机取得标签。

得到点分数图后，按阈值与局部峰值抑制保留点，并可按分数限制最多点数；描述在这些点上插值并单位归一化。若输入缩放为原图的一半，网络点$(100,60)$需还原为原图$(200,120)$，描述向量本身则不乘二。保存点坐标所在网格和图号，后续匹配器才能用同一坐标系做几何验证；没有点时应返回空特征集合，不虚造一对对应。

窗格重复时，检测器可以在许多窗角上都稳定响应，但描述仍可能相似。因此“点能重复检测”与“点能匹配到正确窗格”是两个要求。尺度、旋转、遮挡及保留点数都会影响后续候选，比较匹配器时需固定这些输入条件。

\subsection{精细匹配与结果排序}
\label{sec:vision-recognition-verification}

候选图像中可能混入外观相似的其他实例。局部描述取得候选点对，上下文比较帮助消歧，几何一致性再排除不能共同成立的对应。几何通过也仍需回到查询定义判断相关性。

\Needspace{55mm}
\subsubsection{上下文局部配对与自适应计算：LightGlue}
\label{sec:vision-recognition-lightglue}

独立最近邻匹配只比较两个描述，无法利用邻近窗格怎样共同移动的信息。LightGlue接收已有关键点和描述，交替通过图内自注意力与跨图信息交换更新表示，再预测匹配关系；它不负责从像素检测关键点\scite{vision-2306-13643}

设两图分别有$n,m$个点，将更新后的描述投影为$d$维，再以$1/\sqrt d$缩放两两内积，得到$S_{i,j}$。两侧点可找到对应的分数为$\mu_i,\nu_j\in[0,1]$，匹配得分结合双向归一化与可匹配性：
\begin{equation}
 P_{i,j}=\mu_i\nu_j
 \frac{\exp S_{i,j}}{\sum_{j'=1}^{m}\exp S_{i,j'}}
 \frac{\exp S_{i,j}}{\sum_{i'=1}^{n}\exp S_{i',j}}\text{。}
 \label{eq:vision-lightglue-score}
\end{equation}
第一项归一化要求$i$在对面候选中偏好$j$，第二项要求$j$在本侧候选中偏好$i$。保留互为最高分且超过阈值的点对，有助于避免多个点争抢同一对应。训练使用真实对应监督匹配得分，对没有对应的点监督可匹配性。

一条匹配训练记录含两图的点坐标、描述和真实点对，以及哪些点在另一图不可见。点对可以由有深度、相机的图像重投影获得，也可以由已知图像变换取得；遮挡和越界必须排除，否则“找不到对应”的点会被错误配对。输入先按图像尺寸规范化关键点坐标，为点加入位置表示，再更新描述。若真实匹配集合为$\mathcal M$，两侧无对应点集为$U,V$，一个说明职责的负对数目标为
\begin{equation}
 \begin{aligned}
 L_{\mathrm{match}}={}&-\frac1{|\mathcal M|}
 \sum_{(i,j)\in\mathcal M}\log P_{i,j}\\
 &-\frac1{|U|}\sum_{i\in U}\log(1-\mu_i)
 -\frac1{|V|}\sum_{j\in V}\log(1-\nu_j)\text{。}
 \end{aligned}
 \label{eq:vision-lightglue-training-process}
\end{equation}
空集合对应项应跳过或用规定的非零归一化，不直接除零。匹配项同时监督双方向选择与可匹配性，无对应项监督放弃错误配对；它不是把全部$n m$点对都当同等负例。匹配器训练更新位置编码、信息交换和评分分支，关键点描述来源需固定，不能把为一种特征训练的权重无条件用于另一种描述。

若$\mu_i=0.9,\nu_j=0.8$，两个方向的归一化分数为0.7和0.6，则$P_{i,j}=0.3024$。只有单方向相似度高，或其中一点已被遮挡时，最终得分都可能降低。这是模型的配对评分，不是对几何正确性的独立证明。

不同图像对的难度差异很大。LightGlue在各层预测当前配对是否已稳定，足够比例的点达到置信要求时提前退出；对于已经高置信、但可匹配分数低的点，停止后续计算。尚不能判断的点仍应保留，不能把“不确定”直接当成“无对应”。这些门控调节处理层数和保留点数，比较速度时需固定特征、图像尺寸、硬件与难度分布。

层间置信分支以当前结果是否同最终层的决定相容为训练依据，并限制其梯度路径，避免仅为早退判断任意改动描述。使用时每层完成更新后检查达到置信要求的点比例，未达到则继续；被裁点保存原输入索引，最终对应需映射回完整点表。对保留点计算$P$，取行最大与列最大互相指向且得分过阈值的点对，再输出两幅原图的坐标及匹配得分。若任一图无点，应返回空对应；不能为了让后续几何估计运行而强行填点。置信早退和低可匹配裁减分别调节层数与点数，不改变后续还要做几何核验的职责。

\subsubsection{鲁棒几何估计与候选重排：RANSAC路线}
\label{sec:vision-recognition-ransac}

重复纹理可能产生看似可信但错误的点对，需要检查一组对应能否共同解释视角变化。平面上的点在两图间可由\term{单应变换}（homography）联系，齐次坐标满足$\tilde{\bm{p}}'\sim\bm{H}\tilde{\bm{p}}$。一般三维场景的两视图则遵守\term{极线约束}（epipolar constraint），即$\tilde{\bm{p}}'^\top\bm{F}\tilde{\bm{p}}=0$，其中$\bm{F}$为基础矩阵。后者约束另一图中的一条线，而非单个位置。

\term{随机采样一致性}（random sample consensus, RANSAC）通过反复取少量对应估计模型，寻找被许多点共同支持的解。以单应为例，四个处于非退化配置的点对可以产生候选$\bm{H}$；再把所有对应按重投影残差与阈值分为内点和外点，保存支持较好的模型，并利用内点重估。

几何模型的来源是当前点对，不需训练分类头。对平面对应，消去齐次比例可写$\tilde{\bm p}'\times\bm H\tilde{\bm p}=\bm0$，每对提供两条独立线性约束，四对约束八个尺度以外的自由度；共线、重复或近退化点会破坏估计。将输入坐标平移缩放到较合适的数值范围后求候选，再变回原图$\bm H$，才可用像素残差比较。对候选的对称转移误差可用
\[
 e_i^2=\lVert\bm p_i'-\Pi(\bm H\tilde{\bm p}_i)\rVert_2^2
       +\lVert\bm p_i-\Pi(\bm H^{-1}\tilde{\bm p}_i')\rVert_2^2\text{。}
\]
这里要求$\bm H$可逆且两次投影有限，阈值按像素或上述平方像素明确约定。一般三维两视图则不能用单应的点到点误差替代极线误差；$\bm F$只约束对应点靠近极线，估计时还需满足秩二条件。

若每次独立抽取$s$个对应，内点比例近似为$w$，全内点样本概率约为$w^s$。为了至少以概率$p$抽到一次全内点样本，次数满足
\begin{equation}
 k\geq\frac{\log(1-p)}{\log(1-w^s)}\text{。}
 \label{eq:vision-ransac-iterations}
\end{equation}
这是独立、近似固定内点率下的采样计算，还没有保证样本不退化，也没有保证估计模型正确。取$w=0.5,s=4,p=0.99$，需要至少72次；内点率降低时，所需次数迅速增加。

次数公式来自“$k$次都没抽到全内点”的概率$(1-w^s)^k$，要求它不超过$1-p$后取对数；分母为负数，除法时不等号方向随之改变。$w=0$时不存在有限次数保证，$w=1$时理想采样一步足够，实际仍需排除退化。实现从一个未接受模型状态开始，反复抽样、估计、计内点，若支持更多或在同支持下残差更小就替换当前最佳；到预算或按更新内点率估计的停止次数结束，再用全部内点重估。使用内点率提前停止依赖估计可信，错误高支持模型也可能过早结束。

Graph-Cut RANSAC在局部优化中进一步利用邻近对应的联系。为每个对应设置内外点标签，一元代价来自对当前模型的残差，邻接项倾向于让邻近对应具有一致标签；图割在规定的能量形式下求标记，再据选出的内点重估几何\scite{vision-1706-00984}
图割解决的是给定模型附近的一次局部标记问题，不是对所有模型、样本和标签组合取得全局最优。邻近错误匹配也可能具有一致外观，仍须检查模型假设。

在当前候选模型下，二值标签$\ell_i$表示该对应是内点还是外点，可把局部目标概括为$\sum_iD_i(\ell_i;e_i)+\lambda\sum_{(i,j)\in E}w_{i,j}\mathbf1\{\ell_i\ne\ell_j\}$。一元项依据残差惩罚不合适标记，非负邻接项表达邻近对应倾向于共同支持一个模型；在满足二值可割能量条件时，图割求这一次标记最优。重新用内点估计几何后，一元项改变，仍需继续局部比较。检索重排可以依据内点数量、查询范围覆盖和残差，但应固定排序规则；一小块重复窗格的高内点数不能替代整幅查询的相关性。

建筑正面近似平面时，单应适用；有明显前后视差时，同一个单应不能同时对齐所有表面，应考虑两视图几何或分区结构。纯旋转、近共线点和重叠太少也会改变估计条件。图\ref{fig:vision-retrieval-evidence}将全局筛选、局部配对与几何核验分别列出，避免把一个阶段的高分误作整条证据链完成。

\begin{figure}[htbp]
 \centering
 % 流程节点间箭头均连接左右中心；下方为各阶段尚未排除的错误。
\begin{tikzpicture}[font=\figurefont\small,
 block/.style={draw=BookTeal,fill=BookTeal!8,minimum width=28mm,
   minimum height=12mm,align=center,inner sep=1.5mm},
 flow/.style={-{Stealth[length=2mm]},draw=BookInk,line width=0.9pt},
 note/.style={align=center,text width=28mm,font=\figurefont\footnotesize,
   text=BookMuted}]
 \node[block] (global) at (0,0) {全局描述\\筛选候选};
 \node[block] (local) at (3.7,0) {局部对应\\匹配细节};
 \node[block] (geo) at (7.4,0) {几何支持\\核验重排};
 \draw[flow] (global.east)--(local.west);
 \draw[flow] (local.east)--(geo.west);
 \node[note] at (0,-1.6) {可能只同类\\并非同一实例};
 \node[note] at (3.7,-1.6) {重复窗格\\可能相互错配};
 \node[note] at (7.4,-1.6) {几何匹配仍需\\满足查询定义};
 \foreach \x in {0,3.7,7.4}
   \draw[draw=BookRule,line width=0.7pt] (\x,-0.65)--(\x,-1);
\end{tikzpicture}

 \caption{图像检索中逐步保留更具体的证据。上方沿主流程减少候选，下方说明每步仍可能留下的错误。几何一致可以帮助重排，却仍要检查所得实例是否满足查询的相关性定义。}
 \label{fig:vision-retrieval-evidence}
\end{figure}

Revisited Oxford和Paris修订了实例检索标注与难度协议，并引入大规模干扰图，适合比较候选筛选和几何重排分别改变了哪些结果\scite{vision-1803-11285}
不同协议对困难正例与忽略项的处理不一样，应固定协议后报告排序。更多内点可以支持同一建筑的判断，却不能直接证明其中招牌文字相同，或建筑具有相同的实际尺度。

\section{结果检验}
\label{sec:vision-recognition-evaluation}

组合系统需要让类别、框、掩码、属性和测量指向同一对象。若车辆检测正确，属性却来自旁车，逐模块的总体准确率可能仍然较高。街景应逐实例保留标识并核验空间对应；工业零件应同时检查异常图、整图得分与规定缺陷，避免正常材质变化被业务阈值放大。

检测评价先按类别与IoU阈值对预测和真实对象作一对一匹配，未匹配预测为假阳性，未匹配目标为漏检。按得分变化形成精确率、召回率曲线，再计算平均精度；不同数据集的插值、IoU阈值、尺寸分组与每图检测数量上限必须说明，不能把任意一个“AP”视为同一指标。

区域任务也有不同分母。语义分割按类别比较像素IoU，常再对类别平均；实例分割要先匹配对象。全景质量把匹配区域的IoU与对象错误合在一起，标准任务对同类、IoU超过0.5的区域建立匹配集合$T$，并令$F_{\mathrm{p}},F_{\mathrm{n}}$为多报、漏报区域集合：
\begin{equation}
 \mathrm{PQ}=
 \frac{\sum_{(p,g)\in T}\operatorname{IoU}(p,g)}
 {|T|+\frac12|F_{\mathrm{p}}|+\frac12|F_{\mathrm{n}}|}\text{。}
 \label{eq:vision-pq}
\end{equation}
它同时受到范围与实例错误影响\scite{vision-1801-00868}
例如匹配区域IoU为0.8和0.6，另有一个漏报且无多报，则PQ为$1.4/2.5=0.56$；只报告匹配区域平均IoU会得到0.7，掩盖漏掉的区域。空类别与跨类别平均按数据集协议处理。

检索则按相关性检查排名。若有$R$个相关项，排名$k$处的相关指示为$r_k$，截至该位置的精确率为$P(k)$，全排名上的一种非插值平均精度为
\begin{equation}
 \mathrm{AP}=\frac1R\sum_k P(k)r_k\text{。}
 \label{eq:vision-retrieval-ap}
\end{equation}
三个正例出现在第1、3、5位时，AP为$(1+2/3+3/5)/3\approx0.756$。若只返回前四项，计算截断指标时仍需明示分母和遗漏规则。召回率关注是否找到，AP还关注出现得早不早；无相关项的查询应按协议排除或单独评价，不能直接除以零。

未知发现需在隐藏新类别名称的条件下计未知召回，同时记录未知被误报成已知的数量、背景假阳性，以及后续加类前后旧类的AP。教学中有5个未知物体，正确发现3个、误当已知1个、完全漏掉1个，未知召回为$3/5=0.6$；这不能由已知货车AP代替。文字定位和概念分割还需固定候选短语、正负图像示例及提示成本，使用真实框提示取得的优势不能直接归于无提示系统。

姿态比较固定关节集合、单位和坐标，以平均关节位置误差（mean per-joint position error, MPJPE）计算$(1/j)\sum_k\lVert\hat{\bm J}_k-\bm J_k\rVert_2$。根对齐只去平移，相似变换对齐还去旋转和尺度，二者回答的问题不同。人体表面另比对应顶点误差、轮廓重投影及手足局部，不能由少量关节误差证明全部网格准确。单目深度需区分相对排序、边界和米制误差；允许对齐的结果不能与未对齐度量结果混为一表。

离线场景比较固定视图集合、分辨率、首相机参照和拟合步数，分别报告相机、点图、重投影、完整度及总成本。流式场景在每个时刻只允许已到达图，记录缓存预算、局部误差与长程漂移，回访时检查同一位置是否重合；一次全集对齐后的轨迹误差会隐藏部分累计误差。三维语义需先在有效几何下找到同一表面，再比较各视角的查询归属。RGB渲染、语言相似度和几何一致各有误差来源，三者之一正确不能替其他两项通过。

\begin{table}[htbp]
 \centering
 \small
 \begin{tabularx}{\linewidth}{@{}lXX@{}}
 \toprule
 输出 & 需要固定的条件 & 不能由它单独证明的内容\\
 \midrule
 类别、属性 & 类别表、未知项、对象范围 & 对象存在及属性归属都正确\\
 框、掩码 & 匹配规则、提示、可见范围 & 未观测部分及三维形状正确\\
 数量、姿态 & 计入规则、关节与对齐方式 & 实例清单、身份及原始朝向正确\\
 深度、点图 & 尺度、坐标、有效区域 & 距离带实际单位且适合测量\\
 检索排序 & 正例、图库、忽略项及难度 & 对象属性、文字和距离正确\\
 \bottomrule
 \end{tabularx}
 \caption{不同视觉成果所需的评价条件。右列列出仍需另行核验的推论，不表示这些任务必然无法提供相应能力。}
 \label{tab:vision-recognition-protocols}
\end{table}

\section{本章小结}
\label{sec:vision-recognition-summary}

从街景得到可用结果，需要把预测接回明确的对象和观测。类别与属性规定语义，框和掩码规定范围，数量与姿态规定测量对象，相机与尺度规定三维解释，检索相关性则决定哪些对应值得保留。共享表示可以支持这些输出，监督、恢复步骤和核验关系仍各自不可省略。

本章的方法提供了两类互补依据。学习表示利用样本规律形成候选判断，空间、几何与集合约束检查这些判断能否共同成立。候选高分不等于目标存在，区域正确不等于身份正确，相对深度不等于米制距离，几何支持也不等于任意业务相关性。只有把这些条件保留到最终输出，才能知道一项结果已经回答了什么，以及仍需要哪些观测。

开放词汇是在已给概念下寻找，未知发现是在未命名时保留对象，两者不能以同一查询准确率替代。二维关节、三维骨架和人体表面逐步增加几何约束，却始终受尺度与遮挡限制。离线全集压缩、因果滚动缓存又分别保存不同历史：固定参数容量不免除拟合和图像存储，有限缓存不保证长期回访没有漂移。三维语义只有与位置和可见表面共同核验，才构成可用于后续测量、渲染或问答的结果。

\Needspace{95mm}
这些尚未满足的条件也给出了进一步研究的问题。开放概念增加目标歧义，长尾类别需要更细的覆盖分析，精细边界要求保存小尺度证据，遮挡和单目投影又限制可恢复的空间量。LVIS把大词表和少见实例纳入分割评价，MVTec AD把正常训练与测试缺陷定位分开，为这些问题提供了不同的观察条件\scite{vision-1908-03195,vision-mvtec2019}
医学和遥感等专业用途还要固定采集设备、地理区域或患者划分，核验外观先验能否迁移。研究改进应明确补充了哪种证据、减少了哪类错误，结论仍限于相应任务与划分实际检验的范围。

遮挡人体的提示纠错还可研究如何表达剩余深度歧义，而非只输出一个看似合理的网格；长期地图可研究有限记忆、显式回访校正与尺度一致如何配合；几何语义融合可研究未见表面的预测如何同已观察事实分开保存。这些问题需要固定提示、视图到达、缓存及标注条件，分别检验关节、漂移和跨视角归属，不能用单一视觉分数宣称全部问题已解决。

\chapter{图像复原、增强与压缩}
\label{chap:vision-image-restoration}

夜间拍到的招牌偏暗，字画周围有噪点，放大后又出现模糊和方块。把照片处理得明亮、锐利，并不能单凭外观判断招牌上的字已经恢复正确；把处理结果保存成较小文件，也不能只看重建图像就知道实际节省了多少存储。本章讨论受损图像怎样变得可用，并区分内容估计、呈现调整与表示成本三个目的。

设希望恢复的图像为$\bm{X}$，实际观测为$\bm{Y}$。一个便于分析的成像模型是
\begin{equation}
 \bm{Y}=\mathcal{Q}\bigl(\mathcal{S}(\bm{K}*\bm{X})+\bm{N}\bigr)\text{，}
 \label{eq:vision-image-degradation}
\end{equation}
其中$\bm{K}$描述模糊，$*$表示卷积，$\mathcal{S}$表示采样，$\bm{N}$表示噪声，$\mathcal{Q}$概括量化等后处理。它是一个简化模型：传感器噪声可能发生在颜色形成之前，图像也可能经历多次缩放和压缩，不能把这些操作一律交换顺序。RGB图像通常已经经过颜色插值、白平衡和色调映射；RAW保留更接近传感器读出的数据，两者可供反演的信息不同。

\term{图像复原}（image restoration）依据受损观测估计目标内容；\term{图像增强}（image enhancement）调整亮度、对比或颜色，使呈现适合用途；\term{图像压缩}（image compression）用更少比特表达图像。增强可以接在复原之后，压缩也可以与两者联合设计，但各自允许改变的量需要先规定。例如阅读招牌时允许提亮，档案记录却还需要保留原始观测和处理关系。下文的像素、信号和码流数字均为教学构造。

\section{图像复原}
\label{sec:vision-restoration-restoration}

恢复的困难取决于观测丢失了什么。噪声扰动像素，模糊混合邻域，下采样合并多个位置，缺口则完全没有相应观测。已知退化时，可以显式比较候选图像经过退化后是否接近输入；\term{盲复原}（blind restoration）是在使用时不知道全部退化参数，需要从图像及训练先验中共同估计。这里的“盲”描述可用信息，不能解释为已经覆盖所有真实相机和损坏形式。

\subsection{噪声干扰下的复原}
\label{sec:vision-restoration-denoising}

细小字画与噪声都可能表现为局部亮暗变化。要学会保留前者、去掉后者，训练需要说明噪声怎样产生，以及哪幅图像可作为清晰参考。最简单的模型是$\bm{Y}=\bm{X}+\bm{N}$，各元素$n_{i,j,k}$独立服从均值为零、方差为$\sigma^2$的高斯分布；真实传感器噪声可能随亮度变化，经过成像管线后也可能具有空间和通道相关性。

\subsubsection{DnCNN的噪声残差预测}
\label{sec:vision-restoration-dncnn}

DnCNN让网络预测观测与参考的差，而不是直接输出清晰图像。记这一残差网络为$r_{\bm{\theta}}$，由$n$对带噪图像与清晰参考训练
\begin{equation}
 L(\bm{\theta})=\frac{1}{2n}\sum_{i=1}^{n}
 \bigl\lVert r_{\bm{\theta}}(\bm{Y}_i)
             -(\bm{Y}_i-\bm{X}_i)\bigr\rVert_{\mathrm{F}}^2,
 \qquad
 \widehat{\bm{X}}=\bm{Y}-r_{\bm{\theta}}(\bm{Y})\text{。}
 \label{eq:vision-dncnn}
\end{equation}
这里的Frobenius范数表示所有像素与通道的平方和。网络以同尺寸卷积形成局部特征，中间层组合卷积、批归一化与ReLU，末层回到输入通道数，输出可正可负的残差。零填充维持空间大小，使相减能逐位置执行；它同时规定了图像边缘的处理方式\scite{vision-dncnn2017}

以招牌一行的四个灰度值为例，参考为$(0.2,0.2,0.8,0.8)$，加噪后得到$(0.25,0.17,0.84,0.74)$，残差标签就是$(0.05,-0.03,0.04,-0.06)$。若网络预测$(0.04,-0.02,0.03,-0.05)$，相减得到$(0.21,0.19,0.81,0.79)$。若网络把明暗跳变也解释成噪声，输出会抹掉字画；残差形式本身没有提供“噪声”和“真实细节”的语义区分，这一区分来自数据及损失。

固定噪声版本围绕给定$\sigma$合成配对，盲高斯版本在规定范围内抽取多个噪声强度。训练更新卷积参数与归一化统计，使用时保留这些参数和预处理，不再需要清晰参考。多任务的DnCNN-3还混合去噪、超分辨率和JPEG去块效应，但每类残差来自不同退化。把所有任务写成$\bm{Y}-\bm{X}$只是统一输出接口，不能推得真实噪声与JPEG伪影具有同一分布。

平方误差也决定了输出倾向。对固定观测$\bm{Y}=\bm{y}$和任意候选$\bm{a}$，有
\begin{equation}
 \mathbb{E}\!\left[\lVert\bm{X}-\bm{a}\rVert_{\mathrm{F}}^2
                  \mid\bm{Y}=\bm{y}\right]
 =
 \mathbb{E}\!\left[\lVert\bm{X}-\bm{\mu}\rVert_{\mathrm{F}}^2
                  \mid\bm{Y}=\bm{y}\right]
 +\lVert\bm{a}-\bm{\mu}\rVert_{\mathrm{F}}^2\text{，}
 \label{eq:vision-conditional-mean}
\end{equation}
其中$\bm{\mu}=\mathbb{E}[\bm{X}\mid\bm{Y}=\bm{y}]$，并假设二阶矩有限。展开平方后，交叉项因条件均值为零而消失，所以理想平方误差预测是条件均值。若一个细节存在多个可能位置，平均结果可能较平滑；这为后文比较像素保真与生成细节提供了依据。

\subsection{模糊观测下的复原}
\label{sec:vision-restoration-deblurring}

模糊不是逐像素独立扰动。相机运动或失焦使一个位置的光贡献到多个位置，恢复时要拆开这些混合。已知模糊核提供明确的反演关系；核未知时，学习式映射需要从数据识别可能的清晰结构。两条路线都依赖边界、对齐及退化范围。

\subsubsection{已知模糊核的频域正则化反卷积}
\label{sec:vision-restoration-deconvolution}

先考虑一维信号，假设模糊空间不变，并采用周期边界。卷积可写为$\bm{y}=\bm{H}\bm{x}+\bm{n}$，其中$\bm{H}$是循环矩阵。离散傅里叶变换将它对角化，每个频率满足$Y_\omega=H_\omega X_\omega+N_\omega$。直接除以$H_\omega$时，噪声变为$N_\omega/H_\omega$；弱频率响应会被大幅放大，零响应处则没有唯一反演。

为限制这种放大，可以把观测一致性与相邻位置的平滑要求合在一起：
\begin{equation}
 \widehat{\bm{x}}=\argmin_{\bm{x}}
 \lVert\bm{H}\bm{x}-\bm{y}\rVert_2^2+
 \lambda\lVert\bm{D}\bm{x}\rVert_2^2\text{，}
 \qquad \lambda>0\text{，}
 \label{eq:vision-deconv-objective}
\end{equation}
其中$\bm{D}$计算周期相邻差分。令梯度为零得到
$(\bm{H}^{\top}\bm{H}+\lambda\bm{D}^{\top}\bm{D})\widehat{\bm{x}}=\bm{H}^{\top}\bm{y}$。在二者均可由同一傅里叶基对角化且分母非零时，每个频率独立求解为
\begin{equation}
 \widehat X_\omega=
 \frac{\overline{H_\omega}}
 {|H_\omega|^2+\lambda|D_\omega|^2}\,Y_\omega\text{。}
 \label{eq:vision-regularized-inverse}
\end{equation}
上划线表示复共轭。惩罚较大的高频受到更强抑制，因而输出既减少噪声放大，也可能损失尖锐边缘。若$H_\omega=D_\omega=0$，此目标在该频率仍没有唯一解，需要另外规定约束。

取八点边缘$\bm{x}=(0,0,0,0,1,1,1,1)$，模糊核在中心和左右邻点的权重为$(0.6,0.2,0.2)$，再加入交替的$\pm0.04$噪声，得到观测
\[
 \bm{y}=(0.24,-0.04,0.04,0.16,0.84,0.96,1.04,0.76)\text{。}
\]
最高频处$H_\omega=0.2$，直接求逆把噪声幅值放大到0.2；用相邻差分与$\lambda=0.05$时，该处的逆增益为$0.2/(0.04+0.2)\approx0.833$，噪声分量幅值降为约0.033。

图\ref{fig:vision-deconvolution}显示全部八点结果。正则化后的过冲变小，但跃迁也变缓，且仍有轻微振铃。这里保留了负值与大于1的值以展示反演行为，没有在绘图前裁剪。真实图像若采用反射边界，不能继续照搬循环卷积公式；若模糊核错误或随位置变化，正则化只能缓解数值不稳定，不能修正错误的成像模型。

\begin{figure}[htbp]
 \centering
 % Eight-point circular blur, kernel (.2,.6,.2), alternating noise .04.
% DFT ridge inverse uses first differences and lambda=.05; values are un-clipped.
\begin{tikzpicture}[x=11mm,y=30mm,font=\figurefont\small]
 \draw[BookMuted,line width=.7pt,->] (-.35,-.3)--(7.5,-.3)
   node[right] {$i$};
 \draw[BookMuted,line width=.7pt,->] (-.35,-.3)--(-.35,1.38)
   node[above] {灰度};
 \foreach \y in {0,0.5,1} {
  \draw[BookRule,line width=.5pt] (-.35,\y)--(7.2,\y);
  \node[anchor=east] at (-.45,\y) {$\y$};
 }
 \foreach \i in {0,...,7} {
  \draw[BookMuted,line width=.5pt] (\i,-.3)--(\i,-.34);
  \node[below] at (\i,-.34) {$\i$};
 }
 \draw[BookInk,line width=1.2pt]
  plot coordinates {(0,0)(1,0)(2,0)(3,0)(4,1)(5,1)(6,1)(7,1)};
 \draw[BookAmber,dashed,line width=1pt]
  plot[mark=square*,mark size=1.5pt] coordinates
  {(0,.2)(1,-.2)(2,.2)(3,-.2)(4,1.2)(5,.8)(6,1.2)(7,.8)};
 \draw[BookTeal,line width=1pt]
  plot[mark=*,mark size=1.5pt] coordinates
  {(0,.1997)(1,-.0766)(2,-.01)(3,.133)(4,.867)(5,1.01)(6,1.0766)(7,.8003)};
 \draw[BookInk,line width=1.2pt] (0,1.65)--(.6,1.65);
 \node[anchor=west] at (.7,1.65) {参考};
 \draw[BookAmber,dashed,line width=1pt] (2.1,1.65)--(2.7,1.65);
 \node[anchor=west] at (2.8,1.65) {直接求逆};
 \draw[BookTeal,line width=1pt] (4.8,1.65)--(5.4,1.65);
 \node[anchor=west] at (5.5,1.65) {正则化};
\end{tikzpicture}

 \caption{八点周期边缘的教学反卷积。实线为清晰参考，带方块的虚线为直接求逆，带圆点的实线为差分正则化结果。两端因周期边界相接；纵轴未裁剪，便于观察过冲与平滑的取舍。}
 \label{fig:vision-deconvolution}
\end{figure}
\FloatBarrier

\subsubsection{Restormer的局部上下文与通道注意力复原}
\label{sec:vision-restoration-restormer}

运动模糊招牌的字画可能需要较大范围的结构信息。Restormer以多尺度编码器—解码器组织特征，把注意力矩阵放在通道维度上，控制高分辨率下的计算量。推理时不需输入显式模糊核，约束主要来自退化图像与清晰参考的训练配对\scite{vision-restormer2022}

输入经卷积得到特征，编码阶段逐级降低空间分辨率并增加通道，解码阶段上采样并拼接相应编码特征。在一个注意力头内，设位置数$n=hw$，通道数$d$。对归一化特征先作$1\times1$卷积混合通道，再作$3\times3$深度卷积混合邻域，得到按“通道$\times$位置”排列的$\bm{Q},\bm{K},\bm{V}\in\mathbb{R}^{d\times n}$。沿空间维归一化查询与键后，一种与实现一致的行向量记法是
\begin{equation}
 \bm{A}=\operatorname{softmax}_{\mathrm{row}}
       (\beta\,\bar{\bm{Q}}\bar{\bm{K}}^\top),
 \qquad
 \bm{O}=\bm{A}\bm{V}\text{，}
 \label{eq:vision-restormer-attention}
\end{equation}
其中$\beta$是每头学习的缩放量。$\bm{A}\in\mathbb{R}^{d\times d}$比较的是通道在整幅特征图上的响应关系，随后在每个位置混合通道；它不是每个像素直接对所有其他像素作加权求和。局部深度卷积及多尺度计算负责传播空间信息，全图统计又使通道组合依赖当前输入。

若$n=256\times256$而$d=32$，逐位置注意力需为一个头建立约$4.3\times10^9$个关系，通道矩阵则含1024项。仅考虑上述矩阵乘法，通道注意力的时间量级为$O(nd^2)$，关系矩阵存储为$O(d^2)$；还需要$O(nd)$的特征存储，以及卷积和其余网络的计算。说它对位置数近似线性，必须固定通道和其他结构。

门控前馈模块把归一化特征送入两个含逐点与深度卷积的支路，对一侧施加GELU后逐元素相乘，再投影并残差相加。乘法让一个支路调制另一个支路的响应，但门值不是必须位于$[0,1]$的概率。经过高分辨率细化后，末层预测图像修正量，与输入相加。这个残差符号与式\eqref{eq:vision-dncnn}中预测噪声再相减的约定不同。

训练以配对参考的像素损失优化网络，并逐渐增加裁剪块大小、相应减小批量，使模型在训练后期读取更大范围。对于招牌，小块可约束笔画，较大块还包含字间结构和周围边缘；使用时仍需核查笔画是否真实存在。去噪、去雨、运动模糊与失焦使用各自的数据和模型设置，同一个架构名不能代替检查点的适用范围。

\subsection{分辨率不足下的复原}
\label{sec:vision-restoration-super-resolution}

\term{超分辨率}（super-resolution）根据较低分辨率的输入估计更高分辨率的图像。放大网格增加了输出位置，却没有自动增加独立观测。例如把两个像素平均成一个时，$(0,1)$与$(0.5,0.5)$都会得到0.5。输出中的细线和文字怎样出现，取决于参考监督、退化假设与图像先验，而非输出像素数本身。

\subsubsection{插值输入与残差监督：DnCNN的超分辨率变体}
\label{sec:vision-restoration-dncnn-sr}

复用残差网络前，需要先让输入与目标处在同一网格。由高分辨率参考$\bm{X}$合成低分辨率图像$\bm{Y}=\mathcal{S}(\bm{X})$，再用双三次插值$\mathcal{U}$放大为$\bm{B}=\mathcal{U}(\bm{Y})$。DnCNN的相应变体学习$\bm{B}-\bm{X}$，使用时得到$\widehat{\bm{X}}=\bm{B}-r_{\bm{\theta}}(\bm{B})$。残差现在含插值和平滑带来的结构差异，不再是独立高斯噪声。

为了看清监督的含义，暂用两点平均和常数插值作教学简化。原细线为$(0,1)$，观测为0.5，插值得到$(0.5,0.5)$，残差标签为$(0.5,-0.5)$。若训练先验判断细线在右侧，预测残差能恢复这一排列；若真实原图为$(1,0)$，相同观测需要相反残差。唯一低分辨率观测无法区分两者。

像素监督衡量输出与训练参考的差。若应用还要求输出服从特定退化过程，应计算$\mathcal{S}(\widehat{\bm{X}})$并与$\bm{Y}$比较，这称为对观测的回投影核验。上述两个排列都能通过平均采样的回投影，所以观测一致性是必要依据，却不足以确认缺失细节。比较真实结果时还应固定放大倍数、采样滤波器和边界裁剪。

\Needspace{60mm}
\subsubsection{Real-ESRGAN的高阶退化合成与细节学习}
\label{sec:vision-restoration-real-esrgan}

网上照片可能先缩小、保存成JPEG，再截取、放大和重新压缩。只用一次双三次下采样训练，难以覆盖这种输入。Real-ESRGAN从清晰图像反复采样模糊、缩放、噪声和JPEG压缩，并使用sinc滤波模拟锐利边缘附近的振铃；原方法采用两轮退化构成训练配对\scite{vision-real-esrgan2021}

可以将一对训练样本写成$\bm{Y}=\mathcal{D}_2(\mathcal{D}_1(\bm{X}))$与$\bm{X}$。每轮内部都有随机参数，缩放可以向上或向下，噪声和压缩强度也变化，末尾操作另有顺序安排。高阶指重复退化过程，不表示对图像求高阶导数。这样得到的低清招牌仍与同一张清晰参考对应，监督目标可直接计算。

生成器沿用ESRGAN的残差密集结构。先以像素距离训练复原网络，再加入感知损失和对抗损失；前者比较固定视觉网络中的特征，后者让生成器面对区分真实和生成内容的判别器。U形判别器在多个局部位置给出反馈，并采用谱归一化约束训练。训练要更新生成器和判别器，使用时只保留生成器及相同输入约定，一次前向得到放大图像。

用多次JPEG保存的招牌检查结果时，应分别看方块是否减弱、笔画是否连贯、文字是否改变。重复退化扩大了合成覆盖，感知和对抗目标又鼓励自然细节；这些条件允许模型补出看似合理的纹理。它们既没有穷尽所有相机退化，也没有为补出的字画提供新的观测。

\subsubsection{DiffBIR的退化移除与生成先验补充}
\label{sec:vision-restoration-diffbir}

严重退化下，直接复原容易产生较平滑的结果。DiffBIR将处理拆为减轻退化与补充内容两个阶段：前阶段提供较稳定的图像条件，后阶段用预训练扩散模型补充细节。本节采用2024年3月的第二版预印本描述，区分其中的训练条件构造和实际使用\scite{vision-diffbir2024}

训练生成模块时，作者另训练一个像素平方误差复原器，在较宽的简化退化范围内产生条件图像$\bm{X}_{\mathrm{c}}$。这个预处理器的职责是提供多样、较少退化干扰的训练条件；实际使用时改接各任务已有的复原模块，例如面部复原或盲超分辨率模块。把这两者统称为“第一阶段”会掩盖一个重要差别：训练时的条件产生器不必是交付时处理真实照片的复原器。

生成阶段用固定的图像编码器把$\bm{X}_{\mathrm{c}}$变为条件潜变量$\bm{z}_{\mathrm{c}}$。IRControlNet复制预训练去噪网络的编码及中间模块，读取$\bm{z}_{\mathrm{c}}$与当前带噪潜变量的通道拼接。支路经零初始化卷积输出修正，加入固定去噪网络的中间与跳连特征。训练给清晰参考潜变量加噪，以噪声预测平方误差只更新条件网络及连接参数；主去噪网络、图像编码器和解码器保留预训练参数。共同的潜空间条件生成接口将在\ref{sec:vision-models-media-generation}节集中说明。

推理从随机潜变量开始，反复用条件支路与固定生成先验去噪，再解码成图像。可选的复原引导在每步得到清晰潜变量估计$\tilde{\bm{z}}^{(t)}$后，减少解码图像与前阶段结果的加权距离：
\begin{equation}
 \begin{aligned}
 L_{\mathrm{g}}(\tilde{\bm{z}})
 &=\frac{1}{hwc}
   \left\lVert\bm{W}\odot
    \bigl(\mathcal{D}(\tilde{\bm{z}})-\bm{X}_{\mathrm{c}}\bigr)
   \right\rVert_{\mathrm{F}}^2,\\
 \tilde{\bm{z}}_{\mathrm{g}}^{(t)}
 &=\tilde{\bm{z}}^{(t)}
   -s\nabla_{\tilde{\bm{z}}}L_{\mathrm{g}}
       (\tilde{\bm{z}}^{(t)})\text{。}
 \end{aligned}
 \label{eq:vision-diffbir-guidance}
\end{equation}
这里$\mathcal{D}$是图像解码器，$s$控制引导强度，$\bm{W}$根据前阶段图像的局部梯度统计形成，对较平坦区域赋予更大权重。更新后的清晰估计参与下一采样步。损失中的权重在平方之内，实际像素平方误差的系数是$w_{i,j,k}^2$。

在固定低清人像和随机种子下，增加引导会加强对前阶段轮廓及平坦区域的约束，同时可能压低生成细节。这里被追随的仍是一个复原估计，它若已把眼睛或文字恢复错，引导也可能保留错误。皮肤纹理、身份细节与招牌字形应分别核验，不能把新的高频外观当作相机曾经记录到的事实。

\subsection{缺失区域的复原}
\label{sec:vision-restoration-inpainting}

坏点或缺口与加性噪声不同：对应位置根本没有可信观测。以二值掩码$\bm{M}$标明有效像素，$m_{i,j}=1$表示可用，0表示缺失，并在通道维广播。恢复目标是依据周围内容估计缺口原有内容；若要求删除杯子并创造一只花瓶，则还包含新的语义目标，交由\ref{chap:vision-image-generation}章讨论。

\subsubsection{部分卷积的有效像素归一化与掩码更新}
\label{sec:vision-restoration-partial-convolution}

普通卷积会读取缺口中填入的零或均值，输出因占位数值而变化。\term{部分卷积}（partial convolution）先屏蔽无效位置，再按有效数量调整响应。把一个窗口的输入、掩码和权重分别展平为$\bm{x},\bm{m},\bm{w}\in\mathbb{R}^q$，其输出及新掩码为
\begin{equation}
 x'=
 \begin{cases}
 \dfrac{q}{\sum_i m_i}\bm{w}^{\top}(\bm{x}\odot\bm{m})+b,
   &\sum_i m_i>0,\\
 0,&\sum_i m_i=0,
 \end{cases}
 \qquad
 m'=\mathbb{I}\!\left[\sum_i m_i>0\right]\text{。}
 \label{eq:vision-partial-convolution}
\end{equation}
其中$b$是偏置。归一化因子只乘加权和，不乘偏置；完全无有效输入时输出零，新掩码也为零\scite{vision-partial-convolution2018}

取窗口$\bm{x}=(0.2,0,0.8)$、$\bm{m}=(1,0,1)$、$\bm{w}=(1/3,1/3,1/3)$且$b=0$，屏蔽后的加权和为$1/3$，再乘$3/2$得到0.5，$m'=1$。若把缺口占位改成100，结果仍为0.5。窗口内全缺失时则不能计算这种归一化，必须走零输出分支。图\ref{fig:vision-partial-mask}说明，堆叠窗口使可计算范围逐层进入缺口。

\begin{figure}[htbp]
 \centering
 % Binary masks; kernel width 3, stride 1. No image values are fabricated.
\begin{tikzpicture}[x=1mm,y=1mm,font=\figurefont\small,
 arr/.style={-{Stealth[length=2mm]},BookMuted,line width=.8pt}]
 \foreach \r/\textlabel in {0/输入掩码,1/第一层,2/第二层} {
  \node[anchor=east] at (0,{-\r*18}) {\textlabel};
  \foreach \i in {0,...,6} {
   \draw[BookMuted,line width=.6pt] ({5+\i*11},{-\r*18-4})
     rectangle ++(11,8);
  }
 }
 \foreach \r/\indices in {0/{0,1,5,6},1/{0,1,2,4,5,6},2/{0,1,2,3,4,5,6}} {
  \foreach \i in \indices {
   \fill[BookBlue!12] ({5+\i*11+.4},{-\r*18-3.6}) rectangle ++(10.2,7.2);
   \fill[BookBlue] ({10.5+\i*11},{-\r*18}) circle (.8mm);
  }
 }
 \draw[BookAmber,line width=1pt] (27,-5.5) rectangle (60,5.5);
 \draw[BookAmber,line width=1pt] (27,-23.5) rectangle (60,-12.5);
 \draw[arr] (43.5,-6)--(43.5,-12);
 \draw[arr] (43.5,-24)--(43.5,-31);
 \node[anchor=north] at (43.5,-43) {中心掩码：$0\to0\to1$};
\end{tikzpicture}

 \caption{宽度为3、步幅为1的掩码更新教学例。带点填充的格子表示已有有效输入，白格表示缺失；每层只要窗口内有一个有效位置，该层输出就标为有效。图中省略边界外输入。中间格在第二层获得可计算支持，不代表新取得了原像素观测。}
 \label{fig:vision-partial-mask}
\end{figure}

部分卷积修补网络把特征和掩码同时传过编码器、解码器及跳连。训练从清晰图像随机去掉不规则区域，分别约束已知区和缺失区的像素误差，并加入感知特征、风格统计及缺口周边平滑约束。感知项关心特征结构，风格项关心通道相关统计，它们都不能唯一确定遮挡背后的实际对象。

使用时输入受损图像和原始掩码，网络预测$\widehat{\bm{X}}_{\mathrm{net}}$，再形成
$\widehat{\bm{X}}=\bm{M}\odot\bm{Y}+(1-\bm{M})\odot\widehat{\bm{X}}_{\mathrm{net}}$，使已知像素按要求保留。宽度很小的划痕常有两侧结构可参考，大缺口则允许多种内容。中间掩码变为1只说明特征已有计算来源，原始观测掩码仍须保留，用来区分记录与估计。

\section{图像增强}
\label{sec:vision-restoration-enhancement}

有些图像已经保留了需要的结构，只是亮度或对比不适合观看；另一些暗图则同时缺少可靠观测。两种情况可能有相似外观，但只改变显示曲线与重新估计内容会带来不同结果。增强的目标应说明希望看清什么，以及哪些颜色、测量量或局部关系必须保留。

\subsection{暗处可见性与曝光的改善}
\label{sec:vision-restoration-low-light}

低照度时，传感器收到的光子较少，信号相对噪声较弱。若将像素与加性噪声一起乘上$a$，信号幅值变为$a$倍，噪声标准差也变为$a$倍，这种线性提亮没有改善二者的比例。延长真实曝光能收集更多光子，却可能引入运动模糊；对已经拍下的短曝光照片，事后乘增益不能替代重新曝光。

\subsubsection{Learning to See in the Dark的RAW到RGB映射}
\label{sec:vision-restoration-sid}

Learning to See in the Dark直接读取短曝光RAW，并用同场景长曝光图像形成的RGB作为监督。它把颜色采样整理、去噪和颜色形成放在同一学习流程中，因而可利用RGB图像形成之前的信息\scite{vision-sid2018}

以Bayer传感器为例，每个$2\times2$位置采样一红、两绿和一蓝，原始阵列为$h\times w\times1$。按位置打包后得到$(h/2)\times(w/2)\times4$数组；两个绿色位置保留为不同通道。打包只是重新排列观测，没有提前补齐每个像素的三种颜色。输入减去传感器黑电平，按有效范围归一化，再乘外部给定的放大比例。网络用编码器—解码器预测半分辨率的12通道结果，经子像素重排形成全分辨率三通道RGB。

假设某个RAW读数为530，黑电平为512，可用上界为16383，则归一化信号为$(530-512)/(16383-512)\approx0.00113$。乘100后约为0.113。若忽略黑电平，直接用$530/16383$再乘100会超过3，数值主要来自偏置并很快被裁剪；这不是同一种输入。相机的黑电平、颜色阵列和增益约定必须与训练一致。

训练以短曝光输入与对齐长曝光RGB参考之间的$L_1$距离更新网络。曝光比为放大比例提供依据，使用时放大比例也可作为外部控制，原方法并不自动推断最佳曝光。长曝光图像仍可能有噪声，移动人物和树叶也会破坏配对；若参考错位，逐像素监督就会把位置差异也当作待修正误差。

固定同一夜景比较常规RGB、RGB提亮与RAW映射时，应查看暗处结构、颜色和亮部饱和，并记录三者实际使用的数据。RAW路线拥有更早阶段的传感器信息，不能把收益全部归结为更强的RGB后处理。相机更换和动态场景还需要单独验证，静态配对训练没有自动解决运动。

\subsection{色彩、对比与层次的调整}
\label{sec:vision-restoration-tone}

当图像已有可用内容时，可以直接规定数值映射，不必训练网络。全局调整让相同输入值得到相同输出，局部调整则依据邻域分布改变映射；后者更灵活，也更容易改变区域之间的相对关系。

\subsubsection{白平衡与全局色调曲线}
\label{sec:vision-restoration-white-balance}

\term{白平衡}（white balance）按颜色通道调整增益，使选定照明条件下的中性物体呈现为中性色。在简化的线性RGB模型中，$\bm{c}'=\operatorname{diag}(g_r,g_g,g_b)\bm{c}$。若灰色参考被拍成$(0.20,0.16,0.10)$，取增益$(0.8,1,1.6)$会得到$(0.16,0.16,0.16)$。这依赖参考本来应中性；把一面黄色墙误当灰卡会使整图颜色偏移。色彩管理还包含颜色空间变换等环节，通道增益只表达其中的白平衡近似。

色调曲线把亮度$I\in[0,1]$映射成$T(I)$。为避免不同软件对“伽马”的方向约定混淆，这里直接规定$T(I)=I^\gamma$。当$\gamma=1/2$时，0.04、0.25和0.81依次变为0.2、0.5和0.9，暗部变化较大；$\gamma>1$则压暗中间值。若先作两倍线性提亮再裁剪，0.6和0.9都变为1，高光之间的差异就丢失了。

逐通道应用非线性曲线通常会改变通道比例。若希望尽量保持线性RGB比例，可先计算亮度$I$，再在$I>0$处用$\bm{c}'=\bm{c}\,T(I)/I$统一缩放，随后处理超出色域的数值；接近零时需稳定分母并规定黑色输出。保持比例不等于所有显示条件下的感知色相严格不变，颜色空间和裁剪仍有作用。

这些操作重新分配现有层次。招牌暗部可能更容易辨认，平坦暗区的噪声也可能一起显露；已经被传感器饱和或量化合并的字画，不能仅靠单调曲线恢复。

\subsubsection{受限局部直方图均衡化}
\label{sec:vision-restoration-clahe}

同一全局曲线不容易同时照顾阴影和高光。\term{受限局部直方图均衡化}（contrast limited adaptive histogram equalization, CLAHE）为图像分块建立亮度直方图，限制过高的计数峰，并用重新分配后的累计分布构造局部映射。限制峰值是为了避免把少量灰度差异拉伸得过强；在相邻块的映射结果之间插值，则减少块边界突变。

用四个亮度箱说明计算。某块的直方图为$(0,8,0,0)$，共有8个像素，若直接均衡化，第二箱的累计概率立刻达到1，微小变化会引起很大的映射变化。令计数上限为2，截断后得到$(0,2,0,0)$，需要重新分配6个计数。以保持上限为目标，将这些计数均匀补入其余未满的箱，得到$(2,2,2,2)$，累计概率为$(0.25,0.5,0.75,1)$。这个小例采用容量约束的整数分配；实现中的阈值归一化和余数分配规则须记录，不能只写一个无单位的“阈值2”。

一个像素通常受到周围四个块的映射影响。先将它的亮度分别输入四条曲线，再按位置双线性插值，因而即便亮度相同，不同位置的输出也可能不同。块较小时更关注局部，但噪声及小结构更容易主导直方图；上限放宽又增加局部对比。对招牌可以同时检查细字与平坦墙面，避免只看被增强的字画而忽略噪声和光晕。医学、档案和测量图像还应规定哪些灰度关系允许变化。

\section{图像压缩}
\label{sec:vision-restoration-compression}

压缩的输出是可传输的比特流。一个网络即使能把图像映射到很小的浮点数组，也不等于已经压缩：浮点数需要多少位、哪些辅助信息要传、解码器如何复现概率和符号，都必须明确。以下用尺度超先验方法连接表示学习与实际编解码。

\subsection{表示变换与量化}
\label{sec:vision-restoration-quantization}

分析变换$g_{\mathrm{a}}$把图像变为主潜变量$\bm{y}$，合成变换$g_{\mathrm{s}}$从离散潜变量恢复图像：
\begin{equation}
 \bm{y}=g_{\mathrm{a}}(\bm{X}),\qquad
 \hat y_i=\operatorname{round}(y_i),\qquad
 \widehat{\bm{X}}=g_{\mathrm{s}}(\hat{\bm{y}})\text{。}
 \label{eq:vision-transform-coding}
\end{equation}
\term{量化}（quantization）把连续数映射到有限精度的离散值，舍弃区间内的差别。此处以单位步长取整为例，步长的作用也可由分析与合成变换吸收。学习式变换通过卷积及非线性组织空间、通道信息，训练目标共同决定哪些变化值得保留。

用两个像素作可复算的变换例。对$\bm{x}=(0.2,0.8)$，令$y_1=2(x_1+x_2)=2$、$y_2=2(x_2-x_1)=1.2$。取整得到$(2,1)$，合成按$x_1=(y_1-y_2)/4$、$x_2=(y_1+y_2)/4$得到$(0.25,0.75)$。均方误差为0.0025。这里分析与合成是手工线性变换，学习式网络用数据寻找更适合图像的非线性变换；量化造成不可逆信息损失的原因相同。

取整函数的导数几乎处处为零，不能直接提供有用梯度。尺度超先验方法在训练中以独立均匀噪声替代取整，即$\tilde y_i=y_i+u_i$、$u_i\sim\mathcal{U}(-1/2,1/2)$。带波浪号的连续随机变量用于可微训练，带帽号的整数用于实际编码。训练近似不能代替交付时真正取整、编码并解码的检查。

设$R$为估计码长，$D$为重建失真，训练优化$R+\lambda D$。固定两者的归一化口径时，较大的$\lambda$让失真受到更重惩罚，网络可能保留更多细节并付出更多码率；这不是任意单幅图像上单调改善所有指标的保证。通常需要多个工作点形成码率—失真曲线，而非只调一个数就宣称适合全部用途。

\subsection{概率估计与码流形成}
\label{sec:vision-restoration-bitstream}

量化后还需给整数分配码字。常见符号应使用较短表示，少见符号可以较长，但编码器与解码器必须共享同一概率表。主潜变量在平坦区域常聚集于小幅值，在纹理区域可能更分散；一个全局分布难以充分表达这些变化。

\Needspace{45mm}
\subsubsection{尺度超先验的辅助表示与条件概率}
\label{sec:vision-restoration-hyperprior}

\term{尺度超先验}（scale hyperprior）为主潜变量的局部尺度建立辅助表示\scite{vision-hyperprior2018}
超分析变换产生$\bm{z}=h_{\mathrm{a}}(\bm{y})$，量化为$\hat{\bm{z}}$后，由超合成变换得到尺度$\bm{\sigma}=h_{\mathrm{s}}(\hat{\bm{z}})$。每个$\sigma_i>0$控制一个主符号的概率分布，辅助符号本身也需要编码。

原方法以条件零均值高斯描述主潜变量，并将密度在单位量化区间上积分。对整数$k$，离散概率为
\begin{equation}
 P_i(k\mid\hat{\bm{z}})
 =\int_{k-1/2}^{k+1/2}
   \frac{1}{\sqrt{2\pi}\sigma_i}
   \exp\!\left(-\frac{u^2}{2\sigma_i^2}\right)\,\mathrm{d}u
 =\Phi\!\left(\frac{k+1/2}{\sigma_i}\right)
  -\Phi\!\left(\frac{k-1/2}{\sigma_i}\right)\text{，}
 \label{eq:vision-hyperprior-probability}
\end{equation}
其中$\Phi$是标准正态累计分布函数。不能把密度在整数处的值直接当作符号概率；密度有单位，且在整数点求和不保证等于1。小$\sigma_i$给零附近符号较高概率，大$\sigma_i$为较大幅值留出更多概率。

辅助潜变量采用学习到的因子化分布，编码时也在量化区间上取得离散概率。一次图像的理想估计码长为
\begin{equation}
 R_{\mathrm{est}}=
 -\sum_j\log_2 P_j^{z}(\hat z_j)
 -\sum_i\log_2 P_i(\hat y_i\mid\hat{\bm{z}})\text{。}
 \label{eq:vision-hyperprior-rate}
\end{equation}
训练时相应使用加噪潜变量的连续密度，并与重建损失共同更新分析、合成及概率模型。辅助流的项不可省略，否则网络可以把大量信息转移到$\bm{z}$而不付出代价，目标就不再衡量总压缩成本。

假设共享全局分布编码主符号需要100 bit，传入尺度辅助信息后主流降至72 bit。若辅助流为12 bit，总计84 bit，有16 bit收益；若辅助流需要35 bit，总计107 bit，反而增大。只有节省的主流成本超过辅助成本时，这种条件化才划算。图\ref{fig:vision-hyperprior-codec}标出辅助信息如何同时让两端构造主流概率。

\begin{figure*}[htbp]
 \centering
 % Encoder/decoder dependencies. All ordinary edges are axis-aligned.
\begin{tikzpicture}[x=1mm,y=1mm,font=\figurefont\small,
 b/.style={draw=BookBlue,fill=BookBlue!8,minimum width=13mm,
   minimum height=8mm,inner sep=1mm,line width=.8pt,align=center},
 o/.style={b,draw=BookTeal,fill=BookTeal!8},
 s/.style={draw=BookMuted,minimum width=12mm,minimum height=8mm,
   inner sep=1mm,line width=.8pt},
 a/.style={-{Stealth[length=1.8mm]},BookInk,line width=.9pt},
 c/.style={a,BookAmber,rounded corners=1mm}]
 \node at (5,0) (x) {$\bm{X}$};
 \node[b] at (22,0) (ga) {$g_{\mathrm{a}}$};
 \node[b] at (42,0) (qy) {$Q_y$};
 \node[b] at (64,0) (ey) {E};
 \node[s] at (83,0) (sy) {$S_y$};
 \node[o] at (102,0) (dy) {D};
 \node[o] at (123,0) (gs) {$g_{\mathrm{s}}$};
 \node at (144,0) (xh) {$\widehat{\bm{X}}$};
 \node[b] at (22,-32) (ha) {$h_{\mathrm{a}}$};
 \node[b] at (42,-32) (qz) {$Q_z$};
 \node[b] at (64,-32) (ez) {E};
 \node[s] at (83,-32) (sz) {$S_z$};
 \node[o] at (102,-32) (dz) {D};
 \node at (123,-32) (zh) {$\hat{\bm{z}}$};
 \node[b] at (42,-16) (hse) {$h_{\mathrm{s}}$};
 \node[o] at (123,-16) (hsd) {$h_{\mathrm{s}}$};
 \foreach \u/\v in {x/ga,ga/qy,qy/ey,ey/sy,sy/dy,dy/gs,gs/xh,
   ha/qz,qz/ez,ez/sz,sz/dz,dz/zh} {\draw[a] (\u)--(\v);}
 \draw[a] (ga)--(ha);
 \draw[c] (qz)--(hse);
 \draw[c] (hse.east)--(64,-16)--(ey.south);
 \draw[c] (zh)--(hsd);
 \draw[c] (hsd.west)--(102,-16)--(dy.south);
 \node[anchor=west,text=BookAmber] at (66,-13) {$\bm{\sigma}$};
 \node[anchor=east,text=BookAmber] at (100,-13) {$\bm{\sigma}$};
 \node[anchor=south] at (33,4) {$\bm{y}$};
 \node[anchor=south] at (112.5,4) {$\hat{\bm{y}}$};
 \node[anchor=north,text=BookMuted] at (42,-40) {编码端};
 \node[anchor=north,text=BookMuted] at (112,-40) {解码端};
\end{tikzpicture}

 \caption{尺度超先验的实际编解码依赖。$Q$为取整，E与D分别为熵编码和熵解码，$S_y,S_z$为两条比特流。两端都从同一组整数$\hat{\bm{z}}$经$h_{\mathrm{s}}$产生尺度，主流解码因此必须等待辅助流。图中省略参数与尺寸元数据，未画入不可传输的训练噪声或清晰参考。}
 \label{fig:vision-hyperprior-codec}
\end{figure*}
\FloatBarrier

\subsubsection{量化符号的熵编码与实际码长}
\label{sec:vision-restoration-entropy-coding}

\term{熵编码}（entropy coding）根据共享概率把离散符号无损地转成比特。这里“无损”指整数符号能够原样恢复，前面的量化仍然有损。算术编码的一种解释是不断缩小$[0,1)$区间，用落在最终区间内的二进制前缀标识整个序列。

设三个教学符号$A,B,C$的概率依次为$1/2,1/4,1/4$，基础区间为$[0,1/2)$、$[1/2,3/4)$、$[3/4,1)$。编码序列$ABC$时，$A$把当前区间缩为$[0,1/2)$，$B$选择其中第二段得到$[1/4,3/8)$，$C$再选择最后四分之一得到$[11/32,12/32)$。其宽度为$1/32$，对应理想码长5 bit；二进制前缀$01011$恰好表示这个区间。

若解码端已知长度为3，可用区间内代表值$11/32$恢复序列。它落在基础区间的$A$段，归一化到该段后得到$11/16$，落在$B$段；再次减去$1/2$并除以$1/4$，得到$3/4$，落在$C$段。区间统一左闭右开，因此边界也有确定归属。图\ref{fig:vision-arithmetic-coding}画出逐步缩小的区间；实际实现使用有限精度整数和归一化输出位，无需保存无限精度实数。

\begin{figure}[htbp]
 \centering
 % All intervals use the same coordinate scale; probabilities .5,.25,.25.
\begin{tikzpicture}[x=80mm,y=1mm,font=\figurefont\small]
 \foreach \row/\lab in {0/初始,1/$A$,2/$AB$,3/$ABC$} {
  \draw[BookRule,line width=.7pt] (0,{-15*\row})--(1,{-15*\row});
  \node[anchor=east] at (-.05,{-15*\row}) {\lab};
 }
 \fill[BookBlue!15] (0,-3) rectangle (1,3);
 \draw[BookBlue,line width=.8pt] (0,-3) rectangle (1,3);
 \fill[BookTeal!20] (0,-18) rectangle (.5,-12);
 \draw[BookTeal,line width=.8pt] (0,-18) rectangle (.5,-12);
 \fill[BookTeal!35] (.25,-33) rectangle (.375,-27);
 \draw[BookTeal,line width=.8pt] (.25,-33) rectangle (.375,-27);
 \fill[BookTeal] (.34375,-48) rectangle (.375,-42);
 \foreach \x/\lab in {0/0,.5/{1/2},1/1} {
  \node[anchor=south] at (\x,4) {$\lab$};
 }
 \node[anchor=west] at (.54,-15) {$[0,1/2)$};
 \node[anchor=west] at (.43,-30) {$[1/4,3/8)$};
 \node[anchor=west] at (.43,-45) {$[11/32,12/32)$};
 \draw[BookMuted,dashed,line width=.6pt] (.5,-18)--(.375,-27);
 \draw[BookMuted,dashed,line width=.6pt] (0,-18)--(.25,-27);
 \draw[BookMuted,dashed,line width=.6pt] (.25,-33)--(.34375,-42);
 \draw[BookMuted,dashed,line width=.6pt] (.375,-33)--(.375,-42);
\end{tikzpicture}

 \caption{固定概率下序列$ABC$的算术编码教学例。每行共享$[0,1)$坐标，实色区间为读入当前前缀后仍可能的范围。最终区间宽$1/32$，可由5位前缀$01011$表示；序列长度和概率表在本例中预先共享，不计入这5位。}
 \label{fig:vision-arithmetic-coding}
\end{figure}

在尺度超先验模型中，要把式\eqref{eq:vision-hyperprior-probability}离散化成编码器可用的整数累计概率表，规定精度、尾部符号与逃逸机制，并使两端采用一致的符号顺序。零概率不能分配给实际出现的符号，否则无法编码。有限精度概率、流终止、填充以及尺寸和版本元数据都会使实际字节数偏离负对数概率估计。

真实文件应先提供足以解析的尺寸和模型约定，再存辅助流与主流，或用等价可解析的组织方式记录偏移。若一张$256\times256$图像的全部文件为4096 byte，则实际每像素比特数为$8\times4096/(256\times256)=0.5$，分母不再乘RGB的通道数。模型权重若由收发双方预先共享，常规每图码率不含权重；若每张图传一个专用模型，则这项成本也必须纳入所声明的传输预算。

\subsection{解码重建与用途保持}
\label{sec:vision-restoration-decoding}

解码按依赖执行。用共享的辅助概率模型从$S_z$恢复$\hat{\bm{z}}$，由$h_{\mathrm{s}}$产生主流概率，再从$S_y$恢复$\hat{\bm{y}}$，最后由$g_{\mathrm{s}}$形成$\widehat{\bm{X}}$。编码端也必须用量化后的$\hat{\bm{z}}$计算尺度，不能用未经传输的连续$\bm{z}$，否则两侧的概率表可能不同。

可以用一次独立解码检查这一点：只给解码程序文件、声明的模型参数和必要配置，确认它能复现输出。原图、训练参考和编码中间浮点值都不在它的可用输入中。一个只展示分析网络和合成网络前向结果的实验，尚未完成这个检查。

同样的0.5 bit/像素可以优先保留平滑颜色，也可以保留文本边缘，不同失真目标会形成不同取舍。应把解码图像送入固定的文字读取、识别或测量流程，检查招牌字符串、车牌边界及尺寸。感知上自然的纹理不一定有助于这些用途，较小的平均像素误差也可能掩盖一个关键数字的错误。视频中的运动流、残差流和参考帧依赖将在\ref{chap:vision-video-restoration}章接续。

\section{处理效果的验证与选择}
\label{sec:vision-restoration-evaluation}

判断处理是否值得，需要同时看内容、呈现与成本。比较前固定输入退化、参考对齐、像素范围和颜色约定；对没有可信参考的真实图像，应明确哪些结论只能来自任务结果或人工观察，不能给不存在的原图计算“恢复正确率”。

\subsection{内容保真与感知质量}
\label{sec:vision-restoration-quality}

参考指标逐对比较处理结果与对应图像，感知判断则关心人眼看到的结构或自然程度。两者可以相关，但比较对象不同，尤其在一幅低清观测对应多个可能细节时。

\subsubsection{PSNR与SSIM的像素和局部结构比较}
\label{sec:vision-restoration-psnr-ssim}

对含$n$个有效标量元素的对齐图像，均方误差和\term{峰值信噪比}（peak signal-to-noise ratio, PSNR）为
\begin{equation}
 \operatorname{MSE}=\frac1n\sum_{i=1}^{n}(x_i-\hat x_i)^2,
 \qquad
 \operatorname{PSNR}
 =10\log_{10}\frac{v_{\max}^2}{\operatorname{MSE}}\text{。}
 \label{eq:vision-psnr}
\end{equation}
$v_{\max}$是约定的峰值，归一化图像通常取1，8位整数通常取255；误差为零时PSNR为正无穷。前述两像素压缩例的MSE为0.0025，取$v_{\max}=1$得到约26.02 dB。对一组图像平均各图PSNR，与先合并所有像素误差再计算PSNR不是同一聚合方式。

\term{结构相似性}（structural similarity, SSIM）在对应局部窗口中比较亮度、对比和共同变化。其常用合并形式为
\begin{equation}
 \operatorname{SSIM}(\bm{x},\hat{\bm{x}})
 =
 \frac{(2\mu_x\mu_{\hat x}+c_1)(2\sigma_{x,\hat x}+c_2)}
 {(\mu_x^2+\mu_{\hat x}^2+c_1)
  (\sigma_x^2+\sigma_{\hat x}^2+c_2)}\text{，}
 \label{eq:vision-ssim}
\end{equation}
其中$\mu$是局部加权均值，$\sigma^2$是相应方差，$\sigma_{x,\hat x}$是协方差；$c_1,c_2>0$稳定近常量窗口的计算。整图值由局部结果聚合，窗口大小、权重、边界和颜色处理都应固定\scite{vision-ssim2004}

取参考$(0,0,1,1)$，平滑结果$(0,0.25,0.75,1)$的MSE为0.03125；只把第三个位置改成0的结果$(0,0,0,1)$，MSE为0.25，按固定长度比较时也可视为边缘平移一格。再取$(0,0,0.8,1)$，MSE只有0.01，但若第三位置恰是识别一笔关键字画的依据，任务错误仍可能很大。像素分数描述这些改变的平均幅度，SSIM还考察局部共同变化，二者都不知道这笔画的业务含义。

配准误差会使正确内容与错误位置比较。评价应先确认对应关系，再统一是否裁剪边缘、只评亮度还是评RGB，以及线性值或显示编码值。不能为每个方法单独选择最有利的对齐和颜色口径。

\Needspace{65mm}
\subsubsection{LPIPS的深层特征距离}
\label{sec:vision-restoration-lpips}

像素变化的大小未必与感知变化一致。\term{学习感知图像块相似度}（learned perceptual image patch similarity, LPIPS）使用视觉网络的多层特征，并依据人类对图像块相似性的判断校准组合权重\scite{vision-lpips2018}

设第$\ell$层在位置$p$的特征向量经通道归一化后为$\bar{\bm{f}}_{\ell,p}(\bm{X})$，该层共有$n_\ell$个位置，学习到的通道权重为$\bm{w}_\ell$。其距离写成
\begin{equation}
 d_{\mathrm{LPIPS}}(\bm{X},\widehat{\bm{X}})
 =
 \sum_\ell\frac1{n_\ell}\sum_p
 \left\lVert
 \bm{w}_\ell\odot
 \bigl(\bar{\bm{f}}_{\ell,p}(\bm{X})
      -\bar{\bm{f}}_{\ell,p}(\widehat{\bm{X}})\bigr)
 \right\rVert_2^2\text{。}
 \label{eq:vision-lpips}
\end{equation}
两幅图像都经过同一网络与预处理。先按位置比较归一化特征，再平均空间差异并累加层，不能将整图分类向量的余弦距离直接叫作LPIPS。若某位置的两通道差为$(0.2,0.1)$，权重为$(1,2)$，该位置贡献为$0.2^2+(2\times0.1)^2=0.08$。实际值取决于骨干与校准版本。

平滑砖墙可能产生较小像素误差，却改变深层纹理特征；把招牌“8”改成“3”又可能只影响很小区域，整体感知距离仍然较小。LPIPS衡量对应图像的特征差异，不能独立核验文字、身份或事实，也不能把一项分数直接解释成“恢复了多少真实内容”。

感知与失真的理论取舍采用另一个更具体的定义。设$p_{\bm{X}}$为真实图像分布，复原规则$p(\widehat{\bm{X}}\mid\bm{Y})$诱导输出分布$p_{\widehat{\bm{X}}}$，以$\Delta$度量逐对失真，以$d$度量两种分布的差异。给定失真预算$D$，定义
\begin{equation}
 P(D)=
 \inf_{\substack{p(\widehat{\bm{X}}\mid\bm{Y})\\
  \mathbb{E}[\Delta(\bm{X},\widehat{\bm{X}})]\leq D}}
 d(p_{\bm{X}},p_{\widehat{\bm{X}}})\text{。}
 \label{eq:vision-perception-distortion}
\end{equation}
放宽$D$扩大可选复原规则集合，因此$P(D)$不增；若$d$对第二个分布具有适当凸性，还可得到凸的最优边界。这个框架中的感知差异比较分布，不能直接替换成任意一个单图美观分数\scite{vision-perception-distortion2018}

一个完全没有信息的观测也能说明差别。假设原像素以相同概率为0或1，平方误差最优输出为0.5，平均失真0.25，但输出分布只有一个中间值。若独立地以相同概率输出0或1，边缘分布与真实像素相同，平均平方误差却为0.5。后一规则没有因此猜对每个原像素。实际方法远离最优边界时，改进模型仍可能同时降低多项误差；上述取舍不表示所有评价指标必然互相恶化。

\subsection{下游用途与信息保留}
\label{sec:vision-restoration-downstream}

招牌任务可以固定一套文字读取器，分别输入原照片、处理结果和可信参考，再与人工核实的字符串比较。若原图识别成“18路”、处理后成“13路”、参考为“18路”，锐化后的感知改善没有帮助这次读取。还应检查处理图像本身是否改出了“3”，或字符仍为“8”但读取器产生误判，以区分图像处理与下游系统的误差。

对于测量，可固定边缘定位与标定步骤，比较长度、角点和深度；对于识别，可比较相同对象的类别及拒识结果。阈值需要在验证集上确定，不能反复查看测试标签后为某个处理器调优。若联合训练了处理器与下游模型，还应把这看作一个新系统，与使用固定下游模型的对照分别报告。

当真实参考不可得时，可以保留原始RAW、多次曝光或其他视角作为额外证据，并记录生成补全区域。补出的皮肤、纹理和字符是否可接受，由用途决定；需要记录保真的任务应能追溯处理前后对应关系，不能以模型确信度替代丢失观测。

\subsection{质量与资源的取舍}
\label{sec:vision-restoration-resources}

复原比较应固定输入分辨率、退化、硬件、采样步数和候选数量。前馈网络与扩散流程的成本不只相差参数量，后者还需要多次迭代和可能的引导反向计算。图像压缩则同时统计实际文件码率、编码时间、解码时间与内存，不能只比较训练的估计码率。

表\ref{tab:vision-restoration-choice}把三类目的与所需证据对应起来。一个应用可以同时占据多行，例如先从RAW复原，再调整显示并压缩发送。最终判断仍需回到整个流程的输出，前段留下的错误可能被后段放大。

\begin{table}[htbp]
 \centering
 \small
 \begin{tabularx}{\linewidth}{@{}p{18mm}XX@{}}
  \toprule
  主要目的 & 需要保留的证据 & 共同记录的成本与条件\\
  \midrule
  内容复原 & 对齐参考、退化回投影、文字或测量结果 & 输入范围、模型版本、迭代与候选预算\\
  呈现增强 & 原始图像、亮部裁剪、颜色与局部层次变化 & 颜色空间、曲线、增益和局部参数\\
  存储传输 & 独立解码结果、实际文件、用途误差 & 辅助比特、元数据、编码及解码资源\\
  \bottomrule
 \end{tabularx}
 \caption{按处理目的选择证据。表中各行可以组合；每项判断都同时保留输入条件、处理结果与成本，避免只用一个质量分数作选择。}
 \label{tab:vision-restoration-choice}
\end{table}

\section{本章小结}
\label{sec:vision-restoration-summary}

让夜间招牌变得可用，需要区分看清、看着自然和记录正确。残差监督与显式反卷积利用不同依据估计内容，生成先验为不可辨识区域提供可能细节，RAW映射和色调调整又使用不同阶段的信息。观测丢失的细节不会因为输出更锐利而自动获得真实性。

保存结果还要经过量化、概率分配和真实熵编码，辅助信息必须同时进入码率与解码依赖。选择流程时，应把参考失真、感知变化、任务结果和资源放在同一组条件下比较。

进一步研究需要处理现实退化覆盖不足、生成细节缺少可核验依据，以及低码率下关键文字丢失的问题。它们分别指向数据建模、观测约束和用途相关编码。研究这些问题时，应明确减少了哪种错误、付出了多少成本，以及结论依赖哪类输入。平均感知质量上升若伴随关键字段错误增加，就需要改变目标或应用范围。下一章允许根据用户要求创建和修改内容，届时仍需明确哪些变化是目标、哪些内容必须保持。

\chapter{图像生成与编辑}
\label{chap:vision-image-generation}

“画一只玩具狗”和“把这只玩具狗旁边的杯子换成花瓶”都有许多合理结果，但需要保持的内容不同。前者主要服从创作条件，后者已有一幅原图，玩具狗的身份和未指定背景也属于约束。若要求根据绕玩具狗拍摄的照片生成背面视图，还要说明哪些部分受到实际观测支持。本章按这些变化关系，讨论新图像的创建、已有图像的转换、指定内容的编辑和观测场景的新视角合成。

生成模型通过训练获得图像分布或条件分布，再由采样形成候选；条件约束和原内容保持还需在任务中单独规定。通用训练与采样可参考\BookChapterRef{chap:pgm-generative}{卷六“生成模型”章}，图文表示和媒体条件生成接口已在\ref{chap:vision-multimodal-models}章建立。这里保留使用这些接口所必需的输入、条件、监督与输出过程，把重点放在图像任务特有的要求上。贯穿示例中的候选、分数和相机数值均为教学构造，不表示模型实测结果。

理解本章需要区分三种依据。提示告诉模型希望得到什么，参考图像提供某个外观或结构，相机观测则约束一个确实拍摄过的场景。把照片中的玩具狗放到雪地，属于允许改变背景的创作；从环绕拍摄的数据渲染另一视角，属于同一场景的拟合与呈现。两者都能输出新图像，但不能用后一任务的观测证据替前一任务生成的背面作真实性保证。我们先建立条件与结果之间的关系，再考察有输入对应和历史状态时怎样保持内容，最后把真实观测组织成可渲染的场景。

\section{新图像的创建}
\label{sec:vision-generation-creation}

创建新图像时，需要先确定要求，再选择能够表达这些要求的条件，最后生成并筛选候选。这些是创作过程的职责，不要求模型内部恰好存在三个同名模块。自由采样可以只规定分辨率和大致图像分布，有条件创建则还要核验条件是否真正进入结果。

\subsection{内容与成果要求的确定}
\label{sec:vision-generation-requirements}

考虑“红色玩具狗放在蓝色杯子左边”。其中至少有两个对象、各自的颜色、数量及左右关系。把它写成可检查记录，可以得到“一个玩具狗，红色；一个杯子，蓝色；狗在杯子左边”。这份记录供结果核验，不意味着生成器内部已经显式保存同样的对象表。

对象出现与属性绑定是两层要求。画面同时存在红、蓝两色，不能说明颜色分别属于正确对象；包含狗与杯子，也不能说明数量和空间关系满足。若还要求杯子上写“OPEN”，字符应作为单独条件检查，不能仅看整体图文相似度。

输出用途也会改变要求。插图允许多种合理材质和光照，产品展示可能需要固定主体形状，排版素材还要规定留白与裁剪范围。条件没有指定的部分可以变化，已经明确的部分则不能由“画面自然”替代。对多个可接受结果，应保留变化空间，而非把某一参考图的全部像素当作唯一答案。

\subsection{创作条件的组织}
\label{sec:vision-generation-conditions}

文字容易描述语义，空间图更直接地规定位置，主体参考则提供某个具体玩具的外观。它们约束不同对象，也可能相互冲突。选择接口时，应说明每项条件如何编码、在哪一步影响生成，以及冲突时如何判断结果。

\subsubsection{潜空间条件采样与unCLIP的应用入口}
\label{sec:vision-generation-latent-entry}

潜空间生成先在压缩后的图像表示中采样，再由图像解码器还原像素。使用文字条件时，文字编码器产生词元特征，去噪网络在采样过程中读取这些特征；随机初值、采样器、步数和引导强度共同决定一次输出。图像解码器负责把最终潜变量变成$\bm{X}\in\mathbb{R}^{h\times w\times3}$，它不会自动验证文字中的数量和位置\scite{pgm-generative-rombach2022}

unCLIP将CLIP图像表示作为图像生成的条件。由文字创建时，先验模型根据文字产生可能的CLIP图像向量，再由生成模型形成图像；给定参考图像时，也可以直接编码参考，绕过文字到图像表示的先验，生成具有相关语义的变体。两条路径共享的是图像表示接口，不是必须经过同一个输入过程\scite{vision-unclip2022}

这种全局向量压缩了很多信息。参考中的玩具狗与新结果可能同属相近外观，但精确耳朵形状、徽标和原背景不一定被完整保留。若需要逐区域保持，应增加空间或局部条件；若要求同一主体的新姿态，还需要针对身份进行适配和评价。

实际使用时固定模型、提示、分辨率与总候选数，例如对同一描述生成四个随机种子的结果，比较姿态、背景和细节。相同种子只有在模型、采样过程及实现条件一致时才具有可重复含义。提高步数、改提示或增加候选都会改变总预算，不能只呈现最终选中图片而省略这些条件。

一条实际样本从提示到图像，还要经过尺寸和状态的恢复。先把提示编码成文字序列，再按所选自编码器的压缩倍率确定目标潜网格；若高宽各压缩八倍，$512\times512$目标对应$64\times64$网格，但通道数仍由模型规定。以种子产生同形状高斯噪声，按预定噪声日程逐次调用条件网络，把当前状态更新成下一状态，直到日程终点。最后才解码、按已声明的颜色范围转换，并裁剪为目标尺寸。解码前的潜数组不是可交付RGB图片，改变输出长宽比也不能只在最后拉伸而声称采样条件没有变化。

unCLIP中的全局图像向量又不是上述空间潜网格。由文字得到该向量后，生成解码器从新的随机状态形成低分辨率图像，再由相应上采样阶段得到高分辨率输出；直接编码玩具狗参考时，替换的是先验的输出条件，不是把参考像素原样接到最终图像。使用者需要共同保存先验、生成解码器、上采样器及其预处理版本，不能只保留CLIP编码器便认为已具备生成路径。两条路径的参考信息与随机来源不同，应在候选记录中分别注明。

\subsubsection{ControlNet的空间条件支路}
\label{sec:vision-generation-controlnet}

文字“抬起左前爪”不容易精确指定每个关键点的位置。ControlNet把边缘、深度、分割或关键点图作为空间条件接入已有生成模型，使输入结构在去噪中提供约束\scite{vision-controlnet2023}
预训练主网络保留冻结参数，可训练支路由相应编码及中间模块复制而来，连接处使用零初始化卷积。

用一个模块说明其计算。设预训练模块为$F_{\bm{\theta}}$，复制支路为$F_{\bm{\theta}_{\mathrm{c}}}$，输入特征为$\bm{h}$，条件编码为$\bm{c}$，零初始化卷积为$Z_1,Z_2$，则示意输出为
\begin{equation}
 \bm{h}'=F_{\bm{\theta}}(\bm{h})+
 Z_2\!\left(F_{\bm{\theta}_{\mathrm{c}}}
       \bigl(\bm{h}+Z_1(\bm{c})\bigr)\right)\text{。}
 \label{eq:vision-controlnet}
\end{equation}
训练开始时，$Z_2$的输出为零，支路不扰动原模块结果。零输出不意味着整个支路永远没有学习信号：外层卷积的参数可先从其输入特征得到梯度，变为非零后再向复制模块传递梯度。图\ref{fig:vision-controlnet}用加法节点标明两条路径的关系。

\begin{figure}[htbp]
 \centering
 % Single-module sketch of ControlNet Eq.(2); no complete U-Net is implied.
\begin{tikzpicture}[x=1mm,y=1mm,font=\figurefont\small,
 b/.style={draw=BookBlue,fill=BookBlue!10,line width=.8pt,
  minimum width=16mm,minimum height=9mm,inner sep=1mm,align=center},
 op/.style={circle,draw=BookInk,fill=white,line width=.8pt,
  minimum size=5mm,inner sep=0},
 a/.style={-{Stealth[length=2mm]},BookInk,line width=.9pt,
  rounded corners=1mm}]
 \node at (5,0) (h) {$\bm{h}$};
 \node[b,draw=BookMuted,fill=BookPaper,minimum width=25mm] at (42,0)
  (base) {$F_{\bm{\theta}}$};
 \node[op] at (82,0) (sum) {$+$};
 \node at (101,0) (out) {$\bm{h}'$};
 \node[op] at (30,-25) (cond) {$+$};
 \node[b,minimum width=23mm] at (55,-25) (copy) {$F_{\bm{\theta}_{\mathrm{c}}}$};
 \node[b,minimum width=12mm] at (82,-25) (z2) {$Z_2$};
 \node at (5,-45) (c) {$\bm{c}$};
 \node[b,minimum width=12mm] at (30,-45) (z1) {$Z_1$};
 \draw[a] (h)--(base);
 \draw[a] (base)--(sum);
 \draw[a] (sum)--(out);
 \draw[BookInk,line width=.9pt] (14,0)--(14,-25);
 \draw[a] (14,-25)--(cond);
 \fill[BookInk] (14,0) circle (.6mm);
 \draw[a] (cond)--(copy);
 \draw[a] (copy)--(z2);
 \draw[a] (z2)--(sum);
 \draw[a] (c)--(z1);
 \draw[a] (z1)--(cond);
 \node[anchor=south,text=BookMuted] at (42,7) {冻结};
 \node[anchor=north,text=BookBlue] at (55,-32) {可训练副本};
 \node[anchor=north,text=BookBlue] at (82,-32) {从零连接};
\end{tikzpicture}

 \caption{ControlNet的模块级教学示意。主模块参数$\bm{\theta}$冻结，复制模块及零初始化连接可训练；相同输入进入两侧，条件只通过支路注入。零初始化让控制输出从零开始，训练后经加法调节主路径。完整网络在多个尺度接入控制特征，图中仅展开一处。}
 \label{fig:vision-controlnet}
\end{figure}

训练样本包含目标图像、文字和空间条件。可以从目标图像提取边缘或深度形成配对，给目标的潜变量加噪，再以条件噪声预测损失更新支路与连接。空间条件需要经过编码并对齐相应特征分辨率，不能直接把全尺寸深度图加到任意通道张量上。不同控制类型通常具有各自训练设置，使用时应加载与条件类型匹配的权重。

为把监督落实到一次更新，设目标图像编码为$\bm{z}_i$，条件图为$\bm{C}_i$，提示为$c_i$。在每个小批量中抽取噪声步$k_i$和同形噪声$\bm{\epsilon}_i$，构造$\bm{z}_i^{(k_i)}$，再让冻结主干与控制支路共同预测噪声。若小批量含$m$个样本，目标是
\begin{equation}
 L_{\mathrm{ctrl}}=
 \frac{1}{m}\sum_{i=1}^{m}
 \left\lVert\bm{\epsilon}_i-
 \epsilon_{\bm{\theta},\bm{\psi}}
  (\bm{z}_i^{(k_i)},k_i,c_i,\bm{C}_i)
 \right\rVert_2^2\text{，}
 \label{eq:vision-controlnet-training}
\end{equation}
其中$\bm{\psi}$收集复制支路、条件编码器和零卷积的可训练参数，$\bm{\theta}$不更新。平方和比较同一次抽样噪声的目标与预测；重新抽样噪声步使网络在多种当前状态上接受监督。反向传播得到$\nabla_{\bm{\psi}}L_{\mathrm{ctrl}}$，优化器据此更新$\bm{\psi}$，并在新的小批量上重复；训练预算及留出条件图的检查决定保留哪个检查点，不以某次噪声误差下降宣称所有姿态已学会。

零卷积开始学习的顺序可以手算。将式\eqref{eq:vision-controlnet}的外层连接暂化为标量$a q$，主模块输出为$f$，平方误差为$(f+a q-y)^2/2$。若$a=0,q=2,f=0.2,y=0.6$，则$\partial L/\partial a=(f-y)q=-0.8$；用步长0.1作一次教学更新得到$a=0.08$。该次反传对$q$的梯度是$(f-y)a=0$，下一次则不再必为零。这个计算说明外层连接为何能先启动，并不表示完整网络的优化必然按固定次数收敛。

推理同时给出文字和控制图，在每次去噪中注入相应支路特征；控制强度可对这些修正作缩放。固定玩具狗姿态图，改变文字可探索材质与背景，改变强度则改变对输入结构的依从程度。控制强度为零使这条修正消失，强度增加也不保证任意条件都能满足。

多条控制支路可以把修正相加，但若深度图规定前爪在后、姿态图又规定前爪伸向相机，两种条件并没有一致的几何解。应核验条件提取器本身，并记录模型版本、条件分辨率和选择预算。边缘相近不等于主体身份相同，空间控制还需要与下一节的主体要求配合。

\subsubsection{DreamBooth的主体绑定与类别先验保持}
\label{sec:vision-generation-dreambooth}

普通类别词“玩具狗”无法区分某个具体玩具。DreamBooth用少量同一主体的参考照片和带特殊标识的类别提示微调生成模型，让标识与该主体外观建立联系。例如提示可包含一个选定的稀有标识和“toy dog”，使用时再接“在雪地中”。这个标识是文字条件的一部分，不是逐像素复制指令\scite{vision-dreambooth2023}

只在少量主体照片上训练，模型可能连参考背景一起记住，还可能把一般的“玩具狗”类别词也缩窄成这一只。原方法由微调前模型生成一般类别样本，并用类别先验保持目标约束这一变化。类别先验保持（prior preservation）在这里是继续训练模型生成普通类别实例，减轻类别词只指向单个主体的变化。总目标为
\begin{equation}
 L=L_{\mathrm{subj}}+\lambda_{\mathrm{pr}}L_{\mathrm{class}}\text{，}
 \label{eq:vision-dreambooth}
\end{equation}
其中两项分别来自主体与普通类别样本，$\lambda_{\mathrm{pr}}\geq0$控制类别项的相对代价。原论文以预测干净图像的形式写去噪目标。设主体参考为$\bm{x}$，其提示条件为$c_{\mathrm{subj}}$，加噪状态为$\alpha_k\bm{x}+\sigma_k\bm{\epsilon}$，则主体项为
\begin{equation}
 L_{\mathrm{subj}}=
 \mathbb{E}_{\bm{x},k,\bm{\epsilon}}
 \left[
 w_k\left\lVert
 f_{\bm{\theta}}(\alpha_k\bm{x}+\sigma_k\bm{\epsilon},
                  k,c_{\mathrm{subj}})-\bm{x}
 \right\rVert_2^2\right]\text{。}
 \label{eq:vision-dreambooth-subject}
\end{equation}
$\alpha_k,\sigma_k,w_k$分别是噪声日程系数和训练权重，图像或潜表示的维度须与所选基础模型一致。类别项把$\bm{x}$与$c_{\mathrm{subj}}$替换成预先生成的类别图像$\bm{x}_{\mathrm{pr}}$与普通提示$c_{\mathrm{pr}}$，另抽噪声和噪声步。这里的类别参考来自冻结原模型，不是拍摄到的这一只玩具的新视角。

每次训练抽取主体与类别样本，分别计算两项，再把梯度加权合并后更新模型。参数从已有文字生成模型初始化，而不是从这几张照片训练一个新生成器；特殊标识还需经分词器检查，随机拼一串字符不保证它被当成单个稀有词元。训练结束保留微调权重、标识、类别提示及采样配置。应依据事先留出的主体照片或提示检查身份、背景复制和普通类别生成，并设有限训练预算，避免用无限拟合主体照片作为停止规则。

若基础模型输出噪声而不是干净图像，两种目标需要换算。在$\alpha_k,\sigma_k>0$时，由$\hat{\bm{x}}=(\bm{z}^{(k)}-\sigma_k\hat{\bm{\epsilon}})/\alpha_k$可得
$\lVert\hat{\bm{x}}-\bm{x}\rVert_2^2=(\sigma_k^2/\alpha_k^2)\lVert\hat{\bm{\epsilon}}-\bm{\epsilon}\rVert_2^2$。因此保留相同干净图像目标需要同步变换噪声步权重，不能把式\eqref{eq:vision-dreambooth-subject}机械改写成无权噪声误差。端点处某个系数为零时须采用该实现的有效日程与数值约定。

训练后保留微调参数，输入“该标识的玩具狗在雪地中”便可采样，不再需要同时提供目标雪地图像。类别样本约束普通类别语义，主体样本约束具体外观，两者提供不同训练信息。增大类别项并不会补齐主体照片从未展示的背面，也不能保证特殊标识在所有姿态下都保持细节。

核验时应把新结果与多张参考并排比较耳朵、斑纹和形状，另看是否只是复制了原背景。如果只有正面参考，生成背面的合理纹理仍可能来自类别先验。多个主体各自保持又增加了标识绑定与相互遮挡问题，不能从单主体成功直接推出多主体能力。

\subsubsection{参考图像的角色绑定与条件冲突}
\label{sec:vision-generation-reference-roles}

两张参考同时出现时，生成器还需要知道每张图负责什么。假设第一张展示玩具狗的耳朵和斑纹，第二张展示杯子的蓝釉外观。若只写“参考这两张图”，结果可能把杯子的蓝色放到玩具狗身上，或连第一张的桌面一起复制。我们可以先建立“参考1对应玩具狗身份，参考2对应杯子材质；布局由当前提示指定”的记录，再分别检查对象、属性和位置。这个角色绑定（reference role binding）是来源与目标之间的明确对应，不是假设模型内部已经产生了可靠的对象表。

主体微调、直接参考与角色规定承担不同职责。DreamBooth把主体信息写入适配参数，推理主要读取标识和提示；直接参考接口则在这次生成中读取图像，不必为每个主体更新权重。角色规定告诉系统本次应该取用参考的哪些属性，但只有所选模型支持相应图像接口时才可能执行。它不等于为一个只能接收单图的模型增添任意多图能力。图像与文字怎样编码成共同上下文，以及多轮如何再次取用图像，将在\ref{sec:vision-generation-kontext-history}节完整展开。

冲突也要在取样前暴露。若参考2是蓝杯子，而提示要求同一杯子为红色，应明确本轮允许改变颜色、只保留杯型；若两个参考分别要求杯子在狗的左右两侧，则必须选定最终关系。保留属性与可变属性分别记录后，才能判断输出中的颜色改变是成功编辑还是串换。参考未展示的属性仍来自生成先验，精确文字还需读取核验；不以整体参考相似度代替这些判断。

\subsection{图像生成与候选选择}
\label{sec:vision-generation-selection}

给定相同条件，采样过程可以主动改变生成方向，生成之后也可以从候选中挑选。这两种操作都可能使用图文相似度，却在不同时间、对不同对象起作用，成本也不同。

\Needspace{60mm}
\subsubsection{GLIDE的CLIP过程引导与无分类器引导对照}
\label{sec:vision-generation-glide-guidance}

GLIDE比较了噪声感知CLIP引导与无分类器引导。前者训练可读取当前带噪图像及噪声步的图像编码器，让文字相似度的梯度在每步修正采样；它使用的不是只在清晰图片上训练的原始CLIP\scite{vision-glide2022}

设当前反向转移的均值与协方差为$\bm{\mu}^{(t)},\bm{\Sigma}^{(t)}$，噪声感知图文相似度为$s(\bm{X}^{(t)},c,t)$，则均值修正具有形式
\begin{equation}
 \tilde{\bm{\mu}}^{(t)}
 =\bm{\mu}^{(t)}
  +a\,\bm{\Sigma}^{(t)}
   \nabla_{\bm{X}^{(t)}}s(\bm{X}^{(t)},c,t)\text{。}
 \label{eq:vision-clip-process-guidance}
\end{equation}
$a$为引导强度，梯度针对当前带噪图像，目的在于把后续采样推向文字条件。过程需要额外图文网络与梯度计算；将普通CLIP直接接入相同位置，是另一种分布和方法设置，不能沿用GLIDE对噪声感知模型的比较结论。

无分类器引导不运行独立分类器。训练以一定概率去掉文字，使同一去噪网络能分别预测有条件与无条件噪声；使用时按
$\hat{\bm{\epsilon}}=\bm{\epsilon}_{\varnothing}+b(\bm{\epsilon}_{c}-\bm{\epsilon}_{\varnothing})$
组合，$b=1$退回普通条件预测，$b>1$沿条件差异外推。它仍需要约定的条件与无条件预测计算，并不是没有额外推理成本。

固定玩具狗提示与噪声种子改变引导强度，可以观察主体、背景和多样性怎样变化。更强引导可能使某些条件更突出，也可能改变色彩和构图。对“两个杯子”“蓝色杯子在右侧”等组合条件，仍要逐项核验，不能把相似度上升解释为每项要求都更正确。

\subsubsection{图文相似度重排序与条件筛选}
\label{sec:vision-generation-reranking}

重排序发生在候选图像已经形成之后。对$k$幅图像$\bm{X}_1,\ldots,\bm{X}_k$，计算单位图文向量的相似度$s_i$并排序，它不会修改候选像素，也不能创造候选中不存在的正确解。GLIDE对照中的DALL-E使用了这类生成后CLIP筛选；不同对照还使用了不同总候选量，应连同选择协议阅读。

表\ref{tab:vision-generation-candidates}给出四个构造候选。若仅取最高全局相似度，会选择数量错误的A；若先筛选数量、属性归属与关系，再按约定标准排序，则只有D满足全部要求。这里的分项检查来自人工或独立检测、读取工具，不是CLIP分数已经提供了这些保证。

\begin{table}[htbp]
 \centering
 \small
 \begin{tabularx}{\linewidth}{@{}c c X c@{}}
  \toprule
  候选 & 构造相似度 & 按请求核验的主要结果 & 保留\\
  \midrule
  A & 0.34 & 出现两只玩具狗 & 否\\
  B & 0.32 & 狗为蓝色，杯子为红色 & 否\\
  C & 0.30 & 属性正确，但狗在右侧 & 否\\
  D & 0.28 & 数量、颜色归属与左右关系均满足 & 是\\
  \bottomrule
 \end{tabularx}
 \caption{全局相似度与逐项条件检查的教学对照。分数人为设置，用于展示两种选择规则的区别；不代表任何模型的实测排序。}
 \label{tab:vision-generation-candidates}
\end{table}

如果一个请求生成了100幅图，人工筛出1幅，报告时应保留100幅的生成成本和筛选规则。按最佳候选计分与按一次随机生成计分衡量的是不同系统。没有候选满足硬条件时，应返回条件未满足或继续按明确预算修改，而不是把最高分候选自动称为成功。

\subsection{隐含知识条件与生成后的反思修订}
\label{sec:vision-generation-unireason-refinement}

给出“杯子受侧面光照，红色玩具狗在其左边”，除了画出两个对象，还要落实相对位置和光照造成的阴影。把提示扩写成更长文字，可以明确要求，却看不到实际生成是否遵守。生成后的反思修订（reflection and refinement）增加一个反馈环节：先查看草图与原要求的差异，再把差异写成明确编辑请求，生成修订图像。UniReason 1.0采用这一思路，把知识增强的文字推理与编辑式修订接在同一生成流程中；这里按2026年2月v4描述应用，不重复\ref{chap:vision-multimodal-models}章的联合训练\scite{vision-unireason-2026}

本地恢复这一接口时，需要四类量：用户条件$c$、初次细化的文字$r^{(0)}$、草图$\bm{X}^{(0)}$，以及看到草图后产生的修改说明$r^{(1)}$。模型先依据$c$形成$r^{(0)}$并采样草图，再读取$c,r^{(0)},\bm{X}^{(0)}$形成反馈，把反馈与草图作为编辑条件得到$\bm{X}^{(1)}$。文字输出沿前置章的词元预测接口产生，图像输出沿条件潜空间流生成并解码。图像不是对文字解释的可视化保证；第二次文字必须依据实际草图指出哪个对象、哪个属性或哪种关系需要改变，不能只再说一遍原提示。

这种能力来自相应训练数据，而不是对任意生成器添加一句“反思一下”便可获得。论文的反馈训练样本保存草图、诊断指令与修订目标；在联合训练中，这些量分别充当已给上下文和需预测输出。训练数据的外部诊断器与修订教师负责构造监督，交付模型使用时不必再次调用同一组教师。反思与图像生成各自的目标和参数更新在前置模型章说明，本节使用训练后的接口进行一次修订。

沿玩具狗例子，若草图出现两只狗且杯子在左侧，反馈可写成“仅保留一只红色玩具狗；把蓝杯子移到狗右侧，保持狗的斑纹”。如果草图数量正确，仅阴影与光源矛盾，则不应重新创建全部主体。比较直接生成、加入文字细化以及再加一次修订时，固定原条件并保存三个结果；分别检查数量、颜色归属、左右关系和阴影，再检查自然度。反馈文字全部正确而修订仍画错关系，应计为图像实施失败，不算推理已经解决了任务。

原论文实现采用一次修订，不构成无限自我改善的保证。修订可能纠正数量，又引入主体漂移；发现这种交换时，应拒绝修订或按有限预算局部重做。额外文字、草图生成、图像读取及第二次采样都计入总成本。采用与训练诊断相同的模型裁判验收时，还应加入独立的区域、字符或人工检查，避免只证明结果更迎合这一裁判。

\section{已有图像的转换}
\label{sec:vision-generation-translation}

转换以已有图像为起点，获得另一种风格、呈现或成像形式。它的核心是建立输入与输出的对应关系；训练是否有配对参考，决定了哪些对应可直接监督。空间布局到照片也能使用同样的条件映射工具，本节借此解释机制，具体创作要求仍由\ref{sec:vision-generation-creation}节规定。

\subsection{转换关系与保持要求的规定}
\label{sec:vision-generation-translation-requirements}

把玩具狗照片转成绘画时，纹理和笔触可以改变，姿态和对象数量可能需要保持。把一种成像方式转换为另一种时，还可能要求每个位置与实际结构对应。两种用途即便采用相同网络，允许的误差也不同。

数据应与这种要求相符。对齐的照片—绘画配对可约束具体样本关系，独立的照片集与绘画集只说明两个域中分别有哪些外观。若把没有逐位置对应的图像直接作像素监督，损失会强迫网络拟合无意义的差异。

\subsection{转换映射的获得与使用}
\label{sec:vision-generation-translation-mapping}

现在固定输入和需要保持的对应，考察如何获得映射。对齐图像对可以直接给出每个样本的目标；没有配对时，只能分别约束两个域的外观和往返关系；固定风格图则给出纹理统计而非另一幅逐位置目标。三种监督不能相互替代，训练得到的转换器也有不同的使用条件。

\subsubsection{pix2pix的配对条件映射与局部判别}
\label{sec:vision-generation-pix2pix}

pix2pix使用对齐图像对$(\bm{X},\bm{Y})$学习条件映射。生成器$G$读取$\bm{X}$并输出目标形式，判别器$D$同时读取输入与候选输出，从而判断它们是否像真实配对。原方法的U形生成器用跳连把较细的位置特征传到解码端，局部判别器则在多个图像块上提供真假反馈\scite{vision-pix2pix2017}

以边缘图到玩具狗照片为例，$\bm{X}$提供轮廓，$\bm{Y}$提供该轮廓对应的颜色与纹理。标准条件对抗目标为
\begin{equation}
 \begin{aligned}
 L_{\mathrm{cGAN}}(G,D)
 &=\mathbb{E}_{\bm{X},\bm{Y}}[\log D(\bm{X},\bm{Y})]\\
 &\quad+\mathbb{E}_{\bm{X}}[
       \log(1-D(\bm{X},G(\bm{X})))],\\
 L_1(G)&=\mathbb{E}_{\bm{X},\bm{Y}}
       [\lVert\bm{Y}-G(\bm{X})\rVert_1]\text{。}
 \end{aligned}
 \label{eq:vision-pix2pix}
\end{equation}
训练交替更新判别器与生成器。判别器最大化上述对抗项；生成器在实际实现中改用最大化$\log D(\bm{X},G(\bm{X}))$的非饱和目标，并最小化加权像素误差$\lambda L_1$。这种替换在判别器容易辨认假图时提供更有用的梯度。这里省略随机变量，强调输入—目标对应；原工作发现显式噪声容易被忽略，最终在训练和使用时均保留dropout，但输出变化仍有限，不能据此把生成器当成具有充分多样性的条件采样器。

实际训练先同时裁剪和翻转输入—目标，保证配对位置不变，再把像素变成同一数值范围。U形生成器的下采样形成粗尺度特征，解码时把相同尺度编码特征沿通道拼接，以传递轮廓位置。PatchGAN（局部块判别器）不只输出一个整图真假值，而在多个感受野上输出配对是否像真实样本的分数。原论文主要使用$70\times70$感受野；这个尺寸是判别时看到的邻域，不是先把图像切成互不相交的70像素块。

写经验目标时应在每个局部分数上计算对数损失再平均。设一对图像产生$q$个真实配对分数$d_j^{\mathrm{real}}$和假配对分数$d_j^{\mathrm{fake}}$，则判别损失为
$-\frac1q\sum_j[\log d_j^{\mathrm{real}}+\log(1-d_j^{\mathrm{fake}})]$。它与先平均分数、再取对数不同。更新判别器时把生成结果视为固定输入；更新生成器时冻结判别参数，但保留判别输出对生成图像的梯度。两网从随机权重初始化，用小批量交替执行这两步，按预算与留出转换结果保存检查点；对抗训练没有由该目标自动提供的全局收敛保证。

用两个局部块作教学计算，若真实分数为$(0.8,0.6)$，假分数为$(0.2,0.4)$，判别损失为
$-\{\log0.8+\log0.6+\log0.8+\log0.6\}/2\approx0.734$，这里采用自然对数。生成器的非饱和对抗项为$-(\log0.2+\log0.4)/2\approx1.263$。若两个目标值为$(0.2,0.8)$、输出为$(0.3,0.6)$，未归一化$L_1$和为0.3；教学权重取10时，总生成损失为$1.263+10\times0.3=4.263$。整图实现必须说明像素项求和还是求平均，相应权重不能直接混用。

像素项把输出拉向这幅配对参考，局部判别项鼓励目标域的局部外观。两者都不会自行修正训练对中的错位。若每个真实轮廓总与向右偏移的照片配对，网络可能学到这种偏移；若多个合理纹理共享同一输入，像素目标也可能倾向平均。局部真实还不保证整幅图像的几何和身份正确。

使用时只调用生成器，不再需要判别器或目标参考。对于照片到另一成像形式的转换，应另外检查真实对应与诊断所需结构，不能把训练域中看似逼真的输出解释成新的测量观测。

原pix2pix的使用约定还保留dropout，并用当前测试批次计算批归一化统计，而不是默认读取训练累计统计。因此复现时应保存批量与归一化方式，同一输入的随机性也需记录。输出由解码器恢复到输入网格，逆转数值归一化后得到目标图像；若前处理裁剪了图像，恢复全图坐标或拼接也必须另作说明。不存在目标参考不妨碍转换，却意味着不能在使用时借参考来挑选最佳结果而不计入选择协议。

\subsubsection{CycleGAN的无配对双向映射}
\label{sec:vision-generation-cyclegan}

只有照片域$A$与绘画域$B$的独立样本时，CycleGAN训练$G:A\to B$和$F:B\to A$两个方向，并为两个域分别设置判别器。目标域判别要求输出像该域图像，循环一致性则要求转换后还能回到输入\scite{vision-cyclegan2017}

循环损失为
\begin{equation}
 L_{\mathrm{cyc}}=
 \mathbb{E}_{\bm{x}\sim p_A}
 [\lVert F(G(\bm{x}))-\bm{x}\rVert_1]
 +
 \mathbb{E}_{\bm{y}\sim p_B}
 [\lVert G(F(\bm{y}))-\bm{y}\rVert_1]\text{。}
 \label{eq:vision-cycle-consistency}
\end{equation}
训练把两个域的对抗项和加权循环项结合。照片$\bm{x}$经$G$生成绘画，再经$F$返回照片，输入本身就是往返监督，无需知道它应对应哪幅真实绘画。需要额外颜色保持时，还可在特定设置中加入输入已属于目标域时少改动的约束，但它不等于逐样本跨域配对。

对一次更新，从两个域独立抽$\bm{x},\bm{y}$；它们不需要是同一只玩具，也不能假装配成逐像素目标。计算$G(\bm{x}),F(\bm{y})$及两条往返，再同时更新两个生成器。原方法采用最小二乘对抗损失：以目标域$B$为例，判别器最小化
$\frac12\mathbb{E}_{\bm{y}}[(D_B(\bm{y})-1)^2]+\frac12\mathbb{E}_{\bm{x}}[D_B(G(\bm{x}))^2]$，
生成器相应最小化$\mathbb{E}_{\bm{x}}[(D_B(G(\bm{x}))-1)^2]$，再加循环项。另一个域交换$G,F$与$A,B$执行相同计算。参数随机初始化，用小批量交替更新生成器和判别器；原训练还用近期生成图像的缓冲池更新判别器，以减小它只追随最新生成器的变化。

假设一幅输入的两个像素为$(0.2,0.8)$，某轮往返为$(0.3,0.7)$，该方向未归一化循环误差是0.2。若两个方向的往返都准确，循环项为零，仍要看目标域判别是否接受中间图像；反过来，中间图像像绘画也不能说明回到原输入。两类误差约束不同对象，缺一项都不能由另一项代替。训练结束按所需方向保留生成器、输入范围与域说明，并用未参与训练的输入检查转换，不把训练往返误差最低当成语义保持的充分条件。

往返可恢复并不能唯一确定正确语义。假设两个域各只有一幅狗图和一幅猫图，且概率都为$1/2$。$G$把狗映成猫、猫映成狗，$F$执行逆映射，两个输出分布都正确，循环误差也为零，但对象语义完全交换。图\ref{fig:vision-cycle-ambiguity}将这个反例与保持类别的映射并排显示。

\begin{figure}[htbp]
 \centering
 % Both maps are bijections of two equiprobable categories.
% Target order is changed in the right panel to avoid line crossings.
\begin{tikzpicture}[x=1mm,y=1mm,font=\figurefont\small,
 b/.style={draw=BookBlue,fill=BookBlue!8,line width=.8pt,
  minimum width=12mm,minimum height=8mm,inner sep=1mm},
 a/.style={{Stealth[length=1.7mm]}-{Stealth[length=1.7mm]},
  BookInk,line width=.9pt}]
 \node[font=\figurefont\small\bfseries] at (20,14) {保持类别};
 \node[font=\figurefont\small\bfseries] at (79,14) {交换类别};
 \node[anchor=south,text=BookMuted] at (5,5) {域$A$};
 \node[anchor=south,text=BookMuted] at (35,5) {域$B$};
 \node[anchor=south,text=BookMuted] at (64,5) {域$A$};
 \node[anchor=south,text=BookMuted] at (94,5) {域$B$};
 \node[b] at (5,0) (ad) {狗};
 \node[b] at (35,0) (bd) {狗};
 \node[b] at (5,-17) (ac) {猫};
 \node[b] at (35,-17) (bc) {猫};
 \node[b] at (64,0) (ad2) {狗};
 \node[b] at (94,0) (bc2) {猫};
 \node[b] at (64,-17) (ac2) {猫};
 \node[b] at (94,-17) (bd2) {狗};
 \foreach \u/\v in {ad/bd,ac/bc,ad2/bc2,ac2/bd2} {
  \draw[a] (\u)--node[above=1mm] {$G/F$}(\v);
 }
 \node[anchor=north] at (20,-25) {循环误差为0};
 \node[anchor=north] at (79,-25) {循环误差也为0};
\end{tikzpicture}

 \caption{循环一致性不能唯一确定语义对应。两域各有等概率的狗、猫两项。左侧保持类别，右侧交换类别；每条往返路径都回到原项，且目标域边缘分布相同。反例是离散教学构造，用来说明循环与域判别约束仍有多解。}
 \label{fig:vision-cycle-ambiguity}
\end{figure}

使用照片到绘画方向时，只运行$G$；反向生成器和两个判别器是训练职责，不需要每次转换都执行完整循环。对玩具狗仍应检查姿态、身份和文字，对测量或医学成像则需要更强的配对及事实依据。无配对域转换可以形成视觉风格，不能凭循环重建为未观测结构提供真实性。

\Needspace{45mm}
\subsubsection{感知损失监督的固定风格转换}
\label{sec:vision-generation-perceptual-style}

若目标是一幅指定风格图的纹理，可以用固定视觉网络比较内容与风格，再训练一个前向转换网络\scite{vision-perceptual2016}
Johnson等人的方法区分图像变换网络与损失网络：前者生成图像并更新参数，后者提供特征度量并保持固定。

设变换网络输出$\widehat{\bm{X}}=G_{\bm{\theta}}(\bm{X})$，固定损失网络第$\ell$层的特征为$\bm{F}_\ell(\bm{X})\in\mathbb{R}^{c_\ell\times n_\ell}$。内容项比较对应位置特征，风格项比较通道之间的相关统计。令
\begin{equation}
 \bm{\Gamma}_\ell(\bm{X})
 =\frac{1}{c_\ell n_\ell}
   \bm{F}_\ell(\bm{X})\bm{F}_\ell(\bm{X})^\top\text{，}
 \label{eq:vision-style-gram}
\end{equation}
则可组合输出与输入的特征平方差，以及输出与固定风格图$\bm{S}$的
$\lVert\bm{\Gamma}_\ell(\widehat{\bm{X}})-\bm{\Gamma}_\ell(\bm{S})\rVert_{\mathrm{F}}^2$。归一化方式会影响损失权重，必须与采用的实现一致。

Gram矩阵对空间位置作求和，保留通道共同出现的统计，却不直接记录这些纹理出现在何处。例如单通道两位置特征$(1,0)$与$(0,1)$产生相同Gram值，位置已经交换。内容项因此承担保留较高层结构的职责，风格项则鼓励某种纹理统计；二者的权重决定输出怎样折中。

沿相同玩具狗输入，若某层有两个通道、两个位置，内容特征为
$\bm{F}=\left(\begin{smallmatrix}1&0\\0&2\end{smallmatrix}\right)$，
则$\bm{\Gamma}=\operatorname{diag}(1/4,1)$。把两个位置对调得到
$\bm{F}'=\left(\begin{smallmatrix}0&1\\2&0\end{smallmatrix}\right)$，
Gram矩阵仍相同，但对应内容特征平方差和为10。若内容项按$c_\ell n_\ell=4$归一化，其值为2.5。这一复算同时说明风格项看通道统计，内容项才读取相应位置；它不是在特征中逐像素保存所有细节。

固定风格转换的总体目标可写成内容、风格和总变差三项的加权和。内容项比较输出与输入的指定层特征，风格项在多个指定层比较输出与固定$\bm{S}$的Gram矩阵，总变差项惩罚输出相邻像素的过大跳变。每个小批量只抽内容图，先经可训练变换网络产生输出，再让固定损失网络分别处理输出与内容图；固定风格特征可预先计算。反传仍须经过损失网络对输出的导数，但不更新损失网络参数，再用所得梯度更新变换网络。初始化、学习率和停止预算属于该训练实现，实际采用的层与权重应共同保存。

训练在许多内容图上更新$G_{\bm{\theta}}$，固定同一风格图，之后使用时只需一次生成器前向，不必对每张新图重新优化像素，也不再需要运行损失网络。这个原始固定风格路线不能直接解释成“推理时任意输入一幅风格图即可转换”。玩具狗的笔触可改变，徽标和细小字符是否保留仍需另查。

\subsection{转换结果与对应关系的核验}
\label{sec:vision-generation-translation-check}

转换结果应同时具有目标形式和规定的输入对应。照片仍像照片，说明风格可能未充分改变；绘画自然却把狗变成猫，则说明语义保持失败。把图像恢复成输入的循环误差很小，只能回答往返约束，不能代替中间结果的检查。

固定输入后，可沿对象、轮廓、文字与颜色逐项记录变化，按用途区分允许和不允许的改变。配对训练出现系统性错位时，优先检查数据；没有配对依据却要求严格成像保真时，需要增加相应监督或收窄用途。下一节的编辑将变化要求交给当前请求，而不是预先固定一个输入输出域关系。

\section{指定内容的编辑}
\label{sec:vision-generation-editing}

编辑给定原图，并通过当前请求规定变化。它可以改局部，也可以改全图，但都需要检查修改是否完成、应保留内容是否仍在。语言描述目标，掩码表达范围，参考图像提供期望外观，这几种输入的作用应分别说明。

\subsection{修改目标与范围的确定}
\label{sec:vision-generation-editing-scope}

“把玩具狗旁边的杯子换成花瓶”至少包含待替换对象、替换后的内容和周围保留物。若画面有两个杯子，仅给出对象类别还不足以确定操作范围。用户可以直接给出框、掩码或参考编号，也可从\ref{chap:vision-image-recognition}章的检测与分割候选中选择；这样的位置输入已经足以进入编辑计算。更复杂的自动表达定位由\ref{chap:vision-multimodal-understanding}章展开，不是当前编辑算法必须先解决的前提。

沿用上一章的掩码约定，$\bm{M}=1$表示保留区，$1-\bm{M}$表示允许生成区。若花瓶比原杯子高，只遮住原杯子轮廓可能没有容纳新物体的空间，还可能留下杯子原来的阴影。合理范围应覆盖允许变化的接触区和阴影，并明确玩具狗是否可被遮挡。掩码是当前任务要求，不是生成器自动推断出的正确编辑边界。

文字还可能缺少形状或动作的准确位置。视觉指令（visual instruction）把箭头、圈选、骨架或草图作为请求的一部分，使位置和几何意图不只依赖句子描述。VIBE（Visual Instruction Benchmark for Image Editing）按指示定位、形态修改和因果关系三层组织这类评价；这里采用2026年5月v2。它提供编辑任务和判分材料，不是一个新的编辑网络\scite{vision-vibe-2026}

三层改变的是视觉提示的含义。圈住杯子并写“换成花瓶”，圈选用于选择目标；画出更高花瓶轮廓，则草图用于规定形态；画出侧向光线箭头，则要求图像表现相应的受光与阴影后果，而不是在终图保留箭头。若文字相同而箭头分别指向左右两个杯子，应检查被替换对象是否相应改变。若只看整图出现花瓶，两次请求可能都得高分，却没有检验视觉指令是否发挥作用。

位置也须跟随前处理。原图为$1000\times800$像素，目标点为$(600,400)$，若等比缩小一半，点应变为$(300,200)$；若再从缩小图左侧裁去50像素，模型输入点为$(250,200)$。输出恢复时逆向补回裁剪偏移并还原比例，掩码与草图采用相同变换。只缩放图像而保留原坐标，会把目标送到错误位置。提示标记应保留在条件记录或独立图层，交付时检查它们未被当成场景中的真实纹理。

VIBE使用任务专用问题让多模态模型裁判判断位置、形状与关联变化。裁判读取了哪些提示、原图和输出应明确记录，并用人工区域、关键点或几何检查对照其判断。视觉提示帮助限定请求，不证明生成结果已满足形态，更不证明模型会按真实物理规律计算；最后一类后果将在\ref{sec:vision-generation-physics-check}节单独检验。

\subsection{修改的实施与原内容保持}
\label{sec:vision-generation-editing-methods}

固定原图、文字与允许范围后，单轮方法需要把“利用原内容”和“实施新变化”送入生成过程。掩码修补显式提供可见区，指令编辑则从编辑配对学习应改什么。两者都形成完整候选，但其输入约束不同，不能把没有手工掩码解释为自动严格保持所有未提及像素。

\subsubsection{GLIDE的图像与掩码条件修补}
\label{sec:vision-generation-glide-inpainting}

GLIDE的修补模型在当前带噪RGB之外，还读取未遮挡图像与掩码。训练时随机擦去图像区域，把剩余RGB和一个掩码通道作为额外条件；新增输入通道的权重从零开始，微调模型去噪恢复完整训练图像。于是模型在每步都能读取清晰的保留区域，而不只是当前带噪结果。

训练一条样本时，从图文对抽完整图像$\bm{X}$及提示$c$，抽随机保留掩码$\bm{M}$，构造可见条件$\bm{V}=\bm{M}\odot\bm{X}$，并另给完整目标加噪。当前带噪图像有3通道，可见图像有3通道，掩码有1通道，输入卷积因此读取相应的7通道数组。零初始化的是新增加的输入权重，原有生成权重从文字生成检查点继承；随后以完整目标的噪声预测误差微调。缺口中可见图像的零值与真实黑色由掩码区分，不能只提供$\bm{V}$就假定模型知道哪处缺失。

训练重复抽掩码、噪声步和噪声，使模型读取不同可见范围；使用时则固定用户指定的保留区与文字，噪声状态按采样日程更新，可见条件不随每次噪声状态一起被改写。最终通过生成和上采样得到目标尺寸图像。训练时恢复真实完整图像，为修补建立条件能力；使用时文字要求花瓶，属于改变目标内容的创作应用，不能把它说成训练中观测到了这个花瓶。

这一结构与“按普通生成模型采样，再覆盖已知像素”不同。后一做法可以在输出端强制保留某些像素，但没有因此让模型在训练中学会如何利用清晰上下文。GLIDE通过显式条件训练连接保留区域和目标内容；其上采样模型还读取完整低分辨率图像，以及高分辨率的未遮区域，输入职责也有所区别。

使用时给出原图、杯子及可变周边的掩码、花瓶描述与随机初值，反复去噪得到候选。应检查花瓶是否替换到正确位置，边缘、反光和投影是否与场景协调，玩具狗是否保持。若只要求恢复遮挡前的杯子，则属于\ref{sec:vision-restoration-inpainting}节的复原目标；这里的“花瓶”明确允许创造原图中不存在的内容。

掩码条件提供范围信息，但能否严格保持范围之外的每个像素还要看实现和后处理。硬性保留时可按已声明掩码回填原图，随后重新检查边界；这种回填是额外操作，不能把它造成的逐像素相等归为生成模型的能力。

\Needspace{60mm}
\subsubsection{InstructPix2Pix的合成编辑配对与双条件引导}
\label{sec:vision-generation-instruct-pix2pix}

用户有时只给一句指令，没有手工掩码。InstructPix2Pix学习从原图和编辑指令直接生成结果，监督数据采用合成三元组“原图、指令、编辑后图像”。语言模型根据原描述生成编辑指令与新描述，再用图像生成模型把两份描述转成对应图像\scite{vision-instruct-pix2pix2023}

独立生成两幅图会同时改变姿态、背景等大量内容，因而不适合作为局部编辑配对。原方法借助Prompt-to-Prompt思想，在两次生成中共享噪声和部分注意力过程，鼓励对应；其补充材料所述实现对所有编辑采用早期若干步的自注意力权重替换。随后按图文表示中的变化方向等规则筛选配对。共享计算降低了无关变化，却没有保证每个合成指令都被正确执行。

训练时将目标图像编码并加噪，得到$\bm{z}^{(t)}$，将原图编码为$\bm{z}_{I}$，两者沿通道拼接输入去噪网络，文字编码器则读取指令。噪声预测损失为
\begin{equation}
 L=\mathbb{E}\!\left[
 \left\lVert\bm{\epsilon}-
 \epsilon_{\bm{\theta}}
   \bigl(\bm{z}^{(t)},t,\bm{z}_I,c_T\bigr)
 \right\rVert_2^2\right]\text{。}
 \label{eq:vision-instruct-training}
\end{equation}
$c_T$是指令条件，$\bm{\epsilon}$是加到目标潜变量上的噪声。训练从预训练生成模型初始化，新增图像条件通道的权重初始化为零，并随机丢弃部分图像或文字条件，使模型学会不同条件组合下的预测。

原图条件在训练与使用中都保留干净潜表示，目标潜表示才按当前噪声步变化。若基础编码器产生4通道空间潜变量，输入卷积把两者拼成8通道，文字仍通过文字条件接口进入，不能把指令字符串直接拼进像素通道。每次抽取合成三元组，编码原图和目标，给目标加噪并预测噪声；用小批量平均式\eqref{eq:vision-instruct-training}的误差反传，更新编辑去噪网络。图像自编码器和文字编码器保留预训练参数，训练结束保存编辑权重与这些固定组件及预处理约定，再按留出的编辑请求检查修改和保持。

使用时不需要目标图像，也不需要用户写出完整的新场景描述。输入原图与“把杯子换成花瓶”，从随机目标潜变量开始采样。记同一步的三种噪声预测为
$\bm{\epsilon}_0=\epsilon_{\bm{\theta}}(\bm{z}^{(t)},t,\varnothing,\varnothing)$、
$\bm{\epsilon}_I=\epsilon_{\bm{\theta}}(\bm{z}^{(t)},t,\bm{z}_I,\varnothing)$
和$\bm{\epsilon}_{I,T}=\epsilon_{\bm{\theta}}(\bm{z}^{(t)},t,\bm{z}_I,c_T)$，双条件引导为
\begin{equation}
 \hat{\bm{\epsilon}}
 =\bm{\epsilon}_0
  +s_I(\bm{\epsilon}_I-\bm{\epsilon}_0)
  +s_T(\bm{\epsilon}_{I,T}-\bm{\epsilon}_I)\text{。}
 \label{eq:vision-instruct-guidance}
\end{equation}
第二项强调图像条件，第三项在图像条件已给定时强调指令带来的变化。它不是把三个独立概率相加；这些向量是去噪计算的输出。

以两个分量示意，令$\bm{\epsilon}_0=(0.1,0.2)$、$\bm{\epsilon}_I=(0.2,0.1)$、$\bm{\epsilon}_{I,T}=(0.4,-0.1)$，取$s_I=1.5,s_T=5$，结果为$(1.25,-0.95)$。若两项强度均为1，代数上就恢复$\bm{\epsilon}_{I,T}$。改变强度会改变整条采样轨迹，不能从这个线性组合推出某个对象的修改量也线性变化。

较强图像条件通常有助于保留输入结构，较强指令条件会强化修改，但这种取舍需要对当前输入核验。原方法没有显式编辑掩码，因此未指定区域也可能改变。数量、空间关系和重复编辑的错误，还可能继承自合成数据及其筛选器。

\subsection{连续编辑中的身份保持与累积漂移}
\label{sec:vision-generation-kontext-history}

一次杯子替换之后，用户还可能要求“把花瓶改成蓝色，再移到右边”。这时新变化与已接受状态同时存在。只把当前图像送入下一轮，模型能看见蓝色和位置，却未必能恢复被先前采样削弱的玩具狗斑纹；只重新输入最初图像，又可能丢掉已接受的花瓶。连续编辑（iterative editing）需要在每轮保存原始依据、当前版本和修改记录，说明哪项保持、哪项被本轮覆盖。模型的输入能力与外部版本管理应分别描述。

FLUX.1 Kontext用图像和文字作为共同上下文，训练同一网络创建或编辑图像；本节采用2025年6月v2。它从FLUX.1文字生成检查点开始，让已有图像经冻结自编码器变成空间潜词元，再与待生成图像的潜词元沿序列方向拼接。文字词元进入文字流，与视觉词元通过注意力交换信息。不同于InstructPix2Pix按通道拼接同尺寸两幅潜图，序列拼接允许条件与目标有不同长宽和词元数量\scite{vision-kontext-2025}

设上下文图像为$\bm{Y}$，目标图像为$\bm{X}$，冻结编码器为$E$。编码并按既定规则分块后，目标和上下文分别有$n_x,n_y$个词元，每个词元投影到模型宽度$d$。训练时给目标潜变量加噪得到$\bm{z}(t)$，上下文$E(\bm{Y})$保持干净；视觉流读取长度$n_x+n_y$的序列。三维旋转位置编码为目标词元赋$(0,h,w)$，单张上下文赋$(1,h,w)$，其中首坐标只是区分图像来源的虚拟序号，不是拍摄时间或去噪步$t$。序列中相同空间坐标属于不同图像，由这个序号区分。

为使该网络学习“按指令从上下文得到目标”的关系，训练样本包含$\bm{Y},c,\bm{X}$，文字创建样本令$\bm{Y}=\varnothing$。以展平的目标潜变量$\bm{x}=E(\bm{X})$表示一条样本，抽标准高斯噪声$\bm{\epsilon}$与$t\in(0,1)$，采用
\begin{equation}
 \begin{aligned}
 \bm{z}(t)&=(1-t)\bm{x}+t\bm{\epsilon},\\
 L_{\mathrm{ctx}}&=
 \mathbb{E}\!\left[
 \left\lVert v_{\bm{\theta}}
  (\bm{z}(t),t,E(\bm{Y}),c)
   -(\bm{\epsilon}-\bm{x})\right\rVert_2^2
 \right]\text{。}
 \end{aligned}
 \label{eq:vision-kontext-flow}
\end{equation}
速度目标由该直线路径对$t$求导得到；误差只监督待生成目标的预测，不把上下文也训练为必须重生成的输出。这里$t=0$为干净目标、$t=1$为噪声，与\ref{sec:vision-models-flow-matching}节从噪声向数据递增的路径记法方向相反，因此采样沿$t$递减积分。论文按分辨率调整噪声步抽样分布，不应默认所有尺寸都使用同一个均匀日程。

每次从关系图像对或纯文字样本抽取小批量，编码固定条件与目标，构造加噪状态，平均速度误差并反传更新生成主干；自编码器保持冻结。对于文字创建样本，直接省略上下文图像词元。参数来自已有检查点，学习的是编辑关系，不需为这只玩具再做DreamBooth微调。论文没有公开全部关系数据构造、优化器和训练预算，复现不能仅靠本式推定那些配置。所保留的模型版本还须区别：报告中的pro在流训练后使用潜空间对抗蒸馏，dev使用引导蒸馏且专注图像到图像训练，不能把两个版本的训练安排和推理预算视为相同。

使用时只编码这轮实际给出的上下文，按目标尺寸抽噪声，再反复预测目标速度、沿递减日程更新目标潜变量，最后去掉条件词元并解码目标。以一维教学值说明方向，假设目标潜值$x=0.2$、噪声$\epsilon=1$，$t=0.75$时$z=0.8$，训练速度为0.8。若该处预测速度为0.6，步长$\Delta t=-0.25$给出$z_{\mathrm{next}}=0.8-0.25\times0.6=0.65$，而训练直线在$t=0.5$处为0.6。推理没有未知目标$x$可直接使用，例中的$x$只用于检验速度误差怎样进入状态。

上下文条件保持干净，不等于输出在对应区域逐像素不变。潜表示、全图生成与解码都可能改变细节，模型也没有因为序列拼接就自动保存完整编辑历史。论文的形式可扩展为多个图像序列，并给它们不同来源序号，但该版本实现主要使用单张上下文训练。因而本节逐轮以当前已接受图像为单图上下文；外部保存原始图与各轮结果以核验、回退，需要原始参考时先选定或重做相应轮次，不将任意多参考同时输入归为已验证能力。

沿玩具狗流程，保存$\bm{X}^{(0)}$，接受杯子替换得到$\bm{X}^{(1)}$，再要求花瓶变蓝得到$\bm{X}^{(2)}$，第三轮只移动花瓶得到$\bm{X}^{(3)}$。第三轮需同时检查花瓶位置、仍为蓝色、玩具狗耳朵和斑纹、未要求改变的背景。若第三轮颜色回退，应回到已接受的$\bm{X}^{(2)}$局部重做，而不是悄悄把新结果作为正确起点。每轮保存模型、提示、种子、尺寸和候选数，原始依据留在外部验收记录，不假定当前单图中仍有最初所有信息。

ImgEdit提供连续编辑的数据与评价入口，将早期全局要求的保持、依赖前文指代的新请求，以及回到历史版本的操作分别组织。例如“所有结果保持木质纹理”在后轮仍有效；“把它改蓝”需要明确前文的对象；“回到第一版”需要可回查版本。这些历史要求不等于独立的单轮配对，训练或评价时应保留相应上下文。本章用它解释验收材料，不把ImgEdit-E1的模型结构与Kontext混成一个实现\scite{vision-imgedit-2025}

随着轮数增加，编码、采样和解码误差可能累积，即使某轮没有要求重画主体。Kontext报告也保留了多轮造成伪影和指令遗漏的失败情况；角色或文字保持的示例不构成任意轮数的保证。评估应分别报告当前修改完成率、历次已接受要求的保持及相对原始主体的漂移，不能只给最终自然度。外部回退能限制错误被继续传播，却不能补回原始参考从未提供的细节。

\subsection{修改结果的核验与迭代}
\label{sec:vision-generation-editing-check}

结果比较应同时放入原图、编辑后图像、要求记录和允许变化区域。杯子仍在是目标未改，花瓶出现在另一侧是改错位置，玩具狗斑纹变化是保留失败，物体与桌面之间出现断边则属于边界问题。把这些错误混成一个“编辑质量”分数，不容易确定应修改哪一步。

对明确的保留掩码$\bm{M}$，可计算保留区平均变化
\begin{equation}
 E_{\mathrm{keep}}=
 \frac{\sum_{i,j,k}m_{i,j}
       |x'_{i,j,k}-x_{i,j,k}|}
 {c\sum_{i,j}m_{i,j}}\text{，}
 \label{eq:vision-edit-preservation}
\end{equation}
分母要求至少有一个保留像素。若整图照搬输入，$E_{\mathrm{keep}}=0$，但杯子没有变成花瓶；因此这项必须与修改完成度一起报告。允许全图色调变化时，逐像素保留误差也需要换成与要求相符的判断。

为了强制范围外保持，可合成$\bm{X}_{\mathrm{out}}=\bm{M}\odot\bm{X}+(1-\bm{M})\odot\bm{X}'$。它能保证所选保留像素相等，却可能把新花瓶的阴影切断。若允许羽化掩码，就要说明混合带也会改变，不能同时声称边界带严格不变。

迭代时应持续对照最初的原图和保留要求，避免每轮只与上一轮比较，逐步积累身份漂移。缩小范围、调整条件强度和重新生成都需记录轮数与候选预算；最终成功例应连同得到它的制作成本解释。

\section{观测场景的新视角合成}
\label{sec:vision-generation-novel-view}

\term{新视角合成}（novel view synthesis）根据已拍摄视图生成未拍摄视角的图像。这里的目标是呈现同一个观测场景，所需约束来自相机、跨视图对应和场景表示。若只用一幅图加生成先验构造可能的背面，观测支持更弱，相应空间内容创建将在\ref{chap:vision-multimodal-generation}章继续讨论。

\subsection{观测与相机条件的组织}
\label{sec:vision-generation-cameras}

绕静止玩具狗拍摄多张照片，记录每张图的内参、位置和朝向，并留出不参与场景拟合的测试视图。相机参数把二维像素连接到三维射线，几何估计可复用\ref{chap:vision-image-recognition}章的对应与重建方法。拍摄覆盖决定哪些表面有多视角证据；只有正面照片时，背面并没有因为模型能渲染而变成已观测区域。

原始NeRF和3D Gaussian Splatting都采用逐场景优化：为这一组图像拟合一个场景，再用任意给定相机渲染。换到另一只玩具通常需要重新建立表示。前馈场景预测则把跨场景规律放入预训练模型，使用成本和先验来源不同，不能混用“训练时间”的含义。

静态、曝光一致和相机准确是重要条件。玩具移动、光照变化或标定误差，会让同一空间点需要解释互相矛盾的像素。优化可能通过错误密度或外观来吸收矛盾，训练视图变好不一定意味着场景几何更准确。

\subsection{场景表示的拟合}
\label{sec:vision-generation-scene-fitting}

\Needspace{65mm}
\subsubsection{NeRF的射线采样与体渲染监督}
\label{sec:vision-generation-nerf}

\term{神经辐射场}（neural radiance field, NeRF）用网络表示空间中的体密度与随观察方向变化的颜色。位置$\bm{x}\in\mathbb{R}^3$决定非负密度$\sigma(\bm{x})$，位置和单位方向$\bm{d}$共同决定颜色$\bm{c}(\bm{x},\bm{d})\in[0,1]^3$。密度用于计算射线经过该处时有多少贡献会被截获，颜色则允许同一位置在不同方向呈现不同高光\scite{vision-nerf2020}

相机像素对应射线$\bm{r}(s)=\bm{o}+s\bm{d}$，其中$\bm{o}$是相机中心，$s$是沿射线的距离。给定近远界$s_{\mathrm{n}},s_{\mathrm{f}}$，体渲染颜色为
\begin{equation}
 \begin{aligned}
 \bm{C}(\bm{r})
 &=\int_{s_{\mathrm{n}}}^{s_{\mathrm{f}}}
    T(s)\sigma(\bm{r}(s))
    \bm{c}(\bm{r}(s),\bm{d})\,\mathrm{d}s,\\
 T(s)&=\exp\!\left(
     -\int_{s_{\mathrm{n}}}^{s}\sigma(\bm{r}(u))\,\mathrm{d}u
    \right)\text{。}
 \end{aligned}
 \label{eq:vision-nerf-continuous}
\end{equation}
$T(s)$是到达该位置之前的透射率，近处密度越大，远处颜色的可见贡献越小。若$s$有长度单位，$\sigma$具有其倒数单位，使指数无量纲；实际拟合常把场景归一化。有限远界之外的背景可另加$T(s_{\mathrm{f}})\bm{c}_{\mathrm{bg}}$。

数值计算沿射线选取$n$个有序样本，用每段近似常量的密度和颜色积分。设段长$\delta_i$，得到
\begin{equation}
 \begin{aligned}
 \alpha_i&=1-\exp(-\sigma_i\delta_i),&
 T_i&=\prod_{j<i}(1-\alpha_j),\\
 w_i&=T_i\alpha_i,&
 \widehat{\bm{C}}&=\sum_{i=1}^{n}w_i\bm{c}_i
       +T_{n+1}\bm{c}_{\mathrm{bg}}\text{。}
 \end{aligned}
 \label{eq:vision-nerf-discrete}
\end{equation}
空乘积取1。$w_i$包含“先前未被遮住”和“在当前段产生贡献”两项，故远处颜色不能只按自身$\alpha_i$相加。权重与剩余背景权重之和为1；若不加背景，射线穿过稀疏场景时颜色和可能不足。

取两段$\alpha_1=0.5,\alpha_2=0.75$，颜色分别为红$(1,0,0)$和蓝$(0,0,1)$，背景为白$(1,1,1)$。第一段权重0.5，第二段为$(1-0.5)0.75=0.375$，背景为$0.5\times0.25=0.125$，结果为$(0.625,0.125,0.5)$。如果忘记透射率，直接用0.5与0.75加权，就会重复计算被前景遮住的贡献。图\ref{fig:vision-ray-compositing}同时展示顺序和权重。

\begin{figure}[htbp]
 \centering
 % Two intervals, alpha .5 and .75; background transmittance .125.
\begin{tikzpicture}[x=1mm,y=1mm,font=\figurefont\small]
 \draw[BookMuted,line width=.9pt,->] (5,0)--(100,0) node[right] {$s$};
 \draw[BookInk,line width=.9pt] (5,-3)--(0,-6)--(0,6)--(5,3)--cycle;
 \node[anchor=south] at (4,9) {相机};
 \fill[BookBlue!12] (24,-7) rectangle (41,7);
 \draw[BookBlue,line width=.9pt] (24,-7) rectangle (41,7);
 \fill[BookTeal!15] (58,-7) rectangle (75,7);
 \draw[BookTeal,line width=.9pt] (58,-7) rectangle (75,7);
 \draw[BookMuted,line width=1pt] (91,-7)--(91,7);
 \node[anchor=south] at (32.5,8) {$\alpha_1=0.5$};
 \node[anchor=south] at (66.5,8) {$\alpha_2=0.75$};
 \node[anchor=north] at (32.5,-9) {$T_1=1$};
 \node[anchor=north] at (66.5,-9) {$T_2=0.5$};
 \node[anchor=south] at (91,8) {背景};
 \node[anchor=east] at (7,-31) {$w$};
 \fill[BookBlue!20] (10,-35) rectangle (50,-27);
 \fill[BookTeal!20] (50,-35) rectangle (80,-27);
 \fill[BookPaper] (80,-35) rectangle (90,-27);
 \draw[BookInk,line width=.7pt] (10,-35) rectangle (90,-27);
 \draw[BookInk,line width=.7pt] (50,-35)--(50,-27);
 \draw[BookInk,line width=.7pt] (80,-35)--(80,-27);
 \node at (30,-31) {$0.5$};
 \node at (65,-31) {$0.375$};
 \node[anchor=north] at (30,-37) {第一段};
 \node[anchor=north] at (65,-37) {第二段};
 \node[anchor=north] at (85,-37) {$0.125$};
\end{tikzpicture}

 \caption{两段射线的体渲染教学例。近处段先取得0.5的贡献权重，后段只能使用剩余0.5中的0.75，背景使用最终透射率0.125。下方长度条表示实际权重，三段和为1。颜色值人为设定，用于复算，图中形状不表示真实物体测量。}
 \label{fig:vision-ray-compositing}
\end{figure}

训练从观测图像抽取像素射线，经上述积分生成颜色，再与相应真实像素计算平方误差，对网络参数反向传播。原始NeRF用坐标的多频正余弦编码帮助表示细节，并分别限制密度只依赖位置、颜色还依赖方向。改变这些依赖会改变表示的多视图约束，不能把网络看成直接记忆每个训练像素编号。

空空间与已遮挡区域仍需网络查询，原方法因此使用粗细两级采样。粗网络在分段随机采样点上形成$w_i$，把权重归一化为沿射线的分段分布，再采更多点；细网络在粗点与新增点的并集上重新计算颜色。粗、细颜色都接受像素损失。粗权重全接近零时，需要数值稳定处理以形成有效采样分布；新增样本最终参与非均匀积分，不能只把高权重点的颜色作不带段长的平均。

渲染新相机时，重新为它的每个像素构造射线并查询网络，不是把旧图像中的某个标签搬过来。若一张图有$p$条射线，每条有$n$次网络查询，基本计算随$pn$增长，具体代价还取决于网络和采样层次。未被拍到的表面、镜面反射和错误相机都可能使新视图失败；图像损失本身没有唯一确定真实几何。

\subsubsection{3D Gaussian Splatting的投影、透明度合成与密度调整}
\label{sec:vision-generation-gaussians}

对每条射线反复查询网络的成本较高。3D Gaussian Splatting改用一组显式三维高斯表示场景，把它们投影到图像平面并合成像素。每个高斯保存位置$\bm{\mu}$、协方差$\bm{\Sigma}$、不透明度$o$和方向相关的颜色参数；原方法从相机估计得到的稀疏点云初始化这些对象，再逐场景优化\scite{vision-gaussian2023}

协方差控制高斯的方向与伸展范围。可写$\bm{\Sigma}=\bm{R}_{g}\bm{S}\bm{S}^{\top}\bm{R}_{g}^{\top}$，其中$\bm{R}_{g}$是旋转，$\bm{S}$含正尺度，以保持合法协方差。变换到目标相机后，在高斯中心附近线性化透视投影，二维协方差近似为
\begin{equation}
 \bm{\Sigma}_{2}
 =\bm{J}\bm{R}_{c}\bm{\Sigma}
    \bm{R}_{c}^{\top}\bm{J}^{\top}\text{，}
 \label{eq:vision-gaussian-projection}
\end{equation}
其中$\bm{R}_{c}$将世界方向转到相机坐标，$\bm{J}\in\mathbb{R}^{2\times3}$是该点的投影Jacobian。所得二维椭圆决定哪些像素受到这个高斯影响；这是局部投影近似，不能把三维高斯的透视投影无条件视为精确二维高斯。

若投影中心为$\bm{\mu}_{2}$，像素$\bm{u}$处的有效不透明度可概括为
\begin{equation}
 a_i(\bm{u})=o_i\exp\!\left[
 -\frac12(\bm{u}-\bm{\mu}_{2,i})^\top
 \bm{\Sigma}_{2,i}^{-1}(\bm{u}-\bm{\mu}_{2,i})
 \right]\text{。}
 \label{eq:vision-gaussian-alpha}
\end{equation}
实现还使用有限足迹、数值阈值及稳定处理。将覆盖该像素的高斯按近远顺序合成，颜色形式为
$\sum_i a_i\bm{c}_i\prod_{j<i}(1-a_j)$，必要时加剩余背景。它与式\eqref{eq:vision-nerf-discrete}具有相同的前向透明度合成结构，但$a_i$来自二维投影足迹及学习的不透明度，并非逐段密度积分的直接替换。

原方法按图像块组织高斯并排序，再在GPU上合成，降低空区域遍历与重复操作。排序使用高斯中心的深度来近似可见次序，没有对每个像素执行精确的表面求交。颜色采用球谐系数表示随方向变化的外观，即用一组定义在方向球面上的基函数组合颜色。拟合损失结合像素$L_1$与结构相似性差异，更新位置、尺度、旋转、不透明度和颜色参数。

固定数量的初始高斯未必能表达所有细轮廓。原方法依据投影位置的梯度等信号选择需要加密的区域，小尺度高斯可以克隆，覆盖过大的高斯可以分裂为多个较小对象，并删除贡献很弱或不合适的高斯。玩具耳缘若只有一个大高斯覆盖，局部颜色变化可能被混在一起；分裂提供更多自由度，之后仍需优化位置、颜色和不透明度。分裂动作本身不保证图像误差立刻降低，也不是测得了新的表面点。

新相机仍需重新投影、排序和合成。显式高斯减少了MLP查询，却增加了场景参数存储，像素被大量高斯覆盖时合成成本也会上升。速度比较要固定分辨率、场景规模和硬件。高斯集合没有天然的三角网格拓扑或物理材质，渲染自然不能直接推出资产适合碰撞、制造或精确编辑。

\subsection{目标视角的渲染与核验}
\label{sec:vision-generation-rendering-check}

保留视图应在场景拟合前确定，不用其像素误差调节该场景参数。给定保留相机，比较渲染与真实图像的PSNR、SSIM或LPIPS，同时检查对象轮廓、细节和反光。若相机标定使用了哪些图像信息，也应纳入协议；已知测试相机与未知相机的任务条件不同。

单个视图还可能遗漏不连续现象。沿连续相机路径渲染，可以观察遮挡变化、漂浮斑点和突然外观跳变。连续路径上没有真实图像的位置只能做一致性观察，不能计算缺失参考的真实误差。几何评价则需要点云、深度或其他测量依据，不能从连续且逼真的动画直接推出准确表面。

对于玩具狗背面，训练覆盖不足时，两种场景表示都可能把先验或错误拟合延伸过去。若用途只是某些可见视角的展示，应明确适用视角范围；若需要完整三维资产，则还需要背面采集与几何核验。动态场景的时间变化由\ref{chap:vision-multimodal-generation}章扩展。

\section{生成结果的评价}
\label{sec:vision-generation-evaluation}

创建、转换、编辑和新视角合成有不同参照。自由创建主要看条件和分布中的变化，转换与编辑还要对照输入，新视角则要对照观测场景。比较模型时固定提示、参考、分辨率、候选数量和修订预算，才能知道差异来自模型还是选择过程。

\subsection{条件记录与对象关系核验}
\label{sec:vision-generation-condition-check}

回到玩具狗与杯子的记录，逐项检查存在、数量、颜色归属和左右关系，主体适配还要检查参考外观。可以用检测器统计对象，用区域特征比较主体，用文字读取器检查杯子标识，但这些检查器本身也可能漏检或误读。自动分数应说明输入和输出，必要时抽样人工复核，不把工具的错误都算到生成器头上。

整体图文相似度适合衡量粗粒度相关性，无法独立证明复杂条件全部成立。若只汇总一个平均分，应同时保留关键分项及失败率，例如对象数量正确率和主体保持情况。多个属性绑定错误即使只影响小区域，也可能使整项任务失败。

\subsection{修改范围与输入对应核验}
\label{sec:vision-generation-change-check}

编辑需要同时检查“该改的改了”与“该留的留了”。原样返回玩具狗和杯子可以取得最小图像距离，却没有执行指令；把杯子变成花瓶同时改变狗的斑纹，则完成了局部目标但破坏保留条件。两类错误应分别记录。

转换还应考虑目标形式本身允许的变化。照片变成绘画后，像素距离大是预期现象，应结合内容对应与风格要求解释LPIPS等距离。新视角中，同一对象因投影而移动也不是编辑漂移；评价必须比较对应目标相机，而非直接与原相机图像逐像素相减。

\subsection{观测支持与制作成本}
\label{sec:vision-generation-costs}

一个可交付成果的成本包含生成、检查和修订。主体微调增加前期适配，空间控制增加条件制作，新视角增加采集、标定和场景拟合。一次渲染很快，不表示完成整个场景也只需一次前向；从百幅候选中选出一幅，也不表示单次生成稳定满足要求。

\section{本章小结}
\label{sec:vision-generation-summary}

同一只玩具狗在创建、转换、编辑和新视角合成中承担不同角色：它可能只是一个类别要求，也可能是必须保持的输入主体，或有多视角证据的真实场景。条件控制和主体适配提供不同约束，配对与循环监督建立不同对应，掩码和指令又分别规定可变范围与目标变化。评价需要跟随这些职责，不能仅凭自然外观或全局相似度判断任务完成。

新视角图像由场景表示、相机和渲染共同产生，遮挡与观测覆盖决定哪些内容受到支持。制作过程还要计入采样、筛选、适配与拟合成本。后续连续画面和多模态生成会增加时间、身份与同步约束，本章对条件、保持关系和观测证据的区分仍然适用。

现有约束还留下组合条件冲突、多主体绑定、文字和精确数量控制、连续编辑漂移，以及缺少视角时的观测忠实性等研究问题。改进应定位哪种对应或保持关系失效，再比较新增条件与训练依据是否减少了该类错误。若把生成图像用于训练，还需在真实目标任务上验证收益，不能由生成图像的自然程度推断学习效果。

